%%
% Copyright (c) 2017 - 2021, Pascal Wagler;
% Copyright (c) 2014 - 2021, John MacFarlane
%
% All rights reserved.
%
% Redistribution and use in source and binary forms, with or without
% modification, are permitted provided that the following conditions
% are met:
%
% - Redistributions of source code must retain the above copyright
% notice, this list of conditions and the following disclaimer.
%
% - Redistributions in binary form must reproduce the above copyright
% notice, this list of conditions and the following disclaimer in the
% documentation and/or other materials provided with the distribution.
%
% - Neither the name of John MacFarlane nor the names of other
% contributors may be used to endorse or promote products derived
% from this software without specific prior written permission.
%
% THIS SOFTWARE IS PROVIDED BY THE COPYRIGHT HOLDERS AND CONTRIBUTORS
% "AS IS" AND ANY EXPRESS OR IMPLIED WARRANTIES, INCLUDING, BUT NOT
% LIMITED TO, THE IMPLIED WARRANTIES OF MERCHANTABILITY AND FITNESS
% FOR A PARTICULAR PURPOSE ARE DISCLAIMED. IN NO EVENT SHALL THE
% COPYRIGHT OWNER OR CONTRIBUTORS BE LIABLE FOR ANY DIRECT, INDIRECT,
% INCIDENTAL, SPECIAL, EXEMPLARY, OR CONSEQUENTIAL DAMAGES (INCLUDING,
% BUT NOT LIMITED TO, PROCUREMENT OF SUBSTITUTE GOODS OR SERVICES;
% LOSS OF USE, DATA, OR PROFITS; OR BUSINESS INTERRUPTION) HOWEVER
% CAUSED AND ON ANY THEORY OF LIABILITY, WHETHER IN CONTRACT, STRICT
% LIABILITY, OR TORT (INCLUDING NEGLIGENCE OR OTHERWISE) ARISING IN
% ANY WAY OUT OF THE USE OF THIS SOFTWARE, EVEN IF ADVISED OF THE
% POSSIBILITY OF SUCH DAMAGE.
%%

%%
% This is the Eisvogel pandoc LaTeX template.
%
% For usage information and examples visit the official GitHub page:
% HTTPS://github.com/Wandmalfarbe/pandoc-latex-template
%%

% Options for packages loaded elsewhere
\PassOptionsToPackage{unicode}{hyperref}
\PassOptionsToPackage{hyphens}{url}
\PassOptionsToPackage{dvipsnames,svgnames*,x11names*,table}{xcolor}
%
\documentclass[
  10pt,
  dutch,
  paper=a4,
  ,captions=tableheading
]{scrbook}
\usepackage{amsmath,amssymb}
\usepackage{lmodern}
\usepackage{setspace}
\setstretch{1.2}
\usepackage{ifxetex,ifluatex}
\ifnum 0\ifxetex 1\fi\ifluatex 1\fi=0 % if pdftex
  \usepackage[T1]{fontenc}
  \usepackage[utf8]{inputenc}
  \usepackage{textcomp} % provide euro and other symbols
\else % if luatex or xetex
  \usepackage{unicode-math}
  \defaultfontfeatures{Scale=MatchLowercase}
  \defaultfontfeatures[\rmfamily]{Ligatures=TeX,Scale=1}
\fi
% Use upquote if available, for straight quotes in verbatim environments
\IfFileExists{upquote.sty}{\usepackage{upquote}}{}
\IfFileExists{microtype.sty}{% use microtype if available
  \usepackage[]{microtype}
  \UseMicrotypeSet[protrusion]{basicmath} % disable protrusion for tt fonts
}{}
\makeatletter
\@ifundefined{KOMAClassName}{% if non-KOMA class
  \IfFileExists{parskip.sty}{%
    \usepackage{parskip}
  }{% else
    \setlength{\parindent}{0pt}
    \setlength{\parskip}{6pt plus 2pt minus 1pt}}
}{% if KOMA class
  \KOMAoptions{parskip=half}}
\makeatother
\usepackage{xcolor}
\definecolor{default-linkcolor}{HTML}{A50000}
\definecolor{default-filecolor}{HTML}{A50000}
\definecolor{default-citecolor}{HTML}{4077C0}
\definecolor{default-urlcolor}{HTML}{4077C0}
\IfFileExists{xurl.sty}{\usepackage{xurl}}{} % add URL line breaks if available
\IfFileExists{bookmark.sty}{\usepackage{bookmark}}{\usepackage{hyperref}}
\hypersetup{
  pdftitle={Cyberboswachters},
  pdfauthor={Dams Tim},
  pdflang={nl},
  pdfsubject={security},
  pdfkeywords={cybersecurity, cryptografie, beveiliging, wifi, stuxnet, hacking},
  colorlinks=true,
  linkcolor=default-linkcolor,
  filecolor=default-filecolor,
  citecolor=default-citecolor,
  urlcolor=default-urlcolor,
  breaklinks=true,
  pdfcreator={LaTeX via pandoc with the Eisvogel template}}
\urlstyle{same} % disable monospaced font for URLs
\usepackage[paperheight=9.5in,paperwidth=6.7in,margin=0.7in]{geometry}
\usepackage{color}
\usepackage{fancyvrb}
\newcommand{\VerbBar}{|}
\newcommand{\VERB}{\Verb[commandchars=\\\{\}]}
\DefineVerbatimEnvironment{Highlighting}{Verbatim}{commandchars=\\\{\}}
% Add ',fontsize=\small' for more characters per line
\newenvironment{Shaded}{}{}
\newcommand{\AlertTok}[1]{\textcolor[rgb]{1.00,0.00,0.00}{\textbf{#1}}}
\newcommand{\AnnotationTok}[1]{\textcolor[rgb]{0.38,0.63,0.69}{\textbf{\textit{#1}}}}
\newcommand{\AttributeTok}[1]{\textcolor[rgb]{0.49,0.56,0.16}{#1}}
\newcommand{\BaseNTok}[1]{\textcolor[rgb]{0.25,0.63,0.44}{#1}}
\newcommand{\BuiltInTok}[1]{\textcolor[rgb]{0.00,0.50,0.00}{#1}}
\newcommand{\CharTok}[1]{\textcolor[rgb]{0.25,0.44,0.63}{#1}}
\newcommand{\CommentTok}[1]{\textcolor[rgb]{0.38,0.63,0.69}{\textit{#1}}}
\newcommand{\CommentVarTok}[1]{\textcolor[rgb]{0.38,0.63,0.69}{\textbf{\textit{#1}}}}
\newcommand{\ConstantTok}[1]{\textcolor[rgb]{0.53,0.00,0.00}{#1}}
\newcommand{\ControlFlowTok}[1]{\textcolor[rgb]{0.00,0.44,0.13}{\textbf{#1}}}
\newcommand{\DataTypeTok}[1]{\textcolor[rgb]{0.56,0.13,0.00}{#1}}
\newcommand{\DecValTok}[1]{\textcolor[rgb]{0.25,0.63,0.44}{#1}}
\newcommand{\DocumentationTok}[1]{\textcolor[rgb]{0.73,0.13,0.13}{\textit{#1}}}
\newcommand{\ErrorTok}[1]{\textcolor[rgb]{1.00,0.00,0.00}{\textbf{#1}}}
\newcommand{\ExtensionTok}[1]{#1}
\newcommand{\FloatTok}[1]{\textcolor[rgb]{0.25,0.63,0.44}{#1}}
\newcommand{\FunctionTok}[1]{\textcolor[rgb]{0.02,0.16,0.49}{#1}}
\newcommand{\ImportTok}[1]{\textcolor[rgb]{0.00,0.50,0.00}{\textbf{#1}}}
\newcommand{\InformationTok}[1]{\textcolor[rgb]{0.38,0.63,0.69}{\textbf{\textit{#1}}}}
\newcommand{\KeywordTok}[1]{\textcolor[rgb]{0.00,0.44,0.13}{\textbf{#1}}}
\newcommand{\NormalTok}[1]{#1}
\newcommand{\OperatorTok}[1]{\textcolor[rgb]{0.40,0.40,0.40}{#1}}
\newcommand{\OtherTok}[1]{\textcolor[rgb]{0.00,0.44,0.13}{#1}}
\newcommand{\PreprocessorTok}[1]{\textcolor[rgb]{0.74,0.48,0.00}{#1}}
\newcommand{\RegionMarkerTok}[1]{#1}
\newcommand{\SpecialCharTok}[1]{\textcolor[rgb]{0.25,0.44,0.63}{#1}}
\newcommand{\SpecialStringTok}[1]{\textcolor[rgb]{0.73,0.40,0.53}{#1}}
\newcommand{\StringTok}[1]{\textcolor[rgb]{0.25,0.44,0.63}{#1}}
\newcommand{\VariableTok}[1]{\textcolor[rgb]{0.10,0.09,0.49}{#1}}
\newcommand{\VerbatimStringTok}[1]{\textcolor[rgb]{0.25,0.44,0.63}{#1}}
\newcommand{\WarningTok}[1]{\textcolor[rgb]{0.38,0.63,0.69}{\textbf{\textit{#1}}}}

% Workaround/bugfix from jannick0.
% See HTTPS://github.com/jgm/pandoc/issues/4302#issuecomment-360669013)
% or HTTPS://github.com/Wandmalfarbe/pandoc-latex-template/issues/2
%
% Redefine the verbatim environment 'Highlighting' to break long lines (with
% the help of fvextra). Redefinition is necessary because it is unlikely that
% pandoc includes fvextra in the default template.
\usepackage{fvextra}
\DefineVerbatimEnvironment{Highlighting}{Verbatim}{breaklines,fontsize=\small,commandchars=\\\{\}}

\usepackage{longtable,booktabs,array}
\usepackage{calc} % for calculating minipage widths
% Correct order of tables after \paragraph or \subparagraph
\usepackage{etoolbox}
\makeatletter
\patchcmd\longtable{\par}{\if@noskipsec\mbox{}\fi\par}{}{}
\makeatother
% Allow footnotes in longtable head/foot
\IfFileExists{footnotehyper.sty}{\usepackage{footnotehyper}}{\usepackage{footnote}}
\makesavenoteenv{longtable}
% add backlinks to footnote references, cf. HTTPS://tex.stackexchange.com/questions/302266/make-footnote-clickable-both-ways
\usepackage{footnotebackref}
\usepackage{graphicx}
\makeatletter
\def\maxwidth{\ifdim\Gin@nat@width>\linewidth\linewidth\else\Gin@nat@width\fi}
\def\maxheight{\ifdim\Gin@nat@height>\textheight\textheight\else\Gin@nat@height\fi}
\makeatother
% Scale images if necessary, so that they will not overflow the page
% margins by default, and it is still possible to overwrite the defaults
% using explicit options in \includegraphics[width, height, ...]{}
\setkeys{Gin}{width=\maxwidth,height=\maxheight,keepaspectratio}
% Set default figure placement to htbp
\makeatletter
\def\fps@figure{htbp}
\makeatother
\setlength{\emergencystretch}{3em} % prevent overfull lines
\providecommand{\tightlist}{%
  \setlength{\itemsep}{0pt}\setlength{\parskip}{0pt}}
\setcounter{secnumdepth}{5}

\raggedbottom


% Make use of float-package and set default placement for figures to H.
% The option H means 'PUT IT HERE' (as  opposed to the standard h option which means 'You may put it here if you like').
\usepackage{float}
\floatplacement{figure}{H}

\usepackage{awesomebox}
\providecommand{\pandocbounded}[1]{#1}
\ifxetex
    % See issue HTTPS://github.com/reutenauer/polyglossia/issues/127
  \renewcommand*\familydefault{\sfdefault}
    % Load polyglossia as late as possible: uses bidi with RTL langages (e.g. Hebrew, Arabic)
  \usepackage{polyglossia}
  \setmainlanguage[]{}
\else
  \usepackage[main=dutch]{babel}
% get rid of language-specific shorthands (see #6817):
\let\LanguageShortHands\languageshorthands
\def\languageshorthands#1{}
\fi
\ifluatex
  \usepackage{selnolig}  % disable illegal ligatures
\fi

\title{Cyberboswachters}
\usepackage{etoolbox}
\makeatletter
\providecommand{\subtitle}[1]{% add subtitle to \maketitle
  \apptocmd{\@title}{\par {\large #1 \par}}{}{}
}
\makeatother
\subtitle{De beste digitale stropers zijn ook de beste cyberboswachters}
\author{Dams Tim}
\date{}



%%
%% added
%%

%
% language specification
%
% If no language is specified, use English as the default main document language.
%



%
% for the background color of the title page
%
\usepackage{pagecolor}
\usepackage{afterpage}
\usepackage{tikz}

%
% break urls
%
\PassOptionsToPackage{hyphens}{url}

%
% When using babel or polyglossia with biblatex, loading csquotes is recommended
% to ensure that quoted texts are typeset according to the rules of your main language.
%
\usepackage{csquotes}

%
% captions
%
\definecolor{caption-color}{HTML}{777777}
\usepackage[font={stretch=1.2}, textfont={color=caption-color}, position=top, skip=4mm, labelfont=bf, singlelinecheck=false, justification=raggedright]{caption}
\setcapindent{0em}

%
% blockquote
%
\definecolor{blockquote-border}{RGB}{221,221,221}
\definecolor{blockquote-text}{RGB}{119,119,119}
\usepackage{mdframed}
\newmdenv[rightline=false,bottomline=false,topline=false,linewidth=3pt,linecolor=blockquote-border,skipabove=\parskip]{customblockquote}
\renewenvironment{quote}{\begin{customblockquote}\list{}{\rightmargin=0em\leftmargin=0em}%
\item\relax\color{blockquote-text}\ignorespaces}{\unskip\unskip\endlist\end{customblockquote}}

%
% Source Sans Pro as the de­fault font fam­ily
% Source Code Pro for monospace text
%
% 'default' option sets the default
% font family to Source Sans Pro, not \sfdefault.
%
\ifnum 0\ifxetex 1\fi\ifluatex 1\fi=0 % if pdftex
    \usepackage[default]{sourcesanspro}
  \usepackage{sourcecodepro}
  \else % if not pdftex
    \usepackage[default]{sourcesanspro}
  \usepackage{sourcecodepro}

  % XeLaTeX specific adjustments for straight quotes: HTTPS://tex.stackexchange.com/a/354887
  % This issue is already fixed (see HTTPS://github.com/silkeh/latex-sourcecodepro/pull/5) but the
  % fix is still unreleased.
  % TODO: Remove this workaround when the new version of sourcecodepro is released on CTAN.
  \ifxetex
    \makeatletter
    \defaultfontfeatures[\ttfamily]
      { Numbers   = \sourcecodepro@figurestyle,
        Scale     = \SourceCodePro@scale,
        Extension = .otf }
    \setmonofont
      [ UprightFont    = *-\sourcecodepro@regstyle,
        ItalicFont     = *-\sourcecodepro@regstyle It,
        BoldFont       = *-\sourcecodepro@boldstyle,
        BoldItalicFont = *-\sourcecodepro@boldstyle It ]
      {SourceCodePro}
    \makeatother
  \fi
  \fi

%
% heading color
%
\definecolor{heading-color}{RGB}{40,40,40}
\addtokomafont{section}{\color{heading-color}}
% When using the classes report, scrreprt, book,
% scrbook or memoir, uncomment the following line.
%\addtokomafont{chapter}{\color{heading-color}}

%
% variables for title, author and date
%
\usepackage{titling}
\title{Cyberboswachters}
\author{Dams Tim}
\date{}

%
% tables
%

\definecolor{table-row-color}{HTML}{F5F5F5}
\definecolor{table-rule-color}{HTML}{999999}

%\arrayrulecolor{black!40}
\arrayrulecolor{table-rule-color}     % color of \toprule, \midrule, \bottomrule
\setlength\heavyrulewidth{0.3ex}      % thickness of \toprule, \bottomrule
\renewcommand{\arraystretch}{1.3}     % spacing (padding)


%
% remove paragraph indention
%
\setlength{\parindent}{0pt}
\setlength{\parskip}{6pt plus 2pt minus 1pt}
\setlength{\emergencystretch}{3em}  % prevent overfull lines

%
%
% Listings
%
%


%
% header and footer
%
\usepackage{fancyhdr}

\fancypagestyle{eisvogel-header-footer}{
  %\fancyhead{}
  \fancyfoot{}
  %\lhead[]{Cyberboswachters}
  %\chead[]{}
  %\rhead[Cyberboswachters]{}
  %\lfoot[\thepage]{Dams Tim}
  \lfoot[\thepage]{}
  \cfoot[]{}
  % \rfoot[Dams Tim]{\thepage}
  \rfoot[]{\thepage}
  \renewcommand{\headrulewidth}{0.4pt}
  \renewcommand{\footrulewidth}{0.4pt}
}
\pagestyle{eisvogel-header-footer}

%%
%% end added
%%

\begin{document}

%%
%% begin titlepage
%%
\begin{titlepage}
\newgeometry{top=2cm, right=4cm, bottom=3cm, left=4cm}
\tikz[remember picture,overlay] \node[inner sep=0pt] at (current page.center){
\includegraphics[width=\paperwidth,height=\paperheight]{assets/cover.png}};
\newcommand{\colorRule}[3][black]{\textcolor[HTML]{#1}{\rule{#2}{#3}}}
\begin{flushleft}
\noindent
\\[-1em]
\color[HTML]{5F5F5F}
\makebox[0pt][l]{\colorRule[435488]{1.3\textwidth}{4pt}}
\par
\noindent

% The titlepage with a background image has other text spacing and text size
{
  \setstretch{2}
  \vfill
  \vskip -8em
  \noindent {\huge \textbf{\textsf{Cyberboswachters}}}
    \vskip 1em
  {\Large \textsf{De beste digitale stropers zijn ook de beste
cyberboswachters}}
    \vskip 2em
  \noindent {\Large \textsf{Dams Tim} \vskip 0.6em \textsf{}}
  \vfill
}


\end{flushleft}
\end{titlepage}
\restoregeometry

%%
%% end titlepage
%%

\frontmatter


{
\hypersetup{linkcolor=}
\setcounter{tocdepth}{1}
\tableofcontents
\newpage
}
\mainmatter
\chapter{Introductie}\label{introductie}

Dit handboek wordt gebruikt als basis-cursus binnen de opleidingen
elektronica-ict en toegepaste informatica van de AP Hogeschool. Het
heeft als doel een breed overzicht te geven van de \emph{wereld van de
cybersecurity}. Het doel is niet alomvattend te zijn, maar vooral om
mensen zin te geven om meer te weten te komen over deze belangrijke én
boeiende wereld. Dit boek gaat ervan uit dat de lezer minimale
voorkennis heeft inzake ICT, zo wordt er verwacht dat een basiskenis
netwerk-technologie aanwezig is (denk aan termen zoals IP-adressering,
routering en firewalls).

\section{Boswachters en stropers}\label{boswachters-en-stropers}

De subtitel,\emph{``De beste digitale stropers zijn ook de beste
cyberboswachters''}, van dit handboek verdient een extra woordje uitleg.
Dit boek zal zeker niet de obscure wereld van de hackers en
cybercriminelen verheerlijken, integendeel. We willen echter wel tonen
hoe de stropers, de slechterikken in dit boek, te werk gaan, opdat het
ons zo een beter beeld geeft waar we ons tegen moeten beschermen als
goede cyberboswachters.

\section{Dankwoord}\label{dankwoord}

Dank aan (soms al oud-)studenten Ernie de Magtige, Dimitriy
Vassilchenko, Jasper Van Meel en Erik Van Dyck om als externe fact- en
spellingcheckers aardig wat (gênante) foutjes en typos te ontdekken.

Een oprechte dank aan Stefaan Somerling (RealDolmen) voor de feedback op
het hoofdstuk omtrent GDPR. Geef gerust een sein als u, conform de GDPR
wetgeving, liever uw naam niet in dit document ziet staan ;).

Een dikke fist bump voor collega en \emph{brother from another mother}
Koen Van Eyken, die een deel van z'n paasvakantie opofferde om met een
kritische blik door dit boek te gaan! Uiteraard ook een dikke merci aan
de collega's die feedback op dit document hebben gegeven, en dan zeker
aan Michael Boeynaems en Serge Horsmans.

\section{Waarschuwingen}\label{waarschuwingen}

Dit brengt ons automatisch bij een belangrijke waarschuwing: je zal in
dit handboek geregeld technieken en tools tegenkomen die verregaande
gevolgen voor derden én jezelf kunnen hebben indien ze misbruikt worden.
We tonen deze zaken enkel vanuit het standpunt dat zonet werd besproken:
digitale stropers zijn de beste cyberboswachters. Het is belangrijk dat
je begrijpt dat je deze tools en technieken NOOIT of te NIMMER voor
kwade doeleinden mag gebruiken. Het is zelfs zo dat het gebruik van veel
van deze zaken strafrechtelijke gevolgen kunnen hebben, inclusief boetes
tot zelfs jarenlange opsluiting.

Indien je dus deze tools of technieken wenst te gebruiken dan mag dit
enkel op volgende doelwitten:

\begin{itemize}
\tightlist
\item
  Zaken waar jij de eigenaar van bent. Denk aan servers, webpagina's,
  gebruikers, etc.
\item
  Op doelwitten waar je geen eigenaar van bent, maar waar je wel
  expliciete toestemming voor hebt gekregen.
\end{itemize}

\subsection{Opletten met experimenten}\label{opletten-met-experimenten}

Indien je van plan bent om met bepaalde tools te experimenteren, hou dan
rekening met volgende tips:

\begin{itemize}
\tightlist
\item
  Werk zoveel mogelijk met virtuele machines en test nooit op een
  \emph{live systeem} indien daar geen erg goede reden voor is.
\item
  Hou er rekening mee dat bepaalde tools permanente ``gaten'' slagen in
  je systemen. Deze kunnen dus na (of tijdens!) je experimenten door
  digitale stropers misbruikt worden.
\end{itemize}

\chapter{Wordt het erger?}\label{wordt-het-erger}

Het laatste decennium jaar is de (cyber)security wereld erg veranderd.
Ze doet dit niet omdat ze daar zin in heeft, maar wel als antwoord op
wat er gebeurt in de wereld omtrent cyberaanvallen, geopolitieke
situaties, etc.

Het is altijd een kat en muisspel, waarbij de boswachters helaas bijna
altijd zullen achterlopen op de stropers. De kwaadwillige hackers hoeven
maar één klein gaatje te vinden in je peperdure beveiliging en ze zijn
binnen. Terwijl jij als boswachter wel aan alles moet (proberen te)
denken. Dat het erger wordt, is eigenlijk daarom bijna automatisch een
evidentie. De wereld van de beveiliging gebeurt per opbod en hoe beter
de boswachters het bos kunnen verdedigen, hoe complexer de technieken
zullen worden die de stropers hanteren.

In de volgende secties geven we een klein historisch overzicht van
enkele belangrijke, interessante of spectaculaire gebeurtenissen die
hebben plaatsgevonden in de cyberwereld de voorbije jaren. Dit laat ons
toe om enerzijds enkele begrippen te duiden, anders om de vraag ``wordt
het erger?'' te beantwoorden.

\begin{noteblock}
We gebruiken de term cyber om alles aan te duiden dat zich afspeelt op
de digitale snelweg, namelijk het Internet, de \emph{cloud}, het www, en
alle synoniemen en aanverwanten.
\end{noteblock}

\section{Het voorbije decennium}\label{het-voorbije-decennium}

We zullen geregeld terug in de tijd gaan om te bekijken hoe bepaalde
cyberverdedigings- en aanvalstechnieken zijn ontstaan, maar in dit
hoofdstuk gaan we enkel terug tot 2010. Het jaar waarin Mark Zuckerberg,
CEO van Facebook, door Time Magazine tot persoon van het jaar werd
uitgeroepen en ook waarin Apple de eerste IPad tablet aan het publiek
toonde. Het was helaas ook het jaar waarin het boorplatform, Deepwater
Horizon, voor één van de ergste milieurampen ooit zorgde, alsook het
jaar dat 33 mijnwerkers anderhalve maand lang 700 meter diep opgesloten
zaten in een mijn. Voor ons, als toekomstige cyberboswachters, is echter
de meest cruciale gebeurtenis het ontdekken van \textbf{Stuxnet}.

\subsection{Voor 2010}\label{voor-2010}

\begin{figure}
\centering
\pandocbounded{\includegraphics[keepaspectratio,alt={``Het aardse paradijs met de zondeval van Adam en Eva'' door Peter Paul Rubens en Jan Brueghel (de Oude). Duidelijk gebaseerd op een wereld van voor 2010.}]{intro/eden.jpg}}
\caption{``Het aardse paradijs met de zondeval van Adam en Eva'' door
Peter Paul Rubens en Jan Brueghel (de Oude). Duidelijk gebaseerd op een
wereld van voor 2010.}
\end{figure}

Voor het ontdekken van Stuxnet in 2010 leek de cyberbeveiligingswereld
wel een wereld van (relatieve) peis en vree. Het ergste dat kon gebeuren
was dat je computer een virusje opdeed waardoor je mogelijk wat data, en
vooral veel werkuren, kwijt raakte. Virussen, wormen en spam waren
alomtegenwoordig maar waren eigenlijk niet meer dan luizen in de pels
van de eindgebruikers. Tuurlijk, er werd een grondig gevloekt als een
virus je foto's verwijderde of wanneer een worm zichzelf verspreidde
naar je contacten - denk maar aan het \emph{ILOVEYOU} virus dat als een
soort \emph{social engineer} mensen deed geloven dat ze een potentiële
liefdesmatch hadden waarop ze vol spanning de bijlage openden en zo de
worm \emph{in huis haalden}. Of wat te denken van de \emph{Conficker}
worm die in 2008 meer dan 3 miljoen systemen kon besmetten en telkens de
update-services van het Windows toestel uitschakelde. Uiteraard, dat was
irritant, maar menig mens zou maar al te graag terug naar die tijd
willen gaan als daarmee ransomware, botnets en door overheden
gesponsorde cyberaanvallen onbestaande zouden zijn.

\begin{noteblock}
De termen worm, virus en malware zullen geregeld door elkaar worden
gebruikt in deze cursus. Alle drie betekenen net niet hetzelfde, maar
toch:

\begin{itemize}
\tightlist
\item
  Een \textbf{virus} is een kwaadaardig programma dat, eens het op een
  computer of apparaat staat, digitale schade kan aanbrengen.
\item
  Een \textbf{worm} daarentegen is ook een virus, maar eentje dat
  zichzelf kan \emph{voortplanten} naar andere systemen, iets wat een
  virus niet kan. Een worm kan zichzelf dus verspreiden zonder
  menselijke hulp.
\item
  \textbf{Malware} is een overkoepelende term voor alle programma's en
  code die digitale schade aan een systeem of netwerk brengen. Virussen
  en wormen behoren dus tot deze groep.
\end{itemize}

Naast deze drie termen zullen we ook nog enkele andere malware types
tegenkomen zoals Adware, spyware, spam, trojans, backdoors en rootkit.
Wanneer nodig zullen we deze termen toelichten. Onthoud nu alvast dat
malware de algemene naam is voor alle malafide software.
\end{noteblock}

\subsection{2010: En dan verscheen
Stuxnet}\label{en-dan-verscheen-stuxnet}

Virussen konden computers (tijdelijk) onbruikbaar maken. OK, lastig,
maar niet het einde van de wereld. De Stuxnet worm die in 2010 plots op
de radar verscheen kon dat potentieel wel: de wereld beëindigen! Het was
een worm die op maat was gemaakt om zogenaamde \textbf{PLCs}
(\emph{Programmable Logic Controllers}) te controleren. Héél veel van
onze kritische infrastructuur is geautomatiseerd met behulp van deze
PLCs , krachtige apparaten die machines kunnen aansturen zoals liften,
robotarmen aan assemblagelijnen, zuiveringsinstallaties en zelfs
kernreactors. PLCs vormen de interface tussen machines (hardware) en de
computers (met software op) die de operatoren gebruiken om deze machines
opdrachten te geven. Operators kunnen op hun systemen de machines
bedienen en te allen tijde controleren of deze naar behoren werken.

Wat Stuxnet deed was zich tussen de hardware en de software nestelen.
Als een worm besmette het laptops en computers en zocht het of er
PLC-bedieningssoftware aanwezig was op het toestel. Als dat het geval
was dan plaatste het zichzelf er tussen in zodat het:

\begin{enumerate}
\def\labelenumi{\arabic{enumi}.}
\tightlist
\item
  de machines opdrachten kon geven zonder dat de operator dit wist.
\item
  aan de operator kon vertellen dat alles in orde was, terwijl in
  realiteit de PLC mogelijk desastreuze opdrachten aan de hardware gaf.
\end{enumerate}

\begin{figure}
\centering
\pandocbounded{\includegraphics[keepaspectratio,alt={Werking stuxnet.}]{intro/stuxnet.png}}
\caption{Werking stuxnet.}
\end{figure}

We moeten er geen tekeningetje bij maken wat de effecten kunnen zijn
indien een Stuxnet variant bewust werd ingezet om bijvoorbeeld de
machinerie in een kerncentrale of waterdam te saboteren. Zonder in
details van Stuxnet in te gaan is het duidelijk dat het verschijnen van
dit virus een ommekeer betekende qua potentiële gevolgen van een
cyberaanval: mensenlevens stonden plots op het spel, niet enkel onze
dierbare e-mails en digitale vakantiefoto's.

\begin{noteblock}
Stuxnet bleek een worm te zijn die probeerde een specifiek soort
centrifuge te saboteren: namelijk de centrifuges die Iran gebruikte om
uranium te verrijken. De inlichtingendiensten van de Verenigde Staten en
Israël zouden klaarblijkelijk Stuxnet hebben ontworpen om zo dit
verrijkkingsproces van uranium in Iran te dwarsbomen.
\end{noteblock}

\begin{tipblock}
Wens je meer te weten over Stuxnet? Dan raden we je de uitstekende, in
2016 verschenen, documentaire ``Zero Days'' door Alex Gibney aan. Deze
documentaire geeft een zeer ontluisterend én boeiend beeld over de
zoektocht naar de oorsprong van het Stuxnet virus.
\end{tipblock}

\subsection{2013: Onze privacy te grabbel
gegooid}\label{onze-privacy-te-grabbel-gegooid}

Surfen op het Internet was altijd een risico. We wisten in 2013 al lang
dat onze datapakketjes over tientallen apparaten doorheen het Internet
worden getransporteerd om zo tot bij hun doel te geraken. Cookies
volgden al geregeld onze stappen en ook de aanbevelingen van Amazon
waren soms akelig accuraat. Dat was de prijs die we moesten betalen om
de ontelbare bronnen van het Internet te kunnen gebruiken. En als je
echt wat meer privacy nodig had, dan schakelde je de ``incognito'' modus
in van je browser. Maar ook dan was je je ervan bewust dat zelfs je ISP
(\emph{Internet service provider}) en je doel konden meekijken. Hoe erg
kan dat zijn?

Die gedachte leefde bij veel mensen tot dan: het Internet was een
nuttige tool en je wist dat er kon meegekeken worden door \emph{derden}
als ze dat echt wilden. Maar zou dit nu echt consequent en op grote
schaal gebeuren? Neen toch?!

Dat glazen huisje werd in 2013 hardhandig aan diggelen geslagen door
twee belangrijke personen:

\begin{itemize}
\tightlist
\item
  \textbf{Edward Snowden}: als (contractueel ingehuurde)
  systeembeheerder bij de NSA, de Amerikaanse inlichtingendienst gericht
  op elektronische spionage, had hij toegang tot het doen en laten van
  de dienst. Snowden lekte een grote hoeveelheid top-secret documenten
  naar de bevolking die onder andere het \emph{PRISM}-programma uit de
  doeken deed. Hieruit bleek dat de NSA een complexe samenwerking had
  met enkele grote Amerikaanse Internetbedrijven (o.a. Microsoft,
  Google, Facebook, Apple, etc.) die zogenaamde \emph{taps} op hun
  systemen hebben staan zodat de NSA ``kan meeluisteren'' op de
  netwerken.
\item
  \textbf{Julian Assange}: er zijn altijd klokkenluiders
  (\emph{whistleblowers}) geweest die om persoonlijke overtuigingen
  vonden dat bepaalde informatie met het publiek moesten worden gedeeld.
  Zeker in dictatoriale middens kan dit levensgevaarlijk zijn. Aan het
  eind van de 20e eeuw kregen deze klokkenluiders echter een krachtig
  apparaat om hun nieuws veilig te verspreiden: het Internet. Julian
  Assange besefte dit en richtte daarom in 2006 Wikileaks op. Een site
  waar klokkenluiders op een anonieme manier documenten konden
  \emph{lekken}. Tot 2010 zijn zo enkele erg controversiële documenten
  met de wereld gedeeld die soms verregaande diplomatieke of
  geopolitieke gevolgen hadden. Het lot van Assange is nog steeds
  onduidelijk: momenteel wordt hij in het Verenigd Koninkrijk in
  hechtenis gehouden en voert hij hevig strijd (via de rechtbank) om
  niet uitgeleverd te worden aan de Verenigde Staten.
\end{itemize}

\begin{cautionblock}
Het is niet evident om over mensen als Assange en Snowden te praten
zonder in een controversiële modderpoel te geraken. Voor de één is
Snowden een held, voor de ander een verrader. In dit boek trachten we
een objectief beeld, zonder waardeoordelen, te geven. Wat wel buiten
kijf staat is dat zowel Snowden als Assange de wereld getoond hebben dat
er in de duistere wandelgangen van de veiligheidsdiensten zaken gebeuren
waar ``normale stervelingen'' zelden weet van hebben.
\end{cautionblock}

Het waren vooral de onthullingen van Snowden die bij velen de oogkleppen
deed afvallen. Op het Internet ben je hoegenaamd niét anoniem. Grote (en
kleine) mogendheden en privéfirma's kunnen ons doen en laten op het
Internet bekijken, bewaren en analyseren. Privacy en het Internet zijn
een \emph{oxymoron}: een contradictorische term zoals zwarte sneeuw.
Wanneer je je op het Internet begeeft, op welke manier ook, gooi je een
deel van je privacy te grabbel. En dat is iets dat helaas de voorbije
tien jaar er niet op verbeterd is (alhoewel het \textbf{Tor} netwerk met
z'n \emph{onion routing} toch wel een stevige extra privacy laag kan
aanbieden).

\begin{tipblock}
Een interessant debat dat altijd opduikt bij deze problematiek is de
``Ik heb toch niets te verbergen''-houding. In het kleine maar fijne
boekje ``Je hebt wél iets te verbergen'' van onderzoeksjournalisten
Maurits Martijn en Dimitri Tokmetzis (ISBN 9789082821611) wordt
onherroepelijk brandhout gemaakt met deze stelling. Finaal zijn we
volledig afhankelijk van de online diensten die we gebruiken en wat ze
met onze data nu én belangrijker, in de toekomst zullen doen.
Privé-informatie over jou die nu ogenschijnlijk ongevaarlijk lijkt, kan
dat potentieel in de toekomst wel zijn wanneer normen, waarden of wetten
veranderen.
\end{tipblock}

\subsection{\texorpdfstring{2014: Hoe veilig is \emph{the cloud}
?}{2014: Hoe veilig is the cloud ?}}\label{hoe-veilig-is-the-cloud}

In 2014 vond er een controverse plaats die vooral de puberende tiener
zich zal herinneren: \emph{``The fappening''}. Deze naam laat weinig aan
de verbeelding over en dekt de lading goed. In het jaar dat Oekraïne
Russische legers zag binnenrollen om de Krim (terug) in te palmen,
verschenen er duizenden privéfoto's van een honderdtal
\emph{celebrities} op het Internet. Deze,vaak weinig verhullende,
privéfoto's hadden kwaadwillige hackers van de Apple iCloud accounts van
de slachtoffers gestolen en vervolgens gepubliceerd via Reddit, 4chan,
etc. De foto's verspreidden zich als een lopend vuurtje en de
slachtoffers konden enkel toezien hoe hun in de privésfeer opgenomen
beelden door miljoenen mensen werden gedownload.

De beroemdheden hadden hun foto's in de cloud-opslag van Apple bewaard
zoals op dat moment duizenden gebruikers ook al deden. Deze dienst zorgt
ervoor dat je je als eindgebruiker van eender waar aan je foto's kan,
daar ze ``in the cloud'' staan, wat door de spectaculaire groei van de
smartphones (en iPhones in dit geval) een populair concept was geworden.
Ook diensten zoals Dropbox, Onedrive (toen nog SkyDrive) toonden aan dat
er een grote vraag was naar online opslag van (privé-)informatie.

Om dergelijke diensten te gebruiken dien je natuurlijk een veilig
wachtwoord te hebben, daar eender wie, van eender waar, kan proberen
zich een weg naar je data te verkrijgen door jouw wachtwoord te raden.
In 2014 waren concepten zoals 2-factor-authentication (2FA) en
biometrische beveiliging (meer daarover in hoofdstuk 5) nog niet zo
populair en dus was je geheime wachtwoord je enige bescherming tegen
onrechtmatige toegang tot je data.

De diefstal van de foto's werd echter vergemakkelijkt door twee redenen:

\begin{itemize}
\tightlist
\item
  Sommige slachtoffers gebruikten makkelijk te raden, of snel te
  bruteforcen wachtwoorden.
\item
  Anderen hadden de beveiligingsvragen ter goeder trouw ingediend. Veel
  diensten stellen een extra beveiligingsvraag (zoals ``Wat is de naam
  van je eerste huisdier?'' of ``Op welke school zat je vader?'') die ze
  kunnen gebruiken om jouw identiteit te verifiëren indien je je eigen
  wachtwoord vergeten bent en je dit wenst te resetten. Wanneer jij of
  ik dit soort vragen invullen dan is dit 90\% van de tijd informatie
  die nergens te vinden valt, enkel in de grijze massa in je hoofd en
  misschien in een gênant dagboekje uit je jeugd. Bij de Amerikaanse
  beroemdheden wiens iCloud werd gehackt, is dat niet zo. Hun leven kan
  volledig gereconstrueerd worden aan de hand van de ontelbare
  interviews die ze al hebben gegeven. Gegarandeerd dat ooit een
  interviewer al heeft gevraagd wat de kleur van zijn of haar eerste
  auto was, of in welk dorp hij of zij is opgegroeid.
\end{itemize}

\begin{warningblock}
Het is een verkeerde reflex om in dit soort zaken ogenblikkelijk aan
\emph{victim shaming} te doen (kijk maar naar het recentere voorval
waarbij enkele Bekende Vlamingen het slachtoffer waren van catfishing).
Ieder doet en laat wat hij wenst in z'n privésfeer - daarom heet het ook
privé - maar het is belangrijk te beseffen dat als je zaken \emph{in the
cloud} bewaard, de kans bestaande is dat iets of iemand er ooit
onrechtmatige toegang tot zal krijgen. U weze gewaarschuwd.
\end{warningblock}

\begin{tipblock}
Niemand verplicht je om op de beveiligingsvragen een eerlijk antwoord te
geven. Er bestaat geen leugendetector in die systemen of een mama die op
je vingers komt tikken. Lieg er dus op los, maar onthoud uiteraard je
antwoord. Dankzij die beveiligingsvragen hou je echter een stok achter
de hand moest je je wachtwoord vergeten zijn.
\end{tipblock}

\subsection{Ook in 2014: als landen
vechten}\label{ook-in-2014-als-landen-vechten}

Het was te verwachten dat ook op cyberniveau er ooit een
cyber-wapenwedloop zou starten zoals, helaas, ook in de echte wereld
plaatsvindt. Daar waar landen voorheen nog vooral defensieve
cybermogelijkheden hadden, begon halverwege het vorige decennium dit
arsenaal meer en meer uitgebreid te worden met offensieve cyberwapens.
Het was met andere woorden wachten op de eerste cyberaanvallen door een
soevereine staat op een andere staat.

Het probleem met cyberaanvallen is dat het als aangevallen mogendheid
ongelooflijk moeilijk is om juist te reageren:

\begin{enumerate}
\def\labelenumi{\arabic{enumi}.}
\tightlist
\item
  De bron van de aanval identificeren is soms onmogelijk. Wie weet lijkt
  de aanval te komen vanuit land X, terwijl het eigenlijk land Y is dat
  gewoon de aanval via land X routeert. En het laatste dat je natuurlijk
  wil doen is het verkeerde land beschuldigen (of aanvallen).
\item
  Wat als een hacker, zonder staatsinmengingen, beslist om op eigen
  houtje een cyberaanval te initiëren. Is het land van herkomst van die
  aanvaller dan verantwoordelijk voor de geleden schade?
\item
  De ``schade'' van een cyberaanval is niet altijd duidelijk. Wanneer
  een raket op een stad wordt afgestuurd is dat een duidelijk \emph{act
  of aggression} en is er grote kans op een militair conflict (indien de
  diplomatieke weg geen soelaas brengt). Maar is een cyberaanval, zoals
  een denial-of-service (DoS, zie later), genoeg reden om de oorlog
  buiten het cyberdomein te escaleren? Zijn er eigenlijk \emph{rules of
  engagement}? Is er een conventie van Genève? Neen op dit alles. Daar
  cyberaanvallen zo moeilijk te traceren en identificeren zijn, blijft
  het allemaal het betere nattevingerwerk.
\end{enumerate}

Een bewuste cyberaanval op een soevereine staat is al een enkele keren
gebeurd (denk maar aan de cyberaanvallen in 2007 van Rusland op Estland)
maar in 2014 was er helaas een nieuwe primeur (Stuxnet krijgt eigenlijk
die primeur, maar dat werd pas veel later bevestigd). Ook nu gaan we
naar Hollywood, niet om de beroemdheden en hun privé-kiekjes, maar
omwille van een film waar een andere staat niet om kon lachen.

In 2014 zou Sony een nieuwe film uitbrengen, getiteld ``The Interview'',
met James Franco en Seth Rogan. In deze komedie worden twee journalisten
gevraagd om tijdens hun interview van een \emph{fictieve}
Noord-Koreaanse dictator hem te doden. De vergelijkingen met de
Noord-Koreaanse leider Kim Jong-un waren voor iedereen duidelijk. Dat
Noord-Korea op de tenen zou zijn getrapt, stond dan ook in de sterren
geschreven.

De daaropvolgende gebeurtenissen zijn nog steeds niet uitgeklaard, maar
vast staat wel dat een selecte groep hackers toegang had gekregen tot
confidentiële data op de servers van Sony en deze vervolgens op het
Internet lekte. De financiële schade voor Sony was enorm. De daders
wisten onuitgebrachte films, e-mails, salarissen en andere
privé-informatie te stelen. Heel lang werd gedacht dat de groeperingen
die achter de aanval stonden door Noord-Korea waren aangestuurd. Echter,
tot op de dag van vandaag heeft men dat niet kunnen bevestigen. Zo zijn
er ook experts die denken in het bewijsmateriaal de hand van Rusland
en/of China te zien. Kortom, zoals eerder gezegd, cyberaanvallen zijn
verdomd lastig om in kaart te brengen.

Als het toch Noord-Korea zou zijn geweest - wat nog steeds meer dan
mogelijk is - dan hebben ze hiermee dus een dubieuze primeur: een
soeverein land voert een cyberaanval uit op een bedrijf in een ander
land. En de vraag die vervolgens veel experten zich stelden was: ``Vanaf
wanneer is dit een \emph{act of aggression}?'', en ook: ``Moet een land
op dit soort aanval reageren namens het bedrijf, of namens de hele natie
waar het bedrijf zich bevindt?'' Wie zal het zeggen. Dit boekje helaas
niet, maar het geeft wel voer voor discussie.

\subsection{2015: Stuxnet bewaarheid}\label{stuxnet-bewaarheid}

Keer op keer zullen we dit moeten herhalen: we kunnen nooit met 100\%
zekerheid de origine van een cyberaanval vaststellen. Heel af en toe,
met dank aan klokkenluiders zoals Snowden, of verregaand (al dan niet
journalistiek) onderzoek worden specifiek aanvallen volledig uit de
doeken gedaan (maar helaas vaak jaren na datum).

In 2015, iets meer dan een jaar na de Sony aanval, werd dat waar we voor
vreesden bewaarheid: hackers waren er in geslaagd om de
elektriciteitsinfrastructuur van Oekraïne deels plat te leggen op 23
december, vlak voor kerstavond. Een groepering slaagde er zo in om bijna
een kwart miljoen mensen gedurende meerdere uren zonder stroom te
zetten. Iedereen keek uiteraard ogenblikkelijk in de richting van
Rusland - Oekraïne maakte vroeger deel uit van de voormalige
Sovjetrepubliek - daar de aanvallen kwamen van IP-adressen die waren
toegewezen aan Rusland. Maar zelfs als dat zou zijn, zonder harde
bewijzen dat de hackers ook handelden in naam van de Russische overheid
blijft het koffiedik kijken wie \emph{op de vingers getikt} moet worden
voor de aanval.

Wie het ook waren, één ding staat wel vast: in 2015 waren hackers er
voor het eerst in geslaagd om met één welgemikte aanval honderdduizenden
mensen in problemen te brengen, problemen die verder gingen dan het
verwijderen van enkele bestanden.

Enkele jaren later, in februari 2021, zou het trouwens weer bijna prijs
zijn. Een hacker kreeg toegang tot de systemen die de
waterzuiveringsinstallaties van Oldsmar (Florida, VS) bediende. Een
operator zag op tijd dat z'n muis plots bewoog en de maximum toegelaten
hoeveelheid natriumhydroxide waarde van de filters op een dodelijk
niveau zette. Er werd gelukkig tijdig ingegrepen (en detectoren verderop
in het systeem hadden het vergiftigde water sowieso gedetecteerd), maar
het deed wederom de vraag rijzen of onze kritische systemen wel
voldoende beveiligd zijn tegen dit soort \emph{lone wolves}.

\begin{tipblock}
Sinds de oorlog in Oekraïne is uitgebroken zien we ook terug verregaande
cyberactiviteit vanuit Rusland. De eerste weken van de oorlog (februari
2022) was het opvallend stil en leek het alsof \emph{de angst voor
Russische cyberoorlog} ongegrond was. Ondertussen zijn we helaas meer
dan een jaar verder en verschijnen er meer en meer verhalen, zoals
verwacht, van Russische cyber-inmenging.
\href{https://www.europarl.europa.eu/thinktank/en/document/EPRS_BRI(2022)733549}{Volgende
artikel geeft een heldere tijdslijn hiervan.}

\begin{figure}
\centering
\includegraphics[width=0.7\linewidth,height=\textheight,keepaspectratio,alt={Een stuk van de zonet beschreven tijdslijn uit het door het Europese parlement gepubliceerde rapport.}]{intro/rusukr.jpg}
\caption{Een stuk van de zonet beschreven tijdslijn uit het door het
Europese parlement gepubliceerde rapport.}
\end{figure}

\end{tipblock}

\begin{noteblock}
Microsoft monitort al geruime tijd de vele fake news bronnen die Rusland
rijk is. Er wordt zelfs een \textbf{Russian Propaganda Index} (RPI)
bijgehouden die aangeeft hoe actief de Russische trolls momenteel zijn
in het verspreiden van bewezen onwaarheden (o.a. over de oorlog in
Oekraïne, binnenlands beleid van andere mogendheden, COVID-19 vaccins,
etc.)

\begin{figure}
\centering
\includegraphics[width=0.7\linewidth,height=\textheight,keepaspectratio,alt={Zoals verwacht zag de RPI-grafiek een stevige stijging aan de start van de illegitieme invasie van Oekraïne door Rusland. Bron: https://www.microsoft.com/en-us/security/business/microsoft-digital-defense-report-2022-cyber-influence-operations.}]{intro/rpi.jpg}
\caption{Zoals verwacht zag de RPI-grafiek een stevige stijging aan de
start van de illegitieme invasie van Oekraïne door Rusland. Bron:
https://www.microsoft.com/en-us/security/business/microsoft-digital-defense-report-2022-cyber-influence-operations.}
\end{figure}

\end{noteblock}

\subsection{2015: Scheidingen en
zelfmoord}\label{scheidingen-en-zelfmoord}

In 2015 wordt de schaal van cyberaanvallen steeds groter. De
hoeveelheden informatie die hackers van bedrijven kunnen bemachtigen
kunnen al lang niet meer op één A4'tje afgedrukt worden. Wanneer hackers
toegang krijgen tot de privé-servers van hun doelwit kunnen ze vlotjes
ettelijke gigabytes, tot zelfs terabytes, aan privé informatie stelen.
Tegenwoordig hebben we in Europa de GDPR wetgeving die probeert
bedrijven duidelijk te maken dat zij verantwoordelijk zijn om onze data
op een veilige manier te bewaren (zie appendix) en hen ook te straffen
indien ze dit niet doen. In 2015 was dat veel minder. De datalekken in
die tijd spraken boekdelen: van zodra cybercriminelen toegang hadden tot
de privéservers konden ze de data zonder problemen lezen. Paswoorden,
kredietkaartgegeven, rijksregisternummers, alles stond vaak onbeveiligd
(\emph{ongeëncrypteerd}) op de systemen.

De Ashley Madison website was Tinder voor mensen die een relatie wilden
naast hun ``officiële relatie''. Kortom, de site hielp mensen aan een
affaire. Hun leuze, \emph{``Life is short. Have an affair''}, wond er
geen doekjes rond. En met hun meer dan 20 miljoen ``klanten'' was het
duidelijk dat ze een lucratief idee hadden. Dat ze tegenwind zouden
krijgen was te verwachten\ldots{}

In juli 2015 plaatste een groep hackers een ultimatum op een website aan
het adres van de eigenaars van Ashley Madison: \emph{``Haal jullie
moreel dubieuze website van het Internet, of wij plaatsen meer dan 60
gigabyte aan gestolen data online''}. De site werd niet offline gehaald
en de hackers ``hielden woord.'' De gevolgen waren immens, niet zo zeer
voor Ashley Madison zelf, wel voor de klanten. Plotsklaps kreeg de
wereld een lijst te zien waarin honderdduizenden klanten open en bloot
aan de schandpaal werden genageld. Mensen werden publiekelijk vernederd.
Er is sprake van minstens twee zelfmoorden rechtstreeks als gevolg van
het lek. Mensen werden ontslagen. Kortom, het lekken van wat
uiteindelijk maar een hoop binaire data was, had gevolgen op relaties,
mensenlevens en carrières.

Als positieve noot in dit verhaal halen we hier uit dat dit soort
gigantische lekken mensen heeft doen inzien dat ze tweemaal moeten
nadenken voor ze hun persoonlijke informatie weggeven aan één of andere
Internet-gigant (het is helaas een les die we jaarlijks lijken te
vergeten, kijk maar naar de populariteit van Tik Tok, Facebook, etc.).

\begin{cautionblock}
Een stelregel die bedrijven nu hanteren is de volgende: vraag je niet af
\textbf{of} je gaat gehackt worden maar vraag je af \textbf{wanneer} je
zal gehackt worden. Dit is een heel andere manier van tegen je
beveiligingsprobleem aankijken. Vergelijk het met het in huis halen van
een koffer met daarin 10 miljoen euro aan diamanten. Als je je afvraagt
of er ooit inbrekers zullen binnen geraken en daar naar beveiligt -
waakhonden, videocamera's rond het huis, dubbel slot - dan heb je daar
niets aan als ze vervolgens toch binnen geraken en je koffertje stelen.
Veel beter kan je én je huis beveiligen, én ervan uitgaan dat ze de
koffer gaan bemachtigen: en je dus maar beter ook ervoor zorgt dat de
dieven niets met de diamanten kunnen doen (door bijvoorbeeld je naam er
in te graveren).

Kortom: bedrijven moéten data die ze van ons opslaan minstens
encrypteren en systemen inbouwen die de data onbruikbaar maken als ze
toch gelekt zou worden.
\end{cautionblock}

\subsection{2016: Internet-of-horrors}\label{internet-of-horrors}

In onze huizen verschenen steeds meer apparaatjes die via het, meestal
draadloze, netwerk met elkaar en het Internet konden communiceren.
Internet-of-Things (IoT) bracht een weelde aan nieuwe oplossingen in
onze levens. Domotica, wearables, slimme thermostaten,
beveiligingssystemen, weegschalen, alles werd én slimmer gemaakt én aan
het Internet gehangen.

Een inherent probleem met veel van deze, meestal kleine, toestellen is
dat ze op het gebied van beveiliging ondermaats presteren. Dergelijke
IoT-apparaten werken vaak op batterijen en de makers willen natuurlijk
de batterijduur zo lang mogelijk houden. Iedere extra feature die de
designers in het apparaat willen steken heeft een kost op die
levensduur. Een aspect zoals beveiliging werd dan ook vaak achteraan de
lijst van potentiële features geplaatst. Komt daarbij dat IoT-apparaten
updaten (met bijvoorbeeld nieuwe security patches) soms onmogelijk of
tenminste omslachtig is. Als er dus een beveiligingslek wordt gevonden
in een apparaat, dan is de kans bestaande dat dit lek voor altijd
aanwezig zal blijven én dus ook kan misbruikt worden.

Het was dus wachten tot de eerste aanvallen, specifiek gericht op
IoT-apparaten, zouden plaatsvinden. De meest opvallende aanval gebeurde
eind 2016. Malware, genaamd Mirai, verspreide zich als een lopend
vuurtje over IoT-apparaten waarvan de malware het standaard wachtwoord
en username kende. Gebruikers die vergeten waren het wachtwoord aan te
passen, dat het apparaat heeft wanneer je het uit de doos, haalde zaten
zo plotseling met een ogenschijnlijk perfect werkend, maar besmet
apparaat. Echter, de malware nestelde zich onzichtbaar op het apparaat
en wachtte op commando's van de Mirai makers. De malware creëerde met
andere woorden een zogenaamde \emph{botnet}, een groot netwerk van
besmette apparaten die allemaal commando's kunnen uitvoeren van de
\emph{botnet herder} (i.e.~degene die de malware in de eerste plaats is
beginnen verspreiden). De Mirai-makers hadden zo een leger minicomputers
onder hun bevel die ze konden zeggen ``surf nu allemaal naar die
website''. Een website die plots tienduizenden gebruikers onverwacht
extra te verwerken krijgt zal vaak onder de druk bezwijken en crashen.
Kortom, dit Mirai botnet kon zo grote distributed denial-of-service
(\emph{DDOS}) aanvullen uitvoeren, allemaal omdat gebruikers de
wachtwoorden van hun gloednieuwe apparaatjes niet hadden aangepast. In
hun verdediging, het is vaak een erg omslachtig, technisch, proces om
het wachtwoord van een IoT-apparaat aan te passen.

\begin{tipblock}
De wildgroei van Internet-of-Things apparaten heeft ervoor gezorgd dat
hackers een grote hoeveelheid extra mogelijkheden hebben bijgekregen om
op huis en bedrijfsnetwerken te infiltreren. Of om het in hacker-termen
te zeggen: dankzij de weelde aan IoT-apparaten is de \emph{attack
surface} voor aanvallers exponentieel vergroot.
\end{tipblock}

\begin{warningblock}
Een dubieuze (deels betalende) site, \textbf{shodan.io}, heeft als enige
doel alle Internet-of-Things apparaten in kaart te brengen die zichtbaar
zijn vanop het Internet. Deze ``Google voor IoT'' toont zelfs om wat
voor apparaten het gaat en welke bijvoorbeeld nog steeds het standaard
(default) wachtwoord hebben.
\end{warningblock}

\subsection{2017: Ransomware wordt
gemeengoed}\label{ransomware-wordt-gemeengoed}

De virussen in de vorige eeuw durfden al eens je data te verwijderen.
Lastig, maar erg duidelijk: je data was je kwijt, tenzij je ergens een
back-up had liggen. De nieuwe virussen gingen echter een stapje verder:
ze versleutelden al je data (vaak ook je back-ups als je zo dom was
geweest deze op het zelfde apparaat te hebben) én vroegen vervolgens
losgeld in ruil voor je data. De naam \emph{ransomware} kon, helaas,
niet beter gekozen zijn. Menig particulier betaalde ogenblikkelijk om de
eenvoudige reden dat de ransomware de gebruiker nog een uitweg aanbood
daar ransomware vaak werd ``binnengehaald'' door een gênante actie (bv.
illegale software downloaden, op reclame voor vage pornosites klikken,
etc.).

Voor particulieren was ransomware vooralsnog irritant. Maar wat als de
ransomware bedrijfskritische data begon te encrypteren? Dit is exact wat
er gebeurde in 2017 met de WannaCry en Petya ransomwares. Deze malwares
sloegen er in om banken, hospitalen, havenbedrijven en menig ander groot
(en klein) bedrijf te besmetten. De economische schade werd geschat op
bijna 4 miljard dollar vanwege het lamleggen van de productie (of in het
geval van enkele ziekenhuizen: het redden van mensenlevens).

De schade én de losgeldbedragen werden ook steeds groter. Zo was er het
voorval met Garmin in de zomer van 2020 dat niet alleen de populaire
sport-tracking diensten gedurende meerdere dagen uit de lucht haalde,
maar er ook voor zorgde dat menig vliegtuig niet mocht vliegen omdat de
Garmin Pilot apps niet werkten waardoor de piloten geen up-to-date
aeronautische plannen voorhanden hadden. Het is nooit geweten of Garmin
het losgeld (10 miljoen dollar!) wel of niet heeft betaald, vast staat
wel dat ransomware tegenwoordig een, helaas, erg lucratieve handel is
geworden voor cybercriminelen.

\begin{noteblock}
In december 2022 was Digipolis, de IT-backbone van de stad Antwerpen,
het slachtoffer van een zeer impacterende ransomware aanval door
hackercollectief Play (een soort spin-off van Conti, een Russische
hackersgroep). Quasi alle online-diensten van de stad Antwerpen zijn
meerdere dagen onklaar gemaakt, inclusief politie, brandweer,
bibliotheken, zwembaden, stadsloketten, stedelijke scholen, etc. Quasi
iedere burger heeft in meer of mindere mate last ondervonden van deze
gigantische aanval. Het toont vooral ook weer aan hoe gevoelig ons
digitale leven is en dat er altijd een keerzijde is de digitalisering.
\end{noteblock}

De WannaCry-ransomware maakt gebruik van een kwetsbaarheid in Windows
(EternalBlue) die oorspronkelijk werd ontwikkeld als hackingtool door de
NSA. In 2017 werd deze exploit gelekt door het hackerscollectief
\textbf{The Shadow Brokers}. Hoewel Microsoft kort daarna een
beveiligingspatch uitbracht, hadden veel organisaties deze nog niet
toegepast. Zo konden criminelen WannaCry inzetten met gebruik van een
NSA-wapen, met wereldwijd miljoenen dollars schade tot gevolg \ldots{}
ironisch genoeg ook bij Amerikaanse bedrijven die de NSA net hoort te
beschermen.

\begin{noteblock}
Ondertussen weten we vrij zeker (uiteraard in de cyberwereld zal
attributie nooit 100\% bewezen kunnen worden) dat The Shadow Brokers
Russische staatshackers waren.
\end{noteblock}

\begin{noteblock}
WannaCry was de eerste \textbf{wormable ransomware}. Tot dan verspreidde
ransomware zich voornamelijk via botnets of phishing mails. De
ransomware werkte trouwens \emph{te goed}. Het encrypteerde vaak ook
essentiële bestanden waardoor de besmette computers niet meer
opstartten. Hierdoor kregen de slachtoffers niet de mogelijkheid om
terug te betalen, en kunnen we dus stellen dat de tool veel minder
efficiënt werkte (vanuit het standpunt van de makers) dan verhoopt.
\end{noteblock}

In mei 2025 verschenen gelekte interne chatlogs van de ransomwaregroep
LockBit, wat ongezien inzicht gaf in hun professionele werkwijze. De
groep opereert als een Ransomware-as-a-Service-bedrijf met hiërarchie,
\emph{klantenservice}-achtige onderhandelingen en strikte deadlines. Van
de 208 bekeken gesprekken leidde slechts 18 tot effectieve betalingen,
vaak na forse kortingen, wat aangeeft dat ook cybercriminelen
winstoptimalisatie en reputatie hoog in het vaandel dragen.

Het lek toont tegelijk de kwetsbaarheid van organisaties: betalingen
gebeurden omdat back-ups onbruikbaar waren, besluitvorming traag verliep
of verzekeringen verkeerd werden ingezet.

\begin{figure}
\centering
\includegraphics[width=0.7\linewidth,height=\textheight,keepaspectratio,alt={Een screenshot van één van de chats uit de lek.}]{intro/bitlock.png}
\caption{Een screenshot van één van de chats uit de lek.}
\end{figure}

\subsection{2020: De macht van sociale
media}\label{de-macht-van-sociale-media}

2016 ging er een schokgolf doorheen de wereld. Tegen alle verwachtingen
in won Donald J. Trump de presidentsverkiezingen na een bitsige strijd
tegen Hillary Clinton. Dat social media een belangrijke rol zouden
spelen dit decennium was al lang voorspeld. Facebook, Twitter, Google en
konsoorten hadden miljarden gebruikers die met plezier hun privacy te
grabbel gooiden in ruil voor dagelijkse dopamine-shots dankzij
\emph{likes} en \emph{retweets}. Met dank aan gigantische \emph{troll
farms} (organisaties die duizenden fake social media accounts aanmaken
en zo mee de social media algoritmes beïnvloeden om bepaalde informatie
te \emph{nudgen} in de gewenste politieke richting) kon Trump
honderdduizenden potentiële twijfelaars doen inzien dat hij het juiste
antwoord was tijdens de verkiezingen.

De inmenging van Rusland in een buitenlandse verkiezing (een daad waar
menig land zich reeds schuldig aan heeft gemaakt) was ontluisterend
vanwege de schaal waarop het gebeurde. Het toonde de almacht van de
grote Silicon Valley bedrijven en hoe zij op geopolitiek niveau een
belangrijke speler zijn geworden. Iets dat vervolgens nogmaals bevestigd
werd door het simultaan plaatsgevonden Cambridge Analytica schandaal,
dat niet alleen mede ervoor gezorgd heeft dat Trump president werd, maar
dat hoogstwaarschijnlijk ook een grote groep twijfelaars heeft kunnen
overhalen om een pro-Brexit stem uit te brengen.

Drie jaar later zagen onderzoekers een soortgelijk fenomeen tijdens de
verkiezingen van de Europese Unie in 2019. Een rapport toonde aan dat
Rusland actieve misinformatie campagnes organiseerde om zo de
verkiezingen te beïnvloeden. Ze gebruiken hierbij zogenaamde \emph{bad
actors}: fake social media accounts die (al dan niet fake) nieuws
verspreiden dat ``in de winkel van Rusland past''. Uit het onderzoek
bleek dat de toenmalige belangrijkste EU-mandatarissen (Juncker, King,
Tajani, etc.) soms tot 20\% volgers op Twitter hadden die eigenlijk
\emph{bad actors} waren.

\begin{noteblock}
Wat is het nut van \emph{bad actors} die bekende mandatarissen volgen?
Deze trolls posten hun ``anti-boodschappen'' als reacties op posts van
diegene dat ze volgden. Vervolgens zorgden ze ervoor (via bijvoorbeeld
likes en retweets door mede trolls) dat hun boodschap bovenaan de lijst
van reacties kwam. Hierdoor kreeg iedereen die de mandataris volgde vaak
ook de troll-reactie(s) ogenblikkelijk te zien.
\end{noteblock}

De voorbije jaren is er een (lichte) kentering bezig inzake de macht van
de social media bedrijven, maar het blijft een feit dat momenteel wij
allen grotendeels afhankelijk zijn van door artificiële intelligentie
aangedreven algoritmes die bepalen welke informatie wij
willen/zouden/moeten consumeren, ieder uur van de dag. Zolang echter
data het nieuwe goud is en bedrijven dit vrij kunnen bewaren (GDPR poogt
dit in te perken) zullen zij hun algoritmes kunnen blijven voeden en
trainen om nog griezeliger/knapper te maken.

Het gevolg van die \emph{datahoarding} is echter ook dat datalekken ook
steeds nefastere gevolgen hebben. Daar waar het bij Ashley Madison nog
ging om een dikke 20 miljoen user accounts, medio 2021 zijn het aantal
accounts dat bij een datalek betrokken zijn soms vertienvoudigd - de MGM
Grand Hotels keten zag in de zomer van 2020 plots 142 miljoen van z'n
gast-accounts te koop staan voor een schamele 3000 dollar in bitcoin.

\begin{tipblock}
Twee boeken, die lezen als rasechte thrillers, gaan dieper in op de
zonet beschreven gebeurtenissen: ``Sandworm - A New Era of Cyberwar and
the Hunt for the Kremlin's Most Dangerous Hackers'' van Andy Greenberg
en ``How they tell me the world ends'' van Nicole Perlroth..
\end{tipblock}

\subsection{2021: Oude bibliotheken, nieuwe
problemen}\label{oude-bibliotheken-nieuwe-problemen}

De opmerkelijkste gebeurtenis in 2021 was de lek in de \textbf{Log4J}
Java-bibliotheek. Deze ogenschijnlijk onschuldige bibliotheek wordt al
bijna twee decennia in miljoenen Java-applicaties én servers (waaronder
de Apache webservers!) gebruikt om loginformatie weg te schrijven. Tot
in November 2021 onderzoekers een kritische bug ontdekten waardoor al
deze servers en applicaties plots erg kwetsbaar werden. Het voorval
toonde nog maar eens aan hoe afhankelijk we zijn geworden van onze code
én netwerkinfrastructuur die vaak al jaren oud is en zo goed als zeker
nog ongekende bugs bevatten die misbruikt kunnen worden. In 2012 zagen
we al eens wat de gevolgen kunnen zijn van ``een kleine bug'' in een
veel gebruikte bibliotheek. OpenSSL werd toen al door ontelbare websites
en routers gebruikt om een beveiligde TLS tunnel (zie hoofdstuk 3) op te
zetten wanneer HTTPS werd gebruikt. De bug resulteerde in de
\textbf{Heartbleed}-lek die aanvallers konden gebruiken om data van
servers te stelen die de OpenSSL-bibliotheek gebruikten.

\begin{noteblock}
De Log4j-episode heeft veel bedrijven doen inzien dat ze kritischer
moeten nadenken over hoe ze omgaan met het gebruik van open-source
bibliotheken. De updates van dergelijke bibliotheken worden vaak zonder
nadenken gedownload en geïntegreerd in de eigen software. Veel bedrijven
beginnen daarom nu \emph{policies} op te stellen omtrent het gebruik van
publiekelijk beschikbare stukken software.
\end{noteblock}

\begin{tipblock}
Om de kracht én het gevaar van open-source bibliotheken te bevatten,
lees zeker eens
\href{https://www.bleepingcomputer.com/news/security/dev-corrupts-npm-libs-colors-and-faker-breaking-thousands-of-apps/}{volgende
artikel} waarin de maker van een populaire NPM bibliotheek
(\emph{colors}) het beu was dat grote bedrijven al jaren zijn
bibliotheek(jes) gebruikten zonder hem er ooit voor te bedanken of
betalen (merk op dat zij dit niet moesten doen: de bibliotheken waren
als open-source met de juiste licentie verspreid). Van de één op de
andere dag bracht de maker een update uit die de bibliotheken
``willekeurige output'' liet genereren in de host-applicatie. Wetende
dat zijn bibliotheek wekelijks 20 miljoen keer gedownload wordt, kan je
wel inbeelden dat aardig wat ontwikkelaars, groot en klein, plots met de
handen in het haar zaten.
\end{tipblock}

\subsection{2022-2024: De mens en de
bijna-catastrofe}\label{de-mens-en-de-bijna-catastrofe}

Waar de voorgaande jaren vooral in het teken stonden van ransomware,
zagen we in deze periode twee opvallende trends: de terugkeer naar
\emph{social engineering} om geavanceerde beveiliging te omzeilen én een
wonderbaarlijke ontsnapping aan een wereldwijde IT-ramp.

In 2022 en 2023 bewees de hackersgroep \textbf{Lapsus\$} (bestaande uit
tieners!) dat peperdure beveiliging nutteloos is als je de mens kunt
manipuleren. Ze braken in bij giganten als Uber, Microsoft en Rockstar
Games (waar ze beelden van GTA VI lekten). Hun tactiek? \textbf{MFA
Fatigue}. Ze spamden werknemers 's nachts met honderden
aanmelding-notificaties op hun telefoon, tot het slachtoffer uit
frustratie of vermoeidheid op ``Accepteren'' duwde.

Nog bonter maakte de \textbf{Scattered Spider} groep het in september
2023. Ze legden het \textbf{MGM Resorts} imperium in Las Vegas plat door
simpelweg naar de helpdesk te bellen, zich voor te doen als een
werknemer, en het wachtwoord te resetten. Gevolg: gokkasten op zwart,
sleutelkaarten van hotelkamers werkten niet meer en liften stonden stil.
Schade: meer dan 100 miljoen dollar. Het toonde aan dat \emph{vishing}
(voice phishing) in combinatie met een zwakke helpdesk-procedure
dodelijker kan zijn dan de meest geavanceerde malware.

\subsection{2024: CrowdStrike en de fragiele
verbondenheid}\label{crowdstrike-en-de-fragiele-verbondenheid}

In juli 2024 werd de wereld opgeschrikt door een grootschalige uitval
van Windows-systemen na een foutieve update van beveiligingssoftware
\textbf{CrowdStrike}. Dit was weliswaar geen cybersecurityincident, maar
had het wel kunnen zijn. Het voorval toonde pijnlijk aan hoe gevoelig
onze digitale infrastructuur is vanwege de enorme verbondenheid: één
zwakke schakel in de keten kan wereldwijd miljoenen systemen platleggen.
Het incident geldt als een ernstige waarschuwing voor de potentiële
impact van toekomstige \emph{supply chain attacks}.

\begin{tipblock}
Bij een \textbf{supply chain attack} of ketenaanval richt een aanvaller
zich niet rechtstreeks op het einddoel (bijvoorbeeld een groot bedrijf),
maar op een zwakke schakel in de leveranciersketen, zoals een
softwareleverancier of een externe dienstverlener. Door beveiligde
software van de leverancier te compromitteren (bijvoorbeeld door malware
toe te voegen aan een update), kan de aanvaller ongemerkt toegang
krijgen tot de systemen van alle klanten die deze software gebruiken.

Bekende voorbeelden: * \textbf{SolarWinds (2020)}: Hackers infiltreerden
het netwerkbeheersysteem Orion, waardoor duizenden organisaties
wereldwijd, waaronder de Amerikaanse overheid, werden getroffen. *
\textbf{Kaseya (2021)}: Via beheersoftware voor IT-dienstverleners werd
ransomware verspreid naar duizenden bedrijven.
\end{tipblock}

Begin 2024 ontsnapte de wereld aan een digitale ramp dankzij\ldots{} een
oplettende Microsoft-ingenieur die vond dat zijn systemen ``een halve
seconde te traag reageerden''. Hij ontdekte dat in \textbf{XZ Utils},
een stukje software dat in bijna elke Linux-server ter wereld zit, een
achterdeurtje (\emph{backdoor}) was ingebouwd. Jarenlang had een hacker
het vertrouwen gewonnen van de eenzame beheerder van dit open-source
project, om vervolgens stiekem de backdoor toe te voegen. Was dit niet
ontdekt, dan hadden aanvallers toegang gehad tot miljoenen servers
wereldwijd. Het was de ultieme, langlopende \emph{supply chain attack}
die mislukte op de valreep.

\section{2025: A.I. is here to stay}\label{a.i.-is-here-to-stay}

Sinds de winter van 2022 worden we constant overspoeld met nieuwe
(generatieve) AI toepassingen die zo krachtig zijn dat er zelfs in maart
2023 werd voorgesteld om alle AI onderzoek ``even te pauzeren'', zodat
wij, als gemeenschap, kunnen reflecteren (én bijbenen) over hoe we onze
toekomst met AI willen opbouwen. Onze glazen bol is niet perfect, maar
wees er maar van overtuigd dat we de komende jaren onvoorspelbaar,
bizarre, krachtige cyberaanvallen gaan tegenkomen die door AI worden
ondersteund. Een eerste voorbeeld hiervan zagen we reeds recent in de
zomer van 2023 genaamd \textbf{DarkBert}, een broertje van ChatGPT dat
was getraind op DarkWeb data en dus de ideale chatpartner is voor
conversaties die het daglicht niet mogen zien.

Andere voorbeelden van AI-gedreven incidenten die we recent zagen
opduiken:

\begin{itemize}
\tightlist
\item
  \textbf{Deepfake CFO (2024)}: Een multinational in Hong Kong verloor
  25 miljoen dollar nadat een werknemer geld overmaakte na een videocall
  met wat leek op de CFO en andere collega's. Iedereen in de call,
  behalve het slachtoffer, was een deepfake.
\item
  \textbf{WormGPT \& FraudGPT}: Dit zijn kwaadaardige varianten van
  ChatGPT, specifiek getraind om malware te schrijven, phishing mails op
  te stellen en kwetsbaarheden te vinden, zonder de ethische beperkingen
  die commerciële modellen hebben.
\item
  \textbf{AI-gestuurde Social Engineering}: Aanvallers gebruiken AI om
  in \emph{real-time} stemmen van bekenden na te bootsen (vishing) om zo
  slachtoffers te overtuigen vertrouwelijke gegevens te delen of geld
  over te maken.
\end{itemize}

\section{En de toekomst? AI will rule the
world}\label{en-de-toekomst-ai-will-rule-the-world}

Een glazen bol hebben we niet, maar vast staat dat het er niet op zal
verbeteren. Zoals gezegd gaan bedrijven uit van \emph{when} in plaats
van \emph{if} als het gaat over de vraag of ze al dan niet ooit gehackt
zullen worden. Daarnaast zien we dat hoogtechnologische producten meer
en meer ingeburgerd geraken in het cybercrime-milieu. Zeker door de AI
golf die ons nu al enkele jaren overspoelt aan een razend tempo: deep
fakes video van bekenden zijn nu nog ``grappig'', maar de realistische
beelden (afbeeldingen, video én nu zelfs ook spraak) die ze kunnen
produceren zijn al even niet meer te onderscheiden van het echte en
zullen dus meer en meer kunnen gebruikt worden voor sextortion, spear
phishing en soortgelijke aanvallen.

Dit soort trends zullen nog jaren \emph{fake news} hoogtij laten vieren
waardoor ook toekomstige verkiezingen interessante doelen blijven voor
andere mogendheden om hun stempel te drukken op geopolitieke
tegenstanders. Er gaan zelfs stemmen op dat we leven in een tijdperk van
het ``dead internet'': een tijdperk waarin we niet meer kunnen
vertrouwen op de informatie die we online vinden. En waarin bots met
bots communiceren op sociale media en nieuwe narratieven creëren die de
publieke opinie beïnvloeden.

\section{Het is erger, maar\ldots{}}\label{het-is-erger-maar}

Dus ja, het wordt helaas erger. De wapenwedloop in de cyberwereld gaat
beangstigend snel vooruit en het wordt moeilijker en moeilijker om als
``normale sterveling'' er een antwoord op te geven. Als een IoT botnet,
bestuurd door middel van een krachtige A.I, van 10 miljoen apparaten
morgen beslist om de infrastructuur van jouw KMO plat te leggen, dan
zullen ze daar in slagen, ongeacht de vele euro's die je hebt
geïnvesteerd in firewalls, intrusion detection systems, virusscanners en
honeypots.

Aan de andere kant heeft de wapenwedloop er wel voor gezorgd dat onze
infrastructuur ook steeds complexere aanvallen kan weerstaan. Hierdoor
wordt het voor huis-tuin-en-keuken malware een pak moeilijker om nog
computers van thuisgebruikers te besmetten.

Maar laten we deze cyberstropers geen vrij spel geven! Laten we leren
van hun technieken om zo zelf onze systemen en personeel te
\emph{hardenen} en te beschermen tegen wat niet anders dan het ``wilde
westen van het Internet'' (dixit komiek Steven Wright) kan genoemd
worden.

\begin{tipblock}
In de appendix achteraan dit boek vind je tal van boeiende en nuttige
bronnen om op de hoogte te blijven over het reilen en zeilen in de
schimmige wereld van de defensieve én offensieve cybersecurity wereld.
\end{tipblock}

\chapter{Cybersecurity fundamenten}\label{cybersecurity-fundamenten}

``Iedere goede film begint met een zwart scherm'', zegt Batman in de
``The Lego Batman Movie'' ™ film. Wel, ik vul dit aan met ``\ldots en
ieder goed hoofdstuk begint met een definitie.''

Laten we daarom eerst eens een definitie van cybersecurity neerpennen
dat netwerkbedrijf \emph{Cisco} gebruikt:

\emph{``Cybersecurity is the practice of protecting systems, networks,
and programs from digital attacks. These cyberattacks are usually aimed
at accessing, changing, or destroying sensitive information; extorting
money from users; or interrupting normal business processes.}

\emph{Implementing effective cybersecurity measures is particularly
challenging today because there are more devices than people, and
attackers are becoming more innovative.''}

Alles draait met andere woorden rond het beschermen van informatie, of
dat die nu op een server, in een document of in iemands hoofd zit.

\section{CIA en het security model}\label{cia-en-het-security-model}

Data (of informatie), in welke vorm dan ook (berichten over een netwerk,
bestanden op een harde schijf, tekst in een database), moeten beschermd
worden, dat beseffen we nu. Het doel van onze data is dat deze voldoet
aan het acroniem \textbf{C.I.A} wat staat voor:

\begin{itemize}
\tightlist
\item
  \textbf{C} voor ``\textbf{confidentiality}'': vertrouwelijkheid. De
  data kan enkel door zij die er recht toe hebben gebruikt worden. We
  gaan dit onder andere oplossen met behulp van encryptie en
  wachtwoorden.
\item
  \textbf{I} voor ``\textbf{integrity}'': integriteit. We moeten weten
  of onze data onbeschadigd is en niet werd aangepast door derden (of
  storingen). Bij bestanden gaan we bijvoorbeeld werken met zogenaamde
  (secure) hashes.
\item
  \textbf{A} voor ``\textbf{availability}'': beschikbaarheid. Data die
  niet door rechtmatige gebruikers kan bereikt worden is onbestaande
  data. Zogenaamde ``Denial-of-Service'' (DoS) aanvallen hebben als doel
  deze pijler van CIA aan te vallen. Availability is een breed veld en
  wordt onder andere opgelost door back-ups, redundante servers
  enerzijds, en preventieve maatregelen anderzijds zoals firewalls, load
  balancers, etc.
\end{itemize}

\subsection{McCumber kubus}\label{mccumber-kubus}

De zogenaamde McCumber kubus, ontwikkeld door John McCumber in 1991,
geeft een goed beeld weer waarom het steeds belangrijk is goed te
beseffen binnen welke context we praten. Deze kubus stelt een model voor
dat kan gebruikt worden om je ervan te vergewissen dat je aan alles
denkt wanneer je je informatie wilt beschermen. De kubus bestaat uit
drie dimensies en iedere dimensie bestaat uit een aantal aspecten.
\textbf{Enkel wanneer we alle aspecten van alle dimensies in onze
beveiligingsaanpak voorzien kunnen we hopen dat we onze C.I.A. doelen
hebben bereikt.}

Om ons doel te bereiken (C.I.A.) moeten we ervoor zorgen dat we dit
toepassen op alle vormen die onze data kan hebben (opslag, verzenden,
verwerken). Dit kunnen we bewerkstelligen door technologische
oplossingen (zoals encryptie wat in het volgende hoofdstuk wordt
uitgespit), maar een niet onbelangrijke factor zijn ook de mensen die
met de data moeten werken. Als zij zich niet aan de afspraken
(procedures) houden en hun wachtwoorden gewoon op post-its aan hun
scherm hangen, dan mag je een nog zo'n dure firewall hebben, het zal
niet baten.

\begin{figure}
\centering
\includegraphics[width=0.7\linewidth,height=\textheight,keepaspectratio,alt={De McCumber kubus.}]{crypto/secmodel.png}
\caption{De McCumber kubus.}
\end{figure}

\begin{warningblock}
De wereld van de cyberboswachters houdt van afkortingen, protocollen en
vreemd klinkende standaarden. De wereldwijd gebruikte \textbf{ISO 27001}
standaard (of norm) is er zo eentje, maar wel een erg belangrijke om te
kennen. Wanneer een bedrijf of organisatie wil aantonen dat ze
informatiebeveiliging hoog in het vaandel dragen, dan zullen ze deze
norm trachten te behalen. Om als \emph{ISO 27001 bedrijf} door het leven
te mogen gaan moeten ze aan een hele hoop strenge eisen voldoen (die
alle dimensies van de McCumber kubus omvatten, zou je kunnen zeggen) die
in de standaard beschreven staan. Hierbij zal een externe auditor
vervolgens controleren of je hier aan voldoet als bedrijf (en moet je
dit elke drie jaar opnieuw doen).
\end{warningblock}

\section{Een ongelijke strijd}\label{een-ongelijke-strijd}

De McCumber kubus is een mooi concept, maar het is ook niet meer dan
dat: een theoretisch framework. Het is ideaal om vanuit een high-level
perspectief je ervan te vergewissen dat je aan alles hebt gedacht, maar
in de praktijk komt er uiteraard bij alle aspecten van deze kubus aardig
wat kijken.

De cyberboswachters van de 21e eeuw hebben geen eenvoudige job. Ze was
in de vorige eeuw al pittig, de laatste jaren is de wapenwedloop er
helaas alleen maar grimmiger op geworden.

\begin{itemize}
\tightlist
\item
  De aanvallers kunnen aan bijna supersonische \textbf{snelheid}
  aanvallen op duizenden, of zelfs miljoenen, systemen starten.
\item
  ``Alles is verbonden'': onze huidige IT-netwerken zijn vele malen
  groter en complexer dan circa 20 jaar geleden. Hierdoor is ook de
  zogenaamde \textbf{attack surface} steeds groter. Gedaan zijn de
  tijden dat een middelgroot tot groot bedrijf genoeg had aan één
  cyberboswachter. Er zijn nu zelfs bedrijven die kunnen ingehuurd
  worden om bij problemen (of slimmer: vooraf als audit) te komen helpen
  om de boel te blussen.
\item
  De tools die aanvallers, van welke aard ook, ter hun beschikking
  hebben zijn vaak ongelooflijk \textbf{eenvoudig} geworden. Gedaan is
  de tijd dat een ietwat stevige aanval kon gedaan worden door experts
  met tien jaar script- en netwerkervaring. Sommige vreeswekkende
  aanvallen vereisen niet meer dan het IP-adres van het slachtoffer in
  een invulveld invullen en vervolgens een klik op een grote rode knop
  ``Start attack''.
\item
  Aanvallers schuimen \emph{underground} fora af, op zoek naar de
  nieuwste \emph{zero-day vulnerabilities} die ze in hun arsenaal kunnen
  opnemen. Fabrikanten van besturingssystemen en software kunnen de
  \textbf{snelheid waarmee nieuwe zwakheden} in hun systeem worden
  gevonden niet volgen.
\item
  Door voorgaande snelheid van nieuwe zwakheden, duurt het ook langer en
  langer voor alles kan \textbf{gepatcht} worden door de ontwikkelaars
  van de software.
\item
  Aanvallers kunnen \textbf{gigantische legers van computers en botnets}
  gebruiken om een ijzingwekkende hoeveelheid aan simultane aanvallen op
  één enkel slachtoffer uit te voeren.
\item
  \textbf{Gebruikers} zijn een nog kleiner radertje in dit alles
  geworden en zullen nog sneller fouten maken (klikken op een link in
  een phishing e-mail bijvoorbeeld) dan voorheen, met alle gevolgen van
  dien.
\end{itemize}

\subsection{Zero days en patching}\label{zero-days-en-patching}

Aanvallers hebben voor zero days aardig wat geld over. Een zero day
kopen is een aanval kopen die gegarandeerd zal werken daar deze een
zwakte misbruikt die nog niet bij de maker van de software gekend is.

Een zero day zal quasi gegarandeerd blijven werken tot de ontwikkelaars
een nieuwe patch ervoor maken. Maar zelfs dan blijven zero days nuttig:
het is niet omdat er een patch bestaat, dat de doelwitten deze patch ook
effectief reeds geïnstalleerd hebben. Veel bedrijven hebben nu erg
strenge \emph{patching policies} maar toch blijft het dweilen met de
kraan open: er moet maar één systeem niet gepatcht zijn tegen de zero
day van de aanvallers en het gaatje in de verdedigingslinie is gevonden
en kan misbruikt worden.

De meeste fabrikanten van besturingssystemen (Apple, Microsoft, etc.) en
veelgebruikte softwarepakketten (Adobe, Microsoft, etc.) brengen patches
op welbepaalde dagen uit. Dit zorgt er bijvoorbeeld voor dat
systeembeheerders hier rekening mee kunnen houden in hun wekelijkse
planning. Voor gebruikers van zero days is dit ook nuttig: er ontstaat
een zogenaamde \textbf{window of vulnerability}. Dit is de periode
tussen het ``ontdekken en in gebruik nemen van een zero day'' en de
moment waarop de patch tegen de zero day wordt verspreid. In dit
\emph{window} heeft de aanvaller vrij spel daar geen enkel systeem al
kan gepatcht zijn.

\begin{tipblock}
In 2021 verscheen ``How they tell me the world ends'' van New York Times
journaliste Nicole Perlroth. Dit boek is erg ontluisterend en geeft een
griezelig inzicht in hoe het er momenteel aan toe gaat in de schimmige
wereld van zero days, offensieve cybersecurity, etc. Bekijk bijvoorbeeld
maar eens de twitter-account van Chaouki Bekrar
(\href{https://twitter.com/cbekrar}{twitter.com/cbekrar}), oprichter van
Zerodium, een zero-days \emph{broker} en huiver bij de gigantische
prijzen die zero days waard kunnen zijn (soms meer dan één miljoen
dollar\ldots).
\end{tipblock}

\section{Wie zijn de stropers?}\label{wie-zijn-de-stropers}

Waar moeten we ons tegen beschermen als cyberboswachter? Het zou
verleidelijk zijn om ons te richten op één specifieke doelgroep en ons
daar volledig tegen te wapenen. Helaas is het niet te voorspellen van
wie je last zal hebben. Iedere groep aanvallers heeft eigen motivaties
en middelen en het is niet altijd evident om tegen ieder iets te doen.
Het kan natuurlijk geen kwaad om tenminste te weten met welke groepen je
mogelijk zal geconfronteerd worden:

\begin{itemize}
\tightlist
\item
  \textbf{Hackers}: waarbij we een onderscheid moeten maken tussen
  white, black en grey hat hackers uiteraard.
\item
  \textbf{Scriptkiddies}: een groep die soms te laat beseft dat ook zij
  dingen doen die erg strafbaar zijn.
\item
  \textbf{Werknemers}: een vaak over het hoofd gekeken, maar oh zo veel
  voorkomende groep.
\item
  \textbf{Cybercriminelen}: dé grootste plaag, nog steeds.
\item
  \textbf{Cyberterroristen}: rebel, vrijheidsstrijders, terrorist.
  \emph{It's all in the eye of the beholder} uiteraard.
\item
  \textbf{Spionnen} : \emph{James Bond of the WWW}.
\item
  \textbf{Overheden} (of door overheden gesponsord): een probleem dat
  steeds groter blijkt te worden.
\end{itemize}

\subsection{Hackers}\label{hackers}

Wanneer een digitale stroper niet onder één van de andere noemers kan
gezet worden dan wordt hij ``hacker'' genoemd, een soort catch-all term
die soms een positieve, soms een negatieve connotatie heeft. Om toch wat
onderscheid mogelijk te maken kunnen we zeggen dat er drie soorten
hackers zijn:

\begin{itemize}
\tightlist
\item
  \textbf{Whitehat}: \emph{``de goei''}. Dit zijn hackers die onder
  duidelijke afspraken met hun doelwit de sterktes en zwaktes van een
  systeem zullen testen door deze te proberen te omzeilen. White hat
  hackers werken volledig binnen de krijtlijnen van de wet en zullen
  enkel die zaken testen waartoe zij recht hebben. Denk bijvoorbeeld aan
  audit-bedrijven die de beveiliging van andere bedrijven zullen
  \emph{pentesten} (penetration testing: trachten in een systeem te
  geraken) of bounty hunters die bijvoorbeeld via \emph{intigrity.com}
  op zoek gaan naar nieuwe problemen bij een website of product.
\item
  \textbf{Greyhat}: letterlijk een grijze zone. Grey hat hackers hebben
  meestal een positief doel maar zullen niet altijd volgens ``de regels
  van de wet werken''. Denk maar aan een hacker die ongevraagd een lek
  ontdekt in een bedrijf en dit ook rapporteert. De kans is bestaande
  (maar gelukkig kleiner dan vroeger) dat het bedrijf in kwestie hier
  niet mee opgezet is en alsnog de hacker zal aanklagen. Vergelijk grey
  hat hackers met een soort ``goede dief'': hij gebruikt zijn expertise
  om ongevraagd huizen binnen te breken om dan vervolgens, zonder iets
  te stelen, een briefje achter te laten met uitleg hoe de inbreker net
  is binnen geraakt.
\item
  \textbf{Blackhat}: de ``bad guys'' van de hoop en ook wel
  \emph{crackers} genoemd. Black hat hackers zullen systemen aanvallen
  waar ze geen toestemming voor hebben en meestal met als doel om niet
  legale resultaten (lees: geld, kennis of macht) te bereiken ten koste
  van het doelwit.
\end{itemize}

Recent zijn er nog 3, minder vaak gebruikte, types die aan de hand van
een kleur specifieke hacker-types definiëren:

\begin{itemize}
\tightlist
\item
  \textbf{Redhat}: dit zijn hackers die als doel hebben blackhat hackers
  het leven moeilijk te maken. Ze worden ook soms \emph{vigilantes}
  genoemd. De middelen die de redhat hackers gebruiken zijn echter niet
  noodzakelijk legaal en dus alhoewel hun doel nobel is, moeten we toch
  vraagtekens plaatsen bij hun werkwijzen.
\item
  \textbf{Bluehat}: de wraaklustige hackers. Bluehat hackers hebben maar
  1 doel, en dat is wraak nemen op iets of iemand, gebruik maken van
  alle middelen beschikbaar. Laten we auteur Jodi Picoult citeren die
  het volgende zegt over wraak: \emph{``When you begin your journey of
  revenge, start by digging two graves: one for your enemy, and one for
  yourself''}\ldots{}
\item
  \textbf{Greenhat}: de jonkies, ook wel scriptkiddies genoemd. We
  zullen deze in de volgende sectie uit de doeken doen.
\end{itemize}

\subsection{Scriptkiddies}\label{scriptkiddies}

Iemand die weinig tot niets kent van cybersecurity maar toch bestaande
tools uittest op slachtoffers, wordt ook wel een scriptkiddie genoemd.
Deze groep mensen gebruikt tools die eenvoudig in gebruik zijn, zonder
altijd goed te weten wat de tool doet én wat de gevolgen ervan zijn.
Doordat tools steeds krachtiger worden, zijn ook de gevolgen van
scriptkiddie-aanvallen steeds drastischer. Wat deze gebruikers vaak
vergeten is dat hun acties strafbaar zijn. Zo vergeten ze dat klikken op
een knopje in een programma, vanuit hun gezellige bureaukamer thuis,
erge gevolgen kan hebben voor hun doelwit. Het gebeurt dan ook geregeld
dat een scriptie onzacht én snel in aanraking met het gerecht komt. De
veiligheidsdiensten die op cyberaanvallen reageren kunnen natuurlijk
niet zien wie er achter een aanval zit en zullen dus altijd op dezelfde
manier reageren. Een inval van een team zwaar bewapende agenten omdat
een bank of webshop al dagen wordt lamgelegd is een ``ervaring'' die een
scriptkiddie nog lang zal herinneren (en zien op z'n lege spaarboekje
ten gevolge van de grote boetes die hij of zij nog jaren zal moeten
afbetalen)

\begin{tipblock}
Low Orbit Ion Cannon is zo'n typische scriptkiddie tool: eenvoudig in
gebruik, met potentieel grote gevolgen voor het slachtoffer. Deze tool,
wanneer meerdere gebruikers hem samen gebruiken, voert DDOS-aanvallen op
het doelwit uit. In 2010 werd de tool bijvoorbeeld gebruikt om de
servers van Visa, MasterCard en Paypal uren lang lam te leggen als wraak
op het feit dat deze betalingsdiensten geen betalingen voor WikiLeaks
meer aanvaardden.
\href{https://www.smh.com.au/technology/the-aussie-who-blitzed-visa-mastercard-and-paypal-with-the-low-orbit-ion-cannon-20101209-18qr1.html}{Meer
informatie}.

\begin{figure}
\centering
\includegraphics[width=0.7\linewidth,height=\textheight,keepaspectratio,alt={Low Orbit Ion Cannon UI (Bron Wikipedia).}]{intro/ioncan.jpg}
\caption{Low Orbit Ion Cannon UI (Bron Wikipedia).}
\end{figure}

\end{tipblock}

\subsection{Werknemers}\label{werknemers}

Werknemers die niet tevreden zijn over hun baas of werkomgeving, of net
zijn ontslagen, zijn een veelvoorkomend probleem (volgens sommige
experts zelfs hét belangrijkste cybersecurity probleem). Ze zitten vaak
letterlijk ``aan de binnenkant'' van de beveiligingssystemen en kunnen
daardoor ook veel meer schade aanrichten als ze dat willen. Ontevreden
werknemers die wraak nemen op hun werkgever is een veel voorkomend
probleem. Zeker als die werknemer op de koop toe bijvoorbeeld de
netwerk-administrator was. Er zijn verhalen van ex-admins die voor hun
vertrek de IT-systemen van het bedrijf saboteerden en, als klap op de
vuurpijl, dan ook nog eens losgeld eisten om de systemen terug
up-and-running te brengen. Het hoeft natuurlijk niet zo spectaculair
zijn. Of wat te denken van werknemers die waardevolle documenten stelen,
doorverkopen of aanpassen zonder dat het bedrijf er erg op heeft.

Werknemers zijn ook vaak onbedoeld de oorzaak van veel problemen: ze
gaan misschien slordig om met de manier waarop ze hun wachtwoord
bewaren, waardoor anderen via hun account kunnen inbreken. Ze laten
mensen binnen die gekleed zijn als ``collega's'' zonder te vragen of ze
wel werknemer van het bedrijf zijn. Of ze installeren bijvoorbeeld een
extra draadloos access point om een betere wifi-dekking op de bureau te
hebben, waardoor dit toestel plots een veel minder goed beveiligd
doelwit is dat hackers kunnen misbruiken.

\begin{tipblock}
Veel van dit soort problemen kunnen voorkomen worden door een doordacht
``identity management''-systeem dat er voor zorgt dat de accounts van
recent ontslagen werknemers ogenblikkelijk worden verwijderd of dat
tenminste de toegang tot bedrijfskritische systemen uitschakelt. Voorts
moet personeel (op alle niveaus!) opgeleid en getraind worden zodat ook
dit facet van de McCumber kubus gedekt wordt.
\end{tipblock}

De laatste tien jaar heeft het \textbf{Bring your own device
(BYOD)}-concept in bedrijven voor een extra dimensie gezorgd waar
cyberboswachters rekening mee moeten houden. Vroeger kon je je
verdediging opbouwen als een soort ommuurde burcht waarbij je er van uit
mocht gaan dat alles ``binnen de burcht'' veilig was. Door BYOD gaat dit
concept natuurlijk niet meer op: gebruikers wandelen bedrijven binnen
met hun eigen laptop, tablet, smartwatch en smartphone, etc. Allemaal
apparaten die potentieel door stropers vooraf, bij de gebruiker thuis,
werden geïnfecteerd. Van zodra dit besmette apparaat dan in het bedrijf
wordt geïntroduceerd heeft de aanvaller mogelijks ongelimiteerde toegang
tot het bedrijfsnetwerk. Kortom, we moeten nu ook verdedigingsmuren
rondom de individuele werknemers én hun apparaten bouwen. Denk daarbij
aan on-device firewalls en virusscanners, maar ook netwerk-authenticatie
voor ieder toestel, etc.

\subsection{Cybercriminelen}\label{cybercriminelen}

\emph{Money talks} zegt men wel eens, en dat geldt (geld, snap je'm ;) )
zeker voor deze groep. Cybercriminelen doen wat criminelen al millennia
doen en dat is rijkdom die hen niet toebehoort proberen te pakken
krijgen. Geld, geld, geld is de motivatie van deze groep mensen. Een
onderzoek in 2021
(\href{https://www.iii.org/fact-statistic/facts-statistics-identity-theft-and-cybercrime\#:~:text=There\%20were\%204.8\%20million\%20identity,up\%20from\%20651\%2C000\%20in\%202019.}{bron})
schat dat meer dan 700 miljard dollar verlies werd opgetekend ten
gevolge van online criminaliteit. Doordat steeds meer mensen hun
betalingen en identiteiten (denk maar aan de vele its-me phishing
sms'jes dat je geregeld krijgt) online beheren wordt ook de groep
potentiële slachtoffers steeds groter.

Cybercriminaliteit wordt soms wel eens de motor van de cybersecurity
genoemd omdat ze steeds blijven innoveren en zoeken naar nog betere
manieren om onschuldige slachtoffer hun centjes te stelen. Hierdoor
moeten ook de boswachters steeds blijven vernieuwen.

Cybercriminaliteit is nu zelfs zo ver geëvolueerd dat ze heuse moderne
bedrijfsconcepten overnemen en hun zaakje als echte bedrijven runnen.
Moderne ransomware criminelen hebben zelfs helpdesks die je kan bellen
om je te helpen om de betaling (de \emph{ransom}) te regelen. Of wat te
denken van website die botnets verhuren als waren het legale services.
Hier en daar zie je nu zelfs het ``-as a service'' zinnetje verschijnen
waarbij bijvoorbeeld ``Ransomware as a Service''
(\href{https://zvelo.com/raas-ransomware-as-a-service/}{RaaS}) of ``spam
as a service'' kan gehuurd worden. Het doel hierbij is natuurlijk om de
strafbare feiten zoveel mogelijk te verleggen naar de persoon die de
services inhuurt, en niet naar de aanbieder ervan.

\subsection{Cyberterroristen, spionnen en
overheden}\label{cyberterroristen-spionnen-en-overheden}

De laatste drie groepen bespreken we samen, ook al omdat de termen soms
overvloeien afhankelijk aan wie je vraagt om iets of iemand met dit
label te bestempelen. Zoals reeds in het eerste hoofdstuk aangehaald is
de cyberwereld tegenwoordig ook een belangrijk terrein waar geopolitieke
ruzies op worden uitgevochten. Er bereiken ons steeds meer berichten van
de exploten die hier doorgaan. De financiële en technische middelen die
deze groep voorhanden heeft voor zowel offensieve als defensieve
cyberacties is meestal immens groter dan van alle andere stropers in dit
overzicht. We zagen ooit een presentatie waarin een cybersecurity expert
ietwat lachend sprak over het ``Mossad / Non-Mossad
verdedigingsprincipe'' (Mossad is een Israëlische geheime dienst en
staat in de top van \emph{strafste} cybersecurity expertise). Het
principe gaat uit van de manier waarop je je beveiliging opbouwt: ga er
van uit dat de Mossad in je systemen zal geraken, ongeacht hoeveel geld
en personeel je tegen het probleem aan gooit. Het is met andere woorden
efficiënter dat je een realistische inschatting maakt van je
tegenstanders (qua expertise en middelen) en daar specifiek je op richt,
waarbij je natuurlijk het \emph{low hanging fruit} niet over het hoofd
ziet.

\section{Hoe vallen ze aan?}\label{hoe-vallen-ze-aan}

Alhoewel voorgaande groepen van aanvallers allemaal erg specifieke
redenen hebben, kunnen we toch hun aanvallen generaliseren in vijf
duidelijke stappen:

\begin{enumerate}
\def\labelenumi{\arabic{enumi}.}
\tightlist
\item
  \textbf{Verkennen}: passieve \emph{reconnaissance}. Voor de stropers
  hun aanval effectief starten zullen ze eerst hun doelwit(ten)
  onderzoeken. Ze zoeken zo naar de eenvoudigste, of meest verborgen,
  manier om in een systeem te geraken. Verkennen wordt ook wel passieve
  reconnaissance genoemd omdat in deze fase het doelwit bijna nooit kan
  detecteren dat de stroper actief is. Deze fase gaat erg breed en is
  niet beperkt tot enkel ``tools'' gebruiken. In deze fase zal de
  stroper ook vaak via zoekmachines, social media en \emph{dumpster
  diving} proberen meer te weten te komen over het bedrijf, de
  werknemers, etc.
\item
  \textbf{Scannen}: actieve \emph{reconnaissance}. Van zodra de stroper
  een breed overzicht heeft van z'n doelwit zal hij overgaan op actieve
  scanning. Denk hierbij aan port scanners, vulnerability scanners, etc.
  Uiteraard is deze fase actief én bestaat er dus de kans dat de
  boswachters dit tijdig opmerken en zo de aanvallen in de volgende fase
  kunnen afslaan.
\item
  \textbf{Toegang verkrijgen}: door het scannen weet de stroper nu welk
  systeem hij zal benaderen om toegang tot bijvoorbeeld het
  bedrijfsnetwerk te krijgen. Meestal heeft de aanvaller een systeem
  gedetecteerd in de vorige fase met een gekende kwetsbaarheid, of hij
  heeft bijvoorbeeld via spear phishing een backdoor bij een werknemer
  geïnstalleerd, etc. In deze fase zal de aanvaller ook trachten steeds
  meer rechten te krijgen zodat hij steeds meer kan gedaan krijgen op de
  \emph{gepwnde} systemen.
\item
  \textbf{Toegang bestendigen}: van zodra de stroper \emph{in het
  systeem} zit, begint hij als eerste z'n toegang te bestendigen. De
  aanvaller kan niet voorspellen wanneer het systeem waar hij op zit
  wordt uitgeschakeld, verwijderd, etc. Kortom, de stroper begint nu
  naarstig z'n toegang te garanderen door bijvoorbeeld een (nieuwe)
  backdoor te installeren die ook later actief zal zijn. Voorts zal hij
  ook mogelijk andere systemen in de buurt aanvallen zodat hij niet
  beperkt is tot één systeem (en er dus ook geen \emph{``single point of
  failure''} is voor de stroper).
\item
  \textbf{Sporen wissen}: hoe langer de stroper uit het vizier van de
  boswachters kan blijven, hoe effectiever hij z'n doelen kan behalen.
  Een behendig stroper zal er dan ook voor zorgen dat zijn sporen
  ondetecteerbaar blijven door het aanpassen of wissen van logs, het
  verwijderen van backdoors, etc.
\end{enumerate}

\begin{tipblock}
Om de kracht van de tools die stropers voorhanden hebben aan den lijve
te ontdekken, is het aangeraden om Kali OS te installeren. Deze
Linux-distributie (die je ook virtueel kunt draaien) zit tjokvol
pentest-tools die zowel stropers als boswachters constant gebruiken. Dit
OS laat je toe om alle stappen van de aanvaller te simuleren.

Leer zeker ook werken met Metasploit, dat ook in Kali zit. Het
Metasploit project is een verzameling erg krachtige pentest tools,
inclusief de nieuwste snufjes om bijvoorbeeld malware in een pdf te
embedden, etc.
\end{tipblock}

\begin{noteblock}
Wanneer een netwerk wordt gepentest (legaal) werkt men vaak met twee
teams die tegen elkaar strijden. Het \emph{red team} speelt de rol van
de digitale stropers, terwijl het \emph{blue team} als boswachters zal
proberen de aanvallen te verijdelen.
\end{noteblock}

\section{Classificatie van aanvallen}\label{classificatie-van-aanvallen}

Dit hoofdstuk begon met een definitie, dan moeten we zeker ook eens
enkele zaken classificeren. Alles (moet gedaan worden) om ietwat
wetenschappelijk over te komen, niet waar. In dit geval gaan we eens
kijken hoe we, algemeen gezien, de typische aanvallen kunnen
classificeren die kunnen optreden in de cyber security wereld. Specifiek
zullen we met drie personages werken:

\begin{itemize}
\tightlist
\item
  Alice en Bob: zij zijn de \emph{goeie} en willen op een veilige manier
  met elkaar communiceren.
\item
  Eve: zij is de aanvaller/hacker/crimelord/snoodaard die het leven van
  Alice en Bob zuur wil maken in haar voordeel. Mogelijk wil ze te weten
  komen wat Alice en Bob met elkaar afspreken, misschien wil ze ervoor
  zorgen dat de berichten van Alice niet bij Bob aankomen, etc.
\end{itemize}

\begin{figure}
\centering
\includegraphics[width=0.75\linewidth,height=\textheight,keepaspectratio,alt={Meet the crew.}]{crypto/cast.png}
\caption{Meet the crew.}
\end{figure}

\begin{warningblock}
Herinner je even aan de McCumber kubus: Alice en Bob zijn eender welk
start- en eindpunt van onze informatie, in welke vorm dan ook. Als je
bijvoorbeeld C.I.A. vereist in een computer dan is Alice bijvoorbeeld je
CPU en Bob het RAM-geheugen (om maar iets te zeggen). Kortom, probeer
altijd al deze informatie breed genoeg te plaatsen en niet alleen aan
het klassieke ``netwerkcommunicatie''-systeem te denken waarin Alice
over een netwerk een boodschap naar Bob stuurt.
\end{warningblock}

Eve kan op allerlei manieren \emph{aanvallen} en in de eerste plaats kan
dat \textbf{actief} of \textbf{passief} zijn: bij passieve aanvallen zal
Eve enkel \emph{luisteren} (en wachten) op de communicatie tussen Alice
en Bob. Hierdoor zijn passieve aanvallen veel moeilijker om te
detecteren (soms zelfs onmogelijk), maar uiteraard is Eve 100\%
afhankelijk van wat Bob en Alice doen, daar ze geen dwingende hand heeft
in hun communicatie. Bij actieve aanvallen is dat omgekeerd: de pakkans
is groter (en afhankelijk van het type aanval), maar ze kan ook mogelijk
de communicatie tussen Bob en Alice sturen in de richting die zij nodig
heeft om haar aanval uit te voeren.

\begin{figure}
\centering
\includegraphics[width=1\linewidth,height=\textheight,keepaspectratio,alt={Passieve vs actieve aanvallen.}]{intro/pasact.png}
\caption{Passieve vs actieve aanvallen.}
\end{figure}

\subsection{Passieve aanvallen}\label{passieve-aanvallen}

Bij een \textbf{sniffing} aanval gebruikt Eve een \emph{sniffer}
(bijvoorbeeld Wireshark indien ze netwerk-traffiek wenst te meten) om
alle communicatie tussen twee eindpunten te zien. Alhoewel data steeds
vaker geëncrypteerd wordt, zal Eve toch vaak erg nuttige informatie uit
deze moeilijk te detecteren aanval kunnen halen. Denk maar aan
MAC-adressen, algemene gebruikersinfo, etc. We staan er niet altijd bij
stil hoeveel netwerk-traffiek tegenwoordig constant over netwerken over
en weer vliegt. Daarbij komt ook een iets recenter fenomeen: de minder
veilige third-party apps van bekende merken. Applicaties gemaakt door
derden volgen mogelijk niet altijd de strenge beveiligingscriteria van
het bedrijf waarvoor ze een app hebben gemaakt. Hierdoor bestaat er de
kans dat sommige apps zelfs flagrante fouten maken en bijvoorbeeld user
credentials onbeveiligd opslaan, of erger, over het netwerk sturen. Dit
soort apps maken het werk voor Eve die aan het sniffen is dan ook erg
gemakkelijk.

\begin{figure}
\centering
\includegraphics[width=0.7\linewidth,height=\textheight,keepaspectratio,alt={Passieve aanval, type 1: Sniffing.}]{intro/sniffing.png}
\caption{Passieve aanval, type 1: Sniffing.}
\end{figure}

Soms is geen traffiek kunnen sniffen ook nuttig voor Eve. Later
behandelen we nog side-channel aanvallen, maar we bespreken nu toch al
deze vaak vergeten broer van de sniffing-aanval: de \textbf{trafiek
analyse}. Door het registreren wanneer en hoe een doel communiceert kan
Eve ook erg veel informatie op een passieve manier te pakken krijgen. In
de eerste plaats laat het de stroper toe om te weten wanneer een
gebruiker actief is en wanneer niet. Sommige systemen laten een alarm
afgaan als ze zien dat een legale gebruiker op een onverwacht moment
actief is, iets waar Eve nu rekening mee kan houden. Voorts laat het Eve
ook toe om te ontdekken wat voor activiteiten het slachtoffer gebruikt
(zijn er veel e-mail-gerelateerde berichten? Of net veel VoIP-calls?).

\begin{figure}
\centering
\includegraphics[width=0.7\linewidth,height=\textheight,keepaspectratio,alt={Passieve aanval, type 2: Trafiek analyse.}]{intro/trafanaly.png}
\caption{Passieve aanval, type 2: Trafiek analyse.}
\end{figure}

\subsection{Actieve aanvallen}\label{actieve-aanvallen}

Het domein van de actieve aanvallen is natuurlijk het domein waar Eve de
meeste slaagkansen zal produceren, maar ze heeft ook een veel hogere
kans op gevat te worden. Om die kans te verkleinen zal de stroper bijna
altijd de aanval uitvoeren door zich als iemand anders voor te doen:
\emph{masquerading}. Via \textbf{spoofing} zal de stroper de digitale
identiteit van een legitieme gebruiker overnemen (denk maar aan
MAC-spoofing waarbij Eve het hardware adres van een bedrade of draadloze
netwerk-kaart overneemt). Masquerading heeft een dubbel doel:

\begin{enumerate}
\def\labelenumi{\arabic{enumi}.}
\tightlist
\item
  Het zal de daaropvolgende aanvallen moeilijker kunnen linken aan Eve,
  daar ze onder een pseudoniem actief is.
\item
  Het zal Eve mogelijk toegang verschaffen tot bronnen waar ze onder
  haar eigen \emph{identiteit} niet de juiste rechten toe heeft.
\end{enumerate}

\begin{figure}
\centering
\includegraphics[width=0.7\linewidth,height=\textheight,keepaspectratio,alt={Actieve aanval, type 1: Masquerading.}]{intro/maske.png}
\caption{Actieve aanval, type 1: Masquerading.}
\end{figure}

Het tweede type actieve aanvallen zijn \textbf{replay}-attacks. Hierbij
zal de stroper eerder bewaarde, legitieme, communicatie heruitzenden in
de hoop dat de ontvanger er zich geen vragen bijstelt. Beeld je in dat
Eve een login-pakket heeft gesnift van een erg zwak beveiligd systeem:
als Eve de volgende dag wil inloggen onder de naam van haar slachtoffer
dan hoeft ze enkel dat bewaarde pakket opnieuw te versturen.

\begin{figure}
\centering
\includegraphics[width=0.7\linewidth,height=\textheight,keepaspectratio,alt={Actieve aanval, type 2: Replay attack.}]{intro/replay.png}
\caption{Actieve aanval, type 2: Replay attack.}
\end{figure}

Type 3 is vanuit het standpunt van de aanvaller de interessantste: de
\textbf{man-in-the-middle} of \textbf{MitM}-aanval. Hierbij zal Eve zich
tussenin de communicatie van Bob en Alice nestelen met als doel op een
onzichtbare manier hun communicatie te lezen, aanpassen of blokkeren.
Het laat Eve toe als een soort \emph{puppetmaster} de volledige
communicatie te bepalen en beïnvloeden. Deze aanval is erg krachtig,
maar vereist ook vaak een stevige technische opbouw door Eve daar ze nu
twee eindpunten heeft die ze met behulp van onder andere masquerading
moet aanvallen.

\begin{figure}
\centering
\includegraphics[width=0.7\linewidth,height=\textheight,keepaspectratio,alt={Actieve aanval, type 3: Man-in-the-middle aanval.}]{intro/mitm.png}
\caption{Actieve aanval, type 3: Man-in-the-middle aanval.}
\end{figure}

Als laatste de meest voorkomende aanval: de \textbf{Denial-of-Service}
(DoS). Deze aanval heeft als doel om een systeem of gebruiker \emph{lam
te leggen} zodat deze niet meer voor andere gebruikers of systemen
bereikbaar is. De reden om een DoS uit te voeren zijn velerlei en de
manier waarop deze uitgevoerd kan worden is ook quasi eindeloos: de
stekker uittrekken, gigantische hoeveelheden communicatie versturen, of
het signaal verstoren met een microgolf-oven. Alles is mogelijk en het
hangt vooral van de creativiteit van de aanvaller af hoe effectief de
aanval is.

\begin{figure}
\centering
\includegraphics[width=0.7\linewidth,height=\textheight,keepaspectratio,alt={Actieve aanval, type 4: Denial-of-Service aanval.}]{intro/dos.png}
\caption{Actieve aanval, type 4: Denial-of-Service aanval.}
\end{figure}

\section{Hoe verdedigen}\label{hoe-verdedigen}

Wat en wie je wilt verdedigen in de cyberwereld kan erg gevarieerd zijn.
Toch kunnen we alles herleiden tot 4 + 1 fundamentele principes die je
best hanteert indien je een systeem van welke vorm ook wenst te
beschermen tegen cyberstropers. Deze zijn:

\begin{itemize}
\tightlist
\item
  \textbf{Layering}: bouw je beveiliging zoals de lagen van een ui
  rondom je te beschermen data, gebruikers en services. Hoe meer lagen
  hoe beter, op voorwaarde dat ze natuurlijk verschillend zijn. Het
  voordeel van met meerdere lagen werken is evident: indien één laag om
  welke reden dan ook gecompromitteerd raakt, zijn er nog steeds de
  andere lagen die als verdediging werken.
\item
  \textbf{Limiting}: beperk steeds maximaal wat iedereen binnen je
  systeem kan. Zorg ervoor dat gebruikers en services enkel die zaken
  kunnen doen waartoe ze recht hebben volgens hun rol. Geef dus niet
  iedereen admin-rechten, zet je firewall niet op \emph{allow all}, etc.
\item
  \textbf{Diversity}: zorg ervoor dat je een verscheidenheid aan
  beveiligingen hebt. Op die manier voorkomen we, net als bij layering,
  dat het falen van één systeem je hele verdediging neerhaalt.
\item
  \textbf{Simplicity}: KISS, oftewel \emph{keep it simple, stupid}. Al
  het voorgaande lijkt te doen uitschijnen dat je complexe systemen moet
  bouwen die als het ware een doolhof voor de aanvallers maken. Op zich
  is daar iets van aan, maar zorg er wel voor dat je niet zelf in je
  doolhof verdwaalt en daardoor fouten introduceert zonder het te
  beseffen. Soms is \emph{less more} en dat geldt ook bij beveiliging.
\end{itemize}

\begin{cautionblock}
Het vijfde fundamentele principe krijgt een eigen kadertje omdat deze
voor discussie vatbaar is én geregeld voor de nodige controverse kan
zorgen. Als je dus, om welke reden dan ook, niet genoeg tijd of budget
hebt om alle principes toe te passen, probeer dan de volgende als
laatste ``oplossing'' te gebruiken:

\begin{itemize}
\tightlist
\item
  \textbf{Obscurity}: ga niet aan de grote klok hangen hoe jouw
  verdediging is opgebouwd. Echter, let op met deze leuze. Obscurity wil
  niet zeggen dat je zelf een of ander vaag crypto-algoritme zelf gaat
  ontwikkelen en angstvallig gaat geheim houden. Gebruik standaarden en
  vertrouwde producten in je beveiliging. Eén van de eerste regels in
  security is ``ga niet zelf het wiel heruitvinden''!
\end{itemize}

\end{cautionblock}

\begin{tipblock}
Soms zal je digitale stropers horen spreken over ``Ik heb dat systeem
gepwnd'', uitgesproken als \emph{ge-powned}. De term ``to pwn'' is
\emph{hacker-slang} voor ``to own'' om aan te geven dat je in een
systeem bent binnen geraakt en nu controle over het systeem hebt.
Volgens de
\href{https://www.urbandictionary.com/define.php?term=pwnd}{urbandictionary.com}
is de enige reden dat de \emph{o} een \emph{p} werd het gevolg van een
typfout, daar beide letters vlak naast elkaar staan op een toetsenbord.
\end{tipblock}

\subsection{Je verdediging ontwerpen}\label{je-verdediging-ontwerpen}

Er zijn tal van boeken en papers geschreven over hoe je je vervolgens
moet verdedigen. De voorgaande 4 + 1 principes zijn een start, maar nog
iets te algemeen. Omdat we niet alle design principes kunnen beschrijven
bespreken we hier een \emph{gouwe ouwe}: de \textbf{Saltzer en
Schroeder's design principes} naar twee Amerikaanse
computerwetenschappers die in 1975 dit beschreven in een artikel
getiteld \emph{The Protection of Information in Computer Systems}. Het
artikel wordt ook wel eens het ``meest geciteerde, maar minst gelezen
artikel in het domein'' genoemd.

De principes houden nog steeds stand en zijn de volgende
(\href{https://en.wikipedia.org/wiki/Saltzer_and_Schroeder\%27s_design_principles}{droog
overgenomen van Wikipedia}\ldots want ook wij hebben het originele
artikel niet gelezen):

\begin{itemize}
\tightlist
\item
  \textbf{Economy of mechanism}: Keep the design as simple and small as
  possible.
\item
  \textbf{Fail-safe defaults}: Base access decisions on permission
  rather than exclusion.
\item
  \textbf{Complete mediation}: Every access to every object must be
  checked for authority.
\item
  \textbf{Open design}: The design should not be secret.
\item
  \textbf{Separation of privilege}: Where feasible, a protection
  mechanism that requires two keys to unlock it is more robust and
  flexible than one that allows access to the presenter of only a single
  key.
\item
  \textbf{Least privilege}: Every program and every user of the system
  should operate using the least set of privileges necessary to complete
  the job.
\item
  \textbf{Least common mechanism}: Minimize the amount of mechanism
  common to more than one user and depended on by all users.
\item
  \textbf{Psychological acceptability}: It is essential that the human
  interface be designed for ease of use, so that users routinely and
  automatically apply the protection mechanisms correctly.
\item
  \textbf{Work factor}: Compare the cost of circumventing the mechanism
  with the resources of a potential attacker.
\item
  \textbf{Compromise recording}: It is sometimes suggested that
  mechanisms that reliably record that a compromise of information has
  occurred can be used in place of more elaborate mechanisms that
  completely prevent loss.
\end{itemize}

\section{Social Engineering}\label{social-engineering}

In het eerste hoofdstuk kwam de term ``Social engineering'' al enkele
keren voor. We zagen ook bij de McCumber kubus dat we niet enkel op
technologie mogen rekenen wanneer we onze verdediging opzetten, maar dat
we ook een heel belangrijke schakel moeten trainen en opvolgen: de
mensen. Mensen zijn meestal de zwakste schakel in ons securitymodel en
dus daarom ook een interessante ``aanvalsvector'' voor digitale
stropers.

Social engineering wordt ook wel eens \emph{het hacken van mensen}
genoemd en is een erg laagdrempelige, maar oh zo effectieve
aanvalstechniek die vaak over het hoofd wordt gezien. Bij social
engineering zal de aanvaller \emph{menselijke} interacties gebruiken om
zijn slachtoffers onbewust zaken te laten doen of informatie te geven
dat niet zou mogen. Hierbij misbruikt de aanvaller onze aangeboren
gewoonte om mensen te vertrouwen in plaats van te wantrouwen. We gaan
meestal uit \emph{van het goede van onze medemens} en zullen vaak
meerdere goede redenen kunnen verzinnen waarom iemand iets van jou nodig
heeft. Andere menselijke trekken die we kunnen gebruiken als social
engineer zijn onder andere de nieuwsgierigheid van mensen, hun hebzucht,
onwetendheid of angst.

Enkele voorbeelden:

\begin{itemize}
\tightlist
\item
  Je verkleden als pizzakoerier en een grote stapel (lege) pizzadozen
  het bedrijf binnendragen. Werknemers zullen je willen helpen en de
  deur voor je openhouden, ook al moet je normaal gezien met je
  toegangsbadge het gebouw betreden.
\item
  Bij de rokers aan de achterkant van het gebouw gaan staan, wat met hen
  keuvelen en hen sigaretje aanbieden. Wanneer ze terug binnengaan hen
  volgen (\emph{piggybacking}).
\item
  Bellen en je voordoen als een technieker van Telenet en vervolgens de
  login-gegevens vragen van het slachtoffer.
\end{itemize}

Uiteraard hoeven social engineers zich niet te beperken tot
\emph{fysiek} contact. Ze hanteren ook erg vaak geschreven teksten
(e-mail) om aan social engineering te doen. Hierbij onderscheiden we
twee soorten:

\begin{itemize}
\tightlist
\item
  \textbf{Phishing}: trachten zoveel mogelijk mensen naar een bepaalde
  website of document te lokken waar vervolgens de aanvaller informatie
  van het slachtoffer zal proberen te verkrijgen door een fake login
  scherm te tonen, malware ongezien te installeren, etc.
\item
  \textbf{Spear phishing}: deze vorm van phishing heeft als doel één
  persoon of groep. De e-mail zal dan ook opgesteld worden om specifiek
  voor het slachtoffer te werken, i.p.v. een generieke mail.
\end{itemize}

\begin{tipblock}
De term \emph{phishing} is afgeleid van \emph{fishing} oftewel
vissen/hengelen naar iets. Vroeger had je het concept \emph{phreaking}
dat werd toegepast door hackers om op telefooncentrales in te breken
door de 2600 hertz fluittonen die telefoons gebruiken te imiteren.
\end{tipblock}

\subsection{SET}\label{set}

Een veel gebruikte tool die social engineers hanteren is de
\textbf{Social Engineering Toolkit} (beschikbaar via Kali en
\href{https://github.com/trustedsec/social-engineer-toolkit}{github}) en
zal je helpen om (spear) phishing attacks op te zetten, fake websites
(door bestaande te clonen) te hosten, malware in afbeeldingen te
injecteren, etc. Het kan op de koop toe geïntegreerd samenwerken met
Metasploit waardoor stropers erg complexe aanvallen kunnen opzetten.

\subsection{OSINT}\label{osint}

We zagen reeds de vijf fasen die een aanvaller doorloopt: verkennen,
scannen, etc. Aangezien Social engineering ook een vorm van cyber-aanval
is, zullen ook hier deze fasen worden doorlopen. Zeker bij spear
phishing wil de aanvaller zoveel mogelijk informatie over zijn
slachtoffer(s) te weten komen om zo de meest doeltreffende mail op te
stellen. OSINT, oftewel Open Source INTelligence, is de techniek waarbij
een aanvaller open-source bronnen gebruikt om zijn slachtoffers in kaart
te brengen tijdens deze eerste verkennings- of reconnaissance fase.

Enkele nuttige technieken die worden toegepast:

\begin{itemize}
\tightlist
\item
  Google en andere searchengines: besef dat je veel meer informatie kunt
  vinden indien je ook geavanceerde search-eigenschappen gebruikt (denk
  aan ``site:'' en de ``+''-operator, etc. in Google).
\item
  Reverse image searching: er zijn tegenwoordig erg krachtige tools om
  de oorsprong van een afbeelding te traceren, en met bijvoorbeeld
  Google Lens kan je zelfs objecten en locaties op een foto
  identificeren.
\item
  Metadata van bestanden: media bestanden (.docx, .jpg, .png, .mp4,
  etc.) bevatten ook aardig wat onzichtbare informatie die soms
  onbedoeld meereist met het bestand (denk bijvoorbeeld aan de EXIF data
  in een afbeelding) en dus door kwaadwillige personen misbruikt kunnen
  worden.
\end{itemize}

Het \textbf{OSINT Framework} (via
\href{https://osintframework.com/}{osintframework.com}) is ontwikkeld om
mensen er op te wijzen hoeveel publieke informatie social engineers (en
anderen) over iemand te weten kunnen komen. Het is een griezelig
uitgebreide hoeveelheid online bronnen die mooi gecategoriseerd zijn.
Hierbij is het op de koop toe belangrijk te beseffen dat enkel
open-source bronnen worden gebruikt. Digitale stropers zullen zich
uiteraard niet beperken tot enkel open-source bronnen en ook vaak
betaalde of illegaal verkregen bronnen raadplegen.

\section{Soorten aanvallen}\label{soorten-aanvallen}

\subsection{Software-based aanvallen}\label{software-based-aanvallen}

Software-gebaseerde aanvallen zijn de aanvallen die je in het nieuws
geregeld hoort. In deze groep vinden we de virussen, wormen, backdoors,
trojans en alle andere \textbf{malware} terug. De term malware dekt de
lading erg goed: \emph{``any kind of malicious software designed to
damage or harm a computer system.''} Het doel van malware is altijd
geld, macht (denk aan blackmailen) of pestgedrag. De verschillende
malware-types hebben soms overlappende eigenschappen en het is dus niet
altijd mogelijk om een specifiek stuk malware onder één categorie te
plaatsen. Wat ze allemaal gemeen hebben is dat ze bepaalde gekende of
ongekende (\emph{zerodays}) fouten of bugs (\emph{kwetsbaarheden} of
\emph{vulnerabilities}) misbruiken in software om zo toegang tot een
systeem, data of stuk hardware te verkrijgen.

\begin{tipblock}
\textbf{CVE oftewel ``Common Vulnerabilities and Exposures''} is een
publiek beschikbare lijst waarin alle gekende security kwetsbaarheden
opgelijst staan. Het laat cyberboswachters toe om duidelijk te
communiceren over problemen en oplossingen. Volgens een
\href{https://www.cisa.gov/uscert/ncas/alerts/aa20-133a}{rapport} van de
Amerikaanse ``Cybersecurity \& Infrastructure Security Agency'' (CISA)
was de meest misbruikte CVE tussen 2016 en 2019 een kwetsbaarheid
(CVE-2017-11882) in Microsoft Office 2013 die samengevat: \emph{``allow
an attacker to run arbitrary code in the context of the current user by
failing to properly handle objects in memory, aka `Microsoft Office
Memory Corruption Vulnerability'.''}
\end{tipblock}

We overlopen nu de belangrijkste vormen van malware.

\subsubsection{Virus}\label{virus}

De moeder van de malware. Hiermee is alles begonnen. De eerste virussen
werden geschreven door scriptkiddies en mensen met te veel tijd. Ze
hadden zelden een specifiek doel, en wilden gewoon zien wat de kracht
van een virus was en hoe ambetant het anderen zou maken. De eerste
virussen, in pre-Internet tijd, waren afhankelijk van overdracht via
diskettes en raakten dus niet snel verspreid. Virusscanners waren nog
onbestaande, maar als je een virus te pakken had dan was dat uiteraard
even vloeken: sommige virussen infecteerden je bootsector, zodat je
computer niet meer bootte, andere verwijderden bepaalde
systeembestanden, enz. Ze werden geactiveerd door een uitvoerbaar
bestand uit te voeren (zogenaamde .exe, .bat of .com bestanden) waar het
virus zich onzichtbaar aan had vastgeklampt.

\begin{tipblock}
De auteur van dit boek heeft z'n vader een hoop extra grijze haren
gekost door z'n game-verslaving. In de jaren 90 was de enige manier om
aan games te geraken ofwel via de één of twee computerwinkels in de
provincie, oftewel door diskettes van klasgenoten te kopiëren. Geregeld
bevatte die gekopieerde games echter ook virussen, met alle gevolgen van
dien. Moraal van het verhaal: \emph{don't pirate, kids ;)}.
\end{tipblock}

Heden ten dage komen we nog maar weinig ``klassieke'' virussen tegen. Je
kan zeggen dat ze zijn geëvolueerd naar veel lastigere, soms letterlijke
dodelijke varianten zoals wormen, ransomware, etc.

\subsubsection{Wormen}\label{wormen}

Een virus is een beetje zoals een giftige vis op het droge: het ligt
maar wat in het zand te spartelen en enkel als je zo dom bent om het
ding op te rapen en in je aquarium met zeldzame vissen te plaatsen zal
het schade kunnen toebrengen. Wormen daarentegen zijn de
\emph{sharknados} van de giftige vissen op het droge. Wormen zijn
virussen die zichzelf kunnen voortplanten via één of meerdere
communicatiekanalen. Uiteraard in de eerste plaats denken we dan aan het
Internet as is, maar ook via e-mail, WhatsApp-berichten, Bluetooth,
Facebook, etc.

Als bijkomende handigheid zullen wormen ook vaak zichzelf veranderen
zodat ze moeilijker door virusscanners kunnen gedetecteerd worden.
Voorts hebben wormen geen \emph{hostfile} nodig wat virussen wel hebben.
En als laatste verschil met de virusjes is dat wormen zichzelf ook
kunnen verspreiden zonder dat de gebruiker een (on)bewuste handeling
moet doen. Kortom, wormen zijn een pittig probleem, vooral vanwege de
snelheid (en eenvoud) waarmee ze zich verspreiden.

\begin{figure}
\centering
\includegraphics[width=0.65\linewidth,height=\textheight,keepaspectratio,alt={Worm propagation model. Bron: Network Eye: End-to-End Computer Security Visualization - Scientific Figure on ResearchGate.}]{intro/wormprop.png}
\caption{Worm propagation model. Bron: Network Eye: End-to-End Computer
Security Visualization - Scientific Figure on ResearchGate.}
\end{figure}

Bovenstaande afbeelding toont het \emph{worm propagation model}. Dit
model toont hoe snel een worm zich kan verspreiden en hoe belangrijk het
is om een nieuwe worm zo snel mogelijk te detecteren, in de hoop
voldoende systemen vervolgens te beschermen tegen besmettingen ervan.
Naarmate een worm meer systemen kan besmetten, die op hun beurt dan ook
als uitvalsbasis van de worm dienen, merk je een exponentiële groei van
besmette systemen. Deze groei bereikt een plateau wanneer quasi alle
systemen besmet zijn die konden besmet worden.

\subsubsection{Trojan}\label{trojan}

Trojans zijn malware die zich verstoppen binnenin een legale, al dan
niet nuttige, applicatie die de gebruiker bewust installeert of
downloadt. Van zodra de gebruiker deze host-applicatie installeert of
start zal de onderliggende trojan zijn aanval op het systeem beginnen.
Die aanval kan velerlei zijn, denk maar aan het installeren van een
backdoor, de computer veranderen in een zombie (zie botnet verderop) of
een keylogger activeren.

De naam Trojan is gebaseerd op de legende van Het Paard van Troje uit de
Aeneid van Vergilius.

\subsubsection{Spyware}\label{spyware}

De naam spyware dekt de lading duidelijk: deze malware heeft als doel om
informatie te verzamelen van het doelsysteem, zonder dat de gebruiker
hiervan op de hoogte is. Deze informatie kan gaan van eenvoudige
gebruikersdata zoals usernames en wachtwoorden, maar gaat soms ook
verder tot het in kaart brengen van surfgedrag of op welke manier de
gebruiker een specifieke applicatie gebruikt. We kunnen stellen dat
spyware twee categorieën heeft:

\begin{enumerate}
\def\labelenumi{\arabic{enumi}.}
\tightlist
\item
  Reclamedoeleinden: hoe beter adverteerders de gebruiker kennen, hoe
  gerichter ze reclame kunnen maken en dus hoe groter de kans wordt dat
  die gebruiker uiteindelijk een geadverteerd product koopt.
\item
  Cybercriminelen: de meeste spyware die we heden ten dage tegenkomen
  verzamelen natuurlijk die dingen waar cyberstropers de meeste
  interesse in hebben: bankgegevens (logins, kredietkaartinformatie).
  Het hoofddoel van deze categorie is natuurlijk: centjes!
\end{enumerate}

Spyware wordt meestal als trojan geïnstalleerd nadat de gebruiker een
legitiem programma installeerde. Enkele bekende spyware-dragers waren de
erg populaire (illegale downloader) Kazaa en Messenger Plus!, de plugin
voor Windows Live Messenger.

\subsubsection{Adware}\label{adware}

Daar waar spyware soms door adverteerders wordt gebruikt om ``het
doelwit beter te leren kennen'', heeft adware als doel om effectief
reclame aan deze gebruiker, meestal ongewild te tonen. Ongewild is het
sleutelwoord hier in. Adware is op zich niet noodzakelijk malware. Veel
adware wordt gemaakt zodat de maker ervan extra inkomsten kan genereren
om bijvoorbeeld het hoofddoel van de adware te blijven door ontwikkelen.
Denk maar aan bepaalde gratis \emph{musicplayers} die ook reclamebanners
tonen terwijl je muziek afspeelt. Er is echter ook een categorie adware
die \textbf{ongewild} reclame toont en soms zelfs op een zodanige manier
dat een huis-tuin-en-keuken gebruiker zelfs niet beseft waar de reclame
vandaan komt. Eind jaren 2010 had je veel browser-extensies die in je
menu-balk als flitsende knoppen allerlei handige extra hulpmiddelen
beloofden. Vaak zorgden deze extensies er echter ook voor dat je tal van
extra reclame-pop-ups kreeg die in eerste instantie het gevolg waren van
de website waar je op dat moment naartoe surfte.

\begin{tipblock}
Een bijkomend probleem van sommige malware is dat ze
\emph{rootkit}-achtige (zie hierna) technieken gebruiken waardoor ze erg
moeilijk te verwijderen zijn en je moet opletten dat je niet essentiële
bestanden van je computer verminkt of verwijderd.
\end{tipblock}

\begin{figure}
\centering
\includegraphics[width=0.4\linewidth,height=\textheight,keepaspectratio,alt={Een wat overdreven voorstelling van het soort browsers vol adware die de auteur soms te zien kreeg toen hij computers ging repareren die ``nogal traag waren''. Bron: https://www.directive.com/blog/disabling-those-pesky-browser-toolbars.html.}]{intro/adware.jpg}
\caption{Een wat overdreven voorstelling van het soort browsers vol
adware die de auteur soms te zien kreeg toen hij computers ging
repareren die ``nogal traag waren''. Bron:
https://www.directive.com/blog/disabling-those-pesky-browser-toolbars.html.}
\end{figure}

\subsubsection{Rootkit}\label{rootkit}

Een virus zal zich nestelen op ``gebruikersniveau'', terwijl een rootkit
zich veel dieper in het besturingssysteem zal nestelen (op
\emph{kernel-niveau} of administrator-niveau). Hierdoor zijn rootkits
veel onzichtbaarder, daar ze ook kunnen bepalen welke informatie de
gebruiker van het besturingssysteem te zien krijgt. Een rootkit is
eerder een techniek die kan gebruikt worden door de andere vormen van
malware die we hier beschrijven. De essentie van een rootkit is
natuurlijk dat deze veel robuuster zijn en bijgevolg moeilijker te
verwijderen zijn door het slachtoffer. Vaak zal een rootkit ook
systeembestanden aanpassen (zie ook de infobox bij adware) waardoor je
rootkits niet kunt verwijderen zonder permanente schade aan je
besturingssysteem toe te brengen. De enige manier om een rootkit dan
goed weg te krijgen is door je harde schijf te formatteren. Je OS
opnieuw installeren/resetten, gebruik makend van de aanwezige
installatiebestanden op de computer, is namelijk niet gegarandeerd dat
dit voldoende zal zijn: mogelijk heeft de rootkit zich ook al in de
installatiebestanden op de harde schijf geïnstalleerd!

\subsubsection{Ransomware}\label{ransomware}

De plaag van de laatste jaren! Het einddoel van cybercriminelen is
natuurlijk centjes. Als je virus al je bestanden verwijderd dan heeft
een digitale stroper weinig \emph{leverage} om nog geld van z'n
slachtoffer te pakken te krijgen. Ransomware lost dit probleem op voor
de stropers: het zal de data letterlijk gijzelen en losgeld
(\emph{ransom}) vragen aan de gebruiker. Wanneer een ransomware op een
systeem geraakt (via bijvoorbeeld een trojan of worm) zal het de data
van de gebruiker versleutelen met een sleutel die enkel de maker van de
malware kent. Vervolgens verschijnt er een bericht op de computer met
daarin wat de gebruiker moet doen (betalen) indien deze z'n data terug
wenst. Meestal zal de ransomware betalingen in crypto-coins vragen zodat
het geld niet kan getraceerd worden.

\begin{figure}
\centering
\includegraphics[width=0.7\linewidth,height=\textheight,keepaspectratio,alt={Voorbeeld van een typisch ransomware scherm dat het slachtoffer te zien krijgt (Bron: wikipedia).}]{intro/ransom.jpg}
\caption{Voorbeeld van een typisch ransomware scherm dat het slachtoffer
te zien krijgt (Bron: wikipedia).}
\end{figure}

Ransomware heeft aangetoond dat back-ups maken van je data erg
belangrijk is. Maar ook HOE en WAAR je back-up'd zal invloed hebben op
hoe ransomware-gevoelig je bent. Indien je back-ups maakt op hetzelfde
systeem als waar de originele data staat, dan bestaat de kans dat de
ransomware ook de back-up zal \emph{gijzelen}. Indien je niet op
geregelde tijdstippen een back-up neemt kan het zijn dat je dagen of
weken aan data kwijt bent moest je het slachtoffer van een
ransomware-aanval zijn. We gaan niet verder in op back-up strategieën,
maar het mogelijk duidelijk zijn dat deze \emph{skillset} een essentieel
onderdeel vormt van het \emph{security-beleid} van een bedrijf.

\begin{tipblock}
Herinner je dat de ransomware-aanvallen uit hoofdstuk 1 (WannaCry en
Petya in 2017) ook meer impact hadden dan enkel ``dataverlies'':
ziekenhuizen moesten patiënten de toegang ontzeggen, containerbedrijven
zaten met duizenden tonnen aan vracht die niet verscheept geraakten.
\end{tipblock}

\subsubsection{Botnet}\label{botnet}

\emph{Wat als een stroper toegang had tot een legioen computers?
Duizenden computers die naar het bevel van de cybercrimineel luisteren
en zonder morren doen wat hen gevraagd wordt? Welkom in de wondere
wereld van botnets, zombies en herders.}

Wanneer een cybercrimineel een grote groep computers nodig heeft, dan
zal deze via voorgaande malware-technieken proberen zoveel mogelijk
computers te besmetten met zijn \emph{zombie-virus}. Dit virus zal twee
dingen doen:

\begin{enumerate}
\def\labelenumi{\arabic{enumi}.}
\tightlist
\item
  Het zal zich onzichtbaar nestelen op de computer en ervoor zorgen dat
  het virus ook na heropstart actief is (\emph{maintain access}).
\item
  Het zal een backdoor creëren en terugbellen naar de \emph{command-en
  control-server} (\textbf{C\&C server}) van de originele virusmaker, de
  zogenaamde \emph{herder}.
\end{enumerate}

Via de C\&C-server kan de botnet-herder nu alle zombies benaderen en
bevelen geven. Het kan de zombies bijvoorbeeld vertellen dat ze allemaal
tegelijkertijd naar een bepaalde website moeten pingen, waardoor een
gigantische DDOS (\emph{Distributed DoS})-aanval plaatsvindt. Of het zou
kunnen bevelen dat alle zombies vijf sterren moeten geven aan een
specifieke applicatie in de app-store. Of wat te denken van duizenden
zombies die permanent als cryptominers naar bitcoins delven voor de
herder?

\begin{figure}
\centering
\includegraphics[width=0.6\linewidth,height=\textheight,keepaspectratio,alt={Een echte botnet zal veel meer dan vier zombies bevatten. Maar we moeten ergens mee beginnen, nietwaar.}]{intro/botnet.png}
\caption{Een echte botnet zal veel meer dan vier zombies bevatten. Maar
we moeten ergens mee beginnen, nietwaar.}
\end{figure}

Hoe groter het botnet, hoe krachtiger en machtiger de herder is. Botnets
kunnen zoveel impact hebben op systemen dat er een hele ecologie rond
illegale botnets is verschenen. Zo zijn er websites waar herders de
diensten van hun botnets verhuren aan anderen
(\emph{botnets-as-a-service}) of gewoon geregeld hun nieuwst verzamelde
botnet verkopen aan de hoogste bieder op het darkweb.

Het is in het voordeel van de herder dat botnets zo onzichtbaar mogelijk
blijven. Daarom dat de meeste botnet-software heel subtiel op de
achtergrond werkt. Veel computers maken maanden, soms jaren, deel uit
van een botnet zonder dat ze dat ooit hebben beseft.

Omdat botnets zo'n grote impact kunnen hebben, jagen Microsoft, Cisco,
McAfee, etc. actief op deze zaken. Een botnet uitschakelen door de
zombies te bestrijden is natuurlijk onbegonnen werk. De oplossing ligt
natuurlijk bij de C\&C-servers! Als je die server uit de lucht krijgt
dan zijn de zombies nutteloos en heb je letterlijk het botnet onthoofd.

\subsection{Netwerk-based aanvallen}\label{netwerk-based-aanvallen}

Een groot deel van de aanvallen gebeurt uiteraard via een bedraad of
draadloos netwerk. We spenderen een volledig apart hoofdstuk aan
draadloze aanvallen verderop in dit handboek. De meer \emph{klassieke}
netwerk-gebaseerde aanvallen komen niet in dit handboek voor. Er zijn
ongelooflijk veel mogelijkheden op netwerk/communicatie-niveau om als
cybercriminal toegang tot systemen te verkrijgen. Ieder bekend protocol
(DNS, IP, TCP, CSMA/CD, SNMP, etc.) heeft ontelbare, gekende, bugs. Nog
steeds worden er nieuwe technieken ontwikkeld, zelfs bij protocollen die
al 20 tot 30 jaar bestaan (om je een idee te geven: het IP-protocol werd
in 1977 geschreven).

\subsection{Hardware-based aanvallen}\label{hardware-based-aanvallen}

\emph{In den ouden tijd} leken aanvallen vooral een software gebeuren.
Enkel in de duurdere Hollywood-films werden er ingewikkelde
hardware-apparaten gebruikt door stropers om toegang tot streng
beveiligde systemen te verkrijgen. Tegenwoordig kan je voor enkele
dollars ongelooflijk krachtige Raspberry Pi's, Arduino, etc. kopen,
waardoor hardwaregebaseerde aanvallen steeds vaker voorkomen (je zou
kunnen spreken van een democratisering van \emph{hacking hardware}).

Dit hoofdstuk zou wederom een heel eigen boek kunnen bevatten, we gaan
daarom enkele van de meest voorkomende of interessantste zaken kort
toelichten. Net zoals met malware zal je merken dat er soms overlap is
tussen verschillende type aanvallen en het dus zeker geen zwartwit
classificatie is:

\begin{itemize}
\tightlist
\item
  \textbf{USB sticks}: deze kleine dingen kosten nog geen euro als je ze
  in bulk koopt. Ideaal dus om ze als aanvaller te vullen met malware
  die zichzelf automatisch installeert wanneer een slachtoffer de stick
  goedbedoeld in z'n computer steekt. Of wat te denken van een \emph{USB
  of death} die 220 volt door je computer jaagt als je hem insteekt,
  gegarandeerd dat je daarmee een DoS aanval kunt uitvoeren want de
  stick zal het moederbord, de harde schijven etc. vernietigen door de
  hoge stroomstoot.
\item
  \textbf{BIOS rootkits}: in moderne computers zitten BIOS chips die
  voorkomen dat je zomaar eender welke software kan inladen bij het
  opstarten - met dank aan UEFI dat onze systemen beveiligt met systemen
  zoals Trusted Platform Modules (TPM), Secure Boot, etc. Maar wat als
  je er in slaagt om die chip te hacken? Recent dook \emph{Moonbounce}
  op, UEFI malware die aanvallers als een springplank kunnen gebruiken
  om vervolgens andere malware op het systeem te krijgen. De sterkte van
  Moonbounce? Het kan zichzelf in de chip installeren en blijft
  permanent aanwezig zodat je besturingssysteem verwijderen of harde
  schijf formatteren geen zin heeft.
\item
  \textbf{Raspberry Pi, Arduino, Malduino \emph{and friends}}: computers
  die niet groter dan een dikke duim zijn. De mogelijkheden zijn
  natuurlijk immens wat je kan doen als je dit soort dingen onzichtbaar
  kan inplanten in een bedrijf. Doordat dit soort dingen zo goedkoop
  zijn geworden zien we ook vaker concepten zoals \textbf{warshipping}
  opduiken: een aanvaller verstuurt zijn geautomatiseerde Raspberry Pi
  naar z'n slachtoffer door het toestel in een gewoon postpakket te
  verstoppen in bijvoorbeeld een dubbele bodem. Van zodra het pakket
  aankomt, zal het toestel automatisch proberen verbinding te maken met
  het aanwezig netwerk en \emph{terugbellen} naar de aanvallers. Zolang
  het slachtoffer de Raspberry Pi niet detecteert kan de stroper vanuit
  de veilige haven van z'n huis aanvallen uitvoeren.
\item
  \textbf{Keyloggers en \emph{juice hacking}}: keyloggers bestaan zowel
  in software als hardware, maar hun doel is natuurlijk hetzelfde: de
  toetsaanslagen van het slachtoffer detecteren om er dan bijvoorbeeld
  wachtwoorden en gebruikersnamen uit te filteren. Hardware keyloggers
  worden door stropers vaak aan een systeem gehangen waar ze niet
  permanent toegang tot hebben. Enkele dagen of weken later moet de
  stroper dan enkel de keylogger ophalen. Gerelateerd hieraan is juice
  hacking, een fenomeen dat hier en daar opduikt. In publieke plaatsen
  zijn er meer en meer publieke oplaadpunten waar gebruikers hun
  digitale toestellen via een USB-draad kunnen opladen. Maar hoe zeker
  ben je eigenlijk dat die draad die in de muur verdwijnt niets meer is
  dan een \emph{oplaadkabel}? Wat als aan de andere kant van de muur de
  draad aan het toestel van de stroper hangt?
\item
  \textbf{USB Ninja, RFID cloners en consoorten op lab401.com en
  hak5.org}: er zijn vele websites waar je tegenwoordig gespecialiseerde
  \emph{hacking hardware} kunt kopen. Wat te denken van de USB Ninja?
  Een \emph{gewone} USB-kabel die echter een ingebouwde verzender heeft
  en automatisch alle data via een RF-verbinding naar de stroper
  verderop verstuurt. RFID cloners vind je ook voor een habbekrats en
  laten je toe om de alomgebruikte RFID toegangskaarten te clonen, zodat
  de stroper beveiligde gebouwen kan binnenstappen zonder een alarm te
  laten afgaan.
\end{itemize}

\subsection{Side-channel aanvallen}\label{side-channel-aanvallen}

De persoonlijke favoriet van de auteur vanwege de inventieve aanvallen
die onder deze categorie bestaan. Het idee van een side-channel aanval
bestaat er uit dat je informatie te pakken krijgt uit een protocol of
hardware op onverwachte manier. Een vergelijking in het echte leven zou
het volgende kunnen zijn: je wil inbreken bij een bank verderop in de
straat. Dit gaat echter enkel wanneer de bewaker slaapt. Na observatie
heb je ontdekt dat de bewaker voor het slapengaan altijd een boek leest
en z'n progressie ervan op GoodReads deelt. Je hebt geen toegang tot de
slaapkamer van de bewaker, maar je volgt hem wel op GoodReads. Van zodra
de bewaker een update over z'n voortgang post weet je dat het tijd is.
Dit is een voorbeeld van een side-channel aanval.

Ieder stuk hardware, software of protocol is vatbaar voor side-channel
aanvallen. Het is onmogelijk om 100\% hiertegen beschermd te zijn.
Uiteraard zijn niet alle aanvallen even effectief als de andere, en
alles hangt dus af van wat de stroper juist nodig heeft.

Enkele voorbeelden:

\begin{itemize}
\tightlist
\item
  Militairen deelden hun workouts op Strava. Ze waren echter vergeten
  dat ze in een geheime basis in Afrika werkten. Plots zagen mensen
  workouts \emph{in the middle of nowhere} gedeeld worden, wat deed
  vermoeden dat er op die plaats meer was dan enkel een lege woestijn.
  Oeps.
\item
  Onderzoekers zijn er in geslaagd om data uit een computer te krijgen
  door de \emph{power consumption} op te meten.
\item
  Rowhammer, spectre (en ook drammer en rampage) zijn recente technieken
  waarbij de informatie uit het geheugen kan gelezen worden waar de
  stroper eigenlijk geen rechten toe heeft.
\item
  Heartbleed: deze bespraken we al kort in het eerste hoofdstuk. Deze
  inventieve techniek maakt gebruik van het feit dat computers altijd
  braaf antwoorden als je ze vragen stelt. Maar stropers ontdekten dat
  de snelheid van antwoorden afhing van de vraag.
\end{itemize}

\begin{noteblock}
Wist je dat een oud liedje van Janet Jackson kan gebruikt worden om
(oude) laptops te doen crashen door het gewoon af te spelen? Over een
side-channel DoS aanval gesproken. De aanval heeft zelfs een CVE-nummer
toegewezen gekregen! Lees hier hoe de aanval
werkt:\url{https://www.bleepingcomputer.com/news/security/janet-jacksons-music-video-is-now-a-vulnerability-for-crashing-hard-disks/}
\end{noteblock}

\begin{figure}
\centering
\pandocbounded{\includegraphics[keepaspectratio,alt={De briljante xkcd.com strip legt perfect uit hoe Heartbleed werkt.}]{intro/heartbleed.png}}
\caption{De briljante xkcd.com strip legt perfect uit hoe Heartbleed
werkt.}
\end{figure}

\chapter{Cryptografie}\label{cryptografie}

In dit hoofdstuk duiken we de boeiende wereld van de cryptografie in:
het versleutelen van informatie zodat enkel zender en ontvanger het
bericht kunnen lezen.

\section{Encryptie en decryptie}\label{encryptie-en-decryptie}

Herinner je dat we in het vorige hoofdstuk de termen CIA aanhaalden en
de McCumber kubus? Een grote pijler van CIA, confidentiality, wordt
opgelost met behulp van encryptie, namelijk het versleutelen van onze
data met behulp van een geheime sleutel. Door deze te versleutelen wordt
deze onleesbaar voor personen die de geheime sleutel niet hebben (en
bijgevolg niet geautoriseerd zijn om de data te mogen lezen). De
moeilijkheid van een goed cryptografisch systeem is dat de data op een
zodanige manier moet versleuteld worden dat het quasi onmogelijk is om
zonder sleutel de originele data terug te vinden. We spreken hierbij
over de originele data als de \textbf{plaintext} en de geëncrypteerde
data als \textbf{ciphertext}. De ontvanger van een ciphertext moet deze,
als hij de juiste sleutel heeft, terug kunnen omzetten naar de originele
plaintext.

Er zijn al veel encryptie algoritmes de revue gepasseerd doorheen de
geschiedenis van de mens. Al van in de tijd van de Romeinen werd er aan
cryptografie gedaan. Mensen hebben altijd gevoelige data gehad waar
vertrouwelijk mee moest om gesprongen worden. Naarmate de
\textbf{cryptanalyse} (dat is het proberen ontcijferen van een
ciphertext zonder dat je de geheime sleutel hebt) evolueerde, moesten
ook de cryptografische algoritmes verbeteren. Ook hier zien we weer
diezelfde wedloop tussen digitale stropers en cyberboswachters. Hoe
sterker onze computers worden (met dank aan de wet van Moore), hoe
krachtiger onze algoritmes moeten worden. De eenvoudigste vorm van
cryptanalyse, \emph{bruteforcing}, is rechtstreeks afhankelijk van de
snelheid van de computer. Hoe meer sleutels per seconde een computer kan
testen, hoe sneller de originele sleutel kan gevonden worden.

\begin{figure}
\centering
\includegraphics[width=0.9\linewidth,height=\textheight,keepaspectratio,alt={Alice gebruikt encryptie om een beveiligd bericht naar Bob te sturen, zodat Eve deze niet kan lezen.}]{crypto/basicencrypt.png}
\caption{Alice gebruikt encryptie om een beveiligd bericht naar Bob te
sturen, zodat Eve deze niet kan lezen.}
\end{figure}

\begin{tipblock}
De volledige geschiedenis van de cryptografie hier vertellen zou
ongeveer 1200 pagina's vereisen. Het briljante boek ``The Codebreakers''
van David Kahn is een aanrader voor eenieder die meer willen weten over
deze boeiende geschiedenis. Laat de 1200 pagina's je niet afschrikken,
het boek leest als een echte thriller.
\end{tipblock}

Alle bestaande cryptografische systemen kunnen op verschillende manieren
gekarakteriseerd worden (bron Network Security Essentials, door William
Stallings):

\begin{itemize}
\tightlist
\item
  De \emph{acties} die op de data worden uitgevoerd om deze te
  encrypteren:

  \begin{itemize}
  \tightlist
  \item
    \textbf{Substitutie}: een teken door een ander teken vervangen.
  \item
    \textbf{Transpositie}: een teken op een andere plek in de tekst
    zetten.
  \item
    \textbf{Product}: een combinatie van meerdere substituties en
    transposities.
  \end{itemize}
\item
  Het \emph{aantal sleutels} dat nodig zijn:

  \begin{itemize}
  \tightlist
  \item
    \textbf{1 sleutel}, ook wel ``private encryption'' genoemd.
  \item
    \textbf{2 sleutels}, ook wel ``public encryption'' genoemd.
  \end{itemize}
\item
  De \emph{manier} waarop de data wordt verwerkt:

  \begin{itemize}
  \tightlist
  \item
    Als een \textbf{blok} data, blok per blok.
  \item
    Als een \textbf{stream}, teken per teken.
  \end{itemize}
\end{itemize}

\section{Kerckhoffs principe}\label{kerckhoffs-principe}

We gaan zo meteen bekijken hoe encryptie effectief gebeurt, maar we
willen al even een mythe onderuit halen. Binnen encryptie heb je de
\emph{principes van Kerckhoff}, zes regels die in de 19e eeuw werden
vastgelegd waaraan encryptie-algoritmes moeten voldoen om als veilig
beschouwd te worden. Sommige principes zijn, onder ander door onze
steeds krachtigere computers, niet meer relevant, maar één principe
blijft ongelooflijk belangrijk: \emph{``Het kennen van het gebruikte
encryptie-algoritme door de tegenstander is geen probleem, het is enkel
de geheime sleutel die ten allen tijde uit de handen van de tegenstander
moet blijven.''}

Dit ogenschijnlijk eenvoudige zinnetje betekent heel veel: \textbf{je
sleutel (of wachtwoord) is wat je het beste moet beschermen. Zonder
kennis van de sleutel zou iemand nooit toegang mogen krijgen tot
versleutelde data, ongeacht dat geweten is met welk systeem de data werd
vercijferd.}

We zagen reeds dat \emph{Security through obscurity} een dubbel snijdend
zwaard binnen de security wereld is. Enerzijds is het niet aangeraden om
met veel tamtam aan te kondigen hoe jij je data beveiligt. Anderzijds
geeft het je mogelijk een vals gevoel van veiligheid (\emph{snake's
oil}) daar het geheimhouden van je algoritme en systemen geen garantie
is dat deze ook effectief veilig zijn. In de 21e eeuw zijn de meest
gebruikte encryptie-algoritmes publiek gekende algoritmes die door
duizenden experts aan de tand zijn gevoeld. De kans dat er dus (bewuste)
fouten in dergelijke standaarden zitten, is véél kleiner dan wanneer je
met een zogenaamd proprietary systeem werkt waarvan de werking
angstvallig geheim wordt gehouden.

\textbf{Finaal draait alles op het geheimhouden van je sleutel, iets dat
we telkens weer in dit boek zullen herhalen!}

\subsection{Opletten met reclame}\label{opletten-met-reclame}

Let op met encryptiesystemen die zichzelf verkopen met zinnen zoals
\emph{``10 jaar nodig op een gewone laptop om alle sleutels te
testen''}. Dit zou kunnen doen vermoeden dat je dus voor minstens tien
jaar goed zit (we gaan er even vanuit dat de gemiddelde cryptanalist
maar toegang heeft tot één laptop, wat uiteraard in de echte wereld niet
zo is). De gemiddelde tijd van voorgaande systeem om te bruteforcen is
echter vijf jaar, de helft.

Stel dat je een sleutel hebt die bestaat uit 8 karakters. Een karakter
is een letter van a tot z (geen onderscheid tussen hoofd- en kleine
letters en geen getallen of speciale tekens). Er zijn \(26^8\) mogelijke
sleutels (we gaan ervan uit dat de sleutel exact 8 karakters moet
bevatten). Een computer kan één miljoen sleutels per seconde testen. De
duur om alle mogelijke sleutels te testen is dus
\(\frac{(26^8)}{1 000 000}\), oftewel 208 827 seconden, pakweg 58 uur.
Intuïtief zou je kunnen denken dat je dus meer dan twee dagen ``veilig''
zit, wat niet zo is. De kans dat de eerste sleutel die je test reeds de
juiste is, is even groot als dat het de laatste sleutel is. Kortom,
gemiddeld gezien zal de sleutel in de helft van de maximum tijd gevonden
worden, oftewel 29 uur.

Zouden we de lengte van de sleutel met één karakter verhogen, naar 9,
dan stijgt de totale duur naar pakweg 1500 uur, oftewel 62 dagen. Eén
karakter extra heeft dus wel degelijk een gigantische impact op de
veiligheid van een sleutel!

\section{De eerste algoritmes}\label{de-eerste-algoritmes}

Het doel van ieder encryptie-algoritme is om data zodanig te
versleutelen zodat enkel eigenaars van de gebruikte sleutel de originele
tekst kunnen terugvinden. We vertelden net dat encryptie-algoritmes
kunnen onderverdeeld worden volgens de actie die ze uitvoeren:
substitutie, transpositie of een combinatie. We tonen van iedere variant
nu een historisch voorbeeld.

\begin{tipblock}
Volgende tool, speciaal gemaakt om crypto te leren, is een erg handig
iets om de verschillende cryptografische systemen te visualiseren én
testen: \href{HTTPS://www.cryptool.org/en/ct2/}{link}
\end{tipblock}

\subsection{Substitutie: Caesar
encryptie}\label{substitutie-caesar-encryptie}

De Caesar encryptie (naar Julius Caesar) bestaat uit een eenvoudig
substitutie algoritme. De sleutel is een getal tussen 1 en 25 en geeft
aan door welk element uit het alfabet een teken wordt aangepast, als
volgt:

\begin{itemize}
\tightlist
\item
  Ieder element wordt voorgesteld als een cijfer. A krijgt de waarde 1,
  B wordt 2,\ldots{} Z wordt 26.
\item
  Als de sleutel het getal 3 is, dan zal nu iedere letter A in de tekst
  vervangen worden door het teken 1+3, dus D. Iedere B wordt een E,
  enzovoort.
\item
  Indien er een ``overflow'' is achteraan komen we uiteraard terug naar
  voor in het alfabet. Iedere Z wordt dus een C, iedere Y een B,
  enzovoort.
\end{itemize}

Het Caesarcipher wordt ook wel kortweg \emph{Rot} genoemd, naar het
woord \emph{rotatie}. Een cijfer erachter geeft dan aan welk de te
gebruiken sleutel is. Rot4 wil dus zeggen dat alle elementen vier
plaatsen opgeschoven moeten worden. Merk op dat Rot13 (ook wel
\emph{Caesaralfabet} genoemd) een speciale sleutel is. Als je namelijk
twee maal na elkaar Rot13 toepast op een tekst (eerst op de plaintext,
dan op de resulterende ciphertext) dan verkrijgt men terug de originele
tekst.

Uiteraard kan iedere weldenkende mens in de 21e eeuw een Caesar
encryptie bruteforcen. Het aantal mogelijke sleutels beperkt zich tot 25
mogelijkheden (sleutels 0, 26, etc. zullen resulteren in géén encryptie:
je plaintext en ciphertext zullen identiek zijn) en je kan dit dus snel
testen.

Door \textbf{frequentieanalyse} op de ciphertext toe te passen kan men
ook de plaintext terugvinden zonder te moeten bruteforcen. Indien de
plaintext een tekst in, bijvoorbeeld, het Nederlands is, dan kunnen we
gebruik maken van de statistische eigenschappen van een taal. Zo weten
we dat bepaalde letters in een standaard Nederlandstalige tekst meer of
minder vaak voorkomen. De letter \texttt{e} komt bijvoorbeeld veel vaker
voor dan de \texttt{v}. Daar iedere letter in de encryptie door een
andere wordt vervangen, is het dus voldoende om te ontdekken (a.d.h.v.
frequentieanalyse) welke letter(s) het meest of minst voorkomen om je zo
een vermoeden te geven van de originele letter.

\begin{figure}
\centering
\pandocbounded{\includegraphics[keepaspectratio,alt={De frequentie-analyse van een typisch Nederlandstalige tekst. Merk op dat de letterscore bij het bordspel Scrabble omgekeerd evenredig is met de frequentie dat de letter gemiddeld voorkomt. Dit verklaart ook waarom er per taal een eigen Scrabble-editie bestaat met eigen letterscores (Bron: Wikipedia).}]{crypto/freqNL.jpg}}
\caption{De frequentie-analyse van een typisch Nederlandstalige tekst.
Merk op dat de letterscore bij het bordspel Scrabble omgekeerd evenredig
is met de frequentie dat de letter gemiddeld voorkomt. Dit verklaart ook
waarom er per taal een eigen Scrabble-editie bestaat met eigen
letterscores (Bron: Wikipedia).}
\end{figure}

Dit verklaart ook waarom je best je te encrypteren berichten zo kort
mogelijk houdt. Hoe minder tekens, hoe minder frequentieanalyse zal
werken. Een andere veelvoorkomende fout (in klassiekere encryptie) was
dat de verzender bijvoorbeeld voorspelbare tekst ging encrypteren. Als
je weet dat de verzender altijd begint met ``Geachte'' in z'n berichten,
dan is de kans groot dat de eerste 7 tekens in de ciphertext deze
plaintext voorstellen.

Het principe van Caesar encryptie, de substitutie, blijft echter
overeind staan en zal je nog zien terugkomen in de komende algoritmes.

\subsubsection{Over de modulo operator}\label{over-de-modulo-operator}

De modulo operator (\%) is nuttig bij substitutie-algoritmes zoals bij
Caesar encryptie. De modulo operator geeft de rest weer wanneer we de
linkse door de rechtse operator zouden delen. 19\%5 geeft dus 4 als
resultaat.

Je kan de operator gebruiken om snel te weten wat de waarde van een
teken wordt bij Caesar-encryptie als volgt:

\texttt{(teken\ +\ sleutel)\ \%\ alfabetLengte\ =\textgreater{}\ nieuw\ teken}

De alfabetLengte is 26 bij Caesar-encryptie, namelijk alle letters van A
tot en met Z.

Als je dus een sleutel hebt met waarde 7 en je wilt weten wat de waarde
van \texttt{Y} (element 24, daar we vanaf 0 tellen) wordt dan schrijf
je:

\texttt{(24\ +\ 7)\ \%\ 26\ =\textgreater{}\ 5}

Dit zal dus 5 worden, oftewel een \texttt{F}.

\subsection{Transpositie: scytale
encryptie}\label{transpositie-scytale-encryptie}

Bij transpositie-algoritmen gaan we de positie van de karakters
veranderen. De sleutel kan hierbij bepalen op welke manier dit moet
gebeuren. De Oude Grieken gebruikten een zogenaamde scytale om aan
transpositie-encryptie te doen. Een scytale was een lange stok bestaande
uit drie of meerdere lange zijden. De boodschap werd op een lang lint
geschreven en dit lint werd dan over de scytale gedraaid. De sleutel gaf
aan uit hoeveel vlakken de te gebruiken scytale moest bestaan. Ieder
volgend karakter van de plaintext (op het lint) kwam op een andere zijde
te liggen. Vervolgens werden alle letters op één zijde achter elkaar
gezet, en dit werd herhaald voor iedere zijde: dit werd de ciphertext
die werd doorgestuurd.

Om nu de ciphertext te decrypteren werd een onbeschreven lint over de
juiste scytale gelegd. Vervolgens werd de verkregen ciphertext op dit
lint, zijde per zijde, overgeschreven. Als de ontvanger dan het lint
ontrolde kreeg hij terug de originele tekst te zien.

\begin{figure}
\centering
\includegraphics[width=0.5\linewidth,height=\textheight,keepaspectratio,alt={Een authentieke scytale (Bron: Wikipedia).}]{crypto/skytale.png}
\caption{Een authentieke scytale (Bron: Wikipedia).}
\end{figure}

Stel dat we de tekst ``De perzen komen er nu aan'' (bron: wikipedia)
over een scytale met vier zijden wikkelen, dan krijgen we:

\begin{figure}
\centering
\includegraphics[width=0.4\linewidth,height=\textheight,keepaspectratio,alt={Tijd om de 300 in te roepen!}]{crypto/skypre.png}
\caption{Tijd om de 300 in te roepen!}
\end{figure}

De ciphertext die we vervolgens versturen (wanneer we het lint
afwikkelen) wordt:

\begin{figure}
\centering
\includegraphics[width=0.7\linewidth,height=\textheight,keepaspectratio,alt={Dit gaat Xerxes nooit kunnen lezen\ldots Zeker niet omdat de Nederlandse taal toen nog niet bestond. Muhahahaha!}]{crypto/skyna.png}
\caption{Dit gaat Xerxes nooit kunnen lezen\ldots Zeker niet omdat de
Nederlandse taal toen nog niet bestond. Muhahahaha!}
\end{figure}

Uiteraard zijn er tal van varianten mogelijk om transpositie te doen.
Eerst kan je beslissen om je plaintext in een bepaalde vorm te plaatsen:
bijvoorbeeld in tien kolommen. Vervolgens kan je dan, gebaseerd op de
sleutel, beslissen in welke volgorde je de kolommen achter elkaar
plaatst om de originele tekst te krijgen. Dit is een zogenaamd
\emph{route cipher} wat onder andere werd gebruikt tijdens de
Amerikaanse Burgeroorlog.

\subsection{Combinatie}\label{combinatie}

Het spreekt voor zich dat een combinatie van een transpositiecipher en
een substitutiecipher je encryptie nog versterkt. Veel moderne
algoritmen kunnen nog steeds herleid worden tot een sequentie van
meerdere basisvormen na elkaar.

De \textbf{Advanced Encryption Standard (AES)} is in de 21e eeuw zo'n
beetje de de facto standaard als het aankomt op symmetrische encryptie
(d.w.z. encryptie waar maar één sleutel voor nodig is, verder meer
hierover). Als we echter eens het algoritme opengooien en een enkele
\emph{encryption round} bekijken (AES bestaat uit een sequentie van deze
rondes) dan zien we dat de bits die bovenaan binnenkomen (\emph{state})
vervolgens een combinatie van substituties (\emph{sub}) en transposities
(\emph{mixcolumns} en \emph{shiftrows}) ondergaan.

We gaan AES nog terug zien opduiken wanneer we gaan bekijken hoe
draadloze netwerken worden beveiligd. Als Belg mogen we trouwens erg
fier zijn op deze wereldwijd gebruikte Amerikaanse standaard. Je zal
later ontdekken waarom dat zo is!

\begin{tipblock}
Doel van dit hoofdstuk is ook aantonen dat je geen wiskundig
wondertalent moet zijn om de basisconcepten van cryptografie te
begrijpen. Hier en daar neem ik wat vrijheden om bepaalde stappen te
vereenvoudigen, maar de essentie van de algoritmen blijft wel bewaard en
daarmee, hopelijk, ook de eenvoud (en dus elegantie) ervan.
\end{tipblock}

\section{Cryptanalyse}\label{cryptanalyse}

De term cryptanalyse is nu al enkele keren gevallen: de wereld van de
cryptologie bestaat uit twee delen, die elkaars tegengestelden zijn:

\begin{enumerate}
\def\labelenumi{\arabic{enumi}.}
\tightlist
\item
  \textbf{Cryptografie}: de wetenschap van het versleutelen van
  informatie.
\item
  \textbf{Cryptanalyse}: de wetenschap van het ontcijferen van
  versleutelde informatie, zonder kennis van de gebruikte sleutel.
\end{enumerate}

We gaan in dit boek niet te veel tijd aan de wondere wereld van
cryptanalyse spenderen, daar dit ons te ver zou brengen. We vatten
echter even de belangrijkste concepten hier samen.

\subsection{Sleutellengtes en
bruteforcen}\label{sleutellengtes-en-bruteforcen}

De term \emph{bruteforce} dekt de lading goed. Letterlijk vertaald wordt
het: met brute kracht forceren. Kortom, je gebruikt het wanneer je niet
weet wat doen tijdens de cryptanalyse en gewoonweg de minst efficiënte
manier mogelijk zal toepassen, maar waarvan wel geweten is dat ze altijd
zal werken. Namelijk iedere mogelijke sleutel testen die het cipher
toelaat.

Zoals je je kan inbeelden is de sleutellengte evenredig met de tijd die
cryptanalysten nodig hebben om je sleutel te bruteforcen. De maximale
tijd die nodig is alle sleutels van een bepaalde lengte te berekenen kan
je als volgt vinden:

\(MaximaleTijd = \frac{AantalMogelijkeTekens^{SleutelLengte}}{pogingen/seconde}\)

\subsubsection{Een bruteforce
getalvoorbeeld}\label{een-bruteforce-getalvoorbeeld}

Een sleutel (of wachtwoord) bestaat uit 8 tekens, enkel kleine letters
van a tot en met z zijn toegestaan. De berekeningen worden op een
GeForce GTX 1080 gedaan die ongeveer 30 miljoen pogingen per seconde kan
doen. We krijgen dan:

\(MaximaleTijd = \frac{26^{8}}{30000000}\)

Oftewel ongeveer 6960 seconden, wat neerkomt op ongeveer 1,9 uur tijd
benodigd om alle mogelijke sleutels te testen (herinner je eraan dat
deze tijd gehalveerd moet worden om te weten hoe lang het gemiddeld zal
duren om de juiste sleutel terug te vinden.

Volgende tabel
(\href{HTTPS://uwnthesis.wordpress.com/2020/07/01/brute-force-password-how-long-will-it-take-to-brute-force-a-password/}{bron})
toont nog voorbeelden waarbij telkens dezelfde GeForce 1080 GTX kaart
werd gebruikt. Het getal tussen haakjes geeft aan hoeveel mogelijke
tekens er in dit type mogelijk zijn:

{\def\LTcaptype{none} % do not increment counter
\begin{longtable}[]{@{}
  >{\raggedright\arraybackslash}p{(\linewidth - 8\tabcolsep) * \real{0.2000}}
  >{\raggedright\arraybackslash}p{(\linewidth - 8\tabcolsep) * \real{0.2000}}
  >{\raggedright\arraybackslash}p{(\linewidth - 8\tabcolsep) * \real{0.2000}}
  >{\raggedright\arraybackslash}p{(\linewidth - 8\tabcolsep) * \real{0.2000}}
  >{\raggedright\arraybackslash}p{(\linewidth - 8\tabcolsep) * \real{0.2000}}@{}}
\toprule\noalign{}
\begin{minipage}[b]{\linewidth}\raggedright
\# tekens
\end{minipage} & \begin{minipage}[b]{\linewidth}\raggedright
enkel nummers (10)
\end{minipage} & \begin{minipage}[b]{\linewidth}\raggedright
kleine letters (26)
\end{minipage} & \begin{minipage}[b]{\linewidth}\raggedright
grote \& kleine letters \& nummers (62)
\end{minipage} & \begin{minipage}[b]{\linewidth}\raggedright
eender welk teken (95)
\end{minipage} \\
\midrule\noalign{}
\endhead
\bottomrule\noalign{}
\endlastfoot
4 & 0,3 ms & 15 ms & 490 ms & 2,7 s \\
5 & 3 ms & 400 ms & 31 s & 4,3 min \\
6 & 33 ms & 10 s & 32 min & 6,8 uur \\
7 & 330 ms & 4,5 min & 33 uur & 27 dagen \\
8 & 3,3 s & 1,9 uur & 84 dagen & 7 jaren \\
9 & 33 s & 2,1 dagen & 14 jaren & 670 jaren \\
10 & 5,6 min & 54 dagen & 890 jaren & \(6,3*10^{4}\) jaren \\
11 & 56 min & 3,9 jaren & \(5,5*10^{4}\) jaren & \(6*10^{6}\) jaren \\
12 & 9,3 u & 100 jaren & \(3,4*10^{6}\) jaren & \(5,7*10^{8}\) jaren \\
13 & 3,9 dagen & \(2,6*10^{3}\) jaren & \(2,1*10^{8}\) jaren &
\(5,4*10^{10}\) jaren \\
14 & 39 dagen & \(6,8*10^{4}\) jaren & \(1,3*10^{10}\) jaren &
\(5,1*10^{12}\) jaren \\
15 & 1,1 jaar & \(1,8*10^{6}\) jaren & \(8,1*10^{11}\) jaren &
\(4,9*10^{14}\) jaren \\
16 & 11 jaar & \(4,6*10^{7}\) jaren & \(5*10^{13}\) jaren &
\(4,7*10^{16}\) jaren \\
\end{longtable}
}

\begin{tipblock}
Per verdubbeling van het aantal GeForce-kaarten halveert de tijd.
\end{tipblock}

\begin{noteblock}
Om bovenstaande gigantische getallen wat te duiden: de leeftijd van ons
universum wordt op 13,8 miljard jaar geschat, oftewel \(1,38*10^{10}\)
jaren. Onze mooie blauwe planeet is ongeveer 4,5 miljard jaar oud. De
Tyrannosaurus Rex liep ongeveer 70 miljoen jaar geleden rond, oftewel
\(7*10^{7}\) jaren geleden.
\end{noteblock}

\subsubsection{Dictionary attack}\label{dictionary-attack}

Wanneer de cryptanalist vermoedt dat de te zoeken sleutel iets anders is
dan volledig willekeurige tekens dan kan hij de bruteforce aanval
verbeteren (denk aan ``administrator2022''). In plaats van alle
mogelijke combinaties (permutaties) van de sleutel te testen, zal hij
een woordenboek (\textbf{dictionary}) gebruiken met daarin alle
mogelijke sleutels en woorden die mogelijk de originele sleutel
bevatten.

Tools zoals John The Ripper kan je \emph{voeden} met een dergelijk
woordenboek en dan vragen om sleutels te testen die gebaseerd zijn op
zaken uit dat woordenboek, inclusief bijvoorbeeld door er tekens voor en
na te zetten. Als in het woordenboek het woord \emph{god} staat, dan kan
Jaohn The Ripper bijvoorbeeld ook alle sleutels testen zoals
\emph{god1}, \emph{god2}, etc.

Er zijn tal van woordenboeken online te downloaden die gevuld zijn met
de meest gebruikte wachtwoorden die cryptanalisten (en dus ook de
digitale stropers) kunnen gebruiken om de sleutel sneller te vinden. Het
is aangeraden om zeker geen wachtwoorden (of permutaties ervan) te
gebruiken die in volgende lijsten voorkomen:
\url{https://github.com/danielmiessler/SecLists/tree/master/Passwords/Common-Credentials}

\begin{tipblock}
Dit waren in 2020 de 10 meest gebruikte wachtwoorden:

123456, password, 12345678, qwerty, 123456789, 12345, 1234, 111111,
1234567, dragon

Dit soort lijsten worden opgesteld door gekende datalekken te analyseren
op welke wachtwoorden er in voorkomen.
\end{tipblock}

\subsection{Soorten cryptanalytische
aanvallen}\label{soorten-cryptanalytische-aanvallen}

Geregeld zullen we in dit boek bepaalde zwakheden beschrijven die in
algoritmes misbruikt kunnen worden door een bepaald type
cryptanalytische aanval. Deze aanvallen zijn afhankelijk van de
informatie die de cryptanalist bezit:

\begin{itemize}
\tightlist
\item
  Enkel de ciphertext: vanuit het standpunt van de cyberboswachters is
  dit het beste soort informatie dat de aanvaller bezit. Hij heeft enkel
  een hoop geëncrypteerde informatie en moet proberen daar de originele
  plaintext uit te krijgen. Vanuit het standpunt van de cryptanalist is
  dit dus de minst goede situatie om vanuit te starten.
\item
  Gekende plaintext: de cryptanalist heeft één of meerdere stukken
  informatie waarvan zowel de ciphertext als de bijhorende plaintext
  gekend is.
\item
  Gekozen plaintext: de cryptanalist kan zelf plaintext kiezen waarvan
  de bijhorende ciphertext moet gemaakt worden. Dit zorgt ervoor dat de
  cryptanalist als het ware kan experimenteren.
\item
  Gekozen ciphertext: het zelfde concept als \emph{gekozen plaintext}
  maar deze keer kiest de cryptanalist de ciphertext waarvan hij de
  bijhorende plaintext wil genereren.
\end{itemize}

Er zijn nog enkele meer gespecialiseerde types, maar voor deze cursus
zullen we het bij deze vier basistypes houden.

\begin{noteblock}
Er wordt in deze sectie soms over aanvaller gesproken, alsof de
cryptanalist automatisch van kwade wil is. De wetenschap van de
cryptanalyse is dat uiteraard verre van: enerzijds zorgt het ervoor dat
bestaande en nieuwe cryptografische algoritmes op hun sterkte kunnen
getest worden. Anderzijds helpen ze in tijden van oorlog om boodschappen
van vijanden te onderscheppen en proberen lezen.
\end{noteblock}

\subsection{En wat met
quantum-computers?}\label{en-wat-met-quantum-computers}

Al jaren houdt de crypto-wereld angstvallig de ontwikkelingen in de
quantum-computer wereld in het oog. Alhoewel we nog maar in de babyfase
van quantum-computers zijn, is het toch best mogelijk dat binnen
afzienbare tijd (20, 30 jaar?) we effectief zodanig sterke
quantum-computers zullen hebben die alle bestaande cryptografische
systemen in een handomdraai kunnen ``kraken''.

Hoe dit zal gebeuren snapt de auteur ook (nog) niet en zal dus niet
verder uitgewerkt worden in dit handboek. Besef gewoon dat
quantum-computers van de toekomst potentiële bruteforce aanvallen
drastisch zullen kunnen versnellen.

Het is om deze reden dat er nu reeds onderzoek wordt gedaan naar
cryptografische ciphers die bestand zullen zijn tegen de computers van
de toekomst. Dit soort ciphers worden \emph{post-quantum cryptografische
ciphers} genoemd en zullen niet in dit boek besproken worden.

\begin{noteblock}
Trouwens, ook andere systemen die gebruik maken van cryptografische
concepten zullen in één klap hun nut verliezen. Of zoals
\href{https://www.uclftr.com/post/what-does-the-rise-of-quantum-computers-mean-for-encryption-technology}{dit
artikel} zegt \emph{``And {[}as {]}encryption is everywhere in modern
day life, from e-commerce, to online payments, to passwords, everything
will be vulnerable!''}

Denk daarbij bijvoorbeeld aan \emph{cryptocurrencies} zoals Ethereum en
Bitcoin:

\emph{Cybersecurity specialist Itan Barmes led the vulnerability study
of the Bitcoin blockchain. He found the level of exposure that a large
enough quantum computer would have on the Bitcoin blockchain presents a
systemic risk. ``If {[}4 million{]} coins are eventually stolen in this
way, then trust in the system will be lost and the value of Bitcoin will
probably go to zero,'' he says.}
\href{https://www.investmentmonitor.ai/tech/quantum-computing-bitcoins-crypto-encryption}{Bron}
\end{noteblock}

\section{Symmetrische encryptie}\label{symmetrische-encryptie}

Er zijn twee soorten encryptiesystemen als we kijken naar het aantal
sleutels. Symmetrische systemen zijn systemen waarbij maar één sleutel
nodig is: zowel ontvanger als verzender gebruiken dezelfde sleutel. De
term symmetrisch verwijst naar het feit dat het algoritme exact
hetzelfde doet aan beide zijden. Het enige verschil is dat bij de
verzender de plaintext in het systeem wordt gestoken, wat resulteert in
een ciphertext. Terwijl de ontvanger de ciphertext in het systeem
plaatst om een plaintext te krijgen.

\begin{figure}
\centering
\pandocbounded{\includegraphics[keepaspectratio,alt={Het basismodel van symmetrische encryptie.}]{crypto/basicencrypt.png}}
\caption{Het basismodel van symmetrische encryptie.}
\end{figure}

\begin{itemize}
\item
  De symmetrische encryptiesystemen zijn de oudste vorm: alle klassieke
  algoritmes waren van dit principe. Asymmetrische systemen zijn pas in
  de 20e eeuw ontwikkeld (circa 1970).
\item
  Voorbeelden van bestaande symmetrische encryptiesystemen zijn: AES,
  DES, IDEA, RC4, Blowfish, etc.
\end{itemize}

Dit type encryptie is nog steeds het meest gebruikte en wordt overal
gebruikt waar data op een veilige manier (confidentiality) moet bewaard,
verstuurd of verwerkt worden.

\subsection{Sleuteloverdracht}\label{sleuteloverdracht}

De moeilijkheid bij symmetrische systemen is de sleuteloverdracht. Daar
ontvanger en verzender dezelfde sleutel hanteren is het natuurlijk
belangrijk dat deze de sleutel op een veilige manier kunnen uitwisselen.
Dit probleem wordt niet opgelost door symmetrische cryptosystemen.
Afhankelijk van de context kan deze uitwisseling op verschillende
manieren gebeuren:

\begin{itemize}
\tightlist
\item
  Via een asymmetrisch encryptiesysteem dat wél sleutels op een veilige
  manier kan uitwisselen (zie verder).
\item
  Via een ander beveiligd kanaal, in eender welke vorm (bijvoorbeeld
  fysiek de sleutel aan de andere persoon geven of zeggen, deze opsturen
  via een reeds opgezet symmetrisch encryptiekanaal, etc.).
\end{itemize}

\subsection{Block- en streamciphers}\label{block--en-streamciphers}

Er zijn twee soorten symmetrische encryptieciphers als we kijken naar de
manier waarop ze de te encrypteren data verwerken:

\begin{itemize}
\tightlist
\item
  \textbf{Streamciphers}: hierbij wordt de data letterlijk als een
  stroom (\emph{stream}) van tekens beschouwd. Waarbij teken per teken
  individueel geëncrypteerd wordt (het bekendste voorbeeld is RC4). Voor
  ieder teken dat verwerkt wordt zal er exact één geëncrypteerd teken
  gegenereerd worden. Dit soort algoritmes zijn over het algemeen
  sneller dan blockciphers.
\item
  \textbf{Blockciphers}: de data wordt in blokken (van bijvoorbeeld 128
  tekens) verwerkt. Bekendste voorbeelden die we verderop behandelen
  zijn AES, DES, 3DES, etc.
\end{itemize}

\subsection{Streamciphers}\label{streamciphers}

De werking van een symmetrisch streamcipher is verrassend eenvoudig en
bestaat uit twee delen:

\begin{itemize}
\tightlist
\item
  Een \textbf{pseudorandom keystream generator}: deze zal de sleutel als
  het ware expanderen naar een sleutel met de zelfde lengte als de
  stream, genaamd een \textbf{keystream}. Als je 400 bytes aan data
  wenst te encrypteren, zal je een keystream van 400 bytes moeten
  genereren. Daar we met een stream werken zal deze generator teken per
  teken genereren. Hoe dit gebeurt, leggen we verderop uit.
\item
  De \textbf{XOR} of ``exclusieve of'' functie: deze zal de plaintext
  naar een ciphertext omzetten door de plaintext met de keystream samen
  te voegen. Deze stap is de feitelijke encryptie!
\end{itemize}

\begin{figure}
\centering
\pandocbounded{\includegraphics[keepaspectratio,alt={Het streamcipher proces.}]{crypto/stream.png}}
\caption{Het streamcipher proces.}
\end{figure}

Aan de ontvangerzijde gebeurt exact hetzelfde. \textbf{Enkel indien de
ontvanger dezelfde sleutel gebruikt, zal deze dezelfde keystream kunnen
genereren, en bijgevolg enkel dan de originele plaintext verkrijgen.}

Het hart van een symmetrisch streamcipher is dus enerzijds de
XOR-functie én, belangrijker, de manier waarop de keystream wordt
gemaakt.

\subsubsection{De XOR functie}\label{de-xor-functie}

De waarheidstabel van de XOR-functie is de volgende:

{\def\LTcaptype{none} % do not increment counter
\begin{longtable}[]{@{}lll@{}}
\toprule\noalign{}
Plaintext input & Keystream input & Ciphertext output \\
\midrule\noalign{}
\endhead
\bottomrule\noalign{}
\endlastfoot
1 & 0 & 1 \\
0 & 1 & 1 \\
0 & 0 & 0 \\
1 & 1 & 0 \\
\end{longtable}
}

De XOR-functie wordt in schema's aangeduid door een cirkel met een
plusje in: \(\oplus\)

Beeld je in dat we het bericht \texttt{1010} willen versleutelen, en we
hebben een gegenereerde keystream \texttt{1101}. Als we deze
\emph{XOR'n} dan geeft dit \texttt{0111}. Dit is dus de ciphertext. Als
de ontvanger dezelfde keystream kan genereren en deze XOR'd met de
verkregen ciphertext, dan krijgt deze terug de originele plaintext.

\subsubsection{De keystream generator}\label{de-keystream-generator}

De keystream generator heeft dus als doel om voor ieder karakter dat
moet geëncrypteerd worden een bijhorend keystream karakter te maken.
Deze karaktergeneratie moet onvoorspelbaar zijn (\emph{random})
tegenover de sleutel die wordt gebruikt en het voorgaande karakter dat
werd gemaakt. Echter, dit moet wel PSEUDO (\emph{schijn})-willekeurig
zijn: dezelfde sleutel als beginpunt (\textbf{seed}) moet steeds
dezelfde reeks genereren.

De kracht (en zwakte) van een symmetrisch streamcipher ligt in de
implementatie van de manier waarom deze keystream generator werkt.
Mogelijke zwakheden kunnen bijvoorbeeld zijn dat de gegenereerde stroom
informatie van de sleutel ``lekt'' naar de keystream (wat desastreuze
gevolgen bleek te hebben bij de originele wifi-security (WEP), waarover
later meer) of een voorspelbare ``randomiteit'' van de keystream.

Om aan encryptie te kunnen doen, hebben we systemen nodig die
onvoorspelbaar zijn. Als de aanvaller kan voorspellen wat de uitvoer van
een onderdeel van de encryptie zal zijn, dan kunnen we geen
confidentiality en integrity voorzien. Kortom, we hebben algoritmes
nodig die willekeurige getallen kunnen generen die 100\% onvoorspelbaar
zijn. Net zoals het werpen van een dobbelsteen niet voorspeld kan
worden, zo ook moeten onze algoritmes een (digitale) dobbelsteen hebben.

Digitale systemen die perfect willekeurige getallen genereren noemt men
\textbf{random number generators} (RNG). Uiteraard moet een RNG
geprogrammeerd kunnen worden: dat behelst dus een algoritme. Een
algoritme is per definitie ``voorspelbaar''. Alles hangt daarom af van
de invoer die het algoritme gebruikt om random getallen te beginnen
genereren. We spreken dan van een \textbf{pseudorandom number generator}
(PRNG), pseudo (\textbf{schijnbaar}) omdat de uitvoer afhankelijk is van
het startgetal, de zogenaamde \textbf{seed}. Die seed kan bijvoorbeeld
de encryptiesleutel zijn: enkel met dié sleutel zal het algoritme
dezelfde reeks getallen generen. Er zijn echter ook systemen die
bijvoorbeeld de huidige tijd of de staat van een flipflop als startpunt
gebruiken (wanneer je een flipflop aanzet kan je niet voorspellen of
deze op 1 of 0 zal staan, daar deze staat beïnvloed wordt door de
elektromagnetische straling). Uiteraard is een dergelijke seed voor een
keystream generator nutteloos, daar zowel verzender én ontvanger
dezelfde reeks getallen moeten kunnen genereren.

\subsubsection{RC4 tot op het bot}\label{rc4-tot-op-het-bot}

Laten we eens één van de meest gebruikte streamciphers bekijken, het RC4
cipher. Dit algoritme, ontwikkeld door Ron Rivest (de afkorting staat
trouwens voor \emph{Rons Cipher 4}), wordt gebruikt onder andere om een
beveiligde SSL-tunnel (zie later) op te zetten en zit in het hart van
veel geëncrypteerde communicatiekanalen.

RC4 werkt zoals we eerder verklaarden hoe een streamcipher werkt: het
heeft een keystream generator en zal de keystream vervolgens XOR'n met
de plaintext. Eerst zal de ingevoerde sleutel (die 40 tot 2048 bits lang
mag zijn) omgezet worden naar een compatibele werksleutel met behulp van
een \textbf{Key scheduling algorithm} (KSA). Deze werksleutel zal dan
als seed gebruikt worden om een keystream in het \textbf{Pseudo-random
generator algorithm} (PRGA) te maken.

\begin{figure}
\centering
\pandocbounded{\includegraphics[keepaspectratio,alt={RC4 tot op het bot.}]{crypto/rc4.png}}
\caption{RC4 tot op het bot.}
\end{figure}

\paragraph{KSA}\label{ksa}

De \textbf{KSA} heeft als doel om de ingevoerde 40 tot 2k-bit sleutel om
te zetten naar een compatibele sleutel voor de PRGA. Het doet dit
volgens een eenvoudig algoritme:

\textbf{Stap 1}: plaats de ingevoerde sleutel in een array (\textbf{T})
van 256 karakters. Herhaal de sleutel indien nodig.

\textbf{Stap 2}: maak een sleutelarray (\textbf{S}) aan, ook van 256
karakters, en plaats er de waarden 0 tot en met 255 in.

\textbf{Stap 3}: permuteer (\emph{transpositie}) de elementen in de
array \texttt{S} met behulp van de array \texttt{T} als volgt:

\begin{Shaded}
\begin{Highlighting}[]
\NormalTok{j = 0}
\NormalTok{for i = 0 tot 255}
\NormalTok{\{ }
\NormalTok{  j = (j + S[i] + T[i]) \% 256}
\NormalTok{  Verwissel(S[i],S[j])}
\NormalTok{\}}
\end{Highlighting}
\end{Shaded}

Deze stap zal dus de elementen in de sleutelarray \texttt{S} naar nieuwe
posities in diezelfde array plaatsen en deze ook de hele tijd van plek
wisselen. De index \texttt{j} is afhankelijk van de originele sleutel
uit stap 1 en zorgt er dus voor dat iedere finale sleutel een unieke
array \texttt{S} zal opleveren.

\textbf{Uitgewerkt voorbeeld van KSA}

Stel dat onze sleutel ``ab'' is. De decimale ASCII-waarden van a en b
zijn 97 en 98 respectievelijk. Onze array T zal dus bestaan uit 128 keer
de waarden 97 en 98 na elkaar aan de start:

{\def\LTcaptype{none} % do not increment counter
\begin{longtable}[]{@{}ll@{}}
\toprule\noalign{}
Index & Waarde \\
\midrule\noalign{}
\endhead
\bottomrule\noalign{}
\endlastfoot
T{[}0{]} & 97 \\
T{[}1{]} & 98 \\
T{[}2{]} & 97 \\
etc. & \\
\end{longtable}
}

Stap twee genereert de tabel \texttt{S} die gewoon de waarden 0 tot en
met 255 heeft in de 255 plekjes van de array (de waarde is dus in deze
fase gewoon ook de index van het element)

Als we dan stap drie toepassen dan krijgen we na de eerste iteratie van
de loop (\texttt{i=0}):

\texttt{j\ =\ (0\ +\ S{[}0{]}\ +\ T{[}0{]})\ \%\ 256}

oftewel

\texttt{j=\ (0\ +\ 0\ +\ 97)\ \%\ 256\ \ =\textgreater{}\ j\ wordt\ 97}

De eerste wissel die in \texttt{S} zal plaatsvinden is dan
\texttt{Verwissel(S{[}0{]},S{[}97{]})}

\texttt{S} ziet er dan als volgt uit na de eerste iteratie:

{\def\LTcaptype{none} % do not increment counter
\begin{longtable}[]{@{}ll@{}}
\toprule\noalign{}
Index & Waarde \\
\midrule\noalign{}
\endhead
\bottomrule\noalign{}
\endlastfoot
S{[}0{]} & 97 \\
S{[}1{]} & 1 \\
S{[}2{]} & 2 \\
\ldots{} & \\
S{[}97{]} & 0 \\
etc. & \\
\end{longtable}
}

En dit herhalen we nog 255 keer. \textbf{Finaal hebben we nu in tabel
\texttt{S} een sleutel die bestaat uit de getallen 0 tot en met 255
verdeeld over willekeurige plekken in de array.} Indien we een andere
initiële sleutel zouden hebben gebruikt dan zou deze tabel er totaal
anders uitzien.

\paragraph{PRGA}\label{prga}

Nu we een compatibele sleutel \texttt{S} hebben kan de keystream
generatie van start gaan. Deze bestaat uit een loop die blijft doorgaan
telkens een nieuw plaintext karakter binnenkomt, als volgt:

\begin{Shaded}
\begin{Highlighting}[]
\NormalTok{i = 0}
\NormalTok{j = 0}
\NormalTok{Herhaal telkens keystream karakter nodig is}
\NormalTok{\{}
\NormalTok{  i = (i+1) \% 256}
\NormalTok{  j = (j +S[i]) \% 256}
\NormalTok{  Verwissel(S[i],S[j])}
\NormalTok{  k = S[(S[i]+S[j]) \% 256]}
\NormalTok{  output k naar keystream}
\NormalTok{\}}
\end{Highlighting}
\end{Shaded}

Telkens zal het algoritme één specifieke waarde (tussen 0 en 255) uit de
array teruggeven als keystream karakter \texttt{k}. Zoals je ziet zal de
array \texttt{S} voorts de hele tijd van ``gedaante'' blijven
veranderen. Telkens we een element \texttt{k} uitsturen zal ook de tabel
\texttt{S} weer wat zijn veranderd daar we de waarden van
\texttt{S{[}i{]}} en \texttt{S{[}j{]}} onderling verwisselen.

\paragraph{Uitgewerkt voorbeeld van
PRGA}\label{uitgewerkt-voorbeeld-van-prga}

Als we verder werken met de tabel \texttt{S} van het vorige uitgewerkte
KSA voorbeeld en de waarden ervan gebruiken na één iteratie dan krijgen
we onderstaande berekeningen (\emph{We veronderstellen even dat we de
KSA niet verder hebben uitgevoerd en de tabel \texttt{S} dus dezelfde is
gebleven als op het einde van het voorbeeld, wat in het echt niet zal
zijn.}):

\begin{verbatim}
i = (0+1) % 256  => 1
j = (0+97) % 256 => 97
Verwissel(S[1],S[97]) => Verwissel(1,0)
k = S[ (S[1]+S[97])%256 ] => k = S[ (0 + 1) % 256] 
we outputten de waarde die op S[1] staat
\end{verbatim}

Finaal zal de output, de waarde \texttt{k}, ge-XOR'd worden met het
huidige karakter van de plaintext stream.

\begin{noteblock}
Zo, dat viel nog mee he? Zoals al gezegd, een belangrijke motivatie van
dit boek is aantonen dat je niet bang hoeft te zijn van wat er achter de
schermen van de cyberwereld gebeurt. De hoeveelheid wiskunde die we
bijvoorbeeld nodig hadden, is beperkt gebleven tot onze trouwe modulo
(\%)-operator en meer niet. Wanneer we zo meteen een blockcipher gaan
uitkleden, zal je ook daar ontdekken dat je best in staat bent
schijnbaar complexe technologieën te begrijpen. Hop naar de blockciphers
dus!
\end{noteblock}

\subsection{Blockciphers}\label{blockciphers}

Blockciphers, de naam zegt het al, zullen eerst de plaintext in blokken
karakters opdelen (bv. 128 bits) en vervolgens blok per blok
encrypteren.

\begin{noteblock}
Indien het aantal te encrypteren bytes aan data geen exact veelvoud is
van de blokgrootte dient er \textbf{padding} te gebeuren. Hierbij zal
het laatste blok opgevuld (\emph{gepad}) worden met extra bytes tot het
blok terug de juiste blokgrootte heeft. De manier waarop de padding
gebeurt, is afhankelijk van het cipher dat gehanteerd worden.
\end{noteblock}

\subsubsection{Feistel-structuren}\label{feistel-structuren}

Ook hier zullen we dezelfde soorten operaties (XOR, substituties en
transposities) zien terugkomen. Echter, ook \textbf{Feistel}-structuren
worden hier gebruikt: in deze operatie zal de data steeds in twee
helften worden gesplitst en wordt steeds een specifieke operatie
(aangeduid met \emph{F} van functie in de figuur), zoals een
substitutie, op één helft uitgevoerd dat dan wordt ge-XOR'd met de
andere helft. Dit wordt meerdere keren herhaald, waarbij de linker
(\emph{L}) en rechterzijde (\emph{R}) steeds afwisselend door de
specifieke encryptie-operatie gaan.

\begin{figure}
\centering
\includegraphics[width=0.3\linewidth,height=\textheight,keepaspectratio,alt={Een enkele Feistel-structuur.}]{crypto/Feistelunit.png}
\caption{Een enkele Feistel-structuur.}
\end{figure}

Net zoals bij RC4 zullen we ook vaak met een zogenaamd key scheduling
algoritme werken zodat de sleutel niet constant doorheen het hele proces
dezelfde is en we met \textbf{subkeys} of \emph{round keys} werken.

\subsubsection{DES tot op het bot}\label{des-tot-op-het-bot}

Een van de oudste blockciphers is \textbf{DES} oftewel \textbf{Data
Encryption Standard}. Deze standaard werd in 1977 geboren en vereiste
toen blokken van 64 bits data en gebruikte een 56bit sleutel.
Ondertussen is dit cipher zeer outdated, maar we gaan hem toch bekijken
omdat deze enerzijds erg belangrijk was voor de wereld - grote delen van
het bankwezen beschermden er hun financiële transacties mee (en nadien
met de opvolger 3DES, zie verder) - en de standaard is erg duidelijk om
een goed inzicht in blockciphers te verkrijgen.

Er was wel veel controverse rond deze standaard (gebaseerd op het door
IBM ontwikkelde Lucifer cipher) omdat de originele versie (genaamd
Lucifer) werkte met een dubbel zo grootte sleutel (128 bit) en bijgevolg
dus veiliger was op lange(re) termijn. De Amerikaanse veiligheidsdienst,
NSA, was niet zo happig om een dergelijk goed beveiligd algoritme
commercieel te maken en zorgde er daarom voor dat de sleutellengte
gehalveerd werd.

Data met DES encrypteren (en bijgevolg ook decrypteren daar het een
symmetrisch cipher is) bestaat uit twee onderdelen:

\begin{enumerate}
\def\labelenumi{\arabic{enumi}.}
\tightlist
\item
  \textbf{Versleuteling}: een reeks Feistel-structuren na elkaar
  (\textbf{16 rondes}) die de data blok per blok versleutelen.
\item
  \textbf{Subkeys maken}: een key scheduling algoritme dat 16 subkeys
  genereert (één voor iedere encryptieronde), gebaseerd op de 56 bit
  sleutel.
\end{enumerate}

\paragraph{Versleuteling}\label{versleuteling}

Volgende schema toont de encryptie bestaande uit 16 rondes:

\begin{figure}
\centering
\includegraphics[width=0.8\linewidth,height=\textheight,keepaspectratio,alt={DES encryptie. 16 Feistel-structuren na elkaar.}]{crypto/des1.png}
\caption{DES encryptie. 16 Feistel-structuren na elkaar.}
\end{figure}

De data wordt blok per blok doorheen dit gedeelte gestuurd. Eerst
gebeurt er een \emph{Initiële Permutatie} (\emph{IP} in de figuur)
waarbij iedere bit naar een andere plek wordt gestuurd volgens een vast
patroon. Achteraan gebeurt dit nogmaals in een \emph{Finale Permutatie}
(\emph{FP}).

\begin{figure}
\centering
\includegraphics[width=0.9\linewidth,height=\textheight,keepaspectratio,alt={De Initiële Permutatie.}]{crypto/des2.png}
\caption{De Initiële Permutatie.}
\end{figure}

Na de \emph{IP} gaat de data door 16 Feistel-structuren die telkens
hetzelfde doen. De data wordt in twee helften gesplitst waarbij de
rechterzijde door de F-operatie gaat (die we zo meteen toelichten), het
resultaat hiervan wordt ge-XOR'd met de linkerhelft van de data. Het
resultaat van deze XOR, een 32 bit blok, wordt nu het rechterblok in de
volgende ronde en omgekeerd.

\begin{figure}
\centering
\includegraphics[width=0.6\linewidth,height=\textheight,keepaspectratio,alt={Binnenin de F-operatie.}]{crypto/des3.png}
\caption{Binnenin de F-operatie.}
\end{figure}

In het F-blok wordt eerst het 32-bit blok uitgebreid (\emph{E} in de
figuur, van expansie) naar een 48 bit blok zodat deze even lang is als
de subkey voor deze ronde. De expansie gebeurt, net als de initiële
permutatie, volgens een vast patroon:

\begin{figure}
\centering
\includegraphics[width=0.9\linewidth,height=\textheight,keepaspectratio,alt={De expansie van 32 naar 48 bits.}]{crypto/des4.png}
\caption{De expansie van 32 naar 48 bits.}
\end{figure}

Nu worden deze 48 bits ge-XOR'd met de subkey. Het resultaat wordt in
blokjes van 6 bits door een \emph{S}-blok gestuurd (zogenaamde
\emph{Selection blocks}). In dit blokje wordt 6 bit omgezet naar 4 bit.
In de figuur hieronder zien we bijvoorbeeld hoe de omzetting in blok
\emph{S5} gebeurt. Ieder blokje heeft een soortgelijke tabel, maar met
andere resultaten. De 6 bits bestaan uit de 2 outer bits, namelijk de
eerste en de laatste bit, alsook de 4 innerbits. De figuur toont
bijvoorbeeld dat de output \texttt{1001} zou zijn indien er
\texttt{011011} in het blok wordt geplaatst.

\begin{figure}
\centering
\includegraphics[width=1\linewidth,height=\textheight,keepaspectratio,alt={Waarheidstabel van het S5-blok (Bron wikipedia).}]{crypto/des5.png}
\caption{Waarheidstabel van het S5-blok (Bron wikipedia).}
\end{figure}

Na 16 rondes krijgen we terug een 32 bit datablok dat nog een
\emph{Finale permutatie} ondergaat die weer de bits van plaats verandert
en de output hiervan is een geëncrypteerd blok data dat kan doorgestuurd
worden naar de ontvanger.

\begin{tipblock}
Je kan de DES standaard op
\href{HTTPS://web.archive.org/web/20040410171758/http://www.itl.nist.gov/fipspubs/fip46-2.htm}{web.archive.org/web/20040410171758/http://www.itl.nist.gov/fipspubs/fip46-2.htm}
nalezen en ontdekken dat deze niet zo lang is zoals je zou verwachten
van een wereldwijd gebruikte standaard.
\end{tipblock}

\paragraph{Subkeys maken}\label{subkeys-maken}

Iedere ronde tijdens de versleuteling vereist een sleutel. Om te
voorkomen dat steeds de hoofdsleutel wordt gebruikt (en er zo potentieel
dezelfde keystreams worden gemaakt), wordt deze sleutel doorheen een
\emph{round-key generator algoritme} gestuurd. Dit algoritme bestaat uit
16 rondes waarbij de 56 bit sleutel (8 van de 64 bits in de originele
sleutel zijn zogenaamde pariteits-bits, die dienen om te controleren of
de sleutel geen fouten bevat) steeds in twee helften van 28 bits wordt
\emph{geknipt}.

\begin{figure}
\centering
\includegraphics[width=0.6\linewidth,height=\textheight,keepaspectratio,alt={De 16 rondes die telkens 1 subkey maken.}]{crypto/des6.png}
\caption{De 16 rondes die telkens 1 subkey maken.}
\end{figure}

Iedere ronde wordt de helft van de sleutel 2 bits \emph{geshift} (in
ronde 1,2, 9 en 16 maar 1 bit). Dat wil zeggen dat alle bits twee
plekjes opschuiven en de eerste (of laatste) bits komen dan achteraan
(of vooraan) te staan. Deze twee geshifte helften worden dan enerzijds
doorgestuurd naar de volgende ronde, anderzijds naar een
\emph{compressie P-box} waarvan het resultaat een 48 bits subkey zal
zijn van die ronde.

De naam \emph{compressie P-box} doet al vermoeden wat er gebeurt:

\begin{itemize}
\tightlist
\item
  Compressie: een aantal bits zullen wegvallen (er komen 56 bits in,
  maar we hebben maar 48 bits nodig).
\item
  P-box: een permutatie oftewel transpositie dat alle bits van plek zal
  veranderen.
\end{itemize}

\begin{figure}
\centering
\includegraphics[width=1\linewidth,height=\textheight,keepaspectratio,alt={Er zijn drie types P-Boxes, afhankelijk van wat ze met de data doen.}]{crypto/pboxex.png}
\caption{Er zijn drie types P-Boxes, afhankelijk van wat ze met de data
doen.}
\end{figure}

In de Compression P-Box wordt een aantal bits van de sleutel
``tegengehouden''. Welke bits dat zijn hangt af van de P-Box. Iedere
ronde wordt er een andere Compression P-Box gehanteerd.

\emph{En zo hebben we het einde van de werking van DES bereikt. Dat viel
al bij al nog mee, niet?}

\subsubsection{3DES}\label{des}

Al van bij de start gingen er stemmen op dat de originele sleutellengte
voor DES (56 bits, 48 in effectiviteit vanwege de pariteitsbits)
redelijk snel zou gebruteforced worden. Om die reden werd 3DES in het
leven geroepen in 1995. De oplossing, 3DES, was een mooi staaltje
compromisvorming: het bood een verhoogde beveiliging doordat het een
lange sleutel had (tot 168 bits lang) maar bleef tegelijkertijd
compatibel met de bestaande DES hardware en software.

De werking van 3DES (\emph{tripple DES}) is verrassend eenvoudig: ieder
blok data wordt drie keer doorheen een DES-cipher gestuurd. Hierbij
wordt steeds een andere sleutel gebruikt. Om de bestaande DES hardware
te gebruiken, wordt hierbij de data eerst door de encryptie gestuurd,
dan doorheen de decryptie en terug door de encryptie. Daar we in iedere
fase een andere sleutel gebruiken heeft dit (dankzij de eigenschappen
van symmetrische ciphers) als effect dat we dus effectief drie maal na
elkaar encrypteren met steeds een andere sleutel. Aan de ontvanger zijde
gebeurt dan het omgekeerde: decryptie, encryptie, decryptie én dit dus
allemaal met de bestaande DES hardware!

\begin{figure}
\centering
\includegraphics[width=0.6\linewidth,height=\textheight,keepaspectratio,alt={3DES.}]{crypto/3des.png}
\caption{3DES.}
\end{figure}

\begin{tipblock}
3DES laat dus ook (single) DES encryptie toe. Het enige dat je hiervoor
moet doen is de subsleutels K2 en K3 gelijkstellen waardoor de tweede en
derde fase tijdens de encryptie (en decryptie) eigenlijk niets doen,
daar het gewoon de data encrypteert in ronde twee en dan ogenblikkelijk
in ronde drie terug decrypteert.
\end{tipblock}

\begin{noteblock}
Het bankwezen gebruikt 3DES nog steeds (of varianten die erop gebaseerd)
zijn om financiële transacties van onder andere Visa en Mastercard te
beveiligen.
\end{noteblock}

\subsubsection{Block cipher modes}\label{block-cipher-modes}

Tot hiertoe gingen we data steeds blok per blok in het encryptiecipher
sturen en het resultaat ervan doorsturen. Klaar. Oplettende mensen
hebben hier mogelijk al een hiaat in gezien: wat als twee blokken exact
dezelfde data bevatten? Beiden zullen dezelfde ciphertext als resultaat
genereren, daar we telkens dezelfde sleutel (en dus subkeys) gebruiken.
Twee ciphertexts die identiek zijn, willen we vermijden daar het
potentiële informatie over de plaintext zichtbaar maakt. Voorts laat dit
soort werking ook replay attacks toe: de aanvaller kan een geëncrypteerd
pakket bewaren en op een later moment terug opsturen, zonder dat hij
moet weten wat de plaintext bevat.

Stel dat we een afbeelding van Tux De Pinguïn opsplitsen in ongeveer 100
bij 100 datablokken. Als we nu ieder blok individueel met een
blockcipher encrypten en zouden visualiseren dan krijgen we iets dat
mogelijks toch nog wat informatie van Tux \emph{doorlekt} (zie de tweede
afbeelding) daar blokken van de afbeelding met exact dezelfde informatie
ook dezelfde ciperblock zullen genereren. Vergelijk dit met de derde
afbeelding waarin we een andere modus gebruiken (die we zo meteen gaan
uitleggen) waarin repetities in de plaintext geen invloed hebben op
repetities in de ciphertext.

\begin{tipblock}
Door pakketten te sniffen zou de aanvaller kunnen achterhalen wat de
mogelijke inhoud van een pakket is. Ieder netwerkprotocol volgt de
standaarden die er voor beschreven zijn en zo kan dus de aanvaller heel
veel informatie ontdekken over een ciphertext gewoon ten opzichte van
wanneer een pakket wordt verstuurd tegenover de andere. Als het
bijvoorbeeld het eerste pakket is dat verstuurd wordt, dan is de kans
groot dat dit pakket de typische \emph{``ik wil een verbinding
opzetten''}-request is.
\end{tipblock}

\begin{figure}
\centering
\includegraphics[width=0.75\linewidth,height=\textheight,keepaspectratio,alt={Het volgend voorbeeld toont een (overdreven) manier waarom ECB minder veilig is dan de modes die we nog gaan behandelen (Bron wikipedia).}]{crypto/ecbfail.png}
\caption{Het volgend voorbeeld toont een (overdreven) manier waarom ECB
minder veilig is dan de modes die we nog gaan behandelen (Bron
wikipedia).}
\end{figure}

Voorgaande modus, waarin we ieder blok onafhankelijk van het vorige
encrypteren, noemen we de \textbf{Electronic Codebook (ECB)} modus.
Alhoewel deze modus dus duidelijk een veiligheidsprobleem met zich mee
draagt, heeft deze modus ook één voordeel:

\begin{itemize}
\tightlist
\item
  Ieder blok wordt onafhankelijk van andere blokken gedecrypteerd. Als
  er dus een blok niet gedecrypteerd kon worden door een fout, dan heeft
  dat geen invloed op de daaropvolgende blokken. Dit is dus voor
  streaming-situaties nuttig: beeld je in dat je decryptie faalt
  halverwege het binnenkrijgen van een film die je aan het bekijken
  bent. Je zou helemaal opnieuw moeten beginnen.
\end{itemize}

ECB is een niet zo veilige manier om een blockcipher toe te passen. Veel
interessanter (veiliger) wordt het wanneer we extra informatie gebruiken
om een blok te encrypteren. \textbf{Enkel het huidige blok en dezelfde
sleutel gebruiken is namelijk niét veilig.} Er zijn verschillende modes
om veiliger te encrypteren dan ECB:

\begin{itemize}
\tightlist
\item
  Cipher block chaining (CBC): de output van het vorige blok (de
  ciphertext) wordt mee als input voor de encryptie van het volgende
  blok gebruikt.
\item
  Propagating CBC (PCBC): zelfde als CBC maar bij decryptie van een blok
  zijn ook alle vorige blokken vereist.
\item
  Cipher feedback (CFB): ongeveer hetzelfde als CBC alleen wordt het
  vorige blok iets later in het encryptieproces van het volgende blok
  gebruikt.
\item
  Output feedback (OFB): het blockcipher wordt als een streamcipher
  gebruikt.
\item
  Counter-mode (CTR): een extra teller wordt gebruikt als input, genaamd
  een \emph{Initialisatie vector}, bij de encryptie van een blok. Deze
  teller wordt steeds verhoogd. Eén van de meest gebruikte modes (in
  onder andere WPA2 en IPSEC).
\end{itemize}

\begin{figure}
\centering
\includegraphics[width=0.8\linewidth,height=\textheight,keepaspectratio,alt={CBC encryptie (Bron wikipedia).}]{crypto/cbc.png}
\caption{CBC encryptie (Bron wikipedia).}
\end{figure}

Alle modes uit de doeken doen is hier niet aan de orde maar het moge
duidelijk zijn dat ECB de minst veilige mode voorhanden is en deze wordt
dan ook best vermeden.

\begin{tipblock}
Het concept \textbf{Initialisatie Vector (IV)} zal je veel zien
terugkomen in ciphers. Een IV is een getal dat men als extra seed
meegeeft tijdens de encryptie, naast de sleutel. Op deze manier
voorkomen we dat steeds enkel de sleutel als seed wordt gebruikt en we
dus effectief steeds met een \emph{andere} sleutel werken. Uiteraard zal
ook de andere zijde over dezelfde IV moeten beschikken en zal deze dus
doorgestuurd moeten worden. Dit gebeurt meestal via de header van het
bijhorende pakketje en is ongeëncrypteerd. Dit lijkt contra-intuïtief -
de IV onbeveiligd doorsturen - maar is geen probleem. Uiteraard is het
belangrijk dat er een goed \emph{IV selectie algoritme} wordt gebruikt
dat bepaalt hoe steeds het volgende IV moet worden berekend (bv. steeds
met 1 verhogen, een willekeurig, etc.).
\end{tipblock}

\subsubsection{AES}\label{aes}

Alhoewel 3DES een verbetering op DES was, was er toch nood aan een
nieuwe encryptie-standaard die langere tijd kon bestaan. In 2001 werd
daarom de \textbf{Advanced Encryption Standard (AES)} boven het doopvont
gehouden als de nieuwe de facto encryptiestandaard wereldwijd. Deze
Amerikaanse standaard is gebaseerd op het \textbf{Rijndael}-algoritme
waar we als Belgen fier op mogen zijn: Rijndael is ontwikkeld door twee
Belgische KUL-cryptografen Vincent Rijmen en Joan Daemen.

\begin{figure}
\centering
\includegraphics[width=0.8\linewidth,height=\textheight,keepaspectratio,alt={AES encryptie (Bron wikipedia).}]{crypto/aes.png}
\caption{AES encryptie (Bron wikipedia).}
\end{figure}

AES is een symmetrisch blockcipher dat data in blokken van 128 bits zal
opsplitsen en sleutels tot 256 bits lang toelaat. De volledige werking
van AES gaan we hier niet uit de doeken doen, het voldoet te begrijpen
dat in grote lijnen hetzelfde soort stappen worden doorlopen als DES en
andere symmetrische ciphers:

\begin{itemize}
\tightlist
\item
  Er is een \emph{key expansie} stap om een unieke sleutel pér ronde te
  hebben.
\item
  De data wordt doorheen meerdere rondes gestuurd.
\item
  Iedere ronde gebeuren er zaken zoals substituties en transposities,
  zowel van bytes als van hele rijen of kolommen data.
\item
  Finaal vindt er een XOR-encryptie plaats.
\end{itemize}

Merk op dat ook hier, onderaan, de XOR-functie nog steeds dienst zal
doen als de feitelijke encryptie van de data. Zonder deze XOR-functie
zou al het voorgaande enkel maar resulteren in data die van plek
verandert, volgens een patroon waar de geheime sleutel niet bij van te
pas komt.

\section{Asymmetrische encryptie}\label{asymmetrische-encryptie}

\subsection{Het probleem met symmetrische
encryptie}\label{het-probleem-met-symmetrische-encryptie}

Wat als je bij symmetrische encryptie met meerdere mensen wilt
communiceren zonder dat iedereen elkaars berichten kan zien? Bob kan
onmogelijk dezelfde sleutel gebruiken om met Alfredo te communiceren die
hij al gebruikte met Alice. Kortom, je hebt per \emph{eindpunt} een
aparte sleutel nodig. Het aantal sleutels dat je nodig hebt, zeker als
ook alle gebruikers onderling nog eens willen communiceren wordt snel
erg groot. Je kan het aantal benodigde sleutels berekenen met de formule
\(n * \frac{(n - 1)}{2}\), waarbij \texttt{n} het aantal gebruikers
voorstelt:

\begin{itemize}
\tightlist
\item
  6 gebruikers vereisen 15 sleutels.
\item
  7 gebruikers vereisen 21 sleutels.
\item
  10 gebruikers vereisen er al 45.
\item
  100 gebruikers vereisen er 4950!
\end{itemize}

\begin{figure}
\centering
\includegraphics[width=0.4\linewidth,height=\textheight,keepaspectratio,alt={Bij 6 gebruikers zijn er al 15 sleutels nodig.}]{crypto/keyprob.png}
\caption{Bij 6 gebruikers zijn er al 15 sleutels nodig.}
\end{figure}

Symmetrische encryptie heeft dus een \emph{key distribution problem}
wanneer er encryptie op een grote schaal nodig is. Zeker als we spreken
over online communicatie, over het Internet, wordt de schaal
ogenblikkelijk gigantisch groot en wordt het \emph{sleutelmanagement}
problematisch. Een andere oplossing is dus aan de orde.

\subsection{Publieke cryptografie}\label{publieke-cryptografie}

Publieke crypto oftewel \textbf{asymmetrische encryptie} zal het
probleem met symmetrische encryptie oplossen doordat er gebruikt wordt
gemaakt van twee sleutels:

\begin{itemize}
\tightlist
\item
  Eén publieke sleutel.
\item
  Eén private sleutel.
\end{itemize}

De publieke sleutel kan iedereen, vrij, gebruiken om versleutelde
berichten mee aan te maken. Echter, enkel de eigenaar van de bijhorende
private sleutel zal deze berichten kunnen decrypteren. Kortom, we lossen
nu een deel van het sleutelprobleem op: iedereen heeft z'n eigen private
sleutel (\textbf{en moet deze geheim houden!}) en kan via z'n publieke
sleutel berichten krijgen.

De techniek werd in de jaren 70 door Diffie en Hellman ontwikkeld en had
een grote impact op de manier waarop beveiligde communicatie op het
Internet mogelijk werd maakt.

\begin{figure}
\centering
\pandocbounded{\includegraphics[keepaspectratio,alt={Publieke crypto: Bob gebruikt de publieke sleutel van Alice om haar een beveiligd bericht te sturen.}]{crypto/public.png}}
\caption{Publieke crypto: Bob gebruikt de publieke sleutel van Alice om
haar een beveiligd bericht te sturen.}
\end{figure}

\begin{tipblock}
Je kan publieke encryptie beschouwen als volgt: de publieke sleutel, die
niet geheim is, is een openstaande kist. Iedereen kan de kist gebruiken
om een geheime boodschap in te plaatsen en vervolgens deze op slot te
klikken (\emph{encrypteren}). Enkel de eigenaar van deze kist heeft de
bijpassende private sleutel die echter deze kist kan opendoen en de
originele boodschap kan lezen (\emph{decrypteren}).
\end{tipblock}

\begin{noteblock}
Enkele van de bekendere publieke cryptosystemen zijn onder andere RSA,
DSS en El Gamal.
\end{noteblock}

\subsection{Dualiteit van publieke
cryptografie}\label{dualiteit-van-publieke-cryptografie}

Asymmetrische crypto zal niet alleen het sleutel-probleem oplossen, het
heeft als extra eigenschap dat het publiek/private sleutelpaar ook
dienst kan doen als \textbf{digitale handtekening} om te controleren of
een boodschap wel degelijk afkomstig is van een specifiek persoon.
Hierbij zal een private sleutel gebruikt worden om het digitale bericht
te ondertekenen. Daar de ondertekenaar de enige persoon kan zijn die
deze private sleutel in z'n bezit heeft, kan men zijn identiteit
bevestigen door de bijhorende publieke sleutel te gebruiken. Enkel de
bijhorende publieke sleutel zal hiervoor gebruikt kunnen worden en zo
hebben we een vorm van integriteit of data authenticatie.

\begin{figure}
\centering
\pandocbounded{\includegraphics[keepaspectratio,alt={Het ondertekenen van een document met behulp van je private sleutel.}]{crypto/sign.png}}
\caption{Het ondertekenen van een document met behulp van je private
sleutel.}
\end{figure}

We zullen dit concept verderop uitwerken, maar eerst gaan we bekijken
hoe publieke crypto juist werkt.

\subsection{Diffie-Hellman
sleuteluitwisseling}\label{diffie-hellman-sleuteluitwisseling}

Dankzij public crypto hebben we nu een systeem om sleutels op een
veilige manier uit te wisselen. Het is namelijk zo dat symmetrische
crypto sneller is én dus voor (realtime) communicatie interessanter is.
We weten echter dat het sleutelmanagement bij symmetric crypto een
probleem is als we met grote groepen gebruikers zitten. Het
Diffie-Helman sleuteluitwisselingsconcept helpt ons hierbij: het laat
toe dat twee gebruikers sleutels over een onveilig kanaal kunnen
uitwisselen op een veilige manier.

Zowel Bob als Alice genereren eerst een publiek/private sleutelpaar dat
ze voor deze sessie wensen te gebruiken voor communicatie. Vervolgens
stuurt ieder z'n publieke sleutel naar de ander (wat kan over een
onbeveiligd kanaal). De ontvanger zal deze publieke sleutel combineren
met de eigen private sleutel wat zal resulteren in een nieuw
\emph{shared secret} dat beiden nu kennen en kunnen gebruiken,
bijvoorbeeld, als de symmetrische sleutel om verdere communicatie te
bestendigen.

De reden dat dit werkt, is met dank aan de modulo operator en de
eigenschappen ervan. Een voorbeeld:

{\def\LTcaptype{none} % do not increment counter
\begin{longtable}[]{@{}
  >{\raggedright\arraybackslash}p{(\linewidth - 4\tabcolsep) * \real{0.0606}}
  >{\raggedright\arraybackslash}p{(\linewidth - 4\tabcolsep) * \real{0.4747}}
  >{\raggedright\arraybackslash}p{(\linewidth - 4\tabcolsep) * \real{0.4646}}@{}}
\toprule\noalign{}
\begin{minipage}[b]{\linewidth}\raggedright
Stap
\end{minipage} & \begin{minipage}[b]{\linewidth}\raggedright
Alice
\end{minipage} & \begin{minipage}[b]{\linewidth}\raggedright
Bob
\end{minipage} \\
\midrule\noalign{}
\endhead
\bottomrule\noalign{}
\endlastfoot
1 & Alice kiest een geheim getal A. Ze kiest A= 3 & Bob kiest ook een
geheim getal B = 6. \\
2 & Alice berekent \(7^A\%11\) =\textgreater{} \(343\%11 = 2\) , genaamd
X. & Bob berekent \(7^B\%11\) =\textgreater{} \(11764\%11 = 4\), genaamd
Y. \\
3 & Alice stuurt X=2 naar Bob & Bob stuurt Y=4 naar Alice. \\
4 & Alice berekent \(Y^A\%11 = 4^3\%11 = 9\). & Bob berekent
\(X^B\%11 = 2^6\%11 = 9\). \\
\end{longtable}
}

Zoals je merkt kunnen nu Alice en Bob het berekende getal \textbf{9} als
gedeeld geheim gebruiken. Enkel zij kennen dit getal.

\begin{warningblock}
Uiteraard zullen in de praktijk Bob en Alice véél grotere getallen
kiezen dan 3 en 6.
\end{warningblock}

\begin{noteblock}
Dat Bob en Alice de waarden X en Y naar elkaar kunnen sturen is dankzij
de eigenschappen van de modulo-berekening. X en Y kunnen het resultaat
zijn van een gigantische hoeveelheid berekeningen en een stroper zal dus
veel rekenwerk nodig hebben om alle mogelijkheden te testen.
\end{noteblock}

\subsection{RSA}\label{rsa}

Eén van de oudste, maar nog steeds populairste, publieke cryptosystemen
is het in 1977 ontwikkelde RSA algoritme. RSA, wat staat voor de
achternamen van de drie ontwikkelaars (Rivest, Shamis en Adleman)
gebruikt sleutels van 1536 tot 4096 bits lang. Het systeem is vrij traag
maar heeft als voordeel dat het veilige sleuteltransmissie toestaat over
een onveilig kanaal: we zien daarom vaak RSA gebruikt worden om eerst
sessiesleutels uit te wisselen, vervolgens wordt overgeschakeld op een
sneller symmetrisch cipher.

De exacte berekeningen die gebeuren tijdens encryptie en decryptie
leiden ons iets te ver, maar volgend voorbeeld toont een vereenvoudigde
wijze waarop RSA wordt toegepast:

Data encrypteren met behulp van asymmetrische versleuteling gebeurt op
bijna dezelfde wijze als de Diffie-Hellman sleutel uitwisseling. Ook nu
zullen beide zijden rekenen op de eigenschappen van de modulo-operator
om over een onveilig kanaal veilige communicatie te kunnen doen.

Eerst dient een publieke sleutel aangemaakt te worden:

\begin{itemize}
\tightlist
\item
  Hiertoe dient Bob twee grote priemgetallen, \texttt{q} en \texttt{p}
  te kiezen, bijvoorbeeld \texttt{p=17} en \texttt{q=11}.
\item
  Vervolgens berekent Bob \texttt{N} door deze priemgetallen met elkaar
  te vermenigvuldigen (\texttt{p*q} geeft 17 * 11, \texttt{N=187}).
\item
  Bob kiest nu nog een priemgetal \texttt{e}, bijvoorbeeld \texttt{7}.
\item
  Bob kan nu zijn eigen geheime, private sleutel \texttt{d} maken,
  namelijk \(e*d = 1\%((p-1)*(q-1))\) wat dus \(7*d=1\%(16*10)\) geeft
  of \(7*d=1\%160\) .
\item
  Om nu \texttt{d} te vinden moeten we een getal vinden zodat \(7*d\)
  een veelvoud van \(1\%160\) geeft, dus bijvoorbeeld 1, 161, etc. In
  dit geval vinden we \texttt{d=23}.
\end{itemize}

\begin{tipblock}
Om \(d\) te berekenen maken we gebruik van de zogenaamde
\emph{Uitgebreid Euclidisch algoritme}, een eeuwenoud algoritme
gebaseerd op het \emph{Algoritme van Euclides} dat we in het lager
leerden gebruiken om de grootste gemene deler te berekenen van twee
getallen.
\end{tipblock}

Bob heeft dus nu:

\begin{itemize}
\tightlist
\item
  Publieke sleutel bestaande uit \(N=187\) en \(e=7\).
\item
  Private sleutel \(d=23\).
\end{itemize}

Iedereen die nu naar Bob iets wilt sturen, kan dit via z'n publieke
sleutel (\texttt{N} en \texttt{e}).

Stel dat Alice het ASCII-karakter X naar Bob wil sturen:

\begin{itemize}
\tightlist
\item
  De ASCII-waarde van X is 88.
\item
  De encryptie door Alice gebeurt dan als volgt: \(C = data^e\%N\).
\item
  De te versturen ciphertext C wordt dus: \((88^7)\%187\) oftewel
  \texttt{C=11}.
\end{itemize}

Enkel Bob zal deze ciphertext met zijn private sleutel \texttt{d} kunnen
decrypteren door \(C^d(\%N)\) te doen, oftewel \(11^{23}\%187\) wat
terug de plaintext \texttt{88} geeft!

\begin{tipblock}
De sterkte van publieke crypto stoelt dus op het feit dat ontbinden van
(grote) getallen in factoren computationeel veel moeilijker is dan de
omgekeerde stap, namelijk twee getallen met elkaar vermenigvuldigen.

15621 in z'n factoren ontbinden is veel moeilijker dan de getallen 123
en 127 vermenigvuldigen (wat dus ook 15621 zal geven).
\end{tipblock}

Zonder in detail te treden hoe cryptocoins en blockchains werken, is het
nuttig om te vermelden dat bij cryptocoins ook de public crypto
concepten worden gebruikt. Ook hier is je private sleutel uiterst
belangrijk: enkel de eigenaar van de private sleutel ``bezit'' de
bijhorende cryptocoins in de \emph{blockchain}. Daarom is het belangrijk
dat je \textbf{NOOIT} je private sleutel aan derden geeft, want zo geef
je hen toegang tot jouw \emph{coins} en kunnen ze vervolgens deze stelen
door de private sleutel te vervangen.

\subsubsection{Intermezzo: Hashes}\label{intermezzo-hashes}

We gaan nu even een zijtak inslaan om het concept ``hash'' te bespreken.
Een hash is een concept uit de informatica dat we gebruiken om te
controleren of een digitaal stuk tekst werd aangepast of niet. Door de
tekst in een hashfuntie te steken wordt een hash aangemaakt. Deze hash
is een stuk code met een vaste lengte, ongeacht de originele input.
Wanneer 1 bit of meer wordt aangepast in de originele boodschap dan zal
deze in een totaal andere hash resulteren. Enkel dus wanneer een
identiek stuk tekst als invoer (tot op bitniveau identiek) wordt
gebruikt, zullen twee hashes gelijk zijn.

Voorgaande is uiteraard onmogelijk: daar een hash meestal veel korter is
dan de originele boodschap, is het mathematisch mogelijk dat twee totaal
verschillende teksten toch dezelfde hash geven. Het is de opdracht van
een goede hashfunctie om dit soort \textbf{hash collisions} zo klein
mogelijk te houden.

\begin{figure}
\centering
\includegraphics[width=0.6\linewidth,height=\textheight,keepaspectratio,alt={Het hash proces.}]{crypto/hash.png}
\caption{Het hash proces.}
\end{figure}

Een hashfunctie is niet omkeerbaar: men mag onmogelijk aan de hand van
een hash (ook wel \emph{digest} of \emph{hashcode} genoemd) terug de
originele tekst kunnen achterhalen. Een hashfunctie is dus een
eenrichtingsfunctie, ook wel afbeelding genoemd in wiskundige termen.

Er bestaan veel verschillende hashfuncties. Enkele van de bekendere
zijn:

\begin{itemize}
\tightlist
\item
  MD5, oftewel Message Digest 5: deze zal een 128-bit hashwaarde
  genereren.
\item
  SHA-X, oftewel Secure Hash Algorithms. Zo is er SHA-256 wat een 256
  bits hash zal genereren.
\end{itemize}

De sterkte van een hash algoritme zit hem in de grootte van de kans
waarop hash collisions kunnen optreden. Zo zijn er bij MD5 al veel meer
collisions gevonden dan bijvoorbeeld bij het recenter gepubliceerde
SHA3-512 algoritme.

\textbf{Een digitale hash is dus een ideaal middel om boodschappen
digitaal te ondertekenen.}

\subsection{Boodschappen ondertekenen}\label{boodschappen-ondertekenen}

Een probleem bij online communicatie is dat we geen zekerheid hebben dat
het ontvangen bericht wel degelijk van de persoon komt van wie we
verwachtten dat deze het bericht had opgesteld. We kunnen daarom het
publieke crypto systeem gebruiken om berichten te ondertekenen. Daar
enkel Bob de bijhorende private sleutel kan hebben die bij z'n publieke
sleutel hoort, is het bezit van deze private sleutel hebben het bewijs
dat hij de rechtmatige eigenaar van een bepaalde publieke sleutel is.

Onder andere RSA laat toe om een digitale handtekening (\textbf{digital
signature}) te plaatsen bij een bericht om zo te bewijzen dat de
verzender de rechtmatige eigenaar van een bijhorende publieke sleutel
is.

Een digitale handtekening wordt als extra bericht achteraan de te
versturen boodschap geplaatst. De signature is als het ware een hash
berekend aan de hand van de private sleutel. De ontvanger zal nu met de
bijhorende publieke sleutel van de verzender kunnen verifiëren of de
bijhorende private sleutel werd gebruikt om de handtekening te
genereren.

\begin{figure}
\centering
\pandocbounded{\includegraphics[keepaspectratio,alt={Boodschappen ondertekenen met je private sleutel.}]{crypto/signover.png}}
\caption{Boodschappen ondertekenen met je private sleutel.}
\end{figure}

Om een digitale handtekening te berekenen moeten we eerst een hash van
het bericht berekenen. We gebruiken hier bijvoorbeeld MD5 of één van de
SHA-algoritmes voor. Deze hash gaan we nu ``verpakken'' met de private
sleutel.

Stel dat we willen versturen \texttt{35} is en we hebben reeds volgende
sleutelpaar berekend
(\href{HTTPS://crypto.stackexchange.com/questions/11117/simple-digital-signature-example-with-number}{bron}):

\begin{itemize}
\tightlist
\item
  publieke sleutel: \(e=5\) en \(n=91\).
\item
  private sleutel: \(d=29\).
\end{itemize}

De handtekening wordt berekend door: \(s=m^d\%n\) oftewel
\(s=35^{29}\%91\) wat \texttt{42} geeft.

We versturen dus naar de ontvanger de boodschap zelf en de bijhorende
handtekening: \texttt{35,42}.

De ontvanger kan nu controleren of de ontvangen boodschap klopt of niet.
Hij zal zijn eigen hash genereren van de ontvangen boodschap. Als deze
overeen komt met de hash die bij de boodschap zat, is alles in orde. Om
de handtekening te decrypteren gebruikt de ontvanger \texttt{e} en
berekent: \(42^e == 35^n\), als dit overeenkomt dan weet de ontvanger
dat hij het bericht kan vertrouwen.

\begin{figure}
\centering
\includegraphics[width=0.8\linewidth,height=\textheight,keepaspectratio,alt={Het volledig proces bij een digitale handtekening.}]{crypto/signaturesend.png}
\caption{Het volledig proces bij een digitale handtekening.}
\end{figure}

\subsubsection{Het probleem met digitale
handtekeningen}\label{het-probleem-met-digitale-handtekeningen}

We hebben echter een probleem. Hoe weet je eigenlijk dat je wel de
juiste publieke sleutel gebruikt? Publieke sleutels zijn, wel, publiek.
Iedereen kan jou een publieke sleutel geven en zeggen \emph{``Dit is de
sleutel van persoon X''} zonder dat jij kan controleren of dat zo is.

We kunnen daarom als kwaadwillig persoon bijvoorbeeld een legaal bericht
onderscheppen, aanpassen en dan vervolgens ondertekenen met onze eigen
handtekening. Als we vervolgens aan de ontvanger kunnen wijsmaken dat
jouw publieke sleutel zogezegd bij de originele verzender hoort, dan zal
de ontvanger jouw aangepaste bericht ``geloven''.

\begin{figure}
\centering
\pandocbounded{\includegraphics[keepaspectratio,alt={Eve misbruikt het vertrouwen dat zit ingebouwd in het digitale handtekening proces.}]{crypto/signfout.png}}
\caption{Eve misbruikt het vertrouwen dat zit ingebouwd in het digitale
handtekening proces.}
\end{figure}

Kortom, we hebben een manier nodig om de \textbf{identiteit van de
eigenaar} van een publieke sleutel te verifiëren. Kom binnen:
\textbf{certificaten}.

\section{Digitale certificaten}\label{digitale-certificaten}

Een certificaat is een (digitaal) document dat de identiteit van een
gebruiker bindt aan een publieke sleutel. Dit document werd digitaal
ondertekend door een vertrouwde derde partij (\textbf{trusted third
party}) zodat bij twijfel van de echtheid van het certificaat men altijd
bij deze derde partij terecht kan. Uiteraard is het belangrijk dat we
deze derde partij kunnen vertrouwen, anders kunnen we ook niet het
certificaat vertrouwen die zij onderschrijven.

\begin{tipblock}
Certificaten worden beschreven in de \textbf{X.509} standaard.
\end{tipblock}

Om een certificaat aan te maken dient Bob naar een \textbf{Registration
authority} (RA) te gaan die zijn identiteit zal verifiëren. Dit gebeurt
aan de hand van de typische documenten die ook buiten het Internet
worden gebruikt om iemands identiteit te bewijzen: identiteitskaart,
paspoort, rijbewijs, etc. In sommige gevallen zal de RA zelfs eisen dat
Bob zich naar een fysiek kantoor begeeft om daar z'n identiteit \emph{in
the flesh} te bewijzen. Indien de RA de identiteit heeft bevestigd zal
deze de aanvraag van Bob doorsturen naar een \textbf{Certification
authority} (CA), inclusief Bobs publieke sleutel, die een certificaat
zal aanmaken én ondertekenen.

\begin{figure}
\centering
\pandocbounded{\includegraphics[keepaspectratio,alt={Een certificaat registreren.}]{crypto/certreg.png}}
\caption{Een certificaat registreren.}
\end{figure}

\begin{tipblock}
Het gehele systeem van CA's, RA's, etc. dat bestaat om certificaten uit
te geven, beheren en bewijzen heet een \textbf{Public Key
Infrastructure} (\textbf{PKI}).
\end{tipblock}

De CA zal deze informatie gebruiken om een certificaat, van een bepaalde
levensduur, te genereren. Hierbij zal de echtheid van het certificaat
achteraf bewezen kunnen worden door de CA:

\begin{itemize}
\tightlist
\item
  Het certificaat is een geëncrypteerde hash van Bobs publieke sleutel,
  informatie over de CA en over Bob. De encryptie van de hash gebeurt
  aan de hand van de private sleutel van de CA.
\item
  Om later de echtheid van een certificaat te testen voldoet het om
  dezelfde hash te genereren (publieke sleutel, info over Bob en CA) en
  deze te vergelijken met het certificaat na decryptie met de publieke
  sleutel van de CA. Als deze gelijk zijn weten we dat het certificaat
  door de gegeven CA werd ondertekend (enkel hun publiek/private
  sleutelpaar zal terug de originele hash geven).
\end{itemize}

\begin{figure}
\centering
\pandocbounded{\includegraphics[keepaspectratio,alt={Een certificaat aanmaken.}]{crypto/certcreatie.png}}
\caption{Een certificaat aanmaken.}
\end{figure}

Voorgaande proces zal bijvoorbeeld plaatsvinden wanneer je browser via
een \textbf{HTTPS} verbinding surft naar een website en zo wil
controleren of wel degelijk met de website wordt gecommuniceerd en niet
met een imposter. Indien de browser (of de gebruiker) twijfelt aan de
echtheid van de publieke sleutel van de CA die het certificaat van de
website ondertekent, dan zal het voorgaande proces zich herhalen, maar
deze keer om het certificaat van de CA te controleren met behulp van een
bovenliggende CA. Op die manier kan het dus zijn dat een keten van CA's
ontstaan die telkens CA's onder zich bewijzen. Uiteraard zal er steeds
bovenaan zo'n ketting een \textbf{root CA} staan. Als je die vertrouwt,
dan kan je al de CA's er onder dus ook vertrouwen\ldots maar ook vice
versa!

\begin{figure}
\centering
\includegraphics[width=0.8\linewidth,height=\textheight,keepaspectratio,alt={Het certificaat tijdens het surfen.}]{crypto/webcert.png}
\caption{Het certificaat tijdens het surfen.}
\end{figure}

\begin{noteblock}
Alhoewel \textbf{HTTPS} al sinds 1995 bestond, werd het tot voor kort
amper door websites aangeboden. Nochtans geeft HTTPS een extra
defensielaag tijdens de communicatie van jouw computer met die waar een
website op \emph{gehost} staat. HTTPS zal namelijk je communicatie
versleutelen zodat enkel zender en ontvanger kunnen lezen wat er gezegd
wordt. Met HTTP is dat niet: al je communicatie kan door eender wie
gelezen worden die zich tussen jouw computer en je eindbestemming
nestelt. Het helpt echter niet dat je data versleuteld wordt als je niet
kan bevestigen dat de ontvangende website ook effectief diegene is die
je nodig hebt, vandaar dat dus certificaten en HTTPS in tandem werken om
gebruikers een veiliger Internet aan te bieden.

Pas in 2017 boden meer dan de helft van de websites wereldwijd HTTPS
aan. In 2021 gebruikt ongeveer 70\% van alle websites HTTPS als
standaard communicatiemiddel aan (vroeger waren er al websites met
HTTPS, maar HTTP was de standaard oplossing).
\end{noteblock}

Het ergste dat voor een CA kan voorvallen is dat de betrouwbaarheid van
de CA in het gedrang komt. Als een CA bijvoorbeeld weet heeft van een
potentiële inbraak op zijn systemen dan bestaat er de kans dat
aanvallers de private sleutel van de CA hebben bemachtigd en dus zelf
certificaten \emph{op naam van de CA} kunnen genereren, met alle
gevolgen van dien! Indien dus deze kans bestaat, is er een \emph{breach
of trust} en zullen alle certificaten van deze CA als ongeldig worden
bestempeld, inclusief alle certificaten van sub-CA's! Dit kan
verregaande gevolgen hebben.

\begin{figure}
\centering
\includegraphics[width=0.6\linewidth,height=\textheight,keepaspectratio,alt={De chain-of-trust: oh zo belangrijk bij digitale certificaten.}]{crypto/chaintrust.png}
\caption{De chain-of-trust: oh zo belangrijk bij digitale certificaten.}
\end{figure}

\subsection{Certificaten bekijken}\label{certificaten-bekijken}

In iedere moderne browser kan je snel bekijken hoe zo'n certificaat er
juist uitziet. Als je via een HTTPS verbinding naar een website surft,
dan op het slotje naast de URL in de adresbalk klikt kan je doorklikken
om het certificaat te openen. Als je naar \emph{HTTPS://www.belgium.be}
surft en dit doet dan krijg je eerst wat samenvattende informatie:

\begin{figure}
\centering
\pandocbounded{\includegraphics[keepaspectratio,alt={Het certificaat van België.}]{crypto/belcert0.png}}
\caption{Het certificaat van België.}
\end{figure}

Zo zien we onder andere de geldigheidsduur, alsook de CA die dit
certificaat heeft gegenereerd. Onder details kunnen we onder andere de
publieke sleutel zien van de website alsook de gebruikte algoritmes voor
de hash, e.d.

En op de laatste tab, Certificeringspad, zien we de chain of trust. We
kunnen vervolgens hier de bovenliggende certificaten bekijken.

Het certificaat van Sectigo is uiteraard een \textbf{selfsigned
certificate}, daar zij ``bovenaan de hiërarchie staan''. Als we Sectigo
niet vertrouwen dan kunnen we ook de communicatie met \emph{belgium.be}
niet vertrouwen.

\begin{figure}
\centering
\includegraphics[width=0.4\linewidth,height=\textheight,keepaspectratio,alt={Sectigo heeft een self-signed certificaat wat je herkent aan het feit dat de velden Verleend aan en Verleend door dezelfde waarde hebben.}]{crypto/belcert3.png}
\caption{Sectigo heeft een self-signed certificaat wat je herkent aan
het feit dat de velden \emph{Verleend aan} en \emph{Verleend door}
dezelfde waarde hebben.}
\end{figure}

\subsection{Persoonlijke certificaten}\label{persoonlijke-certificaten}

Naast certificaten voor webservers (zogenaamde \textbf{SSL
certificaten}) kan je ook een persoonlijk certificaat aankopen om je
eigen identiteit aan derden te bewijzen tijdens bijvoorbeeld
e-mail-communicatie. Voorts heb je ook \textbf{code signing}
certificaten die de echtheid van een applicatie bewijzen zodat je zeker
bent dat je geen malware installeert als je programma X hebt gedownload.

Als je in Windows 10 of nieuwer een applicatie of installer probeert uit
te voeren dan zal de ingebouwde \emph{SmartScreen} service
ogenblikkelijk de echtheid (of ontbreken van) het certificaat
controleren, net zoals dit ook in de browser zou gebeuren.

\begin{figure}
\centering
\includegraphics[width=0.4\linewidth,height=\textheight,keepaspectratio,alt={Windows 10 Smartscreen beschermt je van niet digitaal ondertekende software.}]{crypto/smartscreen.png}
\caption{Windows 10 Smartscreen beschermt je van niet digitaal
ondertekende software.}
\end{figure}

Wil dat dan zeggen dat je applicaties niet kunt vertrouwen die door
Smart Screen als onveilig worden aangeduid? Neen, dat niet. Je mag niet
vergeten dat een certificaat geld kost en dat niet alle
software-ontwikkelaars de middelen hebben om een officieel certificaat
te kopen. Het loont dus altijd om extra waakzaam te zijn wanneer Smart
Screen een waarschuwing geeft, maar het is dus niet zo dat de software
automatisch als onveilig moet gehanteerd worden.

\begin{tipblock}
Je kan via de Certification Manager van Windows bekijken welke
certificaten je lokaal hebt geïnstalleerd, welke worden vertrouwd, etc.
Je kan de GUI-versie van deze tool opstarten door ``certlm.msc'' uit te
voeren.

\begin{figure}
\centering
\includegraphics[width=0.4\linewidth,height=\textheight,keepaspectratio,alt={Let er altijd op welke informatie je deelt via screenshots.}]{crypto/certlm.jpg}
\caption{Let er altijd op welke informatie je deelt via screenshots.}
\end{figure}

\end{tipblock}

\section{HTTPS en TLS}\label{https-en-tls}

Certificaten vormen het hart van een veilige manier van surfen. HTTPS,
``HTTP-Secure'', zorgt ervoor dat de communicatie tussen je browser en
de website via een beveiligde, geëncrypteerde tunnel gebeurt.
Tegenwoordig is HTTPS de default manier om een website te benaderen,
maar dat is maar recent. Vroeger gebeurde alles via HTTP, waardoor
iedereen die jouw trafiek kon sniffen, kon zien welke informatie je met
de website uitwisselde.

HTTPS is een protocol dat een beveiligde encryptietunnel opzet tussen
jou en de website waarover vervolgens gewoon HTTP-verkeer kan verlopen
(dit gebeurt over poort 443 in plaats van de klassieke poort 80 waarover
HTTP verloopt). Deze tunnel wordt opgezet door het
\textbf{TLS}-protocol, het \emph{Transport Layer Security} protocol, dat
de opvolger is van \textbf{SSL} (\emph{Secure Sockets Layer}). TLS
gebruikt certificaten om te vergewissen dat de website aan de andere
zijde wel degelijk de website is die de gebruiker verwacht. Het doet dit
door de publieke sleutel van de website eerst te controleren voor het
vervolgens deze sleutel gebruikt om een gemeenschappelijke sleutel af te
spreken die zal dienst doen als de encryptie-sleutel voor de
gemeenschappelijke tunnel. Vanaf dit punt kunnen derden de trafiek van
en naar de website niet meer sniffen.

Samengevat: \textbf{HTTPS is de combinatie van HTTP en een veilige
encryptietunnel die met behulp van het TLS protocol wordt opgezet.}

\begin{figure}
\centering
\pandocbounded{\includegraphics[keepaspectratio,alt={De erg belangrijke TLS tunnel tijdens het surfen met HTTPS.}]{crypto/tlshttp.png}}
\caption{De erg belangrijke TLS tunnel tijdens het surfen met HTTPS.}
\end{figure}

Samengevat zal dus TLS twee zaken doen:

\begin{enumerate}
\def\labelenumi{\arabic{enumi}.}
\tightlist
\item
  Door middel van een certificaat (asymmetrische crypto) wordt de
  identiteit (de publieke sleutel) van de website gecontroleerd.
\item
  Door middel van een afgesproken algoritme (Diffie-Hellman, Forward
  Secrecy, Elliptic Curve, etc.) een gemeenschappelijke sleutel(s)
  afspreken en uitwisselen.
\end{enumerate}

\begin{tipblock}
Zoals reeds eerder vermeld is asymmetrische crypto trager, waardoor het
altijd aanbevolen is om de trafiek tussen 2 punten finaal via een
symmetrische crypto verbinding te laten plaatsvinden. TLS/HTTPS
combineert met andere woorden de sterktes van beide soorten crypto om zo
de zwaktes van beiden te neutraliseren.
\end{tipblock}

De manier waarop een TLS-verbinding wordt opgezet is vrij uitgebreid.
Volgende briljante website
(\href{HTTPS://tls.ulfheim.net/}{tls.ulfheim.net/}) visualiseert de
berichten die server en client uitwisselen om zo'n verbinding te
starten, onderhouden en eindigen.

\begin{warningblock}
Alhoewel HTTPS onze verbinding een pak veiliger maakt, heeft het voor je
ISP (Internet Service Provider, bijvoorbeeld Telenet of Proximus) en de
website ook enkele nadelen. Omdat alle informatie geëncrypteerd wordt
heeft de ISP geen enkel idee wat voor informatie je aan het uitwisselen
bent, waardoor caching ook niet meer mogelijk is. In een normale
HTTP-omgeving kan een ISP trafiek over het Internet uitsparen door een
reeds bewaarde versie van hetgeen jij nodig hebt uit de cache te halen
en naar je te sturen. Ook de website naar waar je surft, ondervindt dit
nadeel: het zal met HTTPS veel meer trafiek genereren dan wanneer de
tussenliggende ISP een deel van het werk via zijn caching overnemen.
\end{warningblock}

\subsection{Mitmproxy}\label{mitmproxy}

\emph{Mitmproxy} is een krachtige linux-tool die een man-in-the-middle
aanval op HTTPS toelaat. Het zal ervoor zorgen dat een aanvaller zich
tussen jou en het Internet kan nestelen en vervolgens doen alsof al je
HTTPS-verbindingen veilig blijven. In de praktijk zorgt mitmproxy ervoor
dat alle HTTPS-verbindingen van de client naar de aanvaller gebeuren,
die op zijn beurt TLS tunnels zal opzetten met de website waar het
slachtoffer naar surft. Hierdoor kan de aanvaller enerzijds alle trafiek
lezen, maar bijvoorbeeld ook ongezien aanpassen.

\begin{figure}
\centering
\includegraphics[width=0.9\linewidth,height=\textheight,keepaspectratio,alt={Een man-in-the-middle aanval met TLS.}]{crypto/mitmproxy.png}
\caption{Een man-in-the-middle aanval met TLS.}
\end{figure}

De aanvaller zal echter nog steeds geen geldige certificaten kunnen
genereren waardoor moderne browsers normaal gezien hier een waarschuwing
zouden moeten geven.

\begin{figure}
\centering
\includegraphics[width=0.7\linewidth,height=\textheight,keepaspectratio,alt={De waarschuwing die Chrome genereert wanneer het een, potentiële, mitm-aanval detecteert.}]{crypto/mitmbrowser.png}
\caption{De waarschuwing die Chrome genereert wanneer het een,
potentiële, mitm-aanval detecteert.}
\end{figure}

\chapter{Wifi security}\label{wifi-security}

\section{De problemen van wifi}\label{de-problemen-van-wifi}

We kunnen draadloze netwerken, specifiek wifi-netwerken, niet meer uit
ons leven wegdenken. De opkomst van de IEEE 802.11b standaard (spreek
IEEE uit als \emph{``Ai-trippel-i''}) in 1999 veroorzaakte een kleine
revolutie in de manier waarop bedrijven en privégebruikers konden
werken. Plots kon je met een laptop van overal in het gebouw - en zelfs
er buiten - op het netwerk geraken. Die vrijheid voor de gebruikers
betekende wel een nachtmerrie voor de cyberboswachters. Een netwerkkabel
heeft een intrinsieke extra beveiliging: enkel daar waar de kabel ligt
kunnen gebruikers op het netwerk geraken. Zolang je dus geen
netwerkkabel naar de publieke parking brengt kan niemand van daar
illegaal het netwerk benaderen. Met wifi leek het alsof plotseling het
hele bedrijfsnetwerk in een straal van tientallen meters rond het gebouw
beschikbaar was, met alle gevolgen van dien.

\begin{figure}
\centering
\includegraphics[width=0.75\linewidth,height=\textheight,keepaspectratio,alt={Een oud voorbeeld van hoe wifi signalen ``uit'' een gebouw veel verder geraken dan verwacht (bron van de afbeelding: onbekend)}]{wifi/afstand.jpg}
\caption{Een oud voorbeeld van hoe wifi signalen ``uit'' een gebouw veel
verder geraken dan verwacht (bron van de afbeelding: onbekend)}
\end{figure}

Al gauw werd een nieuwe sport uitgevonden door hobbyist hackers en
professionele cybercriminelen: \textbf{wardriving}. Het idee is
eenvoudig: je rijdt rond in de stad en laat de laptop naast je in de
auto scannen naar alle draadloze netwerken, met extra aandacht voor die
netwerken die geen of zwakke beveiliging hebben.

\begin{noteblock}

\begin{figure}
\centering
\includegraphics[width=0.3\linewidth,height=\textheight,keepaspectratio,alt={Wargames: een cultklassieker uit 1983.}]{wifi/wargames.png}
\caption{Wargames: een cultklassieker uit 1983.}
\end{figure}

De term wardriving komt van de term \emph{wardialing} die op zijn beurt
gebaseerd is op de klassieke cyber-cult film ``Wargames'' uit 1983. Voor
de geschiedkundigen onder ons, wardialing was het opbellen van
willekeurige telefoonnummers met je modem in de hoop een zogenaamd
\emph{bulletin board system} oftewel \textbf{BBS} (een pre-Internet
forum zeg maar) te vinden.
\end{noteblock}

Draadloze netwerken die gevonden worden hebben dan ook een schare aan
problemen:

\begin{itemize}
\tightlist
\item
  \textbf{Eavesdropping}: iedereen kan ``zien'' wat er door de lucht
  vliegt. Dit heb je niet met een kabel. Bij een bedraad netwerk is de
  fysieke koperdraad het communicatiemedium, bij wifi is dat eigenlijk
  letterlijk de lucht.
\item
  \textbf{Invasion}: je kan eenvoudig verbinden met het netwerk indien
  er geen beveiliging voorzien werd. Iets dat je bij een bedraad netwerk
  enkel kunt als je fysiek toegang hebt tot een netwerkkabel.
\item
  \textbf{Man-in-the-middle aanvallen}: hier gaan we zo meteen dieper op
  in.
\item
  \textbf{Backdoor}: een veelvoorkomend probleem zijn zogenaamde
  \textbf{rogue access points} die worden bijgeplaatst op het netwerk
  door goedbedoelde werknemers die zo het bereik van het netwerk wat
  willen uitbreiden. Vaak zijn de beveiligingsinstellingen van zo'n
  (meestal goedkoop) access point niet zo sterk als die van het bedrijf.
  Bijgevolg is dit de ideale backdoor voor malafide gebruikers die het
  netwerk aan het surveilleren zijn. Wanneer ze een netwerkscan doen
  zullen ze tientallen goed beveiligde access points zien en één zwak
  beveiligd.
\end{itemize}

\begin{figure}
\centering
\includegraphics[width=0.7\linewidth,height=\textheight,keepaspectratio,alt={Een rogue access point is de ideale manier om binnen te geraken voor hackers.}]{wifi/rogue.png}
\caption{Een rogue access point is de ideale manier om binnen te geraken
voor hackers.}
\end{figure}

\begin{figure}
\centering
\includegraphics[width=0.7\linewidth,height=\textheight,keepaspectratio,alt={Het kan iedereen overkomen\ldots{}}]{wifi/backmaga.png}
\caption{Het kan iedereen overkomen\ldots{}}
\end{figure}

\begin{itemize}
\tightlist
\item
  \textbf{Denial-of-service op fysiek niveau}: om een gebruiker toegang
  tot een bedraad netwerk te ontzeggen op fysiek niveau dien je de kabel
  door te knippen. Bij wifi is dit nog eenvoudiger: de ``kabel'' bij
  wifi zijn de frequentiebanden in de lucht waarbinnen de apparaten
  mogen werken (circa 2.4 Ghz bij de oudere wifi-apparaten, nu meestal
  rond de 5 Ghz band). De wifi-apparaten kunnen enkel met elkaar
  communiceren indien zij een signaal naar elkaar over die
  frequentieband kunnen sturen op een moment dat niemand anders in de
  buurt die band gebruikt (zie note hierna). Als een malafide gebruiker
  die frequentieband continue vult met andere signalen, dan zullen de
  legale gebruikers nooit iets kunnen uitsturen. Wil je dus een wifi
  netwerk \emph{DoS'n}, koop dan een signaalgenerator die op de juiste
  frequentieband de nodige ruist uitzendt en klaar is kees.
\end{itemize}

\begin{noteblock}
Bedrade netwerktoestellen werken volgens het CSMA/CD (Carrier Sense
Multiple Access / Collision Detection) om met elkaar over de draad te
communiceren, hierbij detecteren ze wanneer er zich `botsingen' tussen
signalen voordoen en de informatie dus opnieuw moet uitgestuurd worden.
Bij draadloze netwerken is het echter onmogelijk om \emph{botsingen in
de lucht} te detecteren, daarom werken ze met een variant:
\textbf{CSMA/CA} , oftewel CSMA/ \textbf{Collision Avoidance}. Een
wifi-apparaat dat iets wil uitsturen zal eerst controleren of de
frequentieband waarop ze werken vrij is. Vervolgens zal het apparaat
broadcasten dat het iets gaat versturen, gevolgd door de effectieve
informatie.
\end{noteblock}

\subsection{Onbeveiligde
managementframes}\label{onbeveiligde-managementframes}

Veel keuzes die in de IEEE 802.11 standaard werden gemaakt zijn, zijn
vermoedelijk het gevolg van leentje buur spelen bij de reeds bestaande
IEEE standaarden voor bedrade netwerken. Echter, bij bedrade netwerken
had je niet de inherente onveilige omgeving van het draadloze aspect,
waardoor op gebied van beveiliging hier weinig tot geen aandacht aan
werd besteed.

Management frames in wifi zijn frames die ervoor zorgen dat alle
aanwezigen op het netwerk op een ordentelijke manier met elkaar kunnen
communiceren. Deze frames hebben onder andere als doel:

\begin{itemize}
\tightlist
\item
  Om een client verbinding te laten maken met een netwerk.
\item
  Om SSIDs te broadcasten.
\item
  Clients de opdracht te geven het netwerk te verlaten.
\item
  Om te verbinden met een ander access point van hetzelfde netwerk
  (\emph{handover}).
\end{itemize}

Net zoals bij bedrade netwerken, koos men bij de IEEE wifi standaard om
deze managementframes \textbf{zonder enige vorm van beveiliging} te
gebruiken (er werd geen \emph{CIA} voorzien). Dit zorgde ervoor dat niet
legale gebruikers zelf management frames konden uitsturen op een
netwerk, zonder dat de ontvangers ervan konden controleren of de bron
wel een legaal access point of client was.

Dit resulteerde in onder andere volgende scenario's:

\begin{itemize}
\tightlist
\item
  \textbf{Disassociation flooding}: een hacker kan legale gebruikers
  DoS'n door constant zogenaamde \emph{disassociation frames} naar hen
  te sturen. Dit frame, gebruikt door access points, geeft aan clients
  de opdracht dat ze het netwerk moeten verlaten. De hacker kan zo'n
  frame uitsturen en daarbij het ``source'' veld instellen op het
  MAC-adres van het access points. Deze vorm van spoofing kan
  ongecontroleerd gebeuren, waardoor legale gebruikers dit frame altijd
  zullen aanvaarden én vervolgens uitvoeren: de gebruiker kan niet meer
  verbinden met het netwerk zolang de disassociation frames blijven
  verstuurd worden.
\item
  \textbf{Identity spoofing} was ook eenvoudig daar de hacker eender
  welk veld in de frames kan aanpassen. Van zodra hij een legale
  gebruiker met voorgaande techniek van het netwerk heeft geschopt, kan
  hij vervolgens zichzelf voordoen als deze gebruiker. Hiervoor moet hij
  gewoon het MAC-adres spoofen van de legale gebruiker tijdens de
  communicatie met het access point.
\end{itemize}

\begin{figure}
\centering
\includegraphics[width=0.8\linewidth,height=\textheight,keepaspectratio,alt={Eve gebruikt spoofing om de wifi-sessie van Alice ongemerkt over te nemen.}]{wifi/spoofwifi.png}
\caption{Eve gebruikt spoofing om de wifi-sessie van Alice ongemerkt
over te nemen.}
\end{figure}

\begin{itemize}
\tightlist
\item
  En \emph{last but not least} laten de onbeveiligde management frames
  toe dat we eenvoudig een access point kunnen nabootsen
  (\textbf{impersonation}). Vervolgens kunnen we een
  \textbf{man-in-the-middle aanval} uitvoeren daar een hacker zich kan
  plaatsen tussen de gebruiker en het Internet en zo informatie kan
  ontfutselen.
\end{itemize}

\begin{figure}
\centering
\includegraphics[width=0.7\linewidth,height=\textheight,keepaspectratio,alt={Een fake access point opzetten met dank aan de onbeveiligde managementframes.}]{wifi/mitmwifi.png}
\caption{Een fake access point opzetten met dank aan de onbeveiligde
managementframes.}
\end{figure}

\begin{tipblock}
De linux tool \emph{AirSnarf} laat toe om fake hotspots (publiek wifi
netwerk) op te zetten. Hierbij maakt het gebruikt van de onbeveiligde
management frames. Een scenario in België dat gegarandeerd succes had
tot enkele jaren geleden (vanuit het standpunt van de hacker) was een
fake Telenet Wi-Free hotspot op te zetten. Hierbij werd eerst de
inlogpagina van Telenet Wi-Free door de hacker gecloned en via een
lokale webserver aangeboden aan de gebruikers die op het fake access
point met de naam ``Telenet Wi-Free'' verbonden. Vervolgens kon de
stroper nu van iedere gebruiker de gebruikersnaam en het wachtwoord
stelen telkens deze die informatie op de fake loginpagina invoerde.

Nu dat Telenet Wi-Free is overgeschakeld op WPA-Enterprise (zie verder)
kan deze aanval gelukkig niet meer zo eenvoudig uitgevoerd worden.
\end{tipblock}

\section{De 802.11 standaard qua
beveiliging}\label{de-802.11-standaard-qua-beveiliging}

\emph{Om de huidige, en betere, beveiliging van wifi te appreciëren gaan
we terug in de tijd om te kijken hoe de originele IEEE wifi standaard de
beveiliging beschreef. Het zal een ietwat horror-achtige tocht worden
waarin we gaan ontdekken dat enkele stevige hiaten ervoor gezorgd hebben
dat illegale toegang tot bijna ieder wifi netwerk rond de eeuwwisseling
binnen enkele minuten kon gebeuren. Lees verder en huiver mee.}

De originele ``ANSI/IEEE Std. 802.11'' die de wifi specificaties
beschrijft, werd geschreven in 1999. Het had als doel \emph{``to develop
a medium access control (MAC) and physical layer (PHY) specification for
wireless connectivity for fixed, portable, and moving stations within a
local area.''} Hoofdstuk 8 van deze standaard had als titel
``Authentication and privacy'' en was maar tien pagina's lang, in
vergelijking met de totale grootte van het document (528 pagina's) was
dit misschien wel een voorbode hoe weinig aandacht er aan beveiliging
zou worden gegeven.

Voor we dat kleine hoofdstuk gaan openbreken - en ontdekken hoe WEP in
de eerste generatie wifi-apparaten een navenante beveiliging aanbood -
zullen we eerst bekijken hoe gebruikers met een draadloos netwerk
effectief kunnen verbinden. Zoals reeds vermeld, gebeurt dit gebruik
makende van de onbeveiligde management frames.

\subsection{Verbinden met een netwerk}\label{verbinden-met-een-netwerk}

\begin{tipblock}
Vanaf nu zullen we geregeld het woord access point afkorten naar
\textbf{AP}. In deze tekst kan een AP zowel een eenvoudig access point
zijn als een complexe draadloze router of modem.
\end{tipblock}

Voor digitale stropers actieve aanvallen kunnen starten op een netwerk,
dienen ze verbinding te maken met het netwerk. Dit proces bestaat uit
vier stappen die doorlopen moeten worden voor legale en illegale
gebruikers effectief gebruik kunnen maken van het netwerk en
\emph{higher-layer} datatrafiek kunnen verwerken.

\begin{enumerate}
\def\labelenumi{\arabic{enumi}.}
\tightlist
\item
  \textbf{Scannen}: passief of actief zoeken naar de netwerken in de
  buurt.
\item
  \textbf{Verbinden (\emph{joining})}: kiezen met welk netwerk zal
  verbonden worden.
\item
  \textbf{Authenticatie}: bewijzen dat de gebruiker toegang heeft tot
  het netwerk.
\item
  \textbf{Associatie (\emph{association})}: bijhouden met welk AP de
  gebruiker verbonden is van het netwerk. Pas vanaf deze stap kan het
  netwerk gebruikt en misbruikt worden.
\end{enumerate}

\subsubsection{Scannen}\label{scannen}

In deze stap zal de client ontdekken welke netwerken er in de omgeving
zijn. Oftewel zoekt de client specifiek naar een netwerk, oftewel wil
hij gewoon een oplijsting van alle aanwezige netwerken. Ieder draadloze
netwerk heeft een netwerknaam, de \textbf{SSID} (\emph{service set
identifier}) die bestaat uit een 32 byte ASCII karakter string, en zal
deze op geregeld tijdstip broadcasten voor iedereen die hier nood aan
heeft. Naast het SSID zal ieder netwerk ook fysieke parameters uitsturen
die de client in de volgende stap zal nodig hebben.

Er zijn meerdere kanalen beschikbaar in het spectrum waar binnen een
netwerk mag werken. De client zal bij het scannen daarom steeds enkele
milliseconden op een bepaald kanaal luisteren om potentiële SSID
broadcasts op te vangen, om dan op het volgende kanaal hetzelfde te
doen.

\begin{tipblock}
Sommige gebruikers schakelen het broadcasten van hun netwerknaam (het
zogenaamde \textbf{SSID}) uit in de hoop dat zo dat buren,
voorbijgangers en wardrivers hun netwerk niet kunnen zien. Helaas is dit
zogenaamde \emph{snake's oil}: het geeft een vals gevoel van veiligheid.
Je kan weliswaar SSID-broadcasting uitzetten (dit is een pakketje dat je
access point elke paar seconden uitstuurt om aan iedereen die luistert
te zeggen \emph{``Hallo, hier is het netwerk met naam X, en dit zijn de
parameters die je nodig hebt om met mij te verbinden''}) maar het SSID
wordt ook in een hoop andere pakketjes de lucht in gestuurd. Een beetje
wifi hacker kan dus het SSID ogenblikkelijk uit de lucht plukken,
ongeacht dat broadcasting werd uitgeschakeld of niet.
\end{tipblock}

\subsubsection{Verbinden}\label{verbinden}

Wanneer de gebruiker gekozen heeft met welk SSID er moet verbonden
worden zal in deze stap de wifi-netwerkkaart ingesteld worden op de
juiste fysieke parameters (frequentie kanaal, snelheid, etc.) zodat als
het ware AP en client synchroon lopen (en dus op de juiste moment iets
tegen elkaar kunnen zeggen zonder elkaars signalen te storen). In de
scan-stap keken we als het ware goed rond, in deze stap richten we onze
blik nu op een bepaald netwerk.

\subsubsection{Authenticatie}\label{authenticatie}

\begin{noteblock}
Deze en de voorgaande stappen gebeuren volledig automatisch indien de
gebruiker eerder heeft geopteerd om automatisch met een netwerk te
verbinden.
\end{noteblock}

In een bedraad netwerk zit authenticatie impliciet vervat in het fysiek
aspect: wanneer je de kabel in je laptop kan steken dan is de kans groot
dat je een legale gebruiker bent in het gebouw en dus is er geen extra
authenticatie nodig (kort door de bocht gezien, weliswaar). Dat is niet
zo met draadloze netwerken die voorbije de grenzen van het gebouw hun
signaal uitsturen. Er is daarom een extra stap nodig voor we het netwerk
kunnen betreden.

In de 802.11 standaard staan twee mogelijke authenticatie-methoden
beschreven (die we verderop zullen toelichten):

\begin{itemize}
\tightlist
\item
  \textbf{Open-system authentication}
\item
  \textbf{Shared-key authentication}
\end{itemize}

Origineel was deze authenticatie een enkelrichtingstraat. Enkel de
client dient zich te authenticeren. Hierdoor bestaat dus de kans dat de
gebruiker verbindt met een rogue of fake AP (daar er geen \emph{mutual
authentication} bestaat in de originele standaard).

Naast voorgaande twee methoden is er nog een derde, die niet in de
standaard staat beschreven maar die wel door veel fabrikanten werd
aangeboden in hun wifi-apparaten:

\begin{itemize}
\tightlist
\item
  \textbf{MAC Address Authentication gebruik makend van een MAC-ACL}:
  hierbij kan de beheerder van een AP een \emph{access control list}
  (ACL) aanleggen waarin alle MAC-adressen staan van toegelaten
  apparaten.
\end{itemize}

\textbf{Open-system authentication}

Deze vorm van authenticatie werd eerst gezien als de versie om te
zeggen: ``er is geen authenticatie nodig op dit netwerk''.

Om op netwerken met deze modus te authenticeren, worden er exact twee
frames tussen client en AP uitgewisseld:

\begin{enumerate}
\def\labelenumi{\arabic{enumi}.}
\tightlist
\item
  Eerst vraagt de client aan het AP toegang.
\item
  Vervolgens stuurt het AP naar de client een ``Welkom'' frame.
\end{enumerate}

Meer gebeurt er niet in deze modus.

Als er in dit soort netwerk geen encryptie wordt toegepast dan kan dus
eender welk apparaat in de omgeving verbinding maken met dit netwerk.
Als er wél WEP encryptie wordt gebruikt, dan zal het bezit van de
WEP-sleutel als een soort toegangscontrole werken: je kan namelijk wel
authenticeren (daar het AP iedereen toelaat in deze authenticatie-fase)
maar vervolgens kan je geen data lezen en versturen tenzij je een
geldige WEP-sleutel hebt (zie verder).

Verderop zullen we ontdekken dat - oh ironie - deze
authenticatie-methode veiliger is dan de shared-key authentication die
we zo meteen gaan uitleggen.

\textbf{Shared-key authentication}

In deze modus moet je bewijzen dat je in het bezit bent van een geldige
WEP-sleutel voor dit netwerk. Het AP zal daarom een
\textbf{challenge-response} authenticatie opstarten bestaande uit
volgende sequentie van frames:

\begin{enumerate}
\def\labelenumi{\arabic{enumi}.}
\tightlist
\item
  Het AP maakt een random string aan, de \emph{challenge}, en stuurt
  deze in plaintext naar de client.
\item
  De client encrypteert deze string met z'n WEP-sleutel en stuurt dit
  terug naar het AP (de \emph{response}).
\item
  Het AP zal nu de response decrypteren met z'n eigen WEP-sleutel en het
  resultaat vergelijken met de challenge-string. Als beiden gelijk zijn
  weet het AP dat de client een geldige WEP-sleutel heeft en dus
  toegelaten mag worden op het netwerk.
\end{enumerate}

\begin{figure}
\centering
\pandocbounded{\includegraphics[keepaspectratio,alt={De shared-key authentication aan de start van de verbinding.}]{wifi/sharedkey.png}}
\caption{De shared-key authentication aan de start van de verbinding.}
\end{figure}

\begin{tipblock}
Merk op dat aanvallers die deze \emph{handshake} sniffen, twee
interessante frames zien passeren. Eerst zien ze een plaintext,
ogenblikkelijk gevolgd door de bijhorende ciphertext ervan (indien ze
een gebruiker sniffen met een geldige WEP-sleutel). Dit heet in
cryptanalyse een \emph{gekende plaintext} aanval. Dit zal interessante
informatie blijken verderop in dit horrorverhaal, waarin we zullen tonen
waarom WEP niet zo veilig bleek te zijn als gehoopt.
\end{tipblock}

\subsubsection{Associatie}\label{associatie}

Na een succesvolle authenticatie krijgt de client een \emph{association
ID} toegewezen. Dit ID gebruikt het netwerk om te weten waar in het
netwerk de client zich bevindt. Veel draadloze netwerken bestaan
namelijk uit meerdere AP's en via dit ID weet het netwerk met welk AP de
client momenteel verbonden is en zal alle data voor de client dan naar
dat AP sturen.

Vanaf dit punt kan de gebruiker dus de bronnen van het netwerk beginnen
gebruiken en wordt het tijd om deze communicatie te beveiligen (tenzij
het om een publieke hotspot gaat waar iedereen alles van elkaar kan
zien).

\subsection{WEP}\label{wep}

Om potentiële aanvallers ervan te weerhouden dat ze trafiek kunnen
sniffen (of zelf op het netwerk zetten) voorziet de 802.11 standaard
vanaf de associatie de optie om encryptie te voorzien. Dit gebeurt aan
de hand van \textbf{WEP}, wat staat voor \emph{wired equivalent
privacy}. Een naam die veel beloofde maar niet zo goed was: het idee was
dat WEP even veilig zou zijn als een bedraad netwerk.

WEP is optioneel en bevindt zich vlak voor frames naar de fysieke laag
(\emph{PHY}) worden gestuurd (die de frames ``in de lucht'' zal sturen).
Het IEEE werkt met lagen die ongeveer overeen komen met de OSI-lagen
maar met iets andere namen. In volgende figuur zie je waar WEP,
optioneel, zich bevindt.

\begin{figure}
\centering
\includegraphics[width=0.75\linewidth,height=\textheight,keepaspectratio,alt={De 802.11 stack (rechts) ten opzichte van de OSI stack.}]{wifi/osimac.png}
\caption{De 802.11 stack (rechts) ten opzichte van de OSI stack.}
\end{figure}

\subsubsection{Hoe werkt WEP?}\label{hoe-werkt-wep}

\begin{figure}
\centering
\pandocbounded{\includegraphics[keepaspectratio,alt={Een vereenvoudigde voorstelling van WEP in z'n geheel.}]{wifi/wepencr.png}}
\caption{Een vereenvoudigde voorstelling van WEP in z'n geheel.}
\end{figure}

Het hart van WEP is het RC4-algoritme dat we reeds zagen in het
crypto-hoofdstuk. De WEP-sleutel zal dienst doen als de seed voor de
keystream generatie. Deze keystream zal op zijn beurt ge-XOR'd worden
met het te encrypteren frame.

\begin{figure}
\centering
\pandocbounded{\includegraphics[keepaspectratio,alt={Integrity check.}]{wifi/crc.png}}
\caption{Integrity check.}
\end{figure}

Confidentiality en integriteit worden tegelijkertijd afgehandeld in WEP.
Vlak voor dat het frame via de XOR-operatie wordt geëncrypteerd zal het
frame eerst door een \emph{integrity check algoritme} gestuurd worden.
Deze integrity check gebeurt met behulp van CRC-32, een oude getrouwe op
dit gebied. CRC-32 zal een hash genereren die de ontvanger bij ontvangst
kan gebruiken om te zien of het frame werd aangepast na verzenden
(bewust door een aanvaller, of door bijvoorbeeld ruis in het netwerk).
Deze hash, de \textbf{Integrity Check Value (ICV}), zal mee worden
geëncrypteerd door RC4 voor hij verstuurd wordt. Op deze manier kunnen
ordinaire aanvallers het frame niet aanpassen zonder dat ze daarmee ook
de ICV ongeldig maken.

Finaal verkrijgen we dus volgende WEP-frame dat kan verstuurd worden:

\begin{figure}
\centering
\includegraphics[width=0.8\linewidth,height=\textheight,keepaspectratio,alt={WEP frame lay-out.}]{wifi/wepframe.png}
\caption{WEP frame lay-out.}
\end{figure}

\begin{noteblock}
In de originele standaard zat de mogelijkheid om tot vier WEP-sleutels
in een netwerk te gebruiken. Via de \emph{keyid} kon een client of AP
dan aangeven met welk van de geïnstalleerde sleutels een frame werd
geëncrypteerd.
\end{noteblock}

\subsubsection{Sleutellengte en de IV}\label{sleutellengte-en-de-iv}

Origineel ondersteunde WEP enkel 40-bit WEP sleutels. Ondertussen is dat
opgetrokken maar toen de standaard werd geschreven besefte men al dat er
sowieso een probleem met de sleutel zou zijn als die zo zou gebruikt
worden: ieder frame dat met dezelfde sleutel wordt geëncrypteerd zal
dezelfde keystream hebben gebruikt. Dat was natuurlijk geen optie. Om
die reden werd gekozen om te werken met een \textbf{initialisatie vector
(IV)} van 24-bit. Deze IV werd mee als seed aan RC4 gegeven. Door ieder
frame een andere IV te kiezen zorgde men er zo voor dat ieder frame een
andere keystream gebruikte (\textbf{WEP-sleutel+IV werd de nieuwe seed
per frame}). In de originele standaard werd echter niet beschreven hoe
deze IV moest veranderen, wat nefaste gevolgen zal hebben verderop.

Omdat ook de ontvanger dezelfde keystream moet kunnen genereren tijdens
decryptie is het natuurlijk belangrijk dat de IV ook bij de ontvanger
gekend is. De enige manier om dit op te lossen is door de IV mee in de
header van het frame te plaatsen en door te sturen. Uiteraard moet deze
IV als plaintext door het leven gaan.

\begin{tipblock}
Dit concept van een sleutel verlengen met een arbitrair getal heet
\emph{salting} en zullen we in het hoofdstuk omtrent wachtwoorden en
authenticatie verderop in de cursus nog zien terugkomen.
\end{tipblock}

Finaal krijgen we dus de volgende werking, encryptie en decryptie, als
volgt:

\begin{figure}
\centering
\pandocbounded{\includegraphics[keepaspectratio,alt={WEP encryptie.}]{wifi/wepfull.png}}
\caption{WEP encryptie.}
\end{figure}

\begin{figure}
\centering
\pandocbounded{\includegraphics[keepaspectratio,alt={WEP decryptie.}]{wifi/wepdec.png}}
\caption{WEP decryptie.}
\end{figure}

\section{Hoe WEP faalde}\label{hoe-wep-faalde}

Het leek een verbonden vat: hoe populairder wifi werd over de hele
wereld, hoe meer papers er verschenen die fouten identificeerden in WEP.
Al vrij snel werd duidelijk dat WEP hoegenaamd géén CIA kon aanbieden.
De meeste fouten die werden gevonden kunnen gegroepeerd worden in
volgende vier zaken:

\begin{enumerate}
\def\labelenumi{\arabic{enumi}.}
\tightlist
\item
  Het \textbf{RC4 algoritme} is helemaal niet gemaakt voor een
  datagramnetwerk zoals wifi.
\item
  De manier waarop de \textbf{IV} in de standaard is beschreven, is
  ontoereikend en verhoogt de kans op foute implementaties door
  fabrikanten.
\item
  \textbf{CRC-32} kan omzeild worden en kan dus geen integriteit van
  frames garanderen.
\item
  Er is \textbf{geen sleutel-management systeem}.
\end{enumerate}

Er is echter nog een vijfde fout die de voorgaande vier als het ware nog
versterkt:

\begin{enumerate}
\def\labelenumi{\arabic{enumi}.}
\setcounter{enumi}{4}
\tightlist
\item
  Er is \textbf{geen replay protection}.
\end{enumerate}

Hierdoor heeft de aanvaller dus vrij spel en kan hij een draadloos
netwerk als een experimenteertuin gebruiken en duizenden pakketjes te
pas en te pas onpas heruitzenden. We zullen nu de eerste vier grote
problemen beschrijven. De \emph{replay protection} behandelen we niet
apart, maar zullen we bij de andere problemen zien opduiken.

\subsection{Probleem 1: RC4}\label{probleem-1-rc4}

De meeste problemen met WEP komen van een verkeerd gebruik van het RC4
algoritme. RC4 wordt in veel moderne beveiligingsapparaten toegepast
omdat het een sterk én efficiënt algoritme is (het verbruikt weinig
energie omdat er geen dure vermenigvuldiginsoperaties in voorkomen).
Echter, streamciphers in het algemeen, RC4 specifiek, zijn eigenlijk
geen goede keuze voor datagramnetwerken waarin transmissies
onbetrouwbaar zijn.

In een datagramnetwerk worden pakketjes opnieuw verstuurd wanneer er een
fout optrad en de ontvanger om een \emph{retransmission} vraagt (dit
gebeurt bij ongeveer 20\% van de verstuurde data). Dit is volledig
normaal gedrag in zowel bedrade als draadloze netwerken, echter voor een
streamcipher is dit nefast. Specifiek twee eigenschappen van
streamciphers (en dus ook RC4) zorgen voor stevige gebreken in het
WEP-protocol inzake confidentiality:

\begin{enumerate}
\def\labelenumi{\arabic{enumi}.}
\tightlist
\item
  RC4 heeft \textbf{geen \emph{random access} mogelijkheden}.
\item
  RC4 staat \textbf{geen sleutelhergebruik} toe.
\end{enumerate}

Laten we die twee eigenschappen eens bekijken en welke cascade van
problemen ze met zich meebrengen.

\subsubsection{RC4 heeft geen random access
mogelijkheden}\label{rc4-heeft-geen-random-access-mogelijkheden}

Deze eigenschap is niet zo zeer een probleem vanuit
beveiligingsperspectief, maar wel vanuit performantieperspectief. RC4
had eigenlijk nooit gekozen mogen worden door het IEEE om in WEP
gebruikt te worden. Het verlies van één bit van de datastroom zal er
voor zorgen dat alle bits erna met RC4 ook verloren zijn, daar we met
een streamcipher werken waarbij de synchronisatie van de stroom bits
tussen verzender (encryptie) en ontvanger (decryptie) gelijk moet
blijven. Bij het minste dataverlies moeten beide zijden hun
\emph{``RC4-motortje''} resetten en opnieuw beginnen.

AES bijvoorbeeld heeft wél die random access mogelijkheid: hierdoor kan
steeds herbegonnen worden aan het punt van dataverlies en niet helemaal
opnieuw, wat natuurlijk veel efficiënter is (daar we werken in een
datagram omgeving waar bitverlies bijna continue voorkomt).

\subsubsection{RC4 staat geen sleutelhergebruik
toe}\label{rc4-staat-geen-sleutelhergebruik-toe}

Stream ciphers hebben een tweede belangrijke eigenschap: \textbf{het is
uiterst onveilig om een zelfde sleutel twee keer te gebruiken!}

Stel dat je volgende twee plaintext byte sequenties hebt:
\(p_1,p_2,p_3,…\) en \(q_1,q_2,q_3,…\). Beiden worden met dezelfde
keystream \(k_1,k_2,k_3,…\) geëncrypteerd. Dit geeft ons vervolgens
volgende twee ciphertext sequenties:

\(p_1 \oplus k_1, p_2 \oplus k_2, p_3 \oplus k_3\)

\(q_1 \oplus k_1, q_2 \oplus k_2, q_3 \oplus k_3\)

Als we nu veronderstellen dat een aanvaller deze twee ciphertexts
capteert, wat dan volgt is een grove schending van de confidentialiteit
die RC4 moet garanderen (met dank aan de wiskundige eigenschappen van de
modulo operator):

\((p_i \oplus k_i) \oplus (q_i \oplus k_i) = p_i \oplus q_i\)

Of in andere woorden, wanneer we de beide ciphertexts met elkaar XOR'n
krijgen we een sequentie die \textbf{niet afhankelijk is van de
gebruikte sleutel}! Een stevige hoeveelheid informatie over beide
plaintexts wordt zo onthuld. Als één van beide plaintexts gekend is dan
volstaat een eenvoudige XOR-operatie om ook de andere plaintext te
kennen zonder dat hierbij de gebruikte sleutel moet geweten zijn. Deze
fout zullen we verderop misbruiken.

Kortom, streamciphers zijn niet veilig in een datagram omgeving indien
er geen vorm van sleutelmanagement bestaat die de sleutels kan vervangen
voor ze herbruikt worden. WEP probeert dit gebrek aan sleutelmanagement
te omzeilen door met een IV te werken zodat er geen collisions zoals
eerder beschreven kunnen optreden\ldots maar ook dat zal een resem
problemen met zich meebrengen.

\subsection{Probleem 2: IV}\label{probleem-2-iv}

Het IEEE had dus weet van voorgaand probleem met RC4 en introduceerde
daarom de Initialisatie Vector (IV). Vanuit cryptografisch standpunt is
dit een solide oplossing. Echter, door het gebrek aan replay protection
krijgen we helaas een hoop fouten met de IV.

\subsubsection{De IV veroorzaakt weak
keys}\label{de-iv-veroorzaakt-weak-keys}

In een paper van 2001 door Scott Fluhrer, Itsik Mantin en Adi Shamir
werd aangetoond dat het key scheduling algoritme (KSA) van RC4 een hiaat
bevat:

\begin{enumerate}
\def\labelenumi{\arabic{enumi}.}
\tightlist
\item
  Wanneer een deel van de gebruikte RC4 keystream gekend is dan kan een
  groep \emph{RC4 weak keys} gevonden worden.
\item
  Wanneer deze weak keys gebruikt worden om een keystream te genereren
  dan zal er informatie van de gebruikte WEP sleutel in deze keystream
  gelekt worden.
\end{enumerate}

\textbf{Of anders gezegd: sommige IV's zorgen voor een sleutel-lekkage
naar de keystream, wat natuurlijk nefast is voor eender welk
cryptografisch cipher.}

Een gevolg van die weak keys is dat, indien de eerste 2 bytes van genoeg
keystreams (ongeveer 60) geweten is, men de gebruikte WEP sleutel kan
achterhalen door middel van een FMS aanval (de afkorting staat voor de
eerste letters van de 3 onderzoekers uit de paper).

De FMS aanval werkt indien:

\begin{enumerate}
\def\labelenumi{\arabic{enumi}.}
\tightlist
\item
  We ongeveer 60 keystreams kunnen capteren waarvan geweten is dat ze
  \emph{weak} zijn.
\item
  We de eerste 2 bytes van de plaintext van de bijhorende frames kennen
  die met deze weak keystreams zijn geëncrypteerd.
\end{enumerate}

Dat tweede is geen probleem, met dank aan de netwerkspecificaties: de
payload van een met WEP geëncrypteerd pakket bevat de LLC header (de
header van de logical link layer). Volgens de standaard (RFC 2684)
moeten \emph{``IP datagram pakketten altijd zichzelf in de header
identificeren via de SNAP header''}. En laten de eerste 2 bytes van die
header toch wel niet altijd starten met 0xAA. Kortom, quasi alle frames
die over een WEP-netwerk vliegen zullen altijd met de hexadecimale
waarde AA starten. Nu volstaat het om de bijhorende keystream te
achterhalen. En aangezien we de plaintext kennen, kunnen we ook de
eerste 2 bytes van die keystream kennen daar we weten dat:

\(c_i = k_i \oplus p_i\)

(waarbij \(p_i\) gelijk is aan 0xAA)

En dus:

\(c_i = k_i \oplus p_i \Leftrightarrow k_i = c_i \oplus p_i\)

We hebben zo ook deel één van de FMS in onze handen en kunnen nu het
algoritme de gebruikte WEP-sleutel laten berekenen (de manier waarop dat
gebeurt zou ons te ver brengen).

\begin{tipblock}
De FMS aanval is geïmplementeerd in twee populaire Linux tools: Airsnort
\& WEPCrack. Beiden kunnen dus gebruikt worden om de WEP-sleutel van een
draadloos netwerk te achterhalen.
\end{tipblock}

\subsubsection{IV collisions treden op}\label{iv-collisions-treden-op}

Op zich is de FMS aanval al dramatisch, maar helaas stopt het hier niet.
Doordat de IV maar 24 bit groot is, treden er veel sneller collisions op
dan intuïtief wordt verwacht. Van zodra twee pakketjes met dezelfde IV
zijn verstuurd, treedt er een collision op, en die zijn erg interessant
voor aanvallers. Pakketjes met dezelfde IV zijn pakketjes waarvan de
payload met dezelfde keystream werd geëncrypteerd:

Een 24-bit IV kan \(2^{24}\) oftewel 16 777 216 mogelijke waarden
hebben. Een kleine berekening toont hoe snel collisions optreden:

\emph{Gegeven: een eerste generatie AP die aan een miezerige 11 Mbps
werkt en continue 1.500-byte frames uitzendt:}

\begin{itemize}
\tightlist
\item
  \(\frac{11\;Mbps}{(1500\; bytes/pakket)* 8\;bits/byte} =  916.67\;pakketjes/seconde\)
\item
  \(\frac{16.777.216\;IVs}{916.67\;pakketjes/seconde} \approx  18302 \;seconden\)
\end{itemize}

Dat wil dus zeggen dat na ongeveer vijf uur alle IV's opgebruikt zijn en
er dan ten laatste collisions optreden.

Deze fout kunnen aanvallers op twee manieren misbruiken: met een
passieve of met een actieve aanval.

\textbf{Passieve IV aanval}

Een aanvaller kan passief meeluisteren en stilletjes alle trafiek
onderscheppen tot er een IV collision optreedt. Door twee pakketjes met
eenzelfde IV te XOR'n verkrijgt de aanvaller een pakket dat bestaat uit
de XOR van beide plaintexten van de gecapteerde pakketten. Als dus één
van beide plaintexten gekend is, is de inhoud van het andere pakket ook
gekend.

IP trafiek is vaak voorspelbaar en bevat aardig wat redundantie (om
fouten op te vangen). Hierdoor wordt het makkelijker voor een aanvaller
om via cryptanalyse te achterhalen wat de inhoud, of een deel, van het
pakket bevat. Een voorbeeld hiervan toonden we bij de FMS aanval waarbij
steeds de LLC header gekend was van de meeste pakketten.

Omdat collisions redelijk snel optreden is het voor een aanvaller dus
maar een kwestie van lang genoeg te sniffen om zo een grote hoeveelheid
pakketten met gelijke IV's op te vangen, waardoor de cryptanalyse
ongelooflijk vereenvoudigd wordt.

\begin{noteblock}
In al deze voorbeelden gaan we ervan uit dat de gebruiker geen encryptie
toepast op de hogere lagen waar z'n data vandaan komt. Uiteraard wordt
cryptanalyse een pak moeilijker als de payload van gecapteerde pakketten
geëncrypteerd blijkt te zijn.
\end{noteblock}

\textbf{Actieve IV aanval}

De passieve aanval heeft als nadeel dat we als aanvaller:

\begin{enumerate}
\def\labelenumi{\arabic{enumi}.}
\tightlist
\item
  moeten wachten op collisions, en dus bijgevolg op trafiek over het
  netwerk.
\item
  we enkel door weloverwogen gokken (cryptanalyse) kunnen proberen te
  weten te komen wat de originele plaintext juist is.
\end{enumerate}

Beide problemen kunnen we als aanvaller echter te niet doen door een
actieve rol te gaan spelen. Doordat een AP braaf alle trafiek
encrypteert dat het van het bedrade netwerk krijgt om naar een client te
sturen, is het voor een aanvaller een kwestie van ``gekende'' plaintext
van buitenuit naar het slachtoffer te sturen. Als volgt:

\begin{enumerate}
\def\labelenumi{\arabic{enumi}.}
\tightlist
\item
  Een gekende plaintext boodschap (bijvoorbeeld een e-mailbericht of
  ping) wordt naar het AP gestuurd (via het Internet bijvoorbeeld), dat
  vervolgens door de aanvaller in het oog wordt gehouden. Noot: als de
  aanvaller enkel het draadloze netwerk ter beschikking heeft (en niet
  het Internet) dan zal een bitflip-aanval moeten gebruikt worden, wat
  we verderop zullen uitleggen.
\item
  De aanvaller blijft sniffen tot hij de ciphertext ziet passeren waarin
  (vermoedelijk) z'n gestuurde plaintext zit.
\item
  Vervolgens kan de aanvaller een keystream te pakken krijgen door z'n
  plaintext te XOR'n met de gecapteerde cipherhtext:
  \(c_i = k_i \oplus p_i \Leftrightarrow k_i = c_i \oplus p_i\).
\end{enumerate}

\begin{figure}
\centering
\pandocbounded{\includegraphics[keepaspectratio,alt={Known plaintext in het wifi-netwerk krijgen.}]{wifi/inject.png}}
\caption{Known plaintext in het wifi-netwerk krijgen.}
\end{figure}

\subsubsection{Keystreams groeien}\label{keystreams-groeien}

Wanneer een aanvaller met voorgaande IV collisions keystreams kan
capteren kan hij in principe data op het netwerk beginnen plaatsen:
\textbf{aangezien het netwerk ervan uitgaat dat het gebruiken van
geldige keystreams, wil zeggen dat de gebruiker geauthenticeerd is omdat
hij de bijhorende WEP-sleutel heeft.} De aanvaller kan nu plaintext
XOR'n met deze gevonden keystream en op het netwerk zetten. Echter, hij
is beperkt tot pakketten die maximum even groot zijn als de keystream
die gevangen werd. Het zou véél nuttiger zijn als de aanvaller als het
ware een bibliotheekje heeft van geldige keystreams van allerlei
lengtes.

Omdat er geen replay protection aanwezig is, kan de aanvaller eenvoudig
z'n gecapteerde keystreams doen \emph{groeien} en zo byte per byte een
langere keystream genereren. Dit gaat als volgt te werk:

\begin{enumerate}
\def\labelenumi{\arabic{enumi}.}
\tightlist
\item
  De aanvaller maakt een plaintext pakketje aan dat 1 byte langer is dan
  de keystream die hij al heeft. Het ping-commando (ICMP) kan je met de
  ``-l'' optie bijvoorbeeld een ping van eender welke bytesize laten
  genereren. Het voordeel van het ping-commando gebruiken is ook dat we
  exact weten wat voor response er kan verwacht worden.
\item
  De aanvaller plakt nu 1 byte achter de gecapteerde keystream. Hij
  kiest hierbij een willekeurige waarde en heeft dus 1 kans op 256 om de
  juiste te kiezen.
\item
  De aanvaller XOR'd deze keystream met het commando uit stap 1 en
  stuurt dit op het netwerk.
\item
  Indien de aanvaller in stap twee de juiste byte gekozen heeft dan zal
  er een reactie op de ping volgen (daar het pakket werd gedecrypteerd
  door het AP en dan hoger in de OSI-stack door het netwerk werd
  gestuurd). Als de keystream fout is zal er geen reactie komen daar het
  AP het pakketje als foutief heeft weggegooid en dus heeft genegeerd.
\item
  Als een verkeerde byte werd gekozen in stap twee dan zal de gebruiker
  dit proces opnieuw starten en nu een andere byte-waarde kiezen. Hij
  zal dit blijven herhalen tot hij in stap vier reactie krijgt en dus
  weet dat hij nu z'n keystream met succes heeft doen groeien met 1
  byte.
\end{enumerate}

\begin{figure}
\centering
\pandocbounded{\includegraphics[keepaspectratio,alt={Keystreams byte per byte groeien.}]{wifi/grow.png}}
\caption{Keystreams byte per byte groeien.}
\end{figure}

\begin{warningblock}
We hebben bij deze aanval één belangrijk concept genegeerd waardoor deze
aanval op eerste zicht niet mogelijk is: het aanpassen van een payload
resulteert ook in een nieuwe CRC. Deze is echter mee geëncrypteerd
waardoor het niet duidelijk is hoe we dit kunnen omzeilen. Wacht nog
even tot we aan de bitflip aanval komen en alles zal duidelijk worden
(dat dit dus geen probleem is).
\end{warningblock}

\subsubsection{IV selectie}\label{iv-selectie}

De vierde fout met de Initialisatie Vectoren is de manier waarop de
selectie ervan moet gebeuren in de hardware. De 802.11 standaard gaf
enkel aan dat de IV \emph{``geregeld moest geüpdatet''} worden. Dat is
uiteraard te vaag en heeft ervoor gezorgd dat fabrikanten zelf moesten
bepalen welke IV selectie strategie ze in hun hardware zouden
implementeren. Hierdoor waren er drie strategieën die hun weg in de
verschillende apparaten vonden:

\begin{itemize}
\tightlist
\item
  \textbf{Vast IV}: sommige fabrikanten hadden geen flauw benul wat het
  doel van de IV was vanuit cryptografisch standpunt en kozen daarom
  zelfs gewoon om alle pakketten steeds met het zelfde IV te versturen.
  Hierdoor treden er dus collisions op van zodra er een tweede pakketje
  de lucht wordt ingestuurd.
\item
  \textbf{Willekeurig IV}: andere fabrikanten verkozen het om hun
  hardware bij ieder pakketje een willekeurig IV te laten selecteren.
  Alhoewel dit uiteraard veel veiliger is dan een ``vaste
  IV''-strategie, treden er toch veel sneller collisions op dan
  verwacht. Dit valt te verklaren door het zogenaamde
  \textbf{verjaardagenparadox} (zie kader verder) dat verklaart waarom
  er reeds 50\% kans op een collision is na 4823 pakketjes. Dat wil dus
  zeggen dat al na enkele seconden er meestal collisions optreden.
\item
  \textbf{Incrementele IV}: in deze strategie wordt een circulaire
  teller gebruikt waarbij de IV telkens met 1 wordt verhoogd wanneer een
  pakket moet worden verstuurd. Meestal begint deze teller op een vaste
  waarde. Dit zal er dan ook voor zorgen dat er een collisions optreedt
  van zodra een tweede apparaat zich op het netwerk begeeft en dus
  begint uit te zenden met de IV gelijk aan de IV van het allereerste
  pakketje dat het eerste apparaat gebruikte.
\end{itemize}

Kortom, een 24-bit \emph{salt} is véél te klein in een omgeving met erg
hoge data-rates zoals een draadloos netwerk. Dit probleem zou nog
beperkt kunnen worden indien de originele WEP geregeld sleutels kon
verversen, maar door het gebrek aan enig key management was dat dus uit
den boze (want herinner je: de enige reden dat we IV's nodig hadden was
omdat anders steeds dezelfde WEP-sleutel als seed werd gebruikt en dus
alle keystreams gelijk zouden zijn. Door geregeld een andere sleutel te
gebruiken zou onze kleine IV-lengte minder precair zijn, als we maar
tijdig de sleutels verversen).

\begin{noteblock}
Volgende tekst uit Wikipedia legt de \textbf{verjaardagenparadox} uit:
``De verjaardagenparadox is een paradox uit de kansrekening, die een
resultaat toont dat tegen de verwachting ingaat. Het gaat om de vraag
hoe groot de kans is dat in een groep willekeurig gekozen mensen er
(minstens) twee dezelfde verjaardag hebben. Het blijkt dat, onder enkele
lichte veronderstellingen, deze kans al meer dan 50\% is voor een groep
van maar 23 mensen. Bij 57 mensen is de kans zelfs meer dan 99\%.''

Beeld je nu in dat in plaats van mensen, je honderden pakketten hebt,
niet met een verjaardag maar met een eigen IV: de kans op collisions,
ook al is de IV 24 bit, wordt dus 50\% bij reeds een 5000 tal pakketten.
\end{noteblock}

\subsection{Probleem 3: CRC}\label{probleem-3-crc}

WEP gebruikt een \emph{integrity checksum field} om te voorkomen dat een
pakket wordt aangepast tijdens transmissie, namelijk een CRC-32
checksum. Dit was niet zo'n wijze keuze: CRCs zijn in het algemeen
vooral bedoeld om random transmissiefouten te detecteren, niet bewuste
aanpassingen. Daarnaast heeft de CRC-32 ook nefaste gevolgen in
combinatie met RC4 waardoor bitflipaanvallen mogelijk zijn.

CRC-32 is een zogenaamde lineaire functie. Zonder in detail hierover in
te gaan volstaat het te begrijpen dat we als gevolg hiervan het volgende
hebben: \emph{De CRC van een boodschap ge-XOR'd met de CRC van een
andere boodschap is hetzelfde als de CRC nemen nadat beide boodschappen
samen werden ge-XOR'd}.

Of beter gezegd:

\(CRC (boodschap_1) \oplus CRC (boodschap_2) = CRC (boodschap_1 \oplus boodschap_2)\)

Door deze eigenschap kunnen we gecontroleerd aanpassingen aan pakketten
doen, zonder daarmee de CRC te breken. Welgekome, bitflip aanval.

\subsubsection{Bitflip aanval}\label{bitflip-aanval}

\begin{itemize}
\tightlist
\item
  Om een bitflip aanval te doen heeft de aanvaller enkel \textbf{één
  geldig WEP frame} nodig. Hij hoeft zelfs niet te weten wat de inhoud
  ervan is, zolang het maar een geldig frame is.
\item
  Vervolgens maakt de aanvaller een \textbf{bitflip masker}: dit bestaat
  uit een reeks 0'n, met enkele bits op 1 (het masker). Dit zijn de bits
  die geflipt zullen worden in de volgende stap: de XOR nemen het
  bitflip masker met het payload gedeelte van het oorspronkelijke
  pakket. Welke bits geflipt worden doet er niet toe, integendeel: we
  willen net een \textbf{geldig WEP-frame maken maar mét foute data als
  payload}, zoals we zo meteen zullen zien.
\item
  Deze nieuwe payload willen we nu in een geldig frame plaatsen, en dus
  hebben we een \textbf{geldige ICV} nodig. We doen dit door de ICV van
  de nieuwe payload te berekenen en deze te XOR'n met de, geëncrypteerde
  originele ICV.
\item
  Het resultaat is een geldige, geëncrypteerde ICV voor de nieuwe
  payload. Bijgevolg hebben we een geldig frame kunnen maken dat door
  access points op het netwerk aanvaard zal worden.
\end{itemize}

\begin{figure}
\centering
\pandocbounded{\includegraphics[keepaspectratio,alt={De bitflip aanval.}]{wifi/bitflip.png}}
\caption{De bitflip aanval.}
\end{figure}

Wanneer een AP dergelijk pakket krijgt, zal het dit pakket braaf naar de
volgende laag sturen aangezien het pakket een geldig ICV heeft, ook al
bevat de inhoud \emph{rommel}. Op de volgende laag komt nu een pakket
aan dat duidelijk stuk is, en deze laag zal bijgevolg een foutboodschap
genereren en terugsturen. De aanvaller krijgt deze boodschap in een
WEP-frame aan. En alhoewel dit pakket geëncrypteerd is, weet de
aanvaller de inhoud van dit pakket (daar hij heeft opgezocht wat de
foutboodschap zal bevatten op de volgende laag volgens de standaard van
die laag) en kan hij dus een geldige keystream te pakken krijgen:

\(c = k \oplus p \Leftrightarrow k = c \oplus p\)

\begin{figure}
\centering
\pandocbounded{\includegraphics[keepaspectratio,alt={Layer 3 misbruiken door corrupte, geblitflipte pakketjes, een gekende foutboodschap te laten genereren.}]{wifi/respat.png}}
\caption{Layer 3 misbruiken door corrupte, geblitflipte pakketjes, een
gekende foutboodschap te laten genereren.}
\end{figure}

Dankzij de bitflip aanval kan de aanvaller dus zonder problemen
keystreams laten groeien zoals eerder uitgelegd.

\subsection{Probleem 4:
Sleutelmanagement}\label{probleem-4-sleutelmanagement}

Het moge duidelijk zijn: constant dezelfde WEP-sleutel gebruiken, op
meerdere apparaten, gedurende meerdere dagen, is vragen om problemen.
Werknemers die ontslagen worden kunnen een potentiële sleutel-lekkage
veroorzaken met alle gevolgen van dien. Administrators moeten manueel
nieuwe sleutels invoeren bij de werkgevers, etc. Kortom, twee
belangrijke oorzaken zorgen voor een nog lagere beveiligingsgraad van
WEP dan er al was ten gevolge van de voorbije drie problemen (IV, RC4,
CRC):

\begin{enumerate}
\def\labelenumi{\arabic{enumi}.}
\tightlist
\item
  Er is \textbf{geen geautomatiseerd sleutel verversmechanisme}:
  idealiter worden de sleutels binnen één sessie vervangen voordat de
  verzameling van mogelijke IV's is opgeraakt.
\item
  Er is \textbf{geen gecentraliseerd sleutelmanagementsysteem}.
\end{enumerate}

\begin{noteblock}
Bart Preneel van de KUL/Cosic, één van België's meest vooraanstaande
crypto-experts, zei ooit over WEP dat het het perfecte schoolvoorbeeld
is van wat er allemaal kan fout lopen wanneer je beslist om zelf een
nieuw crypto-algoritme te ontwikkelen.
\end{noteblock}

\section{Hoe WEP oplossen?}\label{hoe-wep-oplossen}

Rond 2001, nog geen twee jaar nadat de eerste wifi-standaard de wereld
``het wonder van draadloze netwerken'' bracht, werd duidelijk dat er
dringend een oplossing moest verschijnen voor de vele
veiligheidsproblemen. Wifi was alomtegenwoordig, zowel bij particulieren
als in bedrijven, en dus moest een oplossing gezocht worden voor al die
honderdduizenden bestaande apparaten die reeds in gebruik waren. Het
IEEE kon moeilijk een nieuwe standaard uitschrijven die alle bestaande
gebruikers in de kou zette. Er werd daarom besloten om twee pistes uit
te werken:

\begin{itemize}
\tightlist
\item
  \textbf{WPA1}: een (tijdelijke) oplossing voor bestaande apparaten,
  rekening houdend met enkele duidelijke beperkingen.
\item
  \textbf{WPA2}: de ``ultieme oplossing'' die een volledig nieuwe suite
  aan veiligheidsprotocollen zou bevatten voor toekomstige
  wifi-producten, maar die niet compatibel zou zijn met bestaande
  apparatuur.
\end{itemize}

\begin{tipblock}
WPA staat voor WiFi Protected Access standard.
\end{tipblock}

De beide oplossingen zouden ook de IEEE 802.1X standaard omarmen. Deze
standaard zou sleutelmanagement en gebruikersauthenticatie voorzien in
combinatie met de nieuwe encryptie en integriteits-oplossingen van WPA1
of WPA2.

Omdat sleutelmanagement én de bijhorende 802.1X infrastructuur redelijk
overkill konden zijn voor huis-tuin-en-keukengebruikers van wifi,
besloot het IEEE om twee versies van zowel WPA1 en WPA2 te maken:

\begin{itemize}
\tightlist
\item
  \textbf{Personal}: de thuisversie voor gewone gebruikers waarbij met
  een gedeelde sleutel of passphrase wordt gewerkt (een zogenaamde
  \textbf{PSK} oftewel \emph{pre-shared key}). Hierbij wordt géén 802.1X
  gebruikt.
\item
  \textbf{Enterprise}: de versie voor (grote) bedrijven waar
  sleutelmanagement belangrijk is en dus 802.1X wordt gebruikt
\end{itemize}

We vatten hier alvast de belangrijkste verschillen tussen WEP, WPA1 en
WPA2 samen.

{\def\LTcaptype{none} % do not increment counter
\begin{longtable}[]{@{}
  >{\raggedright\arraybackslash}p{(\linewidth - 6\tabcolsep) * \real{0.2500}}
  >{\raggedright\arraybackslash}p{(\linewidth - 6\tabcolsep) * \real{0.2500}}
  >{\raggedright\arraybackslash}p{(\linewidth - 6\tabcolsep) * \real{0.2500}}
  >{\raggedright\arraybackslash}p{(\linewidth - 6\tabcolsep) * \real{0.2500}}@{}}
\toprule\noalign{}
\begin{minipage}[b]{\linewidth}\raggedright
Standaard
\end{minipage} & \begin{minipage}[b]{\linewidth}\raggedright
Encryptie
\end{minipage} & \begin{minipage}[b]{\linewidth}\raggedright
Integrity
\end{minipage} & \begin{minipage}[b]{\linewidth}\raggedright
Authenticatie \& sleutelmanagement
\end{minipage} \\
\midrule\noalign{}
\endhead
\bottomrule\noalign{}
\endlastfoot
WEP & RC4 & CRC-32 & Geen \\
WPA1-Personal & TKIP & Michael & pre-shared key \\
WPA1-Enterprise & TKIP & Michael & 802.1X \\
WPA2-Personal & CCMP & CBC-MAC & pre-shared key \\
WPA2-Enterprise & CCMP & CBC-MAC & 802.1X \\
\end{longtable}
}

\begin{noteblock}
Ondertussen bestaat er ook WPA3, die we op het einde van dit hoofdstuk
zullen bespreken. WPA1 en WPA2 zijn echter , historisch gezien,
gerelateerd aan WEP en daarom bekijken we deze beiden samen in functie
van WEP.
\end{noteblock}

\section{WPA1}\label{wpa1}

WPA1 (de tussentijdse oplossing) werd ontwikkeld waarbij rekening werd
gehouden met volgende beperkingen:

\begin{itemize}
\tightlist
\item
  Zoals verteld, miljoenen WEP-gebaseerde apparaten waren reeds in
  gebruik. Deze apparaten zouden met behulp van een firmware upgrade
  gepatcht moeten kunnen worden naar WPA1.
\item
  De meeste AP's werkten met processoren die reeds quasi volcontinue
  tegen hun maximum capaciteit werkten. De extra algoritmes (van WPA1)
  die het AP moest draaien mochten maar een beperkte overhead creëren.
\item
  De RC4 encryptie is deels \emph{hardwired} in de hardware van het AP.
  Hierdoor kunnen bepaalde delen van WEP onmogelijk `omzeild' worden en
  hangen we dus inherent vast aan WEP.
\end{itemize}

\subsection{802.1X}\label{x}

De IEEE 802.1X standaard is ondertussen een veelgebruikte standaard in
veel draadloze én bedrade netwerken. Het voorziet volgende zaken aan
WPA1, WPA2 en WPA3 netwerken die in \emph{Enterprise-mode} werken:

\begin{itemize}
\tightlist
\item
  (Mutual) \textbf{authentication}.
\item
  Een \textbf{centraal user management} systeem.
\item
  Een \textbf{veilige manier om geheime sleutels uit te wisselen}.
\end{itemize}

We gaan de volledige werking van de 802.1X standaard hier niet uit de
doeken doen, dat zou ons te ver brengen. Het is echter nuttig om te
begrijpen dat deze standaard werd gekozen omdat hij ervoor zorgt dat de
AP's niet meer zelf de authenticatie moeten doen, maar dat ze gebruik
maken van de bestaande (bedrade) authenticatie-infrastructuur van het
bedrijf. De AP's zullen gewoon als een doorgeefluik aan de start tussen
gebruiker en de authenticatie-server optreden en de boodschappen tussen
beiden uitwisselen. Enkel wanneer het AP toestemming krijgt van de
authenticatie-server (meestal een RADIUS server) zal het AP de
eindgebruiker toegang tot het draadloze netwerk verschaffen.

\begin{figure}
\centering
\includegraphics[width=0.8\linewidth,height=\textheight,keepaspectratio,alt={Port-based authenticatie met 802.1X.}]{wifi/port.png}
\caption{Port-based authenticatie met 802.1X.}
\end{figure}

802.1X zelf beschrijft niet hoe de authenticatie moet plaatsvinden: het
is geen algoritme. Integendeel: het is een \textbf{framework} waar
binnen andere algoritmen en standaarden, op maat van het bedrijf, kunnen
ingeplugd worden. Hierdoor ontstaat een flexibel concept dat bedrijven
(of diehard eindgebruikers) niet verplicht om een bepaalde manier van
authenticatie (en bijhorende soft- en hardware) te omarmen. 802.1X zorgt
voor de vertaling van de authenticatie-boodschappen tussen enerzijds het
netwerkprotocol (bv. Ethernet, wifi, maar ook Token Ring, etc.) en de
\emph{methode laag}. De methode-laag bevat het te gebruiken
authenticatie-protocol en dient \textbf{EAP}-compatibel zijn.

\begin{figure}
\centering
\includegraphics[width=0.5\linewidth,height=\textheight,keepaspectratio,alt={De modulariteit van het 802.1X framework.}]{wifi/8021x.png}
\caption{De modulariteit van het 802.1X framework.}
\end{figure}

EAP oftewel \emph{Extensible Authentication Protocol} is, zoals de naam
doet vermoeden, een uitbreidbaar protocol van te gebruiken
authenticatie-methoden. Afhankelijk van de keuze van het bedrijf kan
voor een bepaald EAP-protocol gekozen worden, het ene is
gebruiksvriendelijker en/of veiliger dan het andere. Uiteraard dienen
zowel de client als de netwerkinfrastructuur compatibel te zijn met de
gekozen EAP-methoden van het netwerk. Koop je dus een AP dat
WPA2-Enterprise compatibel is, moet je nog steeds controleren of het AP
compatibel is met de gekozen EAP-methode van het bedrijf.

De meest gebruikte EAP-methoden zijn:

\begin{itemize}
\tightlist
\item
  \textbf{EAP-TLS}: gebruikt een TLS tunnel om op een beveiligde manier
  te communiceren (we zagen TLS ook reeds aan het einde van crypto waar
  het gebruikt werd om HTTPS-trafiek te beveiligen). Hierbij gebeurt een
  certificaat-gebaseerde authenticatie.
\item
  \textbf{EAP-TTLS} (\emph{Tunneled TLS}): omdat niet alle
  eindgebruikers zich kunnen authenticeren aan de hand van een
  certificaat, voorziet TTLS authenticatie met behulp van een
  username/wachtwoord login. Hierbij wordt wel nog steeds een TLS tunnel
  gebruikt voor veilige communicatie, maar de gebruiker moet geen eigen
  certificaat bezitten. Ter info: EAP-TTLS is quasi hetzelfde als
  \emph{Protected EAP} (PEAP), een ander EAP-protocol dat je soms zal
  zien passeren.
\end{itemize}

\begin{figure}
\centering
\includegraphics[width=0.75\linewidth,height=\textheight,keepaspectratio,alt={Het authenticatie-proces op een wifi-netwerk met 802.1X (Enterprise-mode).}]{wifi/8021X2.png}
\caption{Het authenticatie-proces op een wifi-netwerk met 802.1X
(Enterprise-mode).}
\end{figure}

De voorgaande figuur toont de typische uitwisseling van boodschappen die
plaatsvinden wanneer een client voor het eerst verbinding wil maken op
een WPA1 of WPA2 Enterprise netwerk:

\begin{enumerate}
\def\labelenumi{\arabic{enumi}.}
\tightlist
\item
  Client en access point zoeken samen welke EAP-methoden zij beiden
  kunnen gebruiken om de authenticatie te starten. Enkel indien ze samen
  tot één keuze kunnen komen wordt overgegaan naar stap 2.
\item
  Vanaf nu zal het AP enkel dienst doen als doorgeefluik tussen de
  client en de authenticatie-server.
\item
  Als de server genoeg bewijs heeft om de client te authenticeren, zal
  de server de nodige sleutels genereren en deze aan het AP geven.
\item
  Het AP zal vanaf nu instaan voor het verdere sleutelbeheer en wanneer
  nodig de sleutels verversen tijdens de sessie.
\item
  Wanneer de sleutels tussen client en AP zijn verwerkt, krijgt de
  client toegang tot het netwerk.
\end{enumerate}

In wifi-netwerken met 802.1X worden twee sets van sleutels aangemaakt:

\begin{itemize}
\tightlist
\item
  \textbf{Sessiesleutels}, ook wel ``\emph{pairwise keys}'' genoemd:
  deze zijn uniek per client en zijn enkel gekend door die client en het
  AP. Zoals de naam doet vermoeden zijn deze sleutels enkel geldig
  tijdens de net opgezette sessie.
\item
  \textbf{Groepssleutels}, ook wel ``\emph{group keys}'' genoemd: deze
  worden tussen alle cliënten van hetzelfde AP gedeeld en worden
  gebruikt voor multicast trafiek.
\end{itemize}

Indien het \emph{dynamic key exhange protocol} is geconfigureerd dan zal
de authenticatie de sessiesleutels eenmalig aanmaken en doorsturen.
Daarna zullen specifieke encryptie-sleutels gegenereerd worden bij de
client en AP gebaseerd op deze sessiesleutel. De client (en AP) kan dan
zelf, automatisch, op gepaste momenten nieuwe encryptie-sleutels
genereren.

\subsection{TKIP}\label{tkip}

Bij WPA1, wanneer de nodige sleutels zijn uitgewisseld, is het tijd om
data te encrypteren. Dit gebeurt door middel van het \textbf{Temporal
Key Integrity Protocol} (TKIP), wat een suite van algoritmes is dat als
een wrapper rond WEP wordt gelegd om zo encryptie aan te bieden en
daarbij de fouten van WEP minimaliseert.

TKIP omhult WEP met volgende zaken:

\begin{enumerate}
\def\labelenumi{\arabic{enumi}.}
\tightlist
\item
  Een nieuwe \textbf{integrity check genaamd Michael} dat een
  \emph{message integrity code} (MIC) genereert die wél bestand is tegen
  bitflip aanvallen.
\item
  Een nieuwe manier van \textbf{IV selectie} die replay aanvallen
  voorkomt.
\item
  Een \textbf{per-pakket sleutel mixing} algoritme dat het probleem met
  \emph{weak keys} in RC4 oplost.
\item
  Een \emph{re-keying} mechanisme dat ongeveer elke 10000 pakketjes een
  nieuwe sleutel doet genereren.
\end{enumerate}

\subsubsection{Michael the Mic}\label{michael-the-mic}

Om bewuste aanpassingen aan de payload van een frame te detecteren
gebruikt WPA1 een \emph{message integrity check} (\textbf{MIC}). Dit is
een cryptografisch sterker concept dan de originele CRC checks en wordt
verzorgd door het ``Michael'' algoritme.

\begin{noteblock}
In de literatuur wordt meestal gesproken over ``Message authentication
codes'' of MAC's. Echter, in de IEEE 802 standaarden wordt MAC reeds
gebruikt voor \emph{media access control} en werd er dus gekozen voor
MIC.
\end{noteblock}

``Michael'' berekent de MIC van een payload maar gebruikt hierbij ook de
Michael sleutel (een afgeleide van de authenticatie-sleutel), het adres
van de verzender én ontvanger. Hierdoor wordt het voor een aanvaller
veel moeilijker om een dergelijke MIC na te bootsen, laat staan te
\emph{replayen} (vergelijk dit met de originele CRC-32 die enkel de
payload gebruikt om de checksum te berekenen).

\begin{figure}
\centering
\includegraphics[width=0.6\linewidth,height=\textheight,keepaspectratio,alt={Het aanmaken van een MIC met Michael.}]{wifi/michael.png}
\caption{Het aanmaken van een MIC met Michael.}
\end{figure}

Wanneer TKIP twee foute MICs na elkaar detecteert, gaat het er van uit
dat er een aanval bezig is. Volgende stappen worden dan ogenblikkelijk
uitgevoerd door de client:

\begin{enumerate}
\def\labelenumi{\arabic{enumi}.}
\tightlist
\item
  Alle huidige sleutels worden verwijderd.
\item
  De client verbreekt de verbinding.
\item
  Er wordt één minuut gewacht voor een nieuwe verbinding
  (\emph{association}) wordt opgezet.
\end{enumerate}

\subsubsection{IV selectie verbetering}\label{iv-selectie-verbetering}

Om te voorkomen dat fabrikanten weer naïeve oplossingen voor de IV
selectie implementeerden, legde WPA1 nu de regels op. Een ontvangen
pakket zal pas aanvaard worden indien de IV van het pakket op de IV van
het vorige pakket volgt. Uiteraard zit er een kleine marge om
hertransmissies toe te staan, maar een pakket met bijvoorbeeld IV 1110
zal nooit aanvaard worden als het AP vlak ervoor een pakket met IV 3789
heeft aangekregen.

Daarnaast wordt ook de IV lengte gevoelig vergroot. TKIP hanteert
namelijk een 48-bit IV, genaamd de \emph{TKIP sequence counter} (TSC).
Voor de WEP-encryptie worden de drie minst significante bytes van de TSC
gebruikt als IV, waardoor deze compatibel blijft met de oorspronkelijke
WEP-specificatie. Door de resterende bits van de TSC te benutten voor
replay-bescherming en sleutelmanagement, vermindert TKIP significant het
risico op IV-collisions, een grote zwakte van WEP.

\subsubsection{Key mixing en re-keying}\label{key-mixing-en-re-keying}

Om weak keys te voorkomen gebruikt een TKIP een \emph{key mixing
function} dat zal resulteren in een \textbf{temporal of per-pakket
sleutel} die dienst zal doen als vervanger voor de WEP-sleutel. Deze
sleutel zal geregeld ververst worden (vandaar \emph{temporal} oftewel
tijdelijk) en heeft dus een beperkte levensduur voor actieve aanvallen.

Het mixen van de sleutel gebeurt in twee fases, waarbij iedere fase een
specifieke zwakte van WEP indijkt:

\begin{itemize}
\tightlist
\item
  Fase 1: zorgt ervoor dat alle clients een eigen sleutel hebben doordat
  het verzender adres (transmitter address (TA)). wordt toegevoegd aan
  de temporal key.
\item
  Fase 2: zorgt voor een `per-pakket' sleutel waarbij kennis van de IV
  niet meer door de aanvallers kan misbruikt worden.
\end{itemize}

\begin{figure}
\centering
\pandocbounded{\includegraphics[keepaspectratio,alt={Het mengen van de verschillende sleutels naar een sleutel die ieder pakketje verandert.}]{wifi/mixing.png}}
\caption{Het mengen van de verschillende sleutels naar een sleutel die
ieder pakketje verandert.}
\end{figure}

\textbf{Fase 1 mix}

Het \textbf{transmitter adres} (het MAC-adres van het apparaat dat het
pakket uitzendt) wordt gecombineerd met de \textbf{Temporal Key}. Deze
sleutel is ofwel afgeleid van een PSK-sleutel in de Personal-modus,
ofwel afgeleid van een PMK-sleutel die tijdens de authenticatie via
802.1X is verkregen. Daarnaast worden de vier meest significante bytes
van de TKIP Sequence Counter (TSC) toegevoegd. Het combineren van deze
inputs gebeurt via een iteratief proces dat hashing-achtige technieken
toepast, waaronder XOR-operaties, modulaire optellingen en
bitverschuivingen. Het resultaat is een intermediate sleutel die dient
als input voor de volgende fase van het sleutelbeheerproces.

\begin{tipblock}
Bij Michael spraken we over een destion adres (bron adres) in de
afbeelding. Terwijl we nu over een transmitter adres spreken. Dit is
geen fout, maar een bewuste keuze:

\begin{itemize}
\tightlist
\item
  TA (Transmitter Address): Verwijst naar het apparaat dat het pakket op
  dat moment uitzendt.
\item
  SA (Source Address): Verwijst naar het apparaat dat het pakket
  oorspronkelijk heeft gemaakt.
\end{itemize}

\end{tipblock}

\textbf{Fase 2 mixing}

De intermediate sleutel, verkregen uit de Phase 1 Key Mixing, wordt
gecombineerd met de twee minst significante bytes van de TSC en de
Temporal Key. Dit proces is ontworpen om de tijdelijke sleutel verder te
versterken en een unieke encryptiesleutel te genereren voor elk
datapakket. De combinatie van deze inputs wordt uitgevoerd via een
iteratief proces dat gebruikmaakt van effectieve cryptografische
technieken, zoals XOR-operaties, modulaire optellingen en
bitverschuivingen. Het resultaat van deze fase is de definitieve
RC4-sleutel, ook wel de WEP-seed genoemd. Deze sleutel wordt samen met
de drie minst significante bytes van de TSC (die de verbeterde
Initialization Vector vormen) ingevoerd in het RC4-algoritme om de
keystream te genereren die nodig is voor encryptie in WEP.

\subsection{Alle blokjes samen}\label{alle-blokjes-samen}

Finaal kunnen we vervolgens WPA1 visualiseren, waarbij duidelijk is dat
we vooral een wrapper rond WEP hebben verkregen, maar dat WEP nog steeds
het hart van het systeem is.

\begin{figure}
\centering
\pandocbounded{\includegraphics[keepaspectratio,alt={WPA1 in volle glorie.}]{wifi/wpa1.png}}
\caption{WPA1 in volle glorie.}
\end{figure}

\begin{noteblock}
Het fragmenteren met de TSC is noodzakelijk vanwege de volgende redenen:

\begin{itemize}
\tightlist
\item
  Wanneer een groot pakket wordt opgesplitst in meerdere fragmenten,
  moet elk fragment uniek zijn.
\item
  De TSC fungeert als een unieke identificatie voor elk fragment omdat
  het incrementeert bij elke transmissie.
\item
  Dit voorkomt verwarring of overlap tussen pakketten en maakt het
  mogelijk voor de ontvanger om de fragmenten correct terug samen te
  stellen.
\end{itemize}

\end{noteblock}

\subsection{Are we there yet?!}\label{are-we-there-yet}

In 2009 verschenen er al enkele exploits die WPA1-Personal misbruikten
waardoor aanvallers de WPA passphrase (de PSK) konden achterhalen door
de handshake aan de start van een sessie te capteren. Deze aanvallen
waren echter niet zo efficiënt als de WEP-aanvallen en vereisten een
aanzienlijke hoeveelheid data om de PSK te achterhalen. Het was echter
duidelijk dat WPA1-Personal niet de ultieme oplossing was voor het
WEP-probleem.

\begin{tipblock}
Bekijk coWPAtty en Aircrack om WPA1-Personal te \emph{hacken}.
\end{tipblock}

\section{WPA2}\label{wpa2}

De finale oplossing voor het WEP-veiligheidsprobleem werd gegoten in
WPA2 waarvan AES het hart vormt. Toen de standaard werd ontwikkeld was
AES zonet verschenen en het bleek de ideale toepassing voor WPA2.
Hierbij werden twee onderdelen van AES gebruikt namelijk:

\begin{enumerate}
\def\labelenumi{\arabic{enumi}.}
\tightlist
\item
  AES met \emph{counter mode} voor de encryptie.
\item
  AES in CBC-MAC mode voor integriteit en authentication.
\end{enumerate}

Samen wordt dit encryptieproces \textbf{CCMP} genoemd, wat staat voor
\emph{``Counter Mode Cipher Block Chaining Message Authentication Code
Protocol''}

Voor de sleuteldistributie en management werd, zoals vermeld, gebruik
gemaakt van de 802.1X standaard.

\begin{tipblock}
Toen WPA2 uitkwam ontdekte men dat toch aardig wat bestaande hardware
kon gepatcht worden om ook WPA2-compatibel te zijn. Dit was een opsteker
voor velen, daar de originele premise van WPA2 net was dat ze geen
rekening zouden houden met de bestaande hardware die reeds op de markt
was.
\end{tipblock}

\subsection{CCMP}\label{ccmp}

CCMP gebruikt, net als TKIP, een 48-bit IV die \emph{packet number} (PN)
werd genoemd. Deze PN wordt, samen met andere informatie, gebruikt om de
AES encryptie van een seed te voorzien. Hierbij wordt het frame (en de
header) in blokken van 128 bit verdeeld en zo blok per blok verwerkt. De
encryptie gebeurt parallel met het berekenen van de MIC (de ICV in WEP)
die finaal achteraan als een laatste blok ook mee wordt geëncrypteerd.
Ook nu wordt met een tijdelijke sleutel gewerkt die gebaseerd is op de
hoofdsleutel verkregen via 802.1X (enterprise mode) of de passphrase (in
personal mode).

\begin{warningblock}
Die laatste zin impliceert een ander belangrijk veiligheidsverschil
tussen enterprise en personal mode: in personal modus wordt een
passphrase sleutel gedeeld met alle gebruikers op het netwerk. Hierdoor
kan iemand die in het bezit is van deze passphrase ook de data
decrypteren van de andere gebruikers, iets wat onmogelijk is in
enterprise modus.
\end{warningblock}

\begin{figure}
\centering
\pandocbounded{\includegraphics[keepaspectratio,alt={CCMP: confidentiality en integrity in één.}]{wifi/ccmp.png}}
\caption{CCMP: confidentiality en integrity in één.}
\end{figure}

In de figuur zie je duidelijk de elegantie van het systeem en hoe de
integriteit in parallel wordt berekend met de encryptie. De MIC wordt
berekend met behulp van CBC-MAC, wat staat voor \textbf{cipherblock
chaining - message authentication code}, een veelzeggende naam. Zoals we
uit het hoofdstuk crypto weten, zal er bij CBC een keten van encrypties
starten, waarbij de encryptie van het huidige blok ge-XOR'd wordt met
het resultaat van de encryptie van het vorige blok. Het finale
geëncrypteerde cipherblock zal daarbij dienst doen als MIC. Bij de
minste bitfout ergens in de payload of header zal een totaal andere MIC
gegenereerd worden.

Om de payload te encrypteren (\emph{merk op dat de header niét
geëncrypteerd wordt, het wordt enkel mee betrokken om de MIC te
berekenen}) wordt de \textbf{counter mode} van AES gebruikt. Hierbij
worden de blokken onafhankelijk van de andere blokken geëncrypteerd en
zal een \emph{payload counter} (PC) telkens met één verhoogd worden en
als extra seed voor de AES-motor van het huidige blok dienen (samen met
de te gebruiken sleutel).

\subsection{Helaas, aan alles komt een
einde}\label{helaas-aan-alles-komt-een-einde}

Tot 2017 ging alles goed. De AES standaard was al jaren een robuuste
standaard gebleken en werd op vele plekken nog steeds gebruikt. En dit
zou ook bij wifi zou zijn geweest, waren het niet dat in mei 2017 een
Belgische onderzoeker, Mathy Vanhoef, de bevindingen van z'n onderzoek
publiceerde. Hij had helaas een belangrijke fout gevonden in de WPA2
standaard. Deze had niets te maken met AES - dat blijft een stevige
standaard zijn - maar wel de manier waarop een bepaalde uitwisseling van
bepaalde berichten tijdens de initiële handshake gebeuren. Deze
aanvallen worden beschreven én gedemonstreerd op
\href{HTTPS://www.krackattacks.com/}{krackattacks.com} en verplichtten
de IEEE om te beginnen werken aan een opvolger voor WPA2.

\begin{tipblock}
Een lek zoals krackattack vereist natuurlijk een snelle reactie van de
vendors. Hoe sneller zij een patch uitbrengen, hoe sneller het lek kan
gedicht worden. Maar wat als je een wifi-kaart hebt die al vele jaren
oud is en waarvan de fabrikant misschien niet meer bestaat? Dit probleem
zien we geregeld opduiken en wordt nog groter wanneer we beseffen dat
ook de interne elektronica (de chips) soms gepatcht moeten worden. In
het hoofdstuk rond IoT Security gaan we hier dieper op in.
\end{tipblock}

\begin{tipblock}
De personal mode van WPA1 en WPA2 zal altijd gevoelig zijn voor
dictionary aanvallen. Aangezien de gedeelde sleutel (de pre-shared key)
een gedeeld geheim is, kunnen andere gebruikers proberen dit te
achterhalen. Het is dus belangrijk dat, indien je in personal modus WPA
hanteert, je een sleutel kiest die voldoet aan de typische vereisten van
een goed wachtwoord.
\end{tipblock}

\section{WPA 3 (``Wifi 6'')}\label{wpa-3-wifi-6}

In 2018 kwam de Wifi Alliance uit met de opvolger van WPA2, de titel, je
raadt het nooit, was uiteraard WPA3, maar werd ook wel Wifi 6 genoemd.
De standaard gaan we hier niet zo gedetailleerd bespreken, het volstaat
te begrijpen dat ze

\begin{itemize}
\tightlist
\item
  op veiligheidsniveau state-of-the-art was.
\item
  rekening hield met de noden van 21e-eeuwse draadloze netwerken (denk
  maar aan Internet-of-Things apparaten, beveiligde publieke hotspots,
  etc.).
\end{itemize}

\begin{tipblock}
De Wifi Alliance is in het leven geroepen als een organisatie die ervoor
zorgde dat producten wel officieel konden claimen dat hun producten
conform een IEEE standaard waren. Enkel wanneer hun producten de nodige
tests van de Wifi Alliance aflegden kon het product het label van ``Wifi
Alliance compatibel'' product dragen.
\end{tipblock}

Ook in WPA3 werden twee modes voorzien: een personal en een enterprise
mode. Enkele van de interessantste verbeteringen zijn:

\begin{itemize}
\tightlist
\item
  \emph{Simultaneous Authentication of Equals (SAE)}: een nieuw
  cryptografisch concept waarbij in de personal mode authenticatie veel
  veiliger kan plaatsvinden dan voorheen.
\item
  Resistent tegen offline dictionary attacks.
\item
  \emph{Forward secrecy}: zelfs als de aanvaller de wifi-sleutel van
  oude gecapteerde pakketten vindt zal hij deze toch niet kunnen
  decrypteren. Het aloude ``store now, decrypt later'' is dus niet van
  toepassing op WPA3.
\item
  \emph{Wifi easy connect}: een gebruiksvriendelijke manier om
  Internet-of-Things apparaten met het netwerk te verbinden.
\item
  \emph{Wifi enhanced open}: publieke hotspots blijven publiek, maar
  iedere client heeft z'n eigen veilige kanaal met het AP. Gedaan zijn
  de dagen van je in de Starbucks zetten om zo privé-trafiek van
  omstaanders te sniffen.
\item
  \emph{Geauthenticeerde encryptie} gebruik makend van \emph{``256-bit
  Galois/Counter Mode Protocol (GCMP-256)''} een cryptocipher dat we
  hier niet uit de doeken gaan doen (maar geef toe, met zo'n naam klinkt
  het toch ogenblikkelijk extra veilig!).
\item
  Gebruikt de meest moderne veilige authenticatie- en
  sleuteldistributiemethoden mogelijk binnen 802.1X (HMAC, HMAC-SHA384
  en ECDH)
\end{itemize}

\chapter{Authenticatie}\label{authenticatie-1}

Bewijzen wie je bent om toegang te krijgen tot een website of applicatie
heet \textbf{authenticatie}. Vervolgens, afhankelijk van wie je bent,
zal je bepaalde rechten toegewezen krijgen die bepalen wat je wel en
niet kunt doen op de website of applicatie, dit heet
\textbf{autorisatie}. In dit hoofdstuk gaan we ons toespitsen op het
eerste deel van dit proces: de authenticatie. Hierbij gaan we vooral
kijken hoe we als web- of applicatiebeheerders op een veilige manier
moeten omgaan met de login-informatie van gebruikers.

We weten al dat je geheime sleutel een belangrijk onderdeel is in heel
veel aspecten van cybersecurity. Het \textbf{weten van de geheime
sleutel} is een eerste vorm van authenticatie (maar uiteraard niet de
beste). Je login gegevens die je gebruikt om toegang te krijgen tot een
website bestaan in primaire vorm meestal uit een combinatie van
gebruikersnaam en wachtwoord. Het wachtwoord in dit verhaal is kortom
gewoon een ander woord voor je geheime sleutel. Als beheerder is het dan
ook essentieel dat we uitermate veilig omgaan met de wachtwoorden van
gebruiker. We zouden niet willen dat alle login gegevens van onze
gebruikers in verkeerde handen vallen.

\section{Veilige wachtwoorden}\label{veilige-wachtwoorden}

Ongeacht de veiligheden die we inbouwen als cyberboswachter, veel blijft
afhangen van de manier waarop eindgebruikers omgaan met hun
wachtwoorden. Volgende regels worden continu niét gehanteerd, met alle
gevolgen van dien:

\begin{itemize}
\tightlist
\item
  Hergebruik nooit een wachtwoord. In principe heb je één wachtwoord pér
  service (applicatie, website, etc.).
\item
  Gebruik geen wachtwoorden die in \emph{dictionaries} staan. Maar
  overweeg volledig random gegenereerde wachtwoorden.
\item
  Zorg ervoor dat je wachtwoorden lang genoeg zijn (minimum 16 tekens).
\item
  Zorg ervoor dat je wachtwoorden steeds een combinatie van cijfers,
  letters (grote én kleine) en leestekens zijn.
\end{itemize}

Trouwens, herinner je je de McCumber kubus waarin we benadrukten dat
technologie maar één aspect is om CIA toe te passen op je data in z'n
drie primaire vormen? Het zal je niet verbazen dat cybercriminelen niet
altijd gaan proberen databanken aan te vallen om wachtwoorden van
gebruikers te pakken te krijgen. Als zij een specifiek doelwit hebben
dan gaan ze vaak op andere manieren te werk:

\begin{itemize}
\tightlist
\item
  \textbf{Password spraying}: hierbij gaat de hacker een (beperkte)
  lijst van veelgebruikte wachtwoorden testen op een grote groep
  \emph{useraccounts} van een bepaalde website, in de hoop een
  \emph{hit} te hebben (\emph{``spray and pray''}).
\item
  \textbf{(Spear) phishing}: bij phishing hanteert de aanvaller de
  goedgelovigheid of onoplettendheid van de gebruiker om een
  ogenschijnlijk betrouwbare mail of bericht te sturen met daarin een
  link naar een pagina die malware installeert of een fake login scherm
  toont. Bij spear phishing gebruikt de aanvaller geen massmail, maar
  gaat hij juist gericht één specifiek doelwit een op maat gemaakte mail
  of bericht sturen. Spear phishing is heden ten dage één van dé
  \textbf{social engineering} aanvallen bij uitstek.
\item
  \textbf{Keyloggers}: als de aanvaller toegang heeft tot de computer
  (wat uiteraard van over het netwerk kan) dan kan hij een (permanente)
  keylogger installeren die alle toetsaanslagen op het systeem opneemt.
  Nadien kan de aanvaller deze logs dan analyseren in de hoop zo ook het
  wachtwoord of andere gevoelige informatie terug te vinden.
\end{itemize}

\section{Hoe wachtwoorden opslaan}\label{hoe-wachtwoorden-opslaan}

Hoe moet je nu als cyberboswachter de login gegevens van je gebruikers
bewaren? We gaan een soort \emph{bottom-up} aanpak hanteren, waarbij we
beginnen met de meest naïeve oplossing en telkens verbeteringen zullen
aanbrengen.

\subsection{Paswoorden als plaintext}\label{paswoorden-als-plaintext}

In het prille begin van het Internet gebeurde dit quasi overal: de
login-databank had twee kolommen:

\begin{enumerate}
\def\labelenumi{\arabic{enumi}.}
\tightlist
\item
  gebruikersnaam.
\item
  gebruikerswachtwoord.
\end{enumerate}

De wachtwoorden in kolom 2 stonden er zoals ze waren. Als een gebruiker
wilde inloggen op dit soort websites dan moest hij z'n wachtwoord
verzenden en dan ging de \emph{backend} controleren of het ingezonden
wachtwoord overeen kwam met het wachtwoord in de database. Het spreekt
voor zich dat dit soort databanken van gigantische waarde zijn voor
aanvallers: van zodra ze de databank hebben te pakken hebben ze alle
wachtwoorden van alle gebruikers! Profit!

\textbf{Paswoorden mogen nooit in onbeveiligde, leesbare vorm in een
databank staan!} Wanneer dit wel zo is dan kan je beter ogenblikkelijk
je account bij die service deleten. Want alhoewel deze aanpak al lang
bestaat en er al bijna even lang van geweten is dat deze erg onveilig
is, toch zijn er nog steeds ontelbare websites en applicaties die
hieraan zondigen. Als ze dus jouw wachtwoord zo behandelen, dan is de
kans reëel dat ook hun andere veiligheidsdiensten niet om over naar huis
te schrijven zijn.

Een goede manier om te weten of een service op deze manier werkt is
gebruik maken van de \emph{``Ik ben m'n wachtwoord vergeten''}-knop. Als
je deze knop gebruikt en je krijgt een e-mail met daarin jouw originele
wachtwoord, dan kan je er zeker van zijn dat de service jouw wachtwoord
op deze manier bewaart. In principe zou een service NOOIT jouw
wachtwoord moeten kunnen zien. We gaan zelfs zien dat \textbf{jouw
wachtwoord nooit je computer mag verlaten}, laat staan dat deze
beschikbaar is als plaintext in een database.

\subsection{Paswoord hashing}\label{paswoord-hashing}

Door een wachtwoord te hashen kunnen we de wachtwoorden al iets veiliger
bewaren. Een gebruiker die zich wenst aan te melden bij een systeem dat
met wachtwoord hashes werkt zal nu op zijn lokale systeem z'n hash
moeten genereren en dit over het netwerk doorsturen. De service zal deze
ontvangen hash vergelijken met de waarde die in de database staat, en
indien deze gelijk is dan wordt verondersteld dat de gebruiker het
juiste wachtwoord kende.

We versturen dus niet meer het wachtwoord over het netwerk, maar we zijn
nu wel vatbaar voor een \textbf{pass-the-hash} aanval. Het volstaat om
een geldige combinatie van gebruikersnaam en hash te capteren en deze
vervolgens te gebruiken om ergens in te loggen. De aanvaller heeft
hierbij geen kennis nodig van het originele wachtwoord.

\begin{tipblock}
Door een secure hash van een wachtwoord te genereren creëren we een stuk
tekst dat niet terug naar het originele wachtwoord kan omgezet worden.
In theorie zal ieder wachtwoord een andere hash creëren. Uiteraard
kunnen er toch twee zaken zich voordoen:

\begin{enumerate}
\def\labelenumi{\arabic{enumi}.}
\tightlist
\item
  Twee totaal verschillende wachtwoorden genereren dezelfde hash, een
  zogenaamde \emph{collision}.
\item
  Twee mensen kiezen hetzelfde wachtwoord en zullen dus ook dezelfde
  hash genereren.
\end{enumerate}

Dit probleem gaan we verderop oplossen met behulp van \emph{salting}.
\end{tipblock}

\begin{warningblock}
Veel websites genereren wel degelijk de hash aan serverzijde. Het
verschil hier is echter dat ze eerst een beveiligde TLS tunnel hebben
opgezet waarover het wachtwoord werd verstuurd. Het laat echter de
website/service toe om meer controle te hebben over de hash-generatie.
\end{warningblock}

\subsection{Rainbow table attack}\label{rainbow-table-attack}

Als de database door aanvallers gestolen wordt dan zitten we ook een
tikkeltje veiliger als voorheen (de pass-the-hash aanval zal uiteraard
nu zeker werken) indien de aanvaller de wachtwoorden van gebruikers
nodig heeft (om bijvoorbeeld vervolgens op een ander systeem te
gebruiken). De aanvaller zal een bruteforce of dictionary attack moeten
uitvoeren om te ontdekken welk wachtwoord resulteert in welke hash. Dit
kan een erg tijdrovend proces zijn want enkele veelgebruikte hashing
algoritmen (onder andere \emph{scrypt} en \emph{bcrypt}) zijn \emph{by
design} zodanig geschreven dat deze erg traag werken. Dit zorgt ervoor
dat de tijd om één hash te genereren geen voelbaar verschil geeft, maar
wanneer een aanvaller er duizenden per seconden wil kunnen testen, dan
zal het algoritme als een stevige \textbf{\emph{timebottleneck}}
optreden. De aanvaller zou dan in de plaats vooraf alle hashes kunnen
\emph{precomputen} als alternatief. Dit heeft dan weer voor gevolg dat
zo'n lijst gigantisch groot is en er dus een
\textbf{\emph{memorybottleneck}} optreedt.

De aanvaller zit dus met het dilemma (tijd versus geheugen) tussen
hashen berekenen ter plekke, wat erg traag zal gaan, oftewel alle
mogelijke hashes op voorhand berekenen, wat veel geheugenplek vereist.
Via een \textbf{rainbow table attack} krijgt de aanvaller echter een
handig instrument in handen dat een compromis tussen beide bottlenecks
aanbiedt.

Een rainbow table is een tabel van precomputed hashes, maar waarvan we
ze niet allemaal moeten bewaren om toch over een grotere set te
beschikken dan die dat in de tabel bewaard worden. Je zou het kunnen
vergelijken met een gecomprimeerde lijst van de getallen van 1 tot en
met 101, waarbij we enkel het start (1) en eindgetal (101) bewaren, en
dan erbij zeggen dat ieder volgend getal het vorige +2 is.

\begin{figure}
\centering
\includegraphics[width=0.7\linewidth,height=\textheight,keepaspectratio,alt={Ieder wachtwoord mapt naar exact één hash.}]{auth/rainbow0.png}
\caption{Ieder wachtwoord mapt naar exact één hash.}
\end{figure}

Een rainbow table stel je als volgt op:

\begin{itemize}
\tightlist
\item
  Je kiest een startpunt, zijnde één van de mogelijk wachtwoorden uit de
  set van wachtwoorden waarvoor je een rainbow table wilt opstellen (zie
  figuur hier voor)
\item
  Je genereert de hash van dit gekozen wachtwoord.
\item
  Je past nu op deze hash een \emph{reduction} functie toe. Dit is een
  zelfgekozen mapping van de verkregen hash terug naar een wachtwoord
  uit de set van mogelijk wachtwoorden.
\end{itemize}

\begin{tipblock}
Als reductiefunctie zou je bijvoorbeeld kunnen beslissen om de hash om
te zetten naar een getal (bv. door de som van de ASCII-waarden van de
letters van de hash te nemen) en dit getal dan te gebruiken als index
die bepaalt welk wachtwoord je uit de wachtwoordenset gaat kiezen.
\end{tipblock}

\begin{figure}
\centering
\includegraphics[width=0.7\linewidth,height=\textheight,keepaspectratio,alt={Van de hash via de reductiefunctie terug naar een ander wachtwoord.}]{auth/rainbow1.png}
\caption{Van de hash via de reductiefunctie terug naar een ander
wachtwoord.}
\end{figure}

\begin{itemize}
\tightlist
\item
  Van dit nieuwe wachtwoord genereer je weer een hash en pas je opnieuw
  de reduction functie toe.
\item
  Die combinatie hash+reductie blijf je X aantal keer herhalen tot je
  een lange lijst hebt (een zogenaamde \emph{chain} van bijvoorbeeld
  10000 elementen).
\item
  Finaal hou je nu van deze lijst enkel het startpunt bij (het gekozen
  wachtwoord uit de set) en de allerlaatste gegenereerde hash.
\end{itemize}

\begin{figure}
\centering
\includegraphics[width=0.7\linewidth,height=\textheight,keepaspectratio,alt={Voorbeeld van een lijst opeenvolgende wachtwoorden en hun hashes.}]{auth/rainbow2.png}
\caption{Voorbeeld van een lijst opeenvolgende wachtwoorden en hun
hashes.}
\end{figure}

Wanneer de aanvaller nu van een gestolen hash terug het wachtwoord te
pakken wil krijgen dan zal hij:

\begin{enumerate}
\def\labelenumi{\arabic{enumi}.}
\tightlist
\item
  Deze hash als startpunt gebruiken en hier telkens weer de combinatie
  reductie+hash op toepassen.
\item
  \textbf{Totdat} een hash wordt gevonden die als eindpunt voor een van
  de lijsten werd ingesteld.
\item
  De aanvaller neemt vervolgens het startpunt van deze lijst (een
  wachtwoord) en past daarop opnieuw de combinatie van hashing en
  reductie toe (als het ware opnieuw een rainbow table genereren).
  Uiteindelijk zal hij weer op de gestolen hash uitkomen.
\item
  Als hij dan één stapje terug kijkt dan zal hij daar het wachtwoord
  zien die bij deze gestolen hash hoort.
\end{enumerate}

\begin{figure}
\centering
\includegraphics[width=0.7\linewidth,height=\textheight,keepaspectratio,alt={Ieder wachtwoord mapt naar exact één hash.}]{auth/rainbow3.png}
\caption{Ieder wachtwoord mapt naar exact één hash.}
\end{figure}

\subsection{Salting}\label{salting}

Om bestand te zijn tegen de rainbow attack dienen we de set van
mogelijke wachtwoorden gevoelig te vergroten waardoor het niet meer
realistisch is om voor die set rainbow tables te genereren. We kunnen
helaas niet verwachten van de eindgebruikers dat zij met véél langere,
meer willekeurige, wachtwoorden op de proppen komen en zullen dus een
`oude' truc moeten gebruiken die we ook al bij wifi hebben gezien. Bij
wifi hanteerden we een initialisatie vector (IV) om de WEP-sleutel met
24 bits te verlengen zodat zelfs bij dezelfde sleutel, iedere IV
eigenlijk zorgt voor een unieke seed.

Wel nu, dit concept kan je ook toepassen bij wachtwoorden en heet
\textbf{salting}. Een salt is een extra stuk dat je toevoegt aan het
wachtwoord \textbf{voor je de hash} berekent. Dit extra stukje is een
willekeurig getal dat je uiteraard mee zal moeten opslaan in de
database. Wanneer twee gebruikers hetzelfde wachtwoord zouden hebben,
dan zouden ze (dankzij hun unieke salt) toch beiden totaal verschillende
hashes genereren. Niet alleen dat, maar de salt zorgt er dus ook voor
dat de set van mogelijk wachtwoorden véél groter wordt.

\begin{figure}
\centering
\includegraphics[width=0.9\linewidth,height=\textheight,keepaspectratio,alt={Het salting proces.}]{auth/salting.png}
\caption{Het salting proces.}
\end{figure}

In de database bewaren we nu volgende informatie:

\begin{itemize}
\tightlist
\item
  Gebruikersnaam.
\item
  Gebruikte salt (minimum 32 bits).
\item
  Hash van het wachtwoord en de salt samen.
\end{itemize}

\begin{warningblock}
Merk op dat ook nu we nog steeds niet beschermd zijn tegen pass-the-hash
aanvallen.
\end{warningblock}

\subsubsection{Mimikatz}\label{mimikatz}

Mimikatz werd origineel ontwikkeld als demo om aan te tonen dat de
authenticatieprotocollen van Microsoft onveilig waren. Helaas is de tool
totaal erg snel opgenomen in het arsenaal van de digitale stropers. De
tool laat toe om \emph{authentication tickets} te tonen en hergebruiken.
Zo'n authentication tickets worden door de loginserver aangemaakt na een
geslaagde loginfase door de gebruiker. Dit ticket kan de gebruiker dan
aan een systeem aanbieden om toegang tot het systeem te krijgen (het is
letterlijk een \emph{toegangsticketje}). Wanneer Mimikatz wordt
losgelaten op een Microsoft Windows besturingssysteem zal het deze
tickets op het systeem zoeken zodat de aanvaller vervolgens zonder login
gegevens toch kan inloggen door technieken zoals:

\begin{itemize}
\tightlist
\item
  Pass-the-hash: vroeger werden Windows wachtwoorden als hash (NTLM)
  bewaard op het systeem waardoor deze techniek erg eenvoudig was.
\item
  Pass-the-ticket: zoals zonet beschreven, maar dan met het Kerberos
  ticket (zie ook hierna)
\item
  Kerberos Golden Ticket: \textbf{Kerberos} is een van de meest
  gebruikte authenticatieprotocollen. Veel systemen die Kerberos
  gebruiken hebben echter een verborgen account (\emph{KRBTGT} genaamd)
  wiens ticket domain admin rechten geeft én dat niet vervalt. Kortom,
  een gouden ticket!
\item
  Pass-the-cash: identiek aan pass-the-ticket maar deze aanval werkt ook
  met logindata die zowel op Mac, Unix én Linux kan worden gevonden.
  Kortom, dit is natuurlijk de \emph{motherload}, daar deze niet meer
  afhankelijk is van enkel Microsoft Windows besturingssystemen.
\end{itemize}

Het nadeel van Mimikatz, voor ons als boswachters, is dat de tool erg
goed werkt én kan geautomatiseerd worden. In 2017 onderging Oekraïne een
stevige ransomware aanval van (zo goed als zeker) Russische makkelij,
genaamd \textbf{NotPetya} (een variant op de WannaCry ransomware).
NotPetya gebruikte een aangepaste versie van Mimikatz zodat de
ransomware zichzelf kon verspreiden over het netwerk en op andere
systemen in het domein kon inloggen met hashes en tickets dat de
Mimikatz variant aantrof.

\begin{tipblock}
De oorsprong van Petya en NotPetya werd getraceerd en is vermoedelijk
het resultaat van een Russische hackinggroep genaamd \emph{Sandworm} die
onder de GRU werken, de Russische militaire inlichtingendienst.
\end{tipblock}

\section{CRAM en SCRAM}\label{cram-en-scram}

Om iemand te authenticeren spraken we tot nog toe enkel over een
username/wachtwoord systeem. Echter, er zijn vele andere manieren om
iemand te authenticeren. We spreken over ``\textbf{Challenge-Response
Authentication Mechanism (CRAM)} wanneer de gebruiker een vraag gesteld
krijgt (de \emph{challenge}) en hij hierop een geldig antwoord (de
\emph{response}) moet geven voor hij wordt toegelaten. Authenticeren met
een wachtwoord is dus een vorm van CRAM. Er zijn er echter nog vele
andere, denk maar aan de gehekelde CAPTCHA's - de ambetante vraag om te
bewijzen dat je geen robot bent door alle boten in een afbeelding aan te
duiden - of inloggen met behulp van je irisscan.

\begin{figure}
\centering
\pandocbounded{\includegraphics[keepaspectratio,alt={CRAM.}]{auth/cram.png}}
\caption{CRAM.}
\end{figure}

Het mechanisme van een CRAM werkt als volgt:

\begin{enumerate}
\def\labelenumi{\arabic{enumi}.}
\tightlist
\item
  De gebruiker stuurt z'n username met de vraag om in te loggen.
\item
  De server genereert een random challenge en stuurt deze terug.
\item
  Server en client creëren nu een hash van deze challenge met de hash
  van het wachtwoord (de server heeft dit bewaard, de client genereert
  de hash door z'n wachtwoord in te voeren)
\item
  De client stuurt deze hash, de response, terug naar de server.
\item
  De server vergelijkt of zijn gegenereerde response hash dezelfde is
  als die van de gebruiker.
\end{enumerate}

\begin{figure}
\centering
\pandocbounded{\includegraphics[keepaspectratio,alt={CRAM flow.}]{auth/cramflow.png}}
\caption{CRAM flow.}
\end{figure}

Om het probleem van \emph{pass-the-hash} op te lossen kan je gebruiken
maken van een \textbf{S}CRAM, een \textbf{Salted Challenge-Response
Authentication Mechanism}. We bespreken een vereenvoudigde versie (een
echte SCRAM voorziet ook \emph{mutual authentication}) waarbij we
hoofdzakelijk willen uitleggen waarom een SCRAM systeem veiliger is dan
een klassieke salted wachtwoord login van daarnet. Met dit systeem
zorgen we ervoor dat :

\begin{enumerate}
\def\labelenumi{\arabic{enumi}.}
\tightlist
\item
  De salted hash van de gebruiker NOOIT moet verzonden worden.
\item
  Geen replay aanval m.b.v. pass-the-hash mogelijk is.
\end{enumerate}

Zoals je in de afbeelding kunt zien zal in dit systeem de server ook de
bewaarde salt naar de client sturen, zodat deze geen gebruik kan maken
van een bewaarde password hash die hij niet zelf heeft gemaakt.

\begin{figure}
\centering
\pandocbounded{\includegraphics[keepaspectratio,alt={SCRAM.}]{auth/scramflow.png}}
\caption{SCRAM.}
\end{figure}

\section{Multifactor authentication}\label{multifactor-authentication}

We hebben enkel nog maar over wachtwoorden gesproken in dit hoofdstuk,
maar uiteraard zijn er ook andere zaken die je kan gebruiken om je te
identificeren. Er zijn verschillende \textbf{factoren} die kunnen
gebruikt worden om te controleren of een persoon wel degelijk toegang
mag krijgen tot een systeem:

\begin{itemize}
\tightlist
\item
  Iets wat je \textbf{weet}: je wachtwoord, je pincode, je
  rijksregisternummer, etc.
\item
  Iets wat je \textbf{bent}: een eigenschap die uniek is per persoon en
  onder de noemer ``\emph{biometrics}'' valt, zoals je vingerafdruk,
  irisscan, etc.
\item
  Iets wat je \textbf{hebt}: een stuk hardware zoals een smartphone,
  USBkey, etc.
\item
  \textbf{Waar} of \textbf{wanneer} je bent: je IP-adres , het moment
  van de dag dat je probeert in te loggen.
\end{itemize}

\begin{figure}
\centering
\pandocbounded{\includegraphics[keepaspectratio,alt={Multifactor authenticatie factoren.}]{auth/mfa.png}}
\caption{Multifactor authenticatie factoren.}
\end{figure}

We zien meer en meer systemen verschijnen die aan zogenaamde
\textbf{multifactor authentication (MFA)} doen waarbij het systeem
minstens twee factoren (\emph{2FA}) wil controleren voor je toegelaten
wordt. Hoe meer verschillende factoren er worden gebruikt bij de
authenticatie hoe veiliger het systeem is, maar ook hoe minder
gebruiksvriendelijk het wordt. Het blijft dus een afweging tussen die
twee eigenschappen om in te schatten wat de ideale hoeveelheid
veiligheid en gebruiksvriendelijkheid is die je wenst te hebben.

Systemen die MFA aanbieden doen dit vaak op een gecontroleerde manier:
afhankelijk van de gebeurtenissen zal het systeem beslissen of meerdere
factoren moeten getest worden of niet. Als je bijvoorbeeld in België
woont, maar Google ziet plots dat iemand met jouw wachtwoord probeert in
te loggen vanuit een IP-adres op de Azoren, dan zal Google beslissen dat
``iets wat je weet'' (het wachtwoord) niet genoeg controle is en extra
informatie vragen (andere factoren controleren).

\subsection{Iets wat je weet:
Paswoorden}\label{iets-wat-je-weet-paswoorden}

Deze factor hebben we reeds uitvoerig behandeld. Het grote probleem met
\emph{dingen weten} is dat:

\begin{itemize}
\tightlist
\item
  Je ze kan vergeten en daardoor niet meer kan inloggen.
\item
  Anderen die informatie kunnen te weten komen en zich vervolgens als
  jou voordoen.
\end{itemize}

Kortom, alhoewel deze factor vaak vanuit technologisch standpunt het
eenvoudigst te implementeren is, is het ook de minst veilige vanuit een
social engineering standpunt. Biometrics en hardware zijn moeilijker
door een digitale stroper te stelen dan het wachtwoord en we hoeven niet
bij een fingerprint te vrezen dat de gebruiker een ``zwakke
vingerafdruk'' kiest, iets wat bij wachtwoorden vaak hét primaire
probleem is.

\subsection{Iets wat je bent:
Biometrics}\label{iets-wat-je-bent-biometrics}

``We zijn allemaal uniek''. Iedere mens heeft een hele hoop
eigenschappen die uniek zijn per persoon. Zelfs als bepaalde van onze
eigenschappen gelijkaardig zijn met andere personen dan zal zeker een
combinatie van twee of meerdere eigenschappen dat niet zijn. Door dus
menselijke eigenschappen van je te gebruiken als authenticatie hebben we
een factor gevonden die zeer moeilijk na te bootsen valt: op voorwaarde
dat je een echt unieke eigenschap kiest én deze op deze juiste manier
meet.

Enkel veel gebruikte biometrieken als authenticatievorm zijn:

\begin{itemize}
\tightlist
\item
  Vingerafdruk.
\item
  Iris.
\item
  Stem.
\item
  Gezicht (vaak met behulp van ``stereo camera'' voor 3D beeld).
\end{itemize}

Maar ook andere metrieken kunnen erg interessant zijn zoals de manier
waarop je je wachtwoord invoert, de manier waarop je wandelt
(\emph{gait}) etc.

\begin{warningblock}
Paswoorden van miljoenen mensen opslaan is één ding. De biometrische
gegevens is een heel ander verhaal waarbij ook \textbf{privacy} plots
een erg heikel punt wordt (beeld je even in dat Hitler en zijn trawanten
80 jaar geleden toegang hadden tot biometrische data waarmee met een
bepaalde zekerheid kon vastgesteld worden of iemand van Joodse origine
was of niet.)

In India is de Aadhaar (Indiaas voor ``basis''), hun rijksregisternummer
zeg maar, een unieke code die gebaseerd is op \emph{onder andere de
irisscan en vingerafdrukken} (alle 10!) van de burger. Deze gigantische
database werd in 2018 nog gehackt waardoor mogelijk de informatie van
1.1 miljard geregistreerde burgers werd gestolen.
\end{warningblock}

Om een biometriek in de wachtwoord database te bewaren hebben we een
manier nodig om deze te digitaliseren op een zodanige manier dat de
unieke aspecten ervan bewaard worden. Voorts moet er rekening mee
gehouden worden dat het ``registreren'' van een biometrische eigenschap
nooit 100\% accuraat kan. Denk maar aan een tijdelijk krasje op je
vinger, je baard die anders geschoren is, etc.

De zogenaamde \emph{feature points} van een biometrische eigenschap
worden in de database bewaard: dit zijn de unieke waarden waarvan
geweten is dat deze per persoon anders zijn. We gaan deze niet per
biometrische eigenschap bespreken, het volstaat te begrijpen dat in de
gebruikersdatabase meestal een korte sequentie van getallen (of letters,
denk maar aan een DNA-sample) wordt bewaard die als het ware jouw unieke
wachtwoord voorstelt voor die specifieke biometrische eigenschap van j.
Enkel wanneer je bij het opnieuw inloggen (quasi) dezelfde feature
points genereert als bij de registratie zal deze factor aanvaard worden
als correct.

\begin{tipblock}
Biometrische eigenschappen kunnen niet alleen dienst doen als een extra
factor bij het authenticeren, ze zijn uiteraard ook erg handig voor
identificatie. In principe kan iemand nog steeds de gebruikersnaam van
een ander persoon gebruiken. Als de biometrische eigenschappen als
identificatie dienen kunnen aanvallers dat niet meer doen: ze kunnen
onmogelijk aan het systeem zeggen \emph{``ik ben persoon x''} terwijl de
vingerafdrukscanner duidelijk een vingerafdruk registreert van
\emph{persoon y}.
\end{tipblock}

\subsection{Iets wat je hebt: Hardware}\label{iets-wat-je-hebt-hardware}

Een fysiek object, zeker als het complex is, kan moeilijk nagemaakt
worden en is dus een ideale factor. De elektronica van de 21e eeuw
behoort tot de meest complexe dingen ooit die de mensheid heeft kunnen
vervaardigen. Het is dan ook logisch dat we deze elektronica gebruiken
als extra authenticatiefactor. In essentie zal dit stuk hardware nog
steeds gewoon een wachtwoord bevatten, maar dit zal echter ongelooflijk
veel langer (en dus sterker) zijn dan het gemiddelde wachtwoord dat een
standaard gebruiker kan onthouden.

Er zijn twee grote families van hardware-gebaseerde
authenticatie-vormen:

\begin{itemize}
\tightlist
\item
  Een smartphone, met daarop een \emph{authenticator} app.
\item
  Een USB sleutel.
\end{itemize}

Een nadeel van deze groep is dat het om een fysiek object gaat dat je
kan verliezen of dat stuk kan gaan.

\section{Federation en Single Sign-On
(SSO)}\label{federation-en-single-sign-on-sso}

Bij cryptografie wordt het ten stelligste afgeraden om zomaar op de
\emph{wilde boef} een eigen crypto-algoritme te ontwikkelen. De kans dat
je fouten met verstrekkende gevolgen maakt is te groot. Ook bij het
omgaan van logindata van gebruikers en hoe je ze authenticeert is het
aangeraden om even te bezinnen voor je er zelf aan begint.

Dankzij het concept \textbf{federation} hoef je niet wakker te liggen
van hoe je je gebruikerswachtwoorden gaat opslaan: gebruikers kunnen
inloggen gebruik makend van hun bestaande Google, Facebook en andere
accounts. Via federatie zal de gebruiker op jouw site (of app), de
\emph{service provider}, kunnen inloggen waarbij een \emph{third-party}
- die jij en je gebruiker vertrouwen - voor de eigenlijke authenticatie
zorgt (\emph{de identity provider}), gebruik makend van zogenaamde
\emph{single sign-on} (\textbf{sso}) authenticatie.

\begin{figure}
\centering
\pandocbounded{\includegraphics[keepaspectratio,alt={Een vereenvoudigd single sign-on proces.}]{auth/sso.png}}
\caption{Een vereenvoudigd single sign-on proces.}
\end{figure}

Federation via SSO is een onderdeel van \emph{federated identity
management}, een groep technologieën en concepten die ervoor zorgen dat
de identiteit van een gebruiker over meerdere, onafhankelijke systemen
wordt bewaard en gebruikt. Je zal de termen \textbf{delegation} en
\textbf{federation} soms door elkaar zien tegenkomen wanneer je meer
informatie over SSO opzoekt.

Samengevat gaan we bij \emph{delegation} een gebruiker verplichten in te
loggen met een bepaalde third-party die dit ondersteunt (bv. inloggen
met je Facebook account). Bij \emph{federation} gaat het breder: je
website zal éénder welke third-party account aanvaarden, zolang deze
maar compatibel is met het authenticatie systeem van je website (een
voorbeeld hiervan is OpenID).

\begin{figure}
\centering
\includegraphics[width=0.4\linewidth,height=\textheight,keepaspectratio,alt={Enkele van de vele typische SSO knoppen die je geregeld zal tegenkomen.}]{auth/ssoexamples.png}
\caption{Enkele van de vele typische SSO knoppen die je geregeld zal
tegenkomen.}
\end{figure}

\begin{tipblock}
\textbf{OAuth} (\emph{open authorization}) is een gestandaardiseerde
manier om aan authenticatie te doen.
\end{tipblock}

Uiteraard moeten we bij federatie benadrukken dat ook hier
\textbf{privacy} een belangrijk aspect wordt. De vraag is dan ook in
hoeverre je een bedrijf zoals Google of Facebook/Meta vertrouwt met jouw
(login)data.

\section{WebAuthn en passkeys}\label{webauthn-en-passkeys}

Een nieuwe technologie die in opmars is, is \textbf{WebAuthn}. Deze
technologie laat toe om in te loggen op websites zonder dat je een
wachtwoord moet ingeven. In plaats daarvan gebruik je een
\emph{authenticator} die je bij je hebt, zoals een USB-sleutel of je
smartphone. Deze authenticator zal een digitale handtekening genereren
die de website kan controleren.

Passkeys combineren de kracht van \textbf{asymmetrische/public-key
crypto en hardware authenticatie}, met als grootste pluspunt dat
gebruikers geen complexe wachtwoorden meer moeten onthouden.

Zonder in de ontstaansgeschiedenis te duiken, is het toch nuttig even
enkele termen in vet te zetten die je zeker zal tegenkomen als je meer
over passkeys wilt leren:

\begin{itemize}
\tightlist
\item
  \textbf{FIDO Alliance}: de organisatie die de standaarden voor
  WebAuthn en U2F beheert. FIDO staat voor \emph{Fast IDentity Online}.
\item
  \textbf{WebAuthn} (\emph{Web Authentication API}): een API die
  websites toelaat om met authenticators te communiceren.
\item
  \textbf{FIDO2}: een standaard die WebAuthn ondersteunt, ontwikkel door
  de FIDO Alliance.
\item
  \textbf{U2F} (\emph{Universal 2nd Factor)}: een oudere standaard die
  ook door WebAuthn wordt ondersteund.
\end{itemize}

\begin{tipblock}
Op \href{https://youtu.be/cEhc6vMFTh4}{youtu.be/cEhc6vMFTh4} vind je een
heel duidelijk overzicht van passkeys, inclusief de technische zijde
ervan.
\end{tipblock}

\subsection{Een Passkey aanmaken: de
registratie}\label{een-passkey-aanmaken-de-registratie}

Een passkey aanmaken is een eenvoudig proces. Je hebt een authenticator
nodig, zoals een USB-sleutel of je smartphone. De authenticator zal een
paar sleutels genereren: een \emph{public key} en een \emph{private
key}. De public key wordt naar de website gestuurd, terwijl de private
key (het wachtwoord met andere woorden) op de authenticator blijft.

Wanneer je dus als gebruiker registreert op een website of app, dan zal
de passkey generatie van start gaan, als volgt:

\begin{enumerate}
\def\labelenumi{\arabic{enumi}.}
\tightlist
\item
  Bij het registreren kiest de gebruiker ervoor om een passkey te
  gebruiken (i.p.v. het klassieke wachtwoord).
\item
  De gebruiker zal op zijn eigen toestel zichzelf nu eerst moeten
  identificeren. Dat kan op verschillende manieren: fingerprint scan,
  een PIN-code, een hardwaresleutel (denk aan bijvoorbeeld aan YubiKey),
  etc.
\item
  Het toestel van de gebruiker genereert een sleutelpaar. De private
  sleutel blijft op het toestel en wordt veilig bewaard, de public
  sleutel wordt naar de website gestuurd.
\end{enumerate}

Stap 3 gaan we even verder uit de doeken doen: we gaan natuurlijk deze
sleutel niet zomaar \emph{over den draad} versturen. We gaan natuurlijk
onze kennis van certificaten gebruiken, die ons toelaten om te bewijzen
dat de publieke wel degelijk de onze. De gebruiker zal zijn publieke
sleutel verpakken in een \emph{attestation object}: de publieke sleutel,
samen met een \emph{signed challenge} (zie verder), een \emph{credential
ID} en een certificaat. Dit attestation object wordt naar de andere
zijde gestuurd, die deze zal bewaren.

De tegenpartij, bijvoorbeeld de website waar je wilt registreren, heeft
uiteraard nog bewijs nodig dat het jouw attestation object wel kan
vertrouwen. Tijdens stap 1 van de registratie zal de server daarom een
challenge sturen, die de gebruiker mee in het attestation object moet
plaatsen.

In dit hele proces heeft de website nooit toegang tot de private sleutel
van de gebruiker. De website kan enkel de public sleutel zien en
gebruiken. Het concept ``het wachtwoord verlaat nooit het apparaat''
wordt hier dus erg letterlijk genomen.

\begin{warningblock}
Doordat deze private sleutels niet meer op de website worden bewaard,
wordt het gebruik van password managers nog belangrijker. De private
sleutels worden immers opgeslagen op de authenticator, en als je die
lange, complexe stukken data verliest, ben je al je accounts kwijt.

Een password manager kan je helpen om je accounts te beheren en je
private sleutels (je passkeys, in dit geval zijn dit \textbf{synced
passkeys}, een concept dat ook mee in WebAuthn is ingebouwd) veilig te
bewaren én te synchroniseren naar je andere apparaten. I
\end{warningblock}

\subsection{Inloggen met een Passkey}\label{inloggen-met-een-passkey}

Het inloggen met een passkey is gebaseerd op wat we weten uit public key
crypto: je publieke sleutel kan je aan iedereen geven, enkel de houder
van de bijhorende private sleutel zal de data kunnen lezen die met deze
publieke sleutel werd geëncrypteerd.

De login-fase is dan ook bijna het zelfde als de registratie. Ook nu zal
de gebruiker een challenge krijgen. Deze challenge zal de gebruiker nu
encrypteren met z'n private sleutel. Wanneer de server deze
geëncrypteerde challenge kan decrypteren met de bewaarde publieke
sleutel van de gebruiker, weet deze dat de gebruiker mag toegelaten
worden.

\begin{figure}
\centering
\pandocbounded{\includegraphics[keepaspectratio,alt={Het login proces met een passkey}]{auth/passkey.png}}
\caption{Het login proces met een passkey}
\end{figure}

\chapter{IoT Security}\label{iot-security}

We spraken in het eerste hoofdstuk al over de problemen die inherent
zijn bij Internet-of-Things (IoT) apparaten en netwerken. Zo had je in
2016 het Mirai-botnet dat aantoonde hoe krachtig een botnet van
IoT-apparaten kan zijn, puur door de grote hoeveelheid potentiële
zombies.

Het is niet toevallig dat we dit boek afsluiten met het hoofdstuk
omtrent IoT Security: alles (en meer) wat je in de voorgaande
hoofdstukken hebt geleerd heb je namelijk nodig om IoT-netwerken te
beveiligen. Er zijn zoveel manieren voor digitale stropers om misbruik
te maken van een slecht beveiligd IoT-apparaat of netwerk dat het
complexe probleem van netwerken beveiligen nog ingewikkelder is
geworden, namelijk:

\begin{itemize}
\tightlist
\item
  Van nature werken IoT-apparaten op \textbf{low-power}, daar ze vele
  dagen of weken moeten voortkunnen zonder te veel verbruik. Ingebouwde
  security protocollen mogen dus geen grote percentages van het vermogen
  opsouperen, zeker niet als dat ten koste is van de hoofdbestaansreden
  van het apparaat.
\item
  Meestal, ook weer om een laag energieverbruik te behouden, hebben ze
  ook een \textbf{beperkte bandbreedte} ter beschikking (daarbij: beeld
  je de miserie op een wifi-netwerk in wanneer honderden kleine sensoren
  aan hoge bandbreedtes het netwerk mee gebruiken).
\item
  Zowel de fysieke dimensies en bovenstaande twee eigenschappen zorgen
  er ook voor dat IoT-apparaten meestal maar een \textbf{beperkt aantal
  zaken kunnen}. Het heeft niet altijd de
  netwerkverbindingsmogelijkheden die je nodig hebt (meestal heeft een
  apparaat een draadloze - wifi, Zigbee, LoRaWAN,etc. - of bedrade
  aansluiting, maar zelden beiden) en ook qua opslag en
  processormogelijkheden zijn de apparaten uiteraard vaak veel beperkter
  dan de kracht van computers en mobiele telefoons. Het grootste
  probleem hiervan is echter dat IoT-apparaten \emph{by nature} vaak
  \textbf{niet evident zijn om te updaten}. Sommige apparaten kunnen
  bijvoorbeeld enkel geüpdatet worden door fysiek met het apparaat te
  zijn verbonden (en dan door de firmware te flashen bijvoorbeeld).
  Vanuit een fysieke beveiligingsconcept is dat goed, maar niet als jouw
  honderden apparaten een kritieke beveiligingslek hebben die dringend
  gepatcht moet worden.
\item
  Vaak bevinden IoT-apparaten zich op de koop toe op
  \textbf{onbeschermde locaties} die maar moeilijk fysiek te beveiligen
  zijn. Hierdoor kunnen digitale stropers, zonder schrik om betrapt te
  worden, ongezien met de apparaten knoeien (\emph{tamperen}).
\end{itemize}

\begin{noteblock}
We gaan nogal vrij om met de term IoT-apparaat, waarbij we ook negeren
dat er grote verschillen zijn tussen ``Enterprise IoT'' en
``huis-tuin-en-keuken IoT''. Die eerste categorie heeft meestal
bijvoorbeeld wel zeer doordachte updatestrategieën en dergelijke,
terwijl de slimme koelkast in je huis dat vermoedelijk niet zal hebben.
\end{noteblock}

\section{OWASP IoT Top 10}\label{owasp-iot-top-10}

Het Open Web Application Security Project (\textbf{OWASP}) is een
open-source-project rond computerbeveiliging. Individuen, scholen en
bedrijven delen via dit platform informatie en technieken. Ook de
Belgische tak van OWASP is erg actief en een interessante organisatie
indien je van plan bent om een carrière binnen de cybersecurity wereld
op te bouwen. Ze organiseren geregeld (meestal gratis) workshops en
evenementen.

Het OWASP heeft tal van erg interessante ``Top 10'' overzichten zoals de
meest voorkomende webapp zwakheden, etc. Deze projecten bestaan meestal
uit een reeks tools, \emph{best practices} en gidsen en zijn dus de
ideale manier om je te verdiepen binnen een specifiek cybersecurity
domein. Eén van de actieve projecten is het ``OWASP Internet of Things''
project (beschikbaar via
\href{https://owasp.org/www-project-Internet-of-Things/}{owasp.org/www-project-Internet-of-Things/})
waarvan we de Top 10 hier zullen bespreken, gerangschikt van meest
voorkomend (en dus belangrijkste om aan te pakken) naar minder.

\begin{tipblock}
Bekijk zeker ook eens het IoT Goat project via
\href{https://github.com/OWASP/IoTGoat/}{github.com/OWASP/IoTGoat/} dat
is een \emph{``insecure firmware based on OpenWrt and maintained by
OWASP as a platform to educate software developers and security
professionals with testing commonly found vulnerabilities in IoT
devices. The vulnerability challenges are based on the OWASP IoT Top 10,
as well as `easter eggs' from project contributors.''}
\end{tipblock}

\begin{warningblock}
Om de wel erg algemene term ``IoT'' wat duidelijker te maken, zullen er
hier en daar merknamen en producten worden vermeld. Dit is hoegenaamd
geen verdoken reclame (of kritiek) voor deze producten, maar gewoon een
poging om de lezer een duidelijkere context te geven.
\end{warningblock}

\begin{figure}
\centering
\pandocbounded{\includegraphics[keepaspectratio,alt={Versie 2018 (Bron owasp.org).}]{iot/top10.jpg}}
\caption{Versie 2018 (Bron owasp.org).}
\end{figure}

\subsection{Paswoorden}\label{paswoorden}

Ook bij IoT begint veiligheid bij het hebben van goede, complexe,
moeilijk te raden wachtwoorden. Bij IoT-apparaten is er echter een
stevige drempel, zeker voor huis-tuin-en-keuken gebruikers: het
aanpassen van sommige IoT-apparaten is soms erg omslachtig. Niet alle
apparaten hebben een webpagina langs waar de gebruiker wachtwoorden kan
aanpassen en dus moet dit ofwel via omslachtige, vreemde tools, ofwel
(gruwel) door fysiek het apparaat te bedienen. Deze apparaten zijn
natuurlijk niet gemaakt om complexe wachtwoorden eenvoudig in te voeren.
Denk maar aan een digitale weegschaal: die heeft vaak een eenvoudig
LED-scherm waarop je gewicht wordt getoond. Sommige weegschalen kan je
vervolgens ``bedienen'' door links of rechts op de schaal te klikken,
afhankelijk van wat je wilt doen. De auteur kan je garanderen: je
wachtwoord zal niet lang zijn, als je een goed wachtwoord wilt invoeren
zal je na enkele minuten in het zweet staan van de tapdans die je moet
uitvoeren op je weegschaal.

Omdat het aanpassen van die wachtwoorden vaak omslachtig is, zullen
gebruikers meestal opteren om het \emph{default wachtwoord} te
gebruiken. Ze redeneren dat \textbf{toch niemand dit simpele
IoT-apparaat zal willen hacken}. ``Wat kan een hacker nu doen met
toegang tot mijn digitale thermometer.'' Wat echter vergeten wordt is
dat cyberstropers altijd op zoek gaan naar de \emph{weakest link}, om
die vervolgens te misbruiken om tot het eigenlijk doel te geraken.
Kortom, ieder IoT-apparaat in je omgeving is een potentiële
toegangspoort tot de rest van je netwerk!

\subsection{Ongebruikte of onveilige
netwerkservices}\label{ongebruikte-of-onveilige-netwerkservices}

Als je anno 2023 een printer koopt dan zal dit apparaat bijna altijd een
IoT-apparaat zijn dat je toelaat om vanop afstand te printen. Echter,
dergelijke apparaten hebben vaak tal van netwerkservices draaien waarvan
de gewone huis-tuin-en-keuken gebruiker weinig of niets van afweet.

\begin{tipblock}
Vergelijk dit bijvoorbeeld met je Microsoft Windows besturingssysteem:
heb je al eens gekeken welke services allemaal permanent op je systeem
draaien (Start -\textgreater{} uitvoeren -\textgreater{}
\texttt{services.msc}) ?
\end{tipblock}

Dit soort diensten uitzetten behelst ook weer dat de gebruiker weet
heeft van een webpagina met administrator-instellingen. Soms zijn die
diensten op de koop toe essentieel voor de goede werking van het
apparaat en ze uitzetten is dan geen optie. Maar wat als die service op
zich inherent onveilig is?! Zo zijn er apparaten die enkel
\emph{telnet}-sessies gebruiken om het apparaat in te stellen. Een
telnet-sessie geeft je een remote shell die echter
\textbf{onversleuteld} met het apparaat communiceert (in tegenstelling
tot bijvoorbeeld SSH, Secure Shell, dat wél veilige communicatie
aanbiedt).

Kortom, het is aangeraden om altijd te controleren welke services nu
eigenlijk actief zijn op je apparaat.

\begin{noteblock}
In 2017 was er een grayhat hacker die vermoedelijk 150.000 printers
vanop afstand in Nederland kon benaderen. Vervolgens printte hij
``ludieke'' berichten op de printers om de gebruikers er op te wijzen
dat ze een onveilig apparaat gebruikten. In dit geval gebruikte hij
poort 9100 van de printer die in veel merken wordt gebruikt om vanop
afstand print-opdrachten door te sturen.

\begin{figure}
\centering
\includegraphics[width=0.6\linewidth,height=\textheight,keepaspectratio,alt={Bron afbeelding en artikel: https://www.inktweb.nl/blog/hacker-neemt-150-000-printers-over/}]{iot/printer.jpg}
\caption{Bron afbeelding en artikel:
https://www.inktweb.nl/blog/hacker-neemt-150-000-printers-over/}
\end{figure}

\end{noteblock}

\subsection{Onveilig ecosysteem}\label{onveilig-ecosysteem}

Een IoT-apparaat bestaat niet alleen: meestal maakt het deel uit van een
verzameling apparaten en diensten die de gebruiker via online dashboards
en apps kan bedienen. Denk maar aan de verschillende Google Nest
producten (thermostaat, deurbel, maar ook de Hub, etc.) of lichten van
Philips (\emph{Hue}) die ook met behulp van andere toestellen kunnen
bediend worden (bijvoorbeeld drukknoppen van Niko of Ikea). Kortom,
IoT-apparaten maken meestal deel uit van een heus ecosysteem en dus de
mate van veiligheid van die externe zaken bepaalt ook die van het
IoT-apparaat. Als bijvoorbeeld de dashboards van fitbit.com zouden
gecompromitteerd worden bestaat de kans dat de digitale stropers ook
toegang tot je gegevens of apparaten krijgen die gebruik maken van het
fitbit-portaal.

Gebruikers vergeten ook met welke externe diensten ze hun apparaten
hebben gekoppeld. Het gebeurt vaak dat een gebruiker ``even iets wil
testen'' (denk maar de nuttige website \href{https://ifttt.com/}{``If
this then that'' (ifttt.com)} dat toelaat om services en apparaten te
koppelen die iets moeten doen gegeven zelfgekozen triggers en regels)
maar nadien wel vergeet deze koppeling uit te zetten. Dit geldt ook voor
de vele third-party websites en services die toegang tot uw Google of
Facebook/Meta-accounts wensen. Gebruikers kijken dit zelden na en hebben
mogelijk tien jaar geleden al toegang gegeven aan een service die
ondertussen al door digitale stropers is gecompromitteerd.

\subsection{Geen of onveilig
updatemechanisme}\label{geen-of-onveilig-updatemechanisme}

Dit werd al in de introductie van dit hoofdstuk aangehaald. Het updaten
van fysieke apparaten is niet evident, zeker niet als het gaat om
kleine, low-power-apparaten die mogelijk op de koop toe op een moeilijk
bereikbare plaats hangen en dus enkel praktisch geüpdatet kunnen worden
als ze een remote update-mechanisme ondersteunen.

Daarnaast kan het ook gebeuren dat een update, om welke reden dan ook,
mislukt waardoor het systeem potentieel \emph{gebrickt} geraakt en de
gebruiker of administrator dan alsnog een fysieke rollback moet doen
(als het apparaat die mogelijkheid heeft). Er zijn ondertussen erg dikke
boeken geschreven over de \emph{wondere wereld van het updaten van
IoT-apparaten}. Weet dus dat het belangrijk is dat je als
cyberboswachter van IoT-netwerken goed weet wat de update-strategieën
van je IoT-apparaten zijn en wat je moet doen als er hier problemen
optreden. Merk op dat die problemen zowel het gevolg van digitale
stropers kunnen zijn, maar ook van slechte updates van de fabrikant die
op hun beurt mogelijk nieuwe beveiligingsproblemen in je apparaten
introduceren.

\begin{tipblock}
Het is ook essentieel dat je je als cyberboswachter abonneert op de
juiste informatiekanalen (CVEs, fabrikant-RSS feeds, etc.)
\end{tipblock}

We zijn afhankelijk van de fabrikanten om beveiligingslekken te dichten
in onze software en hardware. Uiteraard zijn er fabrikanten die maar
\emph{x} aantal jaren patches uitbrengen, en er zijn ook al tal van
IoT-apparaten op de markt waarvan de fabrikant al jaren niet meer
bestaat.

Of wat te denken als er een kritieke bug wordt gevonden in een chip die
ontelbare apparaten gebruiken. Wie moet dan de patch voorzien: de
fabrikant van het IoT-apparaat, of die van de chip?

\begin{tipblock}
In september 2021 verscheen een nieuwe Bluetooth kwetsbaarheid, getiteld
\textbf{BrakTooth} die mogelijk op miljoenen IoT-apparaten kan misbruikt
worden. De bug zelf bevindt zich in de Bluetooth-chip van de apparaten.
Aardig wat fabrikanten van dergelijke chips hebben al beloofd een patch
uit te brengen, maar er zijn er ook die simpelweg zeggen ``we zullen
enkel patchen als er genoeg vraag naar is'', nochtans kan de
kwetsbaarheid erg veel schade aanbrengen (DoS-aanvallen en
\emph{arbitrary code execution (ACE)})

\begin{figure}
\centering

\includegraphics[width=0.3\linewidth,height=\textheight,keepaspectratio,alt={BrakTooth: de zoveelste Bluetooth kwetsbaarheid.}]{iot/brak.jpg}
\caption{BrakTooth: de zoveelste Bluetooth kwetsbaarheid.}
\end{figure}

\end{tipblock}

\subsection{Oude of onveilige
componenten}\label{oude-of-onveilige-componenten}

Het voorgaande voorbeeld i.v.m. BrakTooth kan ook in deze sectie
gebruikt worden. Een IoT-apparaat is en blijft een combinatie van
tientallen, soms honderden hardware-onderdelen. De kans dat al deze
onderdelen van de IoT-fabrikant zelf zijn is onbestaande. Gooi eender
welke smartphone of digitale horloge open en je zal tal van onderdelen
vinden van een andere fabrikant. De IoT-fabrikant moet er dus vanuit
gaan dat deze componenten voldoen aan de veiligheidseisen die hij van
z'n eigen product verwacht. Dat is geen evidentie.

Een zelfde fenomeen zien we enkele lagen hoger in de OSI-stack: ook op
software niveau gebeurt het zelden dat de IoT-ontwikkelaars alle
bibliotheken en protocollen \emph{van scratch} hebben ontwikkeld (en dus
ook zelf kunnen instaan voor de robuustheid ervan inzake
cyberveiligheid). Wanneer er een bug (of bewuste backdoor) wordt ontdekt
in een protocol of bibliotheek, dan heeft dit gevolgen voor alles en
iedereen die deze dingen in zijn of haar apparaten gebruikt (denk maar
aan de Log4J en Heartbleed voorvallen die we in hoofdstuk 1 vernoemden).

\begin{tipblock}
In 2020 was er veel heisa omtrent de \emph{SolarWinds-hack} bij
SolarWinds. Duizenden grote bedrijven (en Amerikaanse
overheidsinstellingen) zoals Microsoft gebruiken het Orion Platform van
dit bedrijf. Dit platform laat toe om enterprise netwerken te monitoren
en managen, iets wat niet evident is op de schaal waarmee de klanten van
SolarWinds werken. Deze tool moet natuurlijk waterdicht zijn, wat deze
niet bleek te zijn: hackers sloegen er in om de FTP-server van
SolarWinds te benaderen (het wachtwoord was \emph{solarwinds123}\ldots)
en konden zo een backdoor in het Orion platform inbouwen. Vervolgens,
wanneer de klanten van SolarWinds, hun Orion-installatie updateten,
kregen ze onvrijwillig een versie met een bewust ingebouwde backdoor die
de hackers vervolgens konden misbruiken.

Dit soort aanval heet een \textbf{supply chain attack}: de aanvallers
richten zich op de \emph{weakest link} binnen de supply chain om zo
finaal tot bij de klant \emph{achteraan} de keten te geraken.

\begin{figure}
\centering
\includegraphics[width=0.75\linewidth,height=\textheight,keepaspectratio,alt={Bron dynatrace.com}]{iot/supchain.png}
\caption{Bron dynatrace.com}
\end{figure}

\end{tipblock}

\subsection{Onvoldoende privacy
bescherming}\label{onvoldoende-privacy-bescherming}

Ook IoT-fabrikanten moeten natuurlijk op een veilige, GDPR-compliant,
manier omgaan met de persoonlijke informatie van hun gebruikers. Helaas
wordt dit geregeld over het hoofd gezien. Ook hier hebben we weer het
probleem van de vele potentiële manieren waarmee een digitale stroper
(of de fabrikant zelf!) een IoT-netwerk en apparaat kan misbruiken. De
persoonlijke data staat mogelijk wel op een ultra-beveiligde database
bij de fabrikant \emph{on-premise}, maar wat baat dat als deze
informatie ook als \emph{plaintext} op het apparaat wordt bewaard. Wat
ons ook automatisch bij het volgende punt brengt.

\subsection{Onveilige datatransfer en
opslag}\label{onveilige-datatransfer-en-opslag}

Herinner je je de McCumber kubus aan de start van dit handboek? Er werd
toen benadrukt dat C.I.A. toepassen op je data geen nut heeft als je
geen rekening houdt met alle vormen waarin de data zal voorkomen. Als je
enkel encryptie gebruikt tijdens het versturen van data, maar niet
tijdens het verwerken en bewaren ervan, dan kan je even goed ook de
transmissie beter ongeëncrypteerd doen (zodat we geen \emph{false sense
of security} aan de gebruiker geven).

We vallen in herhaling, maar dus zeker bij IoT-systemen die binnen een
groot ecosysteem werken, met tal van \emph{third-party} modules en
bibliotheken, is het uitermate belangrijk dat ten allen tijde de
IoT-data in al z'n vormen, op alle momenten, aan
\emph{confidentialiteit, integriteit en beschikbaarheid} doet.

\subsection{Gebrek aan apparaat
management}\label{gebrek-aan-apparaat-management}

Spreken over IoT in een enterprise of industriële omgeving is spreken
over IoT netwerken met honderden tot duizenden apparaten die bediend
moeten worden (denk maar aan remote upgrades, instellingen wijzigen,
etc.). Er zijn vele factoren waarom een robuust \emph{IoT device
management} systeem vereist is:

\begin{itemize}
\tightlist
\item
  De hoeveelheid apparaten: te veel om manueel te bedienen.
\item
  De locatie: mogelijk bevinden de apparaten zich op moeilijk te
  bereiken of gevaarlijke plekken.
\item
  Alarmen benodigd: soms wil je op de hoogte gesteld worden wanneer een
  bepaalde IoT sensor iets meet of niet goed werkt.
\item
  Nuttig gebruik van personeel: personeel enkel naar je IoT-apparaten
  sturen als dat écht nodig is.
\item
  Deel van \emph{mission critical applicaties}: als je apparaten deel
  uitmaken van bedrijfskritische systemen is het uitermate belangrijk
  dat deze apparaten maximaal functioneren, daar iedere downtime of
  crash in winstverlies of boetes kan resulteren.
\end{itemize}

Een deftig apparaat management systeem moet natuurlijk zelf ook veilig
en robuust zijn. Er zijn meerdere high-end oplossingen beschikbaar maar
het spreekt voor zich dat dit soort systemen zelden goedkoop zijn.

\subsection{Onveilige standaard
instellingen}\label{onveilige-standaard-instellingen}

Dit item komt nog steeds voor in alle ``10 veiligheidsproblemen die het
meeste voorkomen''-lijstjes. Wanneer een apparaat \emph{uit de doos}
wordt gehaald is deze ingeladen met een hele hoop standaard
instellingen. Instellingen die digitale stropers kunnen opzoeken en
misbruiken. Niet alleen het wachtwoord moet ogenblikkelijk aangepast
worden, vaak staan ook bepaalde services open zonder dat de gebruiker
deze ooit zal gebruiken. Kortom, het is een gezonde gewoonte om, zowel
als professional als huis-tuin-en-keuken IoT-gebruiker, steeds alle
instellingen van je vers geïnstalleerde apparaat te controleren.

\begin{noteblock}
Er ging lang een \emph{urban legend} de ronde dat je nooit je computer
aan het Internet mocht hangen wanneer je Windows XP aan het installeren
was. Deed je dat wel dan bestond er de kans dat je computer
\emph{gepwnd} werd nog voor je goed en wel de hele installatie had
doorlopen.

Ik heb nog geen 100\% bewijs hiervan gekregen, maar volgende artikels
lijken toch te bevestigen dat het om een effectief probleem gaat:

\begin{itemize}
\tightlist
\item
  \textbf{theregister.com/2004/08/19/infected\_in20\_minutes/}
\item
  \textbf{joeykelly.net/hacks/windows/Windows\_XP---Surviving\_The\_First\_Day.pdf}
\end{itemize}

\end{noteblock}

\subsection{\texorpdfstring{Gebrek aan fysieke
\emph{hardening}}{Gebrek aan fysieke hardening}}\label{gebrek-aan-fysieke-hardening}

IoT-apparaten liggen meestal verspreid over grote gebieden, ruimtes of
gebouwen. De kans dat ze op publiek toegankelijk plekken operationeel
zijn is groot, daardoor ook de kans dat digitale stropers toegang tot de
apparaten krijgen. Het is daarom belangrijk dat het apparaat ook op
fysiek niveau \emph{gehardened} is zodat dit apparaat niet als een open
boek uitgelezen of aangepast kan worden. Het is natuurlijk niet evident
om dit te voorkomen, net vanwege de aard van het beestje. IoT-apparaten
werken dagen, weken, maanden zonder dat er iemand naar om ziet.
Computers en telefoons maken deel uit van onze dagelijkse aandacht en
daar is het dus veel moeilijker voor aanvallers om ongezien, fysieke,
toegang tot te krijgen.

Computers zijn tegenwoordig goed beschermd tegen aanvallers die een
laptop stelen en er trachten op in te breken: moderne beschermingen
zoals Secure Boot (hoofdzakelijk om rootkits tegen te gaan) en
\emph{Trusted Platform Module} (TPM) voorkomen toegang door niet
geautoriseerde gebruikers tot de data op het toestel. Dit soort systemen
zit zelden in kleine, goedkope, IoT-apparaten waardoor inbreken op dit
soort apparaten dus vaak eenvoudig is. Zeker als het apparaat op de koop
toe een UART of andere onbeveiligde interface heeft die door de
fabrikant worden gebruikt om problemen op te lossen. Ook aanvallers
kunnen deze hardware interfaces natuurlijk gebruiken om bijvoorbeeld hun
eigen firmware op het apparaat te plaatsen.

\begin{tipblock}
Er is gelukkig een trend gaande waarbij de TPM-chip (of alternatief)
steeds meer z'n weg vindt in nieuwe IoT-apparaten. De initiële kost
ervan vergaat meestal in het niets in vergelijking met de schade die kan
opgelopen worden wanneer onbeveiligde toestellen (zonder TPM) worden
gecompromitteerd.
\end{tipblock}

De Mccumber kubus in gedachte nemende spreekt het natuurlijk voor zich
dat ook de data geëncrypteerd op het IoT-apparaat moet bewaard worden.

\chapter{Besluit}\label{besluit}

En zo hebben we het einde van het prille begin der cyberboswachters
bereikt. Zoals je vermoedelijk ook zelf al aanvoelt, hebben we enkel
maar het tipje van de ijsberg besproken. Als we er echter gezamenlijk in
slagen om dit tipje alvast en zelf toe te passen én aan anderen aan te
leren, dan maken we het de digitale stropers al een pak moeilijker.
Beschouw dit een beetje als een toepassing van de wet van
\emph{diminishing returns.} De primaire, eenvoudige
veiligheidstechnieken en gewoonten (goede wachtwoorden, geen default
instellingen, KISS, etc.) gaan de grootste impact hebben op je algemene
veiligheid. Alle daaropvolgende stappen, hoe belangrijk en nuttig ook,
zullen steeds een kleiner stuk van de potentiële problemen voorkomen.

\emph{Wees veilig en verspreid het woord! ;)}

\emph{Tim Dams}

\emph{2025}

\pandocbounded{
\includegraphics[keepaspectratio]{gotopolice.png}}

\chapter{Appendix 1: GDPR}\label{appendix-1-gdpr}

De General Data Protection Regulation (2018), of in het Nederlands de
\emph{Algemene Verordening Gegevensbescherming} is een Europese
richtlijn om de privacy van Europese burgers, eender waar in de wereld,
te beschermen. Of zoals de EU zelf op haar website beschrijft: ``the
thoughest privacy and security law in the world'' (wat ook effectief zo
is: het is de enige wet momenteel die erin slaagt om wereldwijd de
regels ervan af te dwingen).

De wet zelf bestaat uit 99 verschillende artikels waarin de rechten,
personen en verplichtingen van bedrijven die onder de verordening
vallen, worden uitgelegd. Ieder bedrijf, eender waar in de wereld, die
data bijhoudt van EU-burgers dient zich aan deze wet te houden op risico
van een boete als ze dit niet doet. GDPR gaat dus niet over de
verplichtingen van de personen, wel om die van de bedrijven, namelijk
ter bescherming van personen.

\begin{tipblock}
We gaan in dit hoofdstuk enkel een high-level overzicht geven van GDPR
opdat we niet verzanden in de taal der rechters en advocaten, het
\emph{legalese}.
\end{tipblock}

\section{Wat omvat GDPR?}\label{wat-omvat-gdpr}

\textbf{Het recht om persoonlijke gegevens te wissen of het recht om
vergeten te worden.} GDPR laat je toe, als EU-burger om eender welke
site die `jou kent' te verplichten alle, of specifieke, informatie over
je te verwijderen. Als je dus vindt dat Google niet moet weten waar je
huisadres is, dan heb je het recht Google te verplichten deze informatie
te verwijderen.

\begin{noteblock}
Google kreeg een tijd terug een zeer stevige boete omdat het niet inging
op de vraag van een EU-burger om \emph{vergeten te worden}. Sindsdien
zijn ze een pak behulpzamer in het naleven van de GDPR-wetgeving.
\end{noteblock}

Een ander belangrijk aspect van GDPR is het zogenaamde \textbf{recht op
overdraagbaarheid van gegevens.} Voor GDPR bestond was het soms een
helse opdracht om bijvoorbeeld van Internetprovider te veranderen.
Iedere provider werkte met eigen dataformaten waarin de klantengegevens
werden bewaard. Bij een overdracht moest dit dan altijd omgezet worden
naar het formaat van de nieuwe provider. GDPR verplicht dat dit soort
data overdrachten van persoonsgegevens nu volgens een vast formaat
moeten gebeuren. En belangrijker: deze overdracht moet zo transparant én
zo geautomatiseerd mogelijk gebeuren. Hierdoor kan de data quasi
ogenblikkelijk in de database van de ontvanger geïntegreerd worden en
zal jij als klant dus veel sneller de overstap kunnen maken zonder alle
bijhorende administratieve rompslomp en vertragingen.

\begin{noteblock}
De voorganger van de GDPR was de DPD (Data Protection Directive). Dit
was geen Europese wet (verordening) maar een aanbeveling en was een stap
in de goede richting maar ging nog niet ver genoeg. GDPR heeft ervoor
gezorgd dat individuele burgers geen speelbal meer zijn van de grote
bedrijven die de data van burgers als het nieuwe goud verzamelen en
ermee doen wat ze zelf willen.
\end{noteblock}

Als derde belangrijke aspect voorziet GDPR \textbf{de plicht om niet
noodzakelijke gegevens te laten wissen}.

\section{Hoe beschermt GDPR mij als
burger?}\label{hoe-beschermt-gdpr-mij-als-burger}

Zoals je uit voorgaande 3 aspecten al merkt staat de (Europese) burger
voorop, niet het bedrijf. GDPR is er om de burger te beschermen tegen de
bedrijven die data van hen verzamelen. Het stopt echter niet bij deze 3
aspecten. Volgende rechten worden door GDPR beschermd en zijn minstens
even belangrijk:

\begin{itemize}
\tightlist
\item
  Recht op \textbf{inzage} (\emph{transparantie}): je hebt het recht om
  te weten welke data een bedrijf over jou heeft.
\item
  Recht op \textbf{rectificatie}: je hebt het recht om foute data te
  laten corrigeren.
\item
  Recht op \textbf{beperking} van verwerking: je hebt het recht om te
  eisen dat een bedrijf stopt met het verwerken van jouw data. Ook mogen
  bedrijven je persoonsgegevens niet oneindig lang bewaren. En moeten er
  duidelijke richtlijnen zijn over hoe lang ze de data mogen bewaren.
\item
  Recht op \textbf{verzet}: je hebt het recht om je te verzetten tegen
  het verwerken van je data. Uiteraard is een bedrijf niet verplicht om
  jou dan gebruik te laten maken van hun diensten.
\item
  Recht op \textbf{geautomatiseerde besluitvorming}: je hebt het recht
  om te weten wanneer een bedrijf beslissingen neemt op basis van
  geautomatiseerde processen.
\item
  Recht op \textbf{gegevensportabiliteit}: je hebt het recht om je data
  in een leesbaar formaat te ontvangen.
\item
  Recht op \textbf{gegevensminimalisatie}: bedrijven mogen enkel de data
  bewaren die ze nodig hebben voor hun dienstverlening.
\end{itemize}

\section{Wat beschermt GDPR?}\label{wat-beschermt-gdpr}

De GDPR wet zorgt ervoor dat bedrijven veel bewuster met gebruikersdata
omgaan. Ze kunnen namelijk verantwoordelijk gehouden worden indien
gebruikersdata gestolen of gelekt wordt. De GDPR is een
privacy-wetgeving in de eerste plaats: het wil de privacy van het
individu beschermen. Volgende data valt dan ook onder de GDPR-wetgeving:

\begin{itemize}
\tightlist
\item
  Persoonlijke identificeerbare informatie: namen, adressen,
  geboortedatums en rijksregisternummers.
\item
  Op Internet gebaseerde gegevens: zoals gebruikerslocatie, IP-adres en
  cookies.
\item
  Gezondheids- en genetische gegevens.
\item
  Biometrische gegevens (denk maar aan je irisscan of vingerafdruk).
\item
  Raciale en/of etnische gegevens.
\item
  Politieke meningen.
\item
  Seksuele geaardheid.
\end{itemize}

\section{Persoonlijke data buiten EU}\label{persoonlijke-data-buiten-eu}

Volgende flowchart toont wat bedrijven, binnen en buiten de EU, met
gebruikersdata mogen doen.

\begin{figure}
\centering
\pandocbounded{\includegraphics[keepaspectratio,alt={Bron: privacy.ucsd.edu/gdpr/}]{gdpr/gdprflow.jpg}}
\caption{Bron: privacy.ucsd.edu/gdpr/}
\end{figure}

\section{Meldplicht datalekken}\label{meldplicht-datalekken}

Bedrijven die onder de GDPR-wetgeving vallen (alle bedrijven die data
van Europese burgers bewaren vallen hier onder, ongeacht hun locatie in
de wereld) hebben een meldplicht indien zij een datalek hebben.

\begin{noteblock}
Onder een datalek verstaan we \emph{``Een inbreuk op de beveiliging die
per ongeluk, of op onrechtmatige wijze leidt tot de vernietiging, het
verlies, de wijziging of de ongeoorloofde verstrekking van of de
ongeoorloofde toegang tot doorgezonden, opgeslagen of anderszins
verwerkte gegevens.''}
\end{noteblock}

Van zodra het bedrijf weet heeft van een (mogelijk) datalek dienen zij
dit ogenblikkelijk te melden:

\begin{itemize}
\tightlist
\item
  Aan de \textbf{toezichthoudende autoriteit}: ieder land heeft z'n
  eigen autoriteit en in België is dat de \textbf{CBPL}, de
  \emph{Commissie voor de Bescherming van de Persoonlijke Levenssfeer}.
  Deze commissie zal iedere datalek \emph{case-by-case} beoordelen om te
  zien of het bedrijf een GDPR-overtreding heeft begaan of niet.
\item
  Aan \textbf{ieder individu waar de data betrekking op had}: als de
  gelekte data ook maar op één manier kan gelinkt worden aan een
  individu, dan dient deze van de lek op de hoogte gebracht te worden.
\item
  Dit dient \textbf{binnen de 72 uur na ontdekking van het datalek te
  gebeuren}.
\end{itemize}

Merk op dat tegenwoordig een datalek zelden zo snel wordt ontdekt, dat
beseft de GDPR ook. Daarom is de regel dat 72 uur ingaan vanaf het
moment dat het bedrijf weet heeft van het datalek. Het is dus perfect
mogelijk dat het datalek reeds maanden eerder plaatsvond maar dat het
bedrijf het nu pas ontdekt.

\subsection{Uitzonderingen meldplicht
individu}\label{uitzonderingen-meldplicht-individu}

Bij een datalek zijn er enkele mogelijke uitzonderingen waarom het
bedrijf niet noodzakelijk het individu op de hoogte moet stellen van het
datalek:

\begin{itemize}
\tightlist
\item
  Indien het bedrijf kan aantonen dat de gestolen data zodanig beveiligd
  is dat deze onbruikbaar is. Als dus de data bijvoorbeeld geëncrypteerd
  is dan zou het bedrijf dit als argument kunnen gebruiken om de
  meldplicht te verzaken. Een kanttekening is hier natuurlijk op z'n
  plaats: er bestaat altijd nog de mogelijkheid dat de data finaal nog
  kan gedecrypteerd zal worden en er dus alsnog private gegevens in
  verkeerde handen komen.
\item
  Stel dat het bedrijf de persoonlijke voorkeuren van de gebruikers in
  een database heeft staan en de identificeerbare persoonsgegevens in
  een andere. Als nu blijkt dat enkel de database met persoonlijke
  voorkeuren werd gehackt, dan kan deze data nooit gelinkt worden aan
  personen en is er dus ook geen sprake van een privacy schending.
\end{itemize}

\begin{tipblock}
Het is sowieso altijd een goede gewoonte om het adagio \emph{``don't put
all your eggs in one basket''} te hanteren. Het is wijzer om de data
over verschillende databases te bewaren. Zo voorkom je dat bij een
datalek automatisch steeds alle data wordt gestolen.
\end{tipblock}

\begin{warningblock}
Indien er een erg grote hoeveelheid mensen moeten ingelicht worden na
een datalek, dan mag het bedrijf ook beslissen om in de plaats daarvan
een publieke aankondiging te doen. Uiteraard zijn de bedrijven hier
minder happig op omdat dit sowieso een PR-nachtmerrie is (iedere
aankondiging van een datalek is dat trouwens).
\end{warningblock}

\section{Boetes}\label{boetes}

Als er uiteindelijk toch een inbreuk op de GDPR-wetgeving wordt
vastgesteld dan zal er dus een boete opgelegd worden afhankelijk van de
ernst van de inbreuk:

\begin{itemize}
\tightlist
\item
  Minder ernstige inbreuken: boete tot € 10 miljoen, of 2\% van de
  wereldwijde jaaromzet van het bedrijf uit het voorgaande boekjaar,
  afhankelijk van welk bedrag hoger is.
\item
  Meer ernstige inbreuken (inbreuken die ingaan tegen de principes van
  het recht op privacy en het recht om te worden vergeten, die de kern
  vormen van GDPR): boete van maximaal € 20 miljoen, of 4\% van de
  wereldwijde jaaromzet, afhankelijk van welk bedrag hoger is.
\end{itemize}

Of er een boete zal zijn en hoe groot deze dan is, is gebaseerd op een
aantal criteria:

\begin{itemize}
\tightlist
\item
  Ernst en aard van de overtreding.
\item
  Intentie: is het datalek het gevolg van een slordigheid of opzettelijk
  gebrek aan beveiliging?
\item
  Schadebeperkende omstandigheden.
\item
  Genomen voorzorgen.
\item
  Geschiedenis van overtredingen.
\item
  Medewerking: in hoeverre werkt het bedrijf mee met de toezichthoudende
  autoriteiten?
\item
  Gegevenscategorie.
\item
  Kennisgeving: heeft men zich aan de meldplicht gehouden?
\item
  Certificering: heeft het bedrijf in het verleden reeds via een externe
  audit aangetoond aan de GDPR wetgeving te voldoen?
\item
  Verzwarende of verzachtende factoren.
\end{itemize}

\begin{tipblock}
Merk dus op dat het schenden van de GDPR wetgeving niet automatisch in
een boete resulteert. Als het bedrijf kan aantonen dat ze echt alles in
het werk hebben gesteld om de data te bewaren zoals een goede huisvader
betaamt, dan kan het zijn dat er geen boete zal volgen.
\end{tipblock}

\begin{noteblock}
De voorbije jaren (mei 2018 tot eind 2020) werden reeds 272 miljoen euro
aan boetes geïnd ten gevolge van GDPR overtredingen. De wet is zo
effectief, dat ze ook ondertussen dient als inspiratie wereldwijd
wanneer landen of groeperingen nieuwe privacy-wetgevingen willen
construeren.
\end{noteblock}

\begin{tipblock}
In de lente van 2025 kreeg TikTok een boete van 530 miljoen opgelegd
voor het overtreden van de GDPR-wetgeving. Het bedrijf kon bijvoorbeeld
niet aantonen dat de data van de Europese gebruikers adequaat beveiligd
werd op de servers in China. We moeten er wel bij vertellen dat de boete
gebaseerd is op de situatie voor 2023. Vermoedelijk én hopelijk heeft
TikTok ondertussen zijn beveiliging op orde.
\end{tipblock}

\section{Conclusie}\label{conclusie}

We kunnen dus stellen dat GDPR een zeer mooi waardevol initiatief is,
dat ook verder gaat dan ``niet toegepaste wetten''. De manier waarop
Europese burgers hun data momenteel wordt beschermd is een voorbeeld
voor de rest van de wereld. Het is een wet die de privacy van de burger
voorop stelt en bedrijven verplicht om hiermee rekening te houden.

Het is nu vooral uitkijken of Europa een soortgelijk voorbeeld kan
stellen met de Digital Services Act (DSA) en de Digital Markets Act
(DMA) die momenteel in de maak zijn. Deze wetten zullen de macht van de
grote techbedrijven zoals Google, Facebook en Amazon moeten inperken en
de privacy van de Europese burger nog verder moeten beschermen. Ook de
wereld van artificiële intelligentie zal hierdoor grondig veranderen: de
AI Act die sinds augustus 2024 gefaseerd wordt ingevoerd zal de manier
waarop AI wordt gebruikt in Europa grondig veranderen. Uiteraard, en dat
geldt voor alle wetgevingen, is het de uitvoering en de handhaving van
de wetten die het verschil zullen maken én wat de impact ervan zal zijn
in de concurrentie met de rest van de wereld.

\chapter{Appendix 4: Meer weten}\label{appendix-4-meer-weten}

De smaak te pakken? Volgende online bronnen beveel ik aan om te volgen
zodat je op de hoogte blijft van wat er allemaal gebeurt in de wereld
van digitale stropers en boswachters.

\section{Nieuws en blogartikels}\label{nieuws-en-blogartikels}

\begin{itemize}
\tightlist
\item
  \href{https://www.darkreading.com/}{DarkReading.com}
\item
  \href{https://blog.fox-it.com/}{Fox It}: zeer in-depth en technisch.
\item
  \href{https://grahamcluley.com/}{Graham Cluley blog}
\item
  \href{https://krebsonsecurity.com/}{Krebs on Security}
\item
  \href{https://www.troyhunt.com/}{Troy Hunt}: de blog van ``Have I been
  pwnd''-oprichter Troy Hunt.
\end{itemize}

\section{Podcasts}\label{podcasts}

\begin{itemize}
\tightlist
\item
  \href{https://darknetdiaries.com/}{Darknet Diaries}: waargebeurde
  ``horror-verhalen'' over het darknet, met boeiende, in-depth
  interviews met slachtoffers én aanvallers. Ik kan deze podcast niet
  genoeg aanbevelen!
\end{itemize}

\section{Tutorials}\label{tutorials}

\begin{itemize}
\tightlist
\item
  \href{https://www.hacksplaining.com/lessons}{Hacksplaining}
\item
  \href{https://book.hacktricks.xyz/}{HackTricks book}
\end{itemize}

\chapter{Appendix 2: Dark web}\label{appendix-2-dark-web}

Het internet is vergelijkbaar met de echte wereld: het heeft goed
verlichte plekken die iedereen kan zien, maar ook verborgen
achterbuurten. Het Darkweb is zo'n verborgen gedeelte van het internet.
Het maakt deel uit van het zogenaamde ``Deep Web'', dat gedeelte van het
internet dat niet via standaard zoekmachines zoals Google te vinden is.
Denk hierbij aan databanken, afgeschermde websites en persoonlijke
accounts die niet publiek toegankelijk zijn.

Binnen dat Deep Web bestaat een extra verborgen deel, het Darkweb,
bekend om anonimiteit en privacy. Dit deel van het internet is alleen
toegankelijk via speciale software zoals TOR (The Onion Router).

\begin{figure}
\centering
\includegraphics[width=0.7\linewidth,height=\textheight,keepaspectratio,alt={De internet ijsberg}]{darkweb/ijsberg.jpg}
\caption{De internet ijsberg}
\end{figure}

\section{Kenmerken van het Darkweb}\label{kenmerken-van-het-darkweb}

Het Darkweb onderscheidt zich van het reguliere internet door een aantal
belangrijke kenmerken:

\textbf{Anoniem:} Gebruikers blijven anoniem doordat hun verbinding via
meerdere tussenliggende servers loopt. Hierdoor is het zeer moeilijk te
achterhalen wie een bepaalde website bezoekt of een bepaalde boodschap
heeft verzonden.

\textbf{Beveiligd:} Communicatie en verkeer op het Darkweb zijn
versleuteld met sterke encryptie. Dit voorkomt dat derden kunnen
meekijken of informatie onderscheppen.

\textbf{Gedecentraliseerd:} Er is geen centrale autoriteit of punt van
controle op het Darkweb. Dit maakt het moeilijk voor overheden en
instanties om toezicht te houden, te censureren of sites uit de lucht te
halen.

Zoals elke technologie kan ook het Darkweb gebruikt worden voor goede én
slechte doeleinden:

\begin{itemize}
\tightlist
\item
  Slecht:

  \begin{itemize}
  \tightlist
  \item
    Handel in illegale goederen zoals drugs, wapens en gestolen
    informatie.
  \item
    Verhuur van criminele diensten (``rent-a-hacker'').
  \end{itemize}
\item
  Goed:

  \begin{itemize}
  \tightlist
  \item
    Bescherming van privacy en anonieme communicatie.
  \item
    Platforms voor journalisten, klokkenluiders en activisten om veilig
    te communiceren.
  \end{itemize}
\end{itemize}

\section{Waarom bestaat het Darkweb?}\label{waarom-bestaat-het-darkweb}

Het bestaan van het Darkweb roept veel vragen op over privacy en
veiligheid. Aan de ene kant biedt het Darkweb een veilige ruimte voor
mensen om zich vrij en anoniem te uiten. Dit kan van groot belang zijn
voor journalisten, activisten en klokkenluiders die in landen leven met
sterke censuur en toezicht (zoals China, Iran of Noord-Korea). Zo kan
anonimiteit een cruciaal hulpmiddel zijn voor vrije meningsuiting en
democratie.

Aan de andere kant trekt anonimiteit ook mensen aan die het internet
gebruiken voor illegale activiteiten, zoals handel in drugs, wapens,
gehackte gegevens, of gestolen creditcardgegevens.

Het Darkweb is dus niet per definitie illegaal, maar wat je erop doet
kan dat wel zijn. In sommige landen, zoals Rusland en China, is het
gebruik van het Darkweb zelf al strafbaar.

\section{Toegang tot het Darkweb:
TOR}\label{toegang-tot-het-darkweb-tor}

De meest voorkomende manier om het Darkweb te bezoeken is via het
zogenaamde TOR-netwerk (The Onion Router). Dit netwerk leidt jouw
internetverkeer via meerdere relais-servers om jouw identiteit te
verhullen. Een website op het Darkweb herken je aan het
``.onion''-adres, wat alleen via een TOR-browser te openen is.

\begin{figure}
\centering
\includegraphics[width=0.7\linewidth,height=\textheight,keepaspectratio,alt={Er worden letterlijk lagen, zoals een UI, tussen jou en je doel opgebouwd.}]{darkweb/relais.jpg}
\caption{Er worden letterlijk lagen, zoals een UI, tussen jou en je doel
opgebouwd.}
\end{figure}

\section{TOR heeft ook nadelen:}\label{tor-heeft-ook-nadelen}

\textbf{Snelheid}: Het TOR-netwerk is relatief traag door de vele
tussenstations.

\textbf{Blokkades}: Sommige websites blokkeren verkeer vanaf
TOR-servers.

\textbf{Firewall}: Veel bedrijfsnetwerken en firewalls blokkeren verkeer
van en naar het TOR-netwerk.

Er bestaan ook browsers, zoals Brave Browser, die TOR al geïntegreerd
hebben, waardoor het gebruik eenvoudiger wordt.

\section{Waarom moet je dit nu
weten?}\label{waarom-moet-je-dit-nu-weten}

Het is belangrijk om inzicht te hebben in de risico's en kansen die het
Darkweb biedt. Denk hierbij aan:

\begin{itemize}
\tightlist
\item
  Het begrijpen van criminele fenomenen zoals online drugshandel,
  cybercrime en mensenhandel.
\item
  Het analyseren van bedreigingen en risico's voor organisaties en
  individuen.
\item
  Het begrijpen van het belang van privacy, vrije meningsuiting, en de
  complexiteit van online toezicht en censuur.
\end{itemize}

Het Darkweb maakt dus deel uit van het bredere veiligheidsvraagstuk
binnen de samenleving.

\begin{noteblock}
Vanwege tijdsgebrek werd deze appendix over het Darkweb grotendeels
gegenereerd m.b.v. ChatGPT, gebaseerd op de slidedeck.
\end{noteblock}

\chapter{Appendix 4: Cybersecurity Awareness
creëren}\label{appendix-4-cybersecurity-awareness-creuxebren}

In dit hoofdstuk bekijken we de belangrijke rol die
\textbf{cybersecurity-awareness} speelt binnen organisaties. Als
cybersecrurity-professionals zijn we ons bewust van de risico's, maar
hoe zorgen we ervoor dat ook onze collega's deze risico's begrijpen en
ernaar handelen? Dit hoofdstuk biedt inzicht in de noodzaak van
awareness, de rol van medewerkers, en praktische tips om een cultuur van
cybersecurity te bevorderen.

\section{Waarom awareness?}\label{waarom-awareness}

De voorbije hoofdstukken hebben hopelijk al heel duidelijk gemaakt dat
er zich nog te vaak cybersecurity-incidenten voordoen. Veel van die
incidenten hadden eenvoudig kunnen worden voorkomen als werknemers (en
gevers!) ook maar enige, minimale cybersecurity-awareness-training
hadden gehad. Uit onderzoek blijkt dat ongeveer \textbf{90\% van alle
beveiligingsincidenten veroorzaakt wordt door menselijke fouten}, bewust
of onbewust. Denk hierbij aan het klikken op phishing-links, het
gebruiken van zwakke wachtwoorden, of het verbinden met onveilige
wifi-netwerken. Deze menselijke factor maakt awareness onmisbaar binnen
een robuust beveiligingsbeleid. Die laatste zin willen we nog even
benadrukken: \textbf{bewustzijn is onmisbaar binnen een robuust
beveiligingsbeleid}. Het is niet voldoende om alleen technische
maatregelen te nemen; medewerkers moeten ook begrijpen waarom deze
maatregelen belangrijk zijn en hoe ze deze kunnen toepassen in hun
dagelijkse werkzaamheden.

\subsection{Jouw rol binnen een
organisatie}\label{jouw-rol-binnen-een-organisatie}

Jouw taak als cybersecurity expert gaat verder dan enkel systemen te
configureren en onderhouden. Jij vormt ook de brug tussen IT en
bedrijfsbeleid. Jij bent als het ware een
\textbf{bewustzijnsambassadeur} die complexe technische informatie
vertaalt naar begrijpelijke taal voor collega's.

Daarnaast heb je ook een signalerende rol:

\begin{itemize}
\tightlist
\item
  Identificeren van potentiële risico's.
\item
  Rapporteren van afwijkingen of verdachte gebeurtenissen, en dit op een
  snelle en duidelijke manier.
\end{itemize}

\subsection{Bescherm de hele keten: McCumber
Cube}\label{bescherm-de-hele-keten-mccumber-cube}

In hoofdstuk 2 bespraken we reeds de McCumber cube. Deze visualisatie
helpt ons om in te zien dat we cybersecurity vanuit al z'n dimensies en
facetten moeten bekijken. Hierbij waren de voorbije hoofdstukken vooral
gericht op het technologie-aspect in functie van C.I.A en de staat van
de data (processing, storage, transmission). In dit hoofdstuk gaan we
daarom dieper in op de \textbf{menselijke} dimensie van cybersecurity.
Dit is een cruciaal aspect dat vaak over het hoofd wordt gezien, maar
dat net zo belangrijk is als de technische maatregelen die we nemen. De
menselijke factor is immers de zwakste schakel in de beveiligingsketen.
Daarom is het essentieel dat we ons richten op het creëren van een
cultuur van cybersecurity-awareness binnen onze organisatie.

\section{Awareness-campagnes}\label{awareness-campagnes}

Om de medewerkers in een bedrijf te trainen in cybersecurity, zijn er
verschillende methoden en technieken beschikbaar. Het doel van deze
trainingen is om medewerkers bewust te maken van de risico's en hen te
leren hoe ze veilig kunnen werken. Dit kan door middel van e-learning,
workshops, of zelfs gamification. Het is belangrijk dat deze trainingen
regelmatig worden herhaald, zodat medewerkers op de hoogte blijven van
de laatste ontwikkelingen en trends op het gebied van cybersecurity.

Meestal gaan we dit doen aan de hand van, terugkerende,
\textbf{awareness-campanges}.

\subsection{Gebruikersverwarring
voorkomen}\label{gebruikersverwarring-voorkomen}

Losse technische maatregelen zijn niet voldoende. We moeten ook zorgen
dat medewerkers begrijpen waarom deze maatregelen belangrijk zijn en hoe
ze deze kunnen toepassen in hun dagelijkse werkzaamheden. Dit vereist
een combinatie van technische kennis, communicatievaardigheden en
leiderschap.

Medewerkers hebben namelijk vaak vragen zoals:

\begin{itemize}
\tightlist
\item
  ``Waarom mag ik geen publiek wifi gebruiken?''
\item
  ``Waarom moet ik steeds opnieuw inloggen met multifactor-authenticatie
  (MFA)?''
\item
  ``Waarom moet ik mijn laptop vergrendelen?''
\end{itemize}

Het doel van awareness-campagnes is medewerkers zelfstandig antwoorden
op deze vragen te laten geven. Niet door ze regels op te leggen, maar
door ze inzicht te geven in de risico's en gevolgen van bepaalde
gedragingen. Deze aanpak geldt trouwens eender waar in een omgeving waar
er regels, wetten en procedures zijn. Het is niet voldoende om enkel te
zeggen ``dat mag niet''. We moeten ook uitleggen waarom iets niet mag en
wat de gevolgen zijn van het negeren van deze regels.

Effectieve awareness-campagnes leggen dus niet enkel uit \textbf{wat} je
moet doen, maar vooral ook \textbf{waarom}:

\begin{itemize}
\tightlist
\item
  Gebruikers leggen zelfstandig de link tussen risicogedrag en de
  mogelijke gevolgen.
\item
  Ze maken bewust de keuze voor veilige werkmethoden.
\item
  Veilig handelen wordt een routine, intrinsiek gemotiveerd door begrip.
\end{itemize}

Uiteindelijk willen we dat medewerkers autonoom zijn:

\begin{itemize}
\tightlist
\item
  Ze herkennen en rapporteren nieuwe risico's zelfstandig.
\item
  Ze kunnen geleerde principes toepassen op onbekende situaties.
\end{itemize}

\subsection{Succesvolle
awareness-methoden}\label{succesvolle-awareness-methoden}

De mate waarin een campagne werkt hangt van vele factoren af. Het is
belangrijk dat de campagne dan ook binnen de context van het bedrijf
past: er is geen ``one size fits all''. Hoe beter de campagne is
afgestemd op het bedrijf en de medewerkers, hoe effectiever de campagne
zal zijn. Gebruik met andere woorden realistische scenario's en
voorbeelden die aansluiten bij de dagelijkse werkzaamheden van de
medewerkers. Dit maakt het makkelijker voor hen om de informatie te
begrijpen en toe te passen in hun eigen werk.

Alhoewel er tal van bedrijven zijn die awareness-campagnes aanbieden,
zijn er ook een aantal goede gratis bronnen beschikbaar. Denk hierbij
aan de website van de \textbf{SANS Institute} of de
\textbf{Cybersecurity \& Infrastructure Security Agency (CISA)}. Deze
organisaties bieden gratis materialen en trainingen aan die bedrijven
kunnen gebruiken om hun medewerkers bewust te maken van
cybersecurity-risico's.

Enkele effectieve methodes voor cybersecurity-awareness zijn:

\begin{itemize}
\tightlist
\item
  \textbf{Regelmatige micro-trainingen}: korte, frequente sessies. Dit
  voorkomt dat medewerkers overweldigd worden door informatie en zorgt
  ervoor dat ze de informatie beter onthouden.
\item
  \textbf{Phishing-simulaties}: directe feedback op acties zorgt voor
  snelle leerprocessen. Dit kan bijvoorbeeld door medewerkers te testen
  met nep-phishing-e-mails. Dit helpt hen om verdachte e-mails te
  herkennen en voorkomt dat ze in de val lopen van echte
  phishing-aanvallen.
\item
  \textbf{Gamification}: quizzen en uitdagingen verhogen motivatie en
  betrokkenheid. Je zou bijvoorbeeld via een dashboard kunnen tonen
  welke dienst binnen het bedrijf het beste presteert op het gebied van
  cybersecurity-awareness. Dit kan medewerkers motiveren om beter op te
  letten en hun kennis te vergroten.
\item
  \textbf{Posters en nieuwsbrieven met impact}: visuele reminders in de
  werkomgeving. Dit kan bijvoorbeeld door posters op te hangen met tips
  voor veilig werken of door het gebruik van stickers op laptops en
  andere apparaten. Dit zorgt ervoor dat medewerkers regelmatig worden
  herinnerd aan de belangrijkste principes van cybersecurity. Als deze
  posters op de koop toe grappig, of zelfs uitdagend zijn, zullen
  medewerkers ze sneller opmerken en onthouden.
\end{itemize}

\begin{figure}
\centering

\includegraphics[width=0.3\linewidth,height=\textheight,keepaspectratio,alt={Awareness posters hoeven geen ellenlange teksten te bevatten. Less is more.}]{aware/funnyposter.png}
\caption{Awareness posters hoeven geen ellenlange teksten te bevatten.
Less is more.}
\end{figure}

\begin{itemize}
\tightlist
\item
  \textbf{Leiderschap}: het management moet het goede voorbeeld geven.
  Dit kan bijvoorbeeld door zelf ook regelmatig deel te nemen aan
  trainingen of door openlijk te praten over cybersecurity-risico's. Dit
  laat zien dat cybersecurity een prioriteit is binnen de organisatie en
  dat iedereen verantwoordelijk is voor het waarborgen van de
  veiligheid.
\end{itemize}

\begin{tipblock}
Gedragsverandering is vaak de moeilijkste uitdaging binnen
cybersecurity. Positieve gedragsverandering bereik je niet door
medewerkers te straffen of te bedreigen, maar door ze te motiveren en te
inspireren.
\end{tipblock}

\section{Praktische best practices}\label{praktische-best-practices}

Doorheen het handboek hebben we al tal van praktische ``best practices''
aangereikt. We vatten ze hier nog eens samen, daar zij de
\emph{backbone} zullen vormen van de inhoud van je awareness-campagnes.
Dit zijn de basisprincipes die elke medewerker moet kennen en toepassen
in zijn of haar dagelijkse werkzaamheden. Deze principes zijn niet
alleen van toepassing op cybersecurity, maar ook op andere gebieden
binnen een organisatie. Ze helpen medewerkers om veilig en efficiënt te
werken, ongeacht hun functie of verantwoordelijkheden.

\textbf{Wachtwoordhygiëne}

\begin{itemize}
\tightlist
\item
  Gebruik sterke, unieke wachtwoorden.
\item
  Maak gebruik van een wachtwoordmanager.
\item
  Zet altijd multifactor-authenticatie aan.
\end{itemize}

\textbf{Secuur apparaatbeheer}

\begin{itemize}
\tightlist
\item
  Installeer automatische updates.
\item
  Gebruik antivirussoftware.
\item
  Vergrendel je laptop wanneer je weggaat.
\end{itemize}

\textbf{Veilig werken met documenten en netwerken}

\begin{itemize}
\tightlist
\item
  Vermijd onbekende USB-sticks en kabels.
\item
  Gebruik enkel vertrouwde wifi-netwerken.
\item
  Maak regelmatig backups van documenten, zowel fysiek als digitaal.
\end{itemize}

\textbf{Blijf alert en leer verder}

\begin{itemize}
\tightlist
\item
  Deel kennis actief binnen je team.
\item
  Blijf bij met nieuwe ontwikkelingen op het gebied van cybersecurity.
\end{itemize}

\section{Incident response}\label{incident-response}

Awareness is ook cruciaal in incidentbeheer. Indien je medewerkers
bewust zijn van de risico's en weten hoe ze moeten reageren op
incidenten, kunnen ze sneller en effectiever handelen. Dit kan de impact
van een incident aanzienlijk verminderen.

Incident response is een gestructureerde aanpak voor het afhandelen van
beveiligingsincidenten. Het doel is om de schade te minimaliseren, de
oorzaak van het incident te begrijpen, en herhaling in de toekomst te
voorkomen. Een goed \textbf{incident response-plan} omvat duidelijke
procedures, verantwoordelijkheden, en communicatiekanalen. Dergelijk
plan bestaat uit verschillende stappen, die vaak worden weergegeven in
een cyclus. De stappen zijn als volgt:

\begin{enumerate}
\def\labelenumi{\arabic{enumi}.}
\tightlist
\item
  \textbf{Prepare}: Voorbereiden op incidenten vóór ze gebeuren. In deze
  fase worden alle procedures en verantwoordelijkheden gedefinieerd. Dit
  omvat ook het trainen van medewerkers in hoe ze moeten reageren op
  incidenten. Dit is waar awareness-campagnes een belangrijke rol
  spelen. Het incident response-plan moet meer zijn dan een document dat
  in de kast ligt. Het moet een levend document zijn dat regelmatig
  wordt bijgewerkt en waar medewerkers actief mee aan de slag gaan.
\item
  \textbf{Identify}: Incidenten snel herkennen. Dit kan door middel van
  monitoring, logging, en rapportage. Het is belangrijk dat medewerkers
  weten hoe ze verdachte activiteiten kunnen herkennen en rapporteren.
  Dit kan bijvoorbeeld door het gebruik van een meldingsformulier of een
  speciaal meldpunt binnen de organisatie.
\item
  \textbf{Contain}: Beperk de schade direct na detectie. Dit kan door
  het isoleren van getroffen systemen of het blokkeren van verdachte
  accounts. Het is belangrijk dat medewerkers weten hoe ze snel kunnen
  handelen om verdere schade te voorkomen.
\item
  \textbf{Eradicate}: Verwijder het incident volledig. Vanaf deze, of de
  vorige stap zal meestal de IT-afdeling de verantwoordelijkheid nemen.
  Het is belangrijk dat medewerkers weten dat ze niet alleen
  verantwoordelijk zijn voor het melden van incidenten, maar ook voor
  het helpen bij het oplossen ervan. Dit kan bijvoorbeeld door het
  verzamelen van informatie over het incident of door het uitvoeren van
  bepaalde taken binnen de organisatie.
\item
  \textbf{Recover}: Herstel systemen en data.
\item
  \textbf{Follow Up}: Analyseer wat er is gebeurd en verbeter processen.
\end{enumerate}

\begin{figure}
\centering
\includegraphics[width=0.85\linewidth,height=\textheight,keepaspectratio,alt={Vaak wordt een incident response plan als een lineaire sequentie getoond, zonder te visualiseren dat er altijd na de followup informatie moet doorstromen naar de prepare fase. Op die manier leren we uit onze fouten, in de hoop dat deze volgende keer niet meer zullen voorkomen.}]{aware/incident.png}
\caption{Vaak wordt een incident response plan als een lineaire
sequentie getoond, zonder te visualiseren dat er altijd na de followup
informatie moet doorstromen naar de prepare fase. Op die manier leren we
uit onze fouten, in de hoop dat deze volgende keer niet meer zullen
voorkomen.}
\end{figure}

Awareness helpt vooral bij de \textbf{Identify} stap:

\begin{itemize}
\tightlist
\item
  Goede meldprocedures zijn dan ook cruciaal (wie, wat, hoe snel).
\item
  Maak meldkanalen toegankelijk en laagdrempelig.
\item
  Snelheid van identificatie bepaalt vaak de impact van een incident.
\end{itemize}

\begin{tipblock}
Het heeft geen nut om uiterst complexe incident response-plannen te
maken, als medewerkers niet weten hoe ze deze moeten toepassen. Daarom
is het belangrijk dat medewerkers goed op de hoogte zijn van de
procedures en verantwoordelijkheden binnen het incident response-plan.
\end{tipblock}

\begin{tipblock}
Het internet staat vol van nuttige hulpmiddelen om awareness campagnes
uit te voeren. Gaande van toffe posters, tot \emph{fullblown} phishing
simulators (kijk zeker eens naar \textbf{GoPhish}). Met dank aan de
immer krachtiger wordende mogelijkheden van ChatGPT en friends wordt het
ook steeds eenvoudiger om succesvolle campagnes uit te werken en te
visualiseren.
\end{tipblock}

\section{Social engineering en
phishing}\label{social-engineering-en-phishing}

De eenvoudigste manier om medewerkers te hacken is door middel van
social engineering. De aanvaller heeft hier namelijk geen tools of
technieken voor nodig, enkel een goed verhaal. Social engineering is het
manipuleren van mensen om vertrouwelijke informatie te verkrijgen of
toegang te krijgen tot systemen. Dit kan bijvoorbeeld door middel van
phishing-e-mails, waarbij de aanvaller zich voordoet als een vertrouwde
bron om medewerkers te misleiden. Het is belangrijk dat medewerkers zich
bewust zijn van deze risico's en weten hoe ze verdachte e-mails kunnen
herkennen en rapporteren.

Omdat het leuk is om lijstjes te maken, hebben we hier een lijstje met
de meest voorkomende social engineering technieken. Meestal zullen de
digitale stropers een combinatie van deze technieken gebruiken om hun
doel te bereiken.

\begin{itemize}
\tightlist
\item
  \textbf{Pretexting}: De aanvaller doet zich voor als iemand anders,
  bijvoorbeeld een collega of een IT-medewerker, om vertrouwelijke
  informatie te verkrijgen. Dit kan bijvoorbeeld door middel van een
  telefoongesprek of een e-mail. Meestal kan je dit herkennen aan de
  vraag om gevoelige informatie, zoals wachtwoorden of persoonlijke
  gegevens. Door een \textbf{challenge} te stellen, kan je de aanvaller
  vaak ontmaskeren. Vraag bijvoorbeeld om een bevestiging van de
  identiteit van de persoon die je aan de lijn hebt. Dit kan
  bijvoorbeeld door middel van een terugbelprocedure of door het
  controleren van hun e-mailadres.
\item
  \textbf{Tailgating}: De aanvaller volgt een medewerker naar binnen in
  een beveiligd gebied. Dit kan bijvoorbeeld door zich voor te doen als
  een medewerker of door simpelweg achter iemand aan te lopen. Dit kan
  je voorkomen door altijd je toegangspas te gebruiken en nooit iemand
  anders binnen te laten zonder hun identiteit te verifiëren. Het voelt
  soms lastig aan om een wildvreemde te vragen om zich te identificeren,
  maar het is beter om voorzichtig te zijn dan om een beveiligingsrisico
  te creëren.
\item
  \textbf{Vishing}: Voice phishing, waarbij de aanvaller zich voordoet
  als iemand anders via de telefoon. Dit kan bijvoorbeeld door zich voor
  te doen als een IT-medewerker of een bankmedewerker. Dit kan je
  herkennen aan verdachte vragen of verzoeken om persoonlijke
  informatie. Wees altijd voorzichtig met het delen van gevoelige
  informatie via de telefoon en verifieer altijd de identiteit van de
  persoon die je aan de lijn hebt. Doordat A.I. steeds beter wordt, is
  het ook mogelijk dat de aanvaller je stem kan imiteren. Dit maakt het
  nog moeilijker om te bepalen of je met een echte persoon of een nep
  persoon aan de lijn hebt. Meestal kan je de aanvaller ontmaskeren door
  ook nu weer een \textbf{challenge} te stellen. Je zou bijvoorbeeld een
  vraag uit het verleden kunnen stellen die enkel jij en de andere
  persoon kan weten. Als de persoon per ongeluk het juiste antwoord
  geeft dan kan je nog steeds doen alsof dit niet juist is om te zien of
  de aanvaller z'n correcte antwoord aanpast (``Uw bank is KBC'', ``Neem
  hoor (ook al is het dat wel)'', ``Oh excuses, ik bedoelde Argenta''.)
  Deze vorm van challenge wordt ook wel **reverse social engineering''
  genoemd.
\item
  \textbf{Baiting}: Deze techniek zal het slachtoffer lokken met iets
  dat hij echt wilt. Denk maar aan een USB-stick van het bedrijf, gratis
  een maand toegang tot Spotify Premium, etc. Van zodra je de
  \emph{bait} aanvaardt, zal via deze weg meestal de aanval verder gezet
  worden (denk maar een virus dat op de USB-stick staat, of de link naar
  de gratis Spotify blijkt een link naar een malware server te zijn.)
  Het paard van Troje is een mooi voorbeeld van social engineering met
  behulp van \emph{baiting}.
\item
  \textbf{Quid Pro Quo}: `Voor wat hoort wat'. De aanvaller zal een iets
  aanbieden (\emph{bait}) maar dat enkel geven in ruil voor iets anders.
  Om bijvoorbeeld de bait te krijgen moet je eerst een vragenlijst
  invullen, maar in de vragenlijst staan ook vragen die je gevoelige
  informatie zullen ontfutselen (``Hoe heette je eerste huisdier''). Dit
  kan ook door bijvoorbeeld een gratis softwarepakket aan te bieden in
  ruil voor je e-mailadres of andere persoonlijke informatie.
\end{itemize}

Het is niet altijd eenvoudig om social engineering-aanvallen te
herkennen, maar er zijn een aantal signalen waar je op kunt letten. Dit
zijn enkele tips om verdachte e-mails of telefoontjes te herkennen: *
Controleer altijd het e-mailadres of telefoonnummer van de afzender. Dit
kan je doen door het e-mailadres of telefoonnummer te vergelijken met
eerdere communicatie of door het op te zoeken op de website van het
bedrijf. * Let op verdachte links of bijlagen in e-mails. Dit kan je
doen door de muisaanwijzer boven de link te houden zonder erop te
klikken. Dit toont je waar de link naartoe leidt. Als de link er
verdacht uitziet, klik er dan niet op. * Wees voorzichtig met het delen
van persoonlijke informatie. Dit kan je doen door altijd te vragen
waarom de informatie nodig is en hoe deze zal worden gebruikt. Als je
twijfelt, deel de informatie dan niet. * Let op verdachte verzoeken of
vragen. Dit kan je doen door altijd te vragen waarom de informatie nodig
is en hoe deze zal worden gebruikt. Als je twijfelt, deel de informatie
dan niet.

Bij twijfel kan je altijd een challenge stellen. En als je nog steeds
twijfelt, doe dan een extra challenge, tot je overtuigd bent. Je zal
verrassend vaak merken dat de aanvaller niet voorbereid is op deze extra
vragen. En als blijkt dat de tegenpartij wel degelijk legitiem is, dan
kan je altijd nog excuses aanbieden. Dit is een kleine prijs om te
betalen voor de veiligheid van jezelf en je organisatie.

\chapter{Appendix 5: Bronnen}\label{appendix-5-bronnen}

\begin{itemize}
\tightlist
\item
  802.11 Wireless Network (first edition), Gast, O'Reilly, ISBN:
  0-596-0018-3
\item
  Meerdere boeken van William Stallings (Computer Security,
  Cybersecurity, etc.)
\item
  \href{https://apwt.gitbook.io/software-security/https/001introductioncryptography}{Cursus
  Software Security van M. Boeynaems.}
\item
  Afbeelding '' Worm Propagation Model'' komt uit: Fink, Glenn \& Ball,
  Robert \& North, Chris \& Jawalkar, Nipun \& Correa, Ricardo. (2004).
  Network Eye: End-to-End Computer Security Visualization.
\item
  De uitleg en tekeningen i.v.m. rainbow tables is gebaseerd op de
  uitstekende uitleg op
  \href{http://kestas.kuliukas.com/RainbowTables/}{kestas.kuliukas.com/RainbowTables/)}.
\end{itemize}

\backmatter
\end{document}
